\quellseite{826}
Schatten und in der Wüste geschehen sind, besonders aber, daß der HErr ihn von Angesicht kannte, das heißt, der so vertraut mit GOtt umging und redete. Dies wird nicht für alle Zeiten gesagt, sondern, so lange das Amt Mosis dauerte. Denn alle, die nach ihm kamen, lehrten Mosis Lehre (Mosen), als Kleinere, die unter ihm und von ihm gelehrt waren. Aber zugleich steht das fest, was Cap. 18 gesagt ist, daß nach Mose ein anderer sollte erweckt werden, der Mosi gleich wäre, nämlich Christus, dem dieser große Moses weichen sollte, da er weitaus größer ist an Würde und Kraft, wie wir dort gesehen haben. Es bedeutet auch das, daß nichts Größeres gelehrt und vorgetragen werden kann, soviel Gesetze betrifft, als das Gesetz Mosis. Denn alles ist in ihm das Höchste, ausgenommen, daß das große Gesetz dem größeren Evangelium weichen sollte. Ende.

Ende.

\chapter[{B. D. Martin Luthers Auslegung über etliche Capitel [Cap.}]{B. D. Martin Luthers Auslegung über etliche Capitel [Cap. 1. 4–9] des fünften Buchs Mosis\footnote{Wie wir aus den ersten Worten dieser Schrift ersehen, predigte Luther fortlaufend über alle fünf Bücher Mosis. Am 13. December 1528 stand er bei dem 31. Capitel des vierten Buchs Mosis (vergleiche Col. 1362 dieses Bandes), während der Monate April und Mai 1529 erlitten Luthers Predigten eine Unterbrechung wegen seiner Krankheit (De Wette, Bd. III, S. X), nach seiner Rückkehr vom Gespräch zu Marburg, Mitte October 1529, stand Luther bei dem 7. Capitel des fünften Buchs Mosis (vgl. in dieser Schrift Cap. 6, § 103 und Cap. 7, § 9 ff.), und am vierten Adventsonntage, den 19. December 1529, war er bei dem Schluß des 9. Capitels angelangt (Eislebensche Ausgabe, Th. 1. Bl. 555). Darauf wurden die Predigten über dieses Buch nicht mehr fortgesetzt, „denn (so sagt Murisacher) der Reichstag zu Augsburg Anno 1530 darauf gefolget ist“. Erst nach Luthers Tode wurde unsere Schrift, welche die Auslegung von 5 Mos. 1–9. enthält (mit Ausnahme von Cap. 2 und 3, welche Luther überging, weil sie das im dritten und vierten Buche Gesagte wiederholen), im Jahre 1564 von Murisacher aus Nachschriften des M. Georg Hörer, M. Anton Lauterbach, Superintendenten zu Pirna, und Philipp Fabricius, gewesenen Pfarrherrn zu Ringelheim, veröffentlicht. Sie findet sich zuerst in der Eislebenschen Ausgabe, Bd. I, Bl. 484; darnach in der Altenburger, Bd. IV, S. 707; in der Leipziger, Bd. IV, S. 202 und in der Erlanger, Bd. 36, S. 164. Wir geben den Text nach der Eislebenschen Ausgabe.}}

Gepredigt zu Wittenberg 1529.

\section{Vorrede über dieses Buch}

1. Wir haben nun vier Bücher Mosis, des Propheten, bis anher gehandelt, in welchen wir gesehen, wie er dem alten jüdischen Volke viel hübsche Gesetze gegeben hat, und dasselbige Volk geordnet in zwei Reiche, als, geistlich und leiblich. Darnach haben wir auch von vielen schönen Exempeln [des] Glaubens und Unglaubens, Gehorsams und Ungehorsams der Altväter und Patriarchen gehört, was und wie es denen widerfahren ist, die solche Gesetze und Ordnungen übertreten haben; die denn uns sollen ein Exempel sein, darnach wir uns auch richten, und ihrem Glauben nachfolgen, und vor dem Unglauben uns hüten sollen, auf daß wir auch nicht in das Unglück geraten, darein sie gekommen sind. Denn solches alles uns zum Vorbilde und Exempel vorgeschrieben ist, auf daß wir nicht in ein gottlos Wesen geraten möchten; wie solches St. Paulus zum Corinth ern [1. Ep. 10, 6.] vermahnt.

2. Nun wollten wir gerne auch das fünfte Buch Mosis euch vortragen, auf daß ihr den ganzen Mosen gehört hättet. Das fünfte Buch Mosis aber ist nicht anders, denn als eine lange Predigt, damit er vor seinem Absterben das Volk gesegnet hat, und alles auf das kürzeste gefasset, was zuvor von ihm gelehrt worden, und von allerlei Geschichten, die sich zugetragen hatten, auf daß sie es sollten beschrieben mit sich nehmen in das Land Canaan, darein sie ziehen würden, und solcher Predigten und Geschichte nicht vergessen, sondern ihr Lebenlang lesen,
\quellseite{827}
sonderlich aber alle sieben Jahre dem ganzen Volk Israel auf dem Fest der Laubhütten vorlesen lassen, wie im 31. Capitel [V. 10. 11.] dieses Buchs geschrieben steht. Darum er auch das Beste aus den andern Büchern in dieses zeucht, und gleichsam ein enchiridion oder compendium, einen kurzen Auszug und Summarien macht über die vier vorhergehenden Bücher, und das Uebrige läßt fahren, daß einer den Deuteronomium wohl möchte nennen eine weitläufige Predigt, darin er reichlich ausgestrichen oder ausgelegt habe die zehn Gebote, und darunter vornehmlich das erste Gebot reichlich tractirt. Denn das ist der Brunn und die Hauptquelle aller Gebote, Gesetze und Künste, und was aus diesem Gebote fleußt und wieder hineingeht, das ist recht.

3. Darum darf es keiner Auslegung der zehn Gebote im Alten Testamente, denn das hatte Moses selber mit einem eigenen besondern Buche gethan. So nennen wir nun das fünfte Buch Mosis eine weitläufige Erklärung der zehn Gebote; und weil ihr in den ersten vier Büchern gehört eine kurze, schlechte Erklärung der zehn Gebote, und sonst auch des Jahres vielmal hört die zehn Gebote predigen, so wollen wir jetzt vor uns nehmen diese weitläufige Auslegung, auf daß ihr könnet sehen, wie ferne und weit sich ihr Verstand erstrecke. Das sei nun die Vorrede, daß wir dies Buch wollen vor uns nehmen als eine reichliche Auslegung, sonderlich des ersten Gebots. Wir können doch nichts Nützlicheres predigen, denn von GOtt, es sei nun gleich von seinem Gesetze und Geboten, oder von seinem Evangelio und Verheißungen.

\section{Das erste Capitel}

\begin{schriftzitat}
V. 1. 2. Das sind die Worte, die Moses redete zum ganzen Israel, jenseit dem Jordan, in der Wüste, auf dem Gefilde, gegen dem Sumpf, zwischen Paran und Tophel, Laban, Hazeroth und Disahab, elf Tagereisen von Horeb.
\end{schriftzitat}

1. Alhier zeigt er an, an welchem Orte Moses diese Predigt gethan habe, und ist so viel gesagt: Diese Predigt oder dies Buch ist geredet worden jenseit dem Jordan; denn Moses ist über den Jordan nicht gekommen. Der Jordan theilt aber das Land in zwei Theile: gegen Morgen war Moses, da er die zwei Könige schlug, 4 Mos. 21, 24. 35. 5 Mos. 2, 33. 3, 3. Und das ist der Ort gewesen. Erstlich waren sie am Berge Sinai, da sie das Gesetz empfingen; wie im andern Buch Mosis am 20. Capitel geschrieben steht. Von demselbigen Orte, bis hieher gegen KadesBarnea, sind elf Tagereisen, das sind nach unserer Rechnung bei dreißig Meil Wegs; denn sie machten nicht weite Tagereisen, darum, daß sie mit allem ihrem Hausgesinde, allerlei Habe und Vieh nicht eilen konnten.

2. Dieses ist nun nichts Sonderliches für die Einfältigen, denn uns ist nichts daran gelegen, allein daß man wisse, wo dies Buch gemacht sei. Ueber der ersten Tagereise sind sie gezogen vierzig Jahre. Die Ursache aber, warum sie also lange umgezogen, habt ihr gehört. Sie haben elf Tage zu reisen gehabt, von Sinai bis ins gelobte Land, und haben gleichwohl müssen vierzig Jahre darüber zubringen; nämlich, daß sie murreten wider GOtt, ungehorsam und ungläubig waren, und sich wider setzten Gottes Worte, wollten nicht GOtt folgen, sperreten sich wider Mosen: da wurden sie alle gestraft, und ward ihr Ziel verkehrt, daß aus elf Tagen vierzig Jahre worden, bis alle streitbaren Männer starben, die aus Egypten gezogen waren, ausgenommen Josua und Caleb, so in das gelobte Land kamen von demselben ganzen Haufen; die andern, so mittlerzeit gezeuget und auferzogen worden, mußten also in der Irre laufen, und da in der Wüste bleiben.

3. Nun, was dies Exempel bedeutet, habt ihr gehört. Moses zeucht stets an für ein groß Exempel, daß er vierzig Jahre hat damit umgehen müssen, daß doch hätte können in elf Tagen geschehen, als wollte er sagen: Das soll man nicht in Wind schlagen und vergessen, sondern allezeit in frischem Gedächtnisse behalten, davon singen und sagen, auf daß sich ein jeder daran stoße und fromm sei. Denn es ist ein
\quellseite{828}
schrecklich Exempel wider alle, die GOtt ungehorsam werden und wider ihn murren, daß sich dieselbigen vorsehen, damit sie nicht auch also gestraft werden, als die Israeliter gestraft sind.

\begin{schriftzitat}
V. 3. Und es geschah im vierzigsten Jahr, am ersten Tage des elften Monden, da redete Mose mit den Kindern Israel.
\end{schriftzitat}

4. Diese Worte zeigen an, zu was Zeit Moses dies Buch geschrieben, nämlich nach vierzig Jahren, da sie aus Egypten gekommen waren, nicht über zwei Monden zuvor, ehe denn Moses stirbt; daß also das Buch gleich als Mosis Testament ist. Es sind seine letzten Worte und Predigten, damit er das Volk gesegnet, und ihnen zur Letzte GOttes Gesetz und Gebot wieder vorbildet, so er allhier kürzlich zusammen gefaßt, und gleich als auf einen Kneuel\footnote{Kneuel = Knäuel.} gewunden, oder in ein Büchlein gefaßt hat; item, GOttes Wohlthaten wiederholt und vorhält, welche er ihnen bewiesen, da er sie als ein gnädiger Vater aus Egypten errettet und geführt hatte. Hält ihnen auch vor allerlei Verheißung, so widerfahren sollen denen, die GOttes Gebot halten; dazu ernstliche Dräuung, wie GOtt die Uebertreter seiner Gebote strafen wolle. Sonderlich aber predigt er im 18. Capitel [V. 15.] von Christo, welchen er abmalt als einen Lehrer und Prediger, den alle Welt hören solle etc.

\begin{schriftzitat}
V. 5–8. Jenseit des Jordans, im Lande der Moabiter, fing an Mose auszulegen dies Gesetz, und sprach: Der HErr, unser GOtt, redete mit uns am Gebirge Horeb, und sprach: Ihr seid lange genug an diesem Berge gewesen, wendet euch und ziehet hin, daß ihr zu dem Gebirge der Amoriter kommet, und zu allen ihren Nachbarn im Gefilde, auf Bergen und in Gründen, gegen Mittag, und gegen der Anfurt des Meers, im Lande Canaan, und zum Berge Libanon, bis an das große Wasser etc.
\end{schriftzitat}

5. Da seht ihr, daß dies Buch ist eine Auslegung des Gesetzes und Erzählung der vorigen Geschichte. Moses fähet an die Gesetze zu glossiren, die er gegeben hatte, und spricht: „Der HErr, unser GOtt, redete mit uns am Berge Horeb (das ist Sinai) und sprach: Ihr seid nun lange genug an diesem Berge gewesen.“

6. Das ist der Anfang dieser Predigt; als sollte er sagen: Lieben Kinder, ihr wißt, daß wir am Berge Sinai waren, und empfingen das Gesetz; da gab GOtt Befehl, daß wir straks in das Land Canaan zögen, das er unsern Vätern verheißen hatte, und nach dem Befehl sollten wir bald in das gelobte Land gekommen sein. Das zeigt Mose nicht umsonst an, sondern daß er sie erinnere, was sie so lange habe aufgehalten und gehindert am Einzuge, nämlich ihre eigene Bosheit und Schuld, daß sie GOtt sind ungehorsam gewesen; sonst wären sie längst hinein gekommen. Es hat sich wohl lange verzogen, aber jetzt wird es ein Ende sein; darum ich euch unterrichten will, wie weit und breit das Land sei, darinnen ihr wohnen sollet. Wie er denn hiemit die Landart beschreibt.

7. Libanon ist ein groß Gebirge, und größer denn der Thüringer Wald. So weit soll euer Land sein, bis an den Abbrat, an Libanon, und an die Anfurt des Meers. Das war die erste Anweisung. Aber sie wurden Uebertreter, und schickten Kundschafter hinein, das Land zu besehen, die erschreckten sie, daß sie nicht fort wollten, dieweil sie sagten: Wer will da hinein kommen, es wohnen Giganten drinnen, die haben starke und feste Städte; wir haben Riesen drinnen, gegen welchen wir wie Heuschrecken anzusehen sind; darum wollen wir nicht hinein. [Sie] versäumten sich also selbst an der Verheißung, daß sie mußten noch ganzer vierzig Jahr draußen bleiben. Diese Sünde rührt hier Moses, und straft den Unglauben und Mißtrauen an den Kindern Israel.

\begin{schriftzitat}
V. 9–12. Da sprach ich zu derselben Zeit zu euch: Ich kann euch nicht allein ertragen, denn der HErr, euer GOtt, hat euch gemehrt, daß ihr heutiges Tages seid wie die Menge der Sterne am Himmel. Der HErr, euer Väter GOtt, mache euer noch viel tausend mehr, und segne euch, wie er euch geredet hat. Wie kann ich allein solche Mühe, Last und Hader von euch ertragen?
\end{schriftzitat}

8. Wir haben gesehen, wie Moses anfähet eine lange Erzählung der großen Sünden, um welcher willen sie haben müssen so lange in der Wüste bleiben und umherziehen; und die Ursache setzt er mit vielen Worten, warum sie das gelobte Land nicht haben einnehmen können, nämlich darum, daß GOtt sie geheißen hatte in das Land ziehen, und sie hatten nicht gewollt.

\quellseite{829}
9. Danach setzt er diesen Text. Den müssen wir ansehen, darum, daß er gar eine feine Lehre ist für die weltliche Obrigkeit. Also spricht er: Ehe wir auszogen, faßte ich das Volk mit einem Regiment. Er redet nicht von dem geistlichen Regiment oder vom Priesterthum, davon im vierten Buche geschrieben ist, sondern von der weltlichen Obrigkeit, daß er nahm Häupter aus den Stämmen, und setzte die über tausend, über hundert, über fünfzig, und über zehn. Und da kommen auch die alten Wörterlein her, Dechant, der über Zehn gesetzt zu regieren; wie wir auch sagen: Viermann, Dreimann, Zweimann, dort Zehnmann, Hundertmann,\footnote{In den Ausgaben: „zehn Mann, hundert Mann, tausend Mann und fünfzig Mann“; doch vermuthen wir nach dem Zusammenhange, daß decurio, centurio, etc. gemeint sei. Vergleiche die folgenden Paragraphen.} Hundertmann, Tausendmann und Fünfzigmann etc.

10. Also ordnete er das Volk, und ist eine sehr feine Ordnung. Aus diesem Texte ist das Matth. 2, 6. auch genommen: „Und du Bethlehem bist mit nichten die geringste unter den Fürsten Juda“; da Michäas spricht Cap. 5, 1.: Unter den Tausenden Juda, das ist, unter so vielen, als mag ein Tausendmann\footnote{Erlanger: Tausend Mann.} regieren. Diese Worte des Propheten sind auch hieraus gezogen. So ferne ist\footnote{Im Original: hat.} diese Ordnung gegangen bis auf tausend; darnach aus jedem Stamme ein Fürst, der hat über die tausend Mann regiert; und über alle Fürsten war denn zuletzt Moses. Es ist über die Maßen eine feine Ordnung gewesen.

11. Zum ersten ist ein Dechant, das ist, zehn Mann tragen es an Einen; denn Ein Mann kann wohl auf zehn sehen; darnach der Hundertte\footnote{„Hundertte“ steht hier für: Hundertmann.} über hundert. Es ist aus der Maßen eine schöne Weise, liebliche und feine Ordnung zu regieren. Die Regierung ist damals nicht so schwer gewesen als jetzt, da es alles vermischt und in einander gemengt ist. Es ist dergleichen Ordnung jetzt nicht in der Welt, wird auch nicht wieder angerichtet werden. In den Klöstern war sie nicht. Wenn die Fünfzigmänner nicht haben etwas entschließen können, da ist die Sache vor die Hundertmänner\footnote{Im Original: „mach“ statt: männer.} gebracht. Jetzt soll es ein Fürst allein alles ausrichten. Dies zeigt der Text hier an, da Moses spricht: „Ich kann euch nicht allein ertragen.“ Die Noth trieb ihn dazu, und sein Schwäher Jethro unterrichtete es ihn, 2 Mos. 18, 19. ff., daß er sprach: „Wie kann ich allein solche Mühe, Last und Hader von euch ertragen?“

12. Nimm nun den Text, und laß uns ihn wohl beherzigen und behalten, wider alle stolzen und hoffärtigen Köpfe, die da gerne regieren wollen; denn ich wollte, daß GOtt denselbigen genug zu regieren gäbe. Es ist ein toll Ding um einen solchen, der gerne regieret, denn er gedenkt: Moses ist ein Fürst, sitzt obenan; also auch: Ein Fürst hat ein trefflich, herrlich Schloß, trägt Edelsteine, güldene Ketten, Sammet; und sperrt das Maul auf, hält solches denn gegen seine Armuth, geringes Häuslein oder Strohdach, und denkt, ein Fürst habe es viel besser denn er. Also steht ein Narr in das Regiment.

13. Aber höre, wie Moses drein sieht, der spricht allhier: „Ich kann euch nicht allein ertragen; wie kann ich alleine solche Mühe und Last und Hader von euch ertragen?“ Als sollte er sagen: Es sei der Teufel ein Regent! ist doch nichts mehr denn Mühe und Arbeit darinnen. Je mehr Gewalt, je mehr Sorge und Bekümmerniß. Ich rede jetzund vom Regiment an ihm selbst. Moses sagt, es sei Mühe, Last und Hader darinne, er habe nichts davon, denn daß er von einem jeglichen müsse getrieben und geplagt werden, will er anders ihrer los werden vom Halse. Sieht aber ein Regent durch die Finger, und läßt das Böse und die Laster ungestraft, so ist er lieb gehalten. Derhalben sind das große Narren, die nur ansehen die große Pracht und Ehre, schöne Kleider und güldene Ketten der Fürsten und Herren.

14. Wenn sie aber betrachteten, was an denselben Ketten geschrieben wäre von Sorge, Arbeit, Gefahr und Mühe, so würden sie sagen: Und wenn sie schon eitel Gold hätten, so wollen wir unser Hüttlein oder Strohdächlein dafür nehmen, und jenen sein großes Schloß behalten lassen, und sagen: Behüte mich GOtt vor deinen güldenen Ketten und Perlen. Denn sie sind also in der Obrigkeit Stand gesetzt, daß ihnen der Kitzel wohl vergeht. Denn wer ein frommer Fürst sein will, wenn ihm die Welt gleich viel Geschenke gäbe, so ist doch seine Arbeit nicht bezahlt; es gehört ein größerer\footnote{Im Original: groß.} Lohn
\quellseite{830}
dazu. Es ist eine große Sache, regieren, und kein Weiser begehrt es. Viel ein größer Ding ist es um das geistliche Regiment oder um die Seelsorge.

15. Dies sehen die Bauern auch nicht also an; sondern wie droben [§ 12] gesagt, so thun sie die Augen zu gegen die Last und Mühe, so weltliche und geistliche Regenten ausstehen müssen, meinen schlecht, regieren sei nichts anders, denn marmorne Schauben tragen; aber die Last und Unlust von den Leuten und dem Teufel, welchen man muß zum Feinde haben; item, GOttes Gerichte, so sie fürchten müssen, sehen sie nicht. Darum sagten die Bauern in dem Aufruhr: Wir wollen auch marmorne Schauben und güldene Ketten tragen, und Rebhühner fressen. Aber es ist kein gleich gemacht auf Erden. Sage mir, wie kann einer ein besser, ruhlicher\footnote{Im Original: rüglicher.} und friedlicher Leben haben und führen, denn eines Bürgers oder Bauers Leben ist?

16. Also pflegte unser Churfürst, Herzog Friedrich, kurz vor seinem Tode zu sagen, der auch aller Stände Leben wohl erfahren hatte, und sprach: das bürgerliche Leben oder das Bauerleben wäre das beste Leben, denn die Könige und Fürsten hätten nur den Glanz und Schein, die Unterthanen aber das Gold. Wie so? Die Bürger und Bauern sitzen daheim sicher, hinter dem Ofen, bei ihrem Acker, Haus und Hof, das Ihre ist wohl in gutem Frieden bewahrt. Was, meinst du, daß [es] ein kleiner Schatz sei, einen Gulden im Beutel haben, item, einen grauen Rock, oder ein schweißig Hemde antragen, oder in einem frohen Hause mit Frieden wohnen? Es ist ja viel besser denn des Fürsten Schatz, und seine sammeten Schauben oder güldenen Ketten.

17. Herzog Johannes, Churfürst, noch kein Fürst hat einen Gulden in seiner Verwahrung mit Frieden; wenn ein Krieg ist in andern Landen, da muß er sorgen, daß solcher Krieg nicht auch in sein Fürstenthum komme. Was helfen ihn seine güldenen Becher, wenn er ihrer nicht brauchen oder genießen kann? Das darfst du gar nicht, sondern lebst, als wärest du Herr im Lande, und Herzog Johannes muß dein Knecht sein und dich beschützen, daß du Frieden und gut Gemach habest, und dein Brod mit Frieden essest. Ja, er muß (mit Urlaub zu reden) noch wohl einem jeden den Hintern wischen. Das sieht oder erkennt niemand, sondern allein sieht man dahin: Du ich habe ein Haus mit Stroh gedeckt, muß einen leinenen Kittel tragen, auf dem Heu liegen etc., und siehst dagegen nicht, daß du hinter dem Ofen liegst, bist sicher, und hast dein geruhiges Leben im Frieden, das doch kein Fürst hat.

18. Da solltest du bedenken, daß dein grauer Rock köstlicher ist denn seine marmornen Schauben und dergleichen köstliche Tracht; so kannst du auch stille sitzen, und nicht ein jeder dich unter die Füße treten und unterdrücken. Hier schreibt irgendes dem Churfürsten ein anderer Fürst einen unnützen Brief; dort gibt ihm ein Höfling einen Stich; da mag er sonst nicht einen guten Bissen essen mit Frieden. Er hat viel Sorge, ja, viel mehr Sorgen und Kümmerniß in seiner marmornen Schaube, als du graue Faden in deinem Rocke hast. Du bist in guter Ruhe mit deinem Weibe und Kindern, und trinkst deinen Kofent sicherer, denn er seinen Malvasier trinkt: hast du nicht Frieden, so ist es deine eigene Schuld.

19. Darum hat auch GOtt geboten, daß man die Obrigkeit, Fürsten und Herren ehren solle, und für sie beten, wie St. Paulus 1 Tim. 2, 1. 2. sagt: „So ermahne ich euch nun, daß man vor allen Dingen zuerst thue Bitte, Gebet, Fürbitte und Danksagung für alle Menschen, für die Könige und für alle Obrigkeit, auf daß wir ein geruhlich und stilles Leben führen mögen, in aller Gottseligkeit und Redlichkeit“ etc. Und Jeremias, der Prophet, vermahnt die Juden zu Babel in dem\footnote{Im Original: „der“ statt: „in dem“, was wir aus der alten Ausgabe herüber genommen haben.} Gefängniß, daß sie für den König von Babel fleißig bitten sollen; denn wenn es ihm wohl ginge, so würden es die Juden auch gut haben, Jer. 29, 7. Darum soll man wider sie nicht murren, sie schmähen und lästern; denn wir haben mehr Güter denn sie, wiewohl es nicht scheint. Denn wir haben mehr Frieden in unserm Armuth, denn sie in ihrem Reichthum. Und wähle allhier, ob du nicht lieber wolltest einen Gulden mit Frieden haben, denn hundert mit Unfrieden.

20. Nem, sage mir, ob du nicht lieber wolltest Eine Kuh haben, und eine Schüssel voll Milch hinter dem Ofen essen, oder Birnen braten mit
\quellseite{831}
gutem Frieden, als neun und neunzig Kühe, die nicht deine wären, und müßtest derselben halben allzeit in Jahr stehen? Wirfst du nicht sagen: Ich ließe sie ein gut Jahr haben, ich will lieber mit Frieden genießen, daß mir Eine Kuh gibt. In der Fürsten Schaube steckt so manche Sorge, so manches Haar daran ist, denn er muß sich eine jegliche böse Sache daran wischen lassen. Da muß er fürchten, daß er den Fürsten nicht zum Feinde habe, oder jenen nicht irgend erzürne; dagegen zu rechnen, so ist dein grauer Rock eitel gülden Stück, um des Friedens und Ruhe willen. Die drei Stücke\footnote{Hier scheinen die Nachschriften mangelhaft zu sein. „Die drei Stücke“ sind wohl (§ 16) „ein Gulden im Beutel, ein grauer Rock, ein fröhlich Haus“ mit Frieden.} sind besser denn Gold und alle Schätze. Derhalben so schreibe über dein Haus: Das Haus des Friedens, in Stille.

21. Ich bin sehr zornig auf die Bauern, die da selbst wollen regieren, und die solchen ihren Reichthum nicht erkennen, daß sie im Frieden sitzen durch der Fürsten Hülfe und Schutz. Ihr ohnmächtigen großen Bauern und Esel, wollt ihr es nicht vernehmen? Daß euch der Donner erschlage! Ihr habt das Beste, nämlich Nutzen, Brauch, Saft aus der Weintraube, und laßt den Fürsten die Hülfen und Körner. Das Mark habt ihr, und sollt noch so undankbar sein, und nicht beten für die Fürsten, und ihnen nur nichts geben wollen? Welche aber verständige Bauern und Bürger sind, das sind seine Leute. Ich habe einen Bauern gehört, der sagte: wenn er zwo Kühe hätte, so wollte er die eine darum geben, daß er die andere möchte mit Frieden behalten und nützen.

22. Denn wenn Krieg ist, da gibst du gerne einem Herrn zehn Gulden um Frieden; jetzt aber wurfst du, wenn du sollst einen Groschen zu Zins und Geschoß geben. Darum wird euch GOtt strafen um solcher Undankbarkeit willen. Jetzt regieren und machen es die Bauern auf dem Markte, wie sie nur selber wollen, lassen sich nicht genügen, daß GOtt sie in Frieden und Ruhe läßt sitzen. Aber kommt einmal ein Unglück über euch, das wird euch lehren, was Friede sei. Gedenket daran. Jetzt sammelt nur getrost auf das Register. Alle Schuld rostet nicht. Merket dieses wohl.

23. Je geringer einer wäre, je fröhlicher er sein sollte, denn er hat das Beste, er hat das Mark und den Saft von den Gütern, die Fürsten haben nur die Hülfen. Also ist es auch mit einem Knechte, Magd, oder sonst mit einem Dienstboten. Der Hausvater ist allewege im gefährlichern und mühseligern Stande als das Gesinde, denn er muß famulus communis sein. Der Herr muß sein selbst Knecht, will er’s im Hause finden recht; die Frau muß sein selbst Magd, will sie im Hause schaffen Rath. Sie müssen die Sorge tragen, und so fortan. Je höher einer ist, je größere Unlust er hat. Wie allhier Moses solches auch sieht, und darum spricht: Ich kann nicht alle solche Mühe, Last und Hader von euch ertragen.

24. Ei, so sollte Moses nicht reden, sondern sagen: Ei, ich bin froh, daß ich vornangehe und güldene Ketten trage. Aber er spricht: Es ist eine Mühe, Last und Hader, und klagt sehr darüber. Wer hat immermehr das gehört, daß man einem Regimente oder Königreiche solche Titel solle haben gegeben? Wer hat das Buch gemacht, daß einer sich stecken soll in alle Sachen und damit umgehen? Darum ist es ein groß und gefährlich Ding um die Obrigkeit, und GOtt hat befohlen, daß man den Herren dienen soll mit allem Fleiß, und aller Gewalt oder Obrigkeit gehorsam und unterthänig sein; wie denn St. Paulus sagt, Röm. 13, 1. 2. 4.: „Jedermann sei unterthan der Obrigkeit und Gewalt; denn die Gewalt, die allenthalben ist, die ist von GOtt verordnet“, also, daß „wer sich wider Gewalt setzet, der widerstrebet GOttes Ordnung, denn sie ist GOttes Dienerin, eine Rächerin zur Strafe über den, der Böses thut“. Item, dergleichen gehet St. Petrus auch in seiner Epistel, daß man aller Obrigkeit, auch der unschlachtigen und bösen, solle gehorsam sein, 1 Petr. 2, 13. 18.

\subsection{Last}

25. Diese „Last“ hat den Papst nicht gedrückt. „Mühe“ ist im Regiment, daß einer viel zu thun hat in einer großen Gemeinde oder Haufen; da hat einer den andern gleichgeladen, der hat einem sein Weib geschändet etc. Uebel das ist auch „Last“, daß Einer muß auf sich nehmen alle Sorge, Arbeit, Mühe und Noth des Regiments, auf daß die Unterthanen Gerichte und Gerechtigkeit, Schutz und Schirm haben. Das kommt die Regenten nicht leicht noch süß an; sie sitzen da nicht auf einem sammeten Pfuhle oder Kissen; sie gehen nicht auf Rosen, sondern müssen manchen Schweiß darüber lassen.

\quellseite{832}
26. Da drehst du dich mittlerweile hinter dem Ofen um, und wendest dem Teufel einen Faulbraten\footnote{Faulbraten = Faulenzer. (Dietz.)}; gedenkst nicht an die Sorge deiner Obrigkeit; lebst, als wärest du ein Junker; gehst auf deinen Acker, und in deinen Garten; indes muß Moses und andere Obrigkeit sorgen, rathschlagen, wie er dich beschütze, und dein Knecht werde. Sonst solltest du wohl zehn Wächter haben, die dich und dein Haus bewahreten, und dir deiner Schafe hüteten, und solltest dennoch wohl keines behalten; ja, du solltest wohl hundert Gulden für ein gut Geleite geben über Land. Aber das muß der Fürst alles versorgen. Siehe, also ist „Mühe“ und „Last“ bei einander; nicht allein Mühe vorhanden, die Sachen zu bestellen, sondern darnach auch Last, dieselben auszurichten.

27. Zum dritten kommen die Haderer. Das sind eitel Drachen; unter denen ist keine Lust, sondern Herzeleid und Hölle. Das fühlen alle Obrigkeiten und frommen Hausväter wohl, was es für ein groß Ding sei, andern dienen oder regieren. Ich wollte wohl ein Exempel geben, und euch erklären, was die drei Dinge, „Last“, „Mühe“ und „Hader“, heißen; dann würdet ihr in einen Fürsten sehen, gleichwie in die Hölle; wie Moses auch allhier thut.

28. Dieses habe ich darum gehandelt und ausgelegt, daß ich euch doch einmal abziehen möchte von der leidigen Undankbarkeit wider die Regenten und Oberherren, und ihr nicht allein den Nutzen und das Gute an ihnen schauet, sondern auch betrachtet, was für eine Gefahr Regieren in sich habe. Drehe das Blatt um, und siehe an das Gute, das du hast, und das Unglück, das ein Regent hat, dann wirst du die Hände aufheben, und GOtt dafür danken, und dir gerne an deinem Glücke und Stande genügen lassen, und das Schieferdach gegen deinem strohernen Häuslein nichts achten. Tolle Heiligen sind mir, die wir über einem kleinen Geschwär oder Schwulst an unserm Leibe ungeduldig werden, und gedenken nicht, daß dagegen unser ganzer Leib stark und gesund ist. Ja, man sollte dir es bestellen, GOtt wird dir es nach deinem Willen machen.

29. Es hat’s GOtt also gemacht, daß ein jeglicher Stand sein Ungemach soll tragen; und sollen nicht die Unlust, sondern auch das Gute ansehen in unsern Ständen. Und merke das: Je geringer ein Amt oder Stand ist, je lieber und besser ist derselbige Stand. Darum spricht der Heide Sophocles: Wenn aller Menschen Unglück würde zusammen getragen auf einen Haufen, daß der Richter käme, und es zugleich austheilete, so würde ein jeder doch sprechen: Gib mir mein Unglück wieder, ich will lieber das meine tragen, denn daß es sollte gleich ausgetheilet werden. Es ist gleich genug getheilt; es mangelt nur an dem, daß man es recht\footnote{Im Original und in der Erlanger: „nicht“ statt: „recht“.} erkenne. Aber darum, daß [es] nicht geschieht, so sagt uns die heilige Schrift, daß, wenn die Obrigkeit, Fürsten und Herren schon Hals und Bauch dahin gäben für die Unterthanen, so hätte es doch der unvernünftige, undankbare Esel, der gemeine Mann, bald vergessen. Also thut man: Der Bauer sieht an das Kleine, daß er Zins geben muß etc., und bedenkt nicht den großen Frieden, so er hat, und die schwere Last und Mühe, so die Obrigkeit trägt für seinen Zins und Geschoß. Die aufrührischen Bauern wollten es machen, daß kein Bauer sollte kein Ungemach haben; aber wäre das Ungemach aus einem Orte gewichen, so wäre es doch an tausend Orten wieder angerichtet worden.

30. Also hat Moses geklagt, daß er nicht könne allein die Last und Mühe im Volk mit der Regierung ertragen. Daran wir lernen, welch ein groß, schweres und gefährliches Amt es ist, wenn ein Mensch das andere\footnote{Hier ist „Mensch“ sächlichen Geschlechts; in der Erlanger unrichtig: „das Andere“.} soll regieren, es sei gleich im geistlichen oder weltlichen Regimente, daß es des Fürsten, Vater und Mutter Stand, oder mein, des Predigers Amt ist, es sei nun, in welchem Stande es wolle. Darum will GOtt die Obrigkeit auch geehrt, und nicht beleidigt noch geschmäht haben, sondern, daß man ihnen gehorsam und unterthan ist.

\begin{schriftzitat}
V. 13. Schaffet her weise, verständige, erfahrene Leute unter euren Stämmen, die will ich über euch zu Häuptern setzen.
\end{schriftzitat}

31. Es sind noch alles Vorreden, so vor der Auslegung der zehn Gebote hergehen, wie wir hören werden im dritten und vierten Capitel, die da ordnen das Volk in zwei Reihe. Nun,
\quellseite{833}
da steht der Text: „Gebet mir her weise, verständige und erfahrne Männer“ etc. Wo nimmt man sie? Wo will man sie finden, wie er allhier Leute zum Regiment erfordert? Drei Stücke setzt er, so zum Regenten gehören. Solche Regenten und Häupter sollen unter ihnen sein.

32. Es ist dieser Text ein Muster und Regel aller derer, die da sollen zu Aemtern geboren werden. Wo es aber also gethan ist, daß die Leute geboren sind zum Amte, und die Regierung ihnen angeerbt wird, wie es geschieht bei Fürsten und Herren, und wie Vater und Mutter den Kindern eine Haushaltung und Regiment hinterlassen, da muß man Acht haben, daß sie recht und wohl auferzogen werden; da liegt die größte Kunst an dem, daß ein junger Regent wohl erzogen werde. An diesem Orte aber redet er von denen, die man erkieset und erwählt, daß sie sollen Fürsten, Häupter, und also fortan Regenten sein. Die sollen weise, verständige und erfahrne Leute sein. Und hernach [V. 17.] wird er ihnen eine Regel setzen, daß sie richten ohne Ansehen der Person, und daß sie recht richten. Ursache, denn das Gericht ist GOttes, 2 Chron. 19, 6. 7.

33. So gehört nun zu einem Amtmanne, erstlich, daß er „weise“ sei, das ist, gottesfürchtig, gottselig. Denn der 111. Psalm, V. 10., sagt: „Die Furcht des HErrn ist der Weisheit Anfang.“ Und darnach, daß er auch wisse, was sein Amt betreffe, und habe die Rechte und Gesetze, die dazu gehören, daß er es nicht erst lernen müsse, und seinem Amte nicht vorstehe, daß er also sitze als ein Stock oder Klotz. Denn Narren soll man nicht über Eier setzen, sie zerbrechen dieselbigen. Darum gehören zum Fürsten- und Regierstande weise und kluge Leute. Das ist stracks beschlossen, welche die Welt sollen regieren. Gewalt und Macht will es nicht thun, sondern Weisheit.

34. Es ist noch nie kein Reich mit Gewalt bestanden, sondern es muß sich mit Weisheit schützen. Wird man das römische Reich mit Gewalt regieren, so wird es ein Wesen sein, wie im wilden Walde unter den unvernünftigen Thieren, da ein Thier das andere frißt; da frißt der Fuchs den Hasen, der Wolf den Fuchs, den Wolf pfleget dann der Bär zu zerreißen. Also würde es auch mit gewaltsamer Regierung unter den Menschen zugehen. Denn so würde der, so stärker wäre, sich erheben mit seiner Macht über den Schwächern, und würde ihn unterdrücken; alsdann wäre ein ewiges Blutvergießen, Zank, Hader und Aufruhr in der Welt. Der Fürst ist eine einzelne Person, und hat einen großen tollen Haufen, bei welchem größere Macht ist, denn bei dem Fürsten, seinem Reich, oder Regimente. Darum soll das Regiment nicht sein, wo die größte Macht ist, wie die aufrührischen Bauern meinten. So spricht nun allhier Moses: Schaffet mir Leute her etc., die weise, verständige und erfahrne Leute sind; nicht die starke Fäuste haben, nicht die da Rotten sind, die das Schwert allein führen wollen, sondern es gehören weise Leute dazu, die da regieren nach Gesetzen und Ordnung, die ihnen vorgestellt sind.

35. Dergestalt hatte auch Herr Antonius Teucher, Losungsherr zu Nürnberg, Kaiser Maximiliano, hochlöblicher Gedächtniß, geantwortet, als er daselbst einmal war eingeritten, und eine große Welt Volks zugelaufen, des Kaisers Einzug zu sehen, und der Kaiser hernachher gefragt: wie man doch dies große Volk im Gehorsam erhalten könnte? Da hat Antonius Teucher, als ein weiser Mann, darauf gesagt: Allergnädigster Herr Kaiser, durch gute Worte und harte Strafe. Und ist weißlich geredet, denn mit dem Worte, mit Gesetzen und Weisheit muß man regieren. Darum auch im geistlichen Regimente herrscht Christus durch das Wort, wie im 110. Psalm, V. 2., gesagt wird: „Das Scepter deines Reichs wird ausgehen aus Sion“, das ist, das Evangelium soll aus Sion in die Welt gepflanzet und gepredigt werden. Also muß man das Wort auch gebrauchen im weltlichen Regimente, denn da geht durch die Weisheit alles besser fort, denn durch Gewalt; die Weisheit ist Kaiserin im Regimente, daß sie ihre Rechte und Gesetze kenne, und nach denselben\footnote{Im Original: demselben.} richte und urtheile.

36. Darnach sollen sie auch „verständig“ sein, das ist im Hebräischen Nabon, das heißt ein fleißiger Mann, der Acht darauf hat, daß [es] so gehe, wie ihn seine Weisheit lehrt. Denn also geht es zu im Regimente, daß einer kann übereilt werden durch seine Affecten, als, mit Zorn, Neid, Hoffahrt, Gunst, und nach seinem Kopfe fahren; und wenn das Geblüt so wüthet, kommt es oft, daß einer der Gesetze und Rechte
\quellseite{834}
vergißt, regieret nach seinem Kopf, und sagt darnach, er richte nach den Gesetzen, wie zornige, grimmige und nachgierige Richter thun, bei denen Fleisch und Blut oben schwebet. In unserm geistlichen Regimente spüren wir es auch fein, da die Secten und Rottengeister auch predigen gleich sowohl als wir; aber sie predigen nach ihrem Kopf und Affecten, und verbergen und verhehlen doch ihre Affecten mit etlichen Sprüchen der heiligen Schrift. Dazu gehört denn ein Verstand, der ihn nicht lasse die Weisheit nehmen durch Hoffart, Neid, Gunst und dergleichen, sondern bleibt bei der Weisheit, und nach derselben Recht richtet er.

37. Weisheit und Verstand unterscheide ich also: Weisheit ist eine gemeine Erkenntniß des Rechten, wie man regieren soll. Klugheit aber, oder Verstand, ist ein eigener Hüter der Weisheit, daß man sich nicht durch Affecten lasse einnehmen, dadurch man oft betrogen wird. Denn wenn ein weiser Regent über das auch scharf ist, genaue Achtung auf die Händel gibt, so sind so viel Einreden, Fälle, Klagen, Bitten etc. nicht so sorglich, und geht einem Amtmanne nimmer ein Wort zum Ohr ein, darauf er nicht zum fleißigsten mercke, und mit seiner Weisheit scharf darauf sehe. Dieweil ein jeglicher, der vor den Richter kommt, schmückt seine Sache aufs beste, als er nur kann, und will eine jede Partei recht haben, und können also den Richter leichtlich betrügen. Darum ist es verordnet, daß ein Regent die Parteien fleißig verhöre, wie man im Sprichwort sagt: Eines Mannes Rede, eine halbe Rede. Derhalben soll man sie verhören beide, denn ein jedes Theil wollte den Richter gern einnehmen, und auf seine Seite bringen.

38. Derhalben bedarf ein weiser Mann wohl gutes Verstandes, daß er sich durch Klugheit vorsehe. Das ist mir zum öfternmal begegnet, und widerfährt gewißlich auch denen, welchen Sachen zu verhören befohlen, daß die Parteien den Richter betrügen. Es fehlt uns schier am meisten an der Klugheit, Witz und Schärfe, daß einer nicht seine Weisheit lasse zur Thorheit machen durch seine eigenen Feinde und böse Mäuler. Darum steht im Alten Testamente fein geschrieben: Man soll Eines Mannes Wort allein nicht hören. Und ein Mensch hat nur Einen Mund und zwei Ohren, daß ein Regent mit einem Ohr den Kläger, und mit dem andern den Verklagten höre; wie der mächtige König Alexander Magnus mit einem Ohr allezeit den Kläger gehört, und das andere Ohr zugehalten, des Beklagten Entschuldigung und Verantwortungen auch anzuhören.

39. Denn die Welt ist so gar verzwieselt arg und böse. So einer eine Sache hat, der kommt gemeiniglich also, daß er des Richters beide Ohren einnehme, schmückt und rechtfertigt sich und seine Sache; wenn aber das Gegentheil kommt, so befindet sich’s viel anders. Viel mehr geschiehet es aber, wenn sich einer selbst betrügt durch seine Weisheit. Darum ist vonnöthen, daß er verständig sei, und nicht sicher hingehe in seiner Weisheit; niemand glaube, und verlasse sich zu weit auf seine Weisheit, sondern examinire wohl alle Umstände. So ist nun Klugheit eine wackere Weisheit, die auf die Schanze siehet.

\subsection{Die unter euren Stämmen bekannt sind}

40. „Bekannte Leute“ sind die, so ein gut Gerücht haben in der Stadt, und feine ehrbare Leute sind, eines guten Wandels; denn er will nicht böse Buben haben. Und dieser Punkt ist so gut, als die vorigen zwei. Es wäre gut, daß solch Regiment in der ganzen Welt wäre, daß man keinen im Regimente litte, auch nicht in einerlei\footnote{einerlei = einigerlei, irgend welche.} Gemeinschaft aufnähme, es wäre denn, daß er ein gut Zeugniß habe. Denn da kommt sonst alle Plage her, daß man dieser Regel nicht gefolgt; daß man, sonderlich unter dem Pabstthum, alle Fremden und Unbekannten in die Städte zu Bürgermeister[n] und Rathsherren aufgenommen. Niemand hat an das Wort gedacht, daß hier Moses braucht, „bekannt sein“. Man sollte keinen in einer Stadt leiden, er wäre denn in der Stadt bekannt, und habe ein gut Zeugniß von seinen Nachbarn; und wenn er das schon alles hätte, darf es GOttes und seiner Gnade noch, daß es wohl gerathe.

41. Man darf nicht Unglück suchen, und einen Unbekannten zum Regenten wählen, man nehme dazu die Klügsten und Bekanntesten, es soll denn noch werden, wie GOtt will, und darf dennoch wohl Glücks. Man soll den Teufel nicht an die Wand malen, oder zu Gewatter bitten, er kommt von sich selber. Regieren ist ein solch Amt, das GOtt angeht; derhalben soll ein Regent bekannt sein, und ein gut Zeugniß haben. Denn der Teufel ist ein Vater der Lügen [Joh. 8, 44.], und wo er einen Regenten kann bethören
\quellseite{835}
, so thut er es nur gern, wenn er auch der Klügste ist. Es greife nur ein jeglicher in seinen Buhlen, der im Regiment ist, und fühle, ob er nicht oft zum Narren ist worden.

42. Er soll ein gut Zeugniß haben, daß die Leute von ihm können zeugen, er sei ein kluger, weiser Mann, 1 Tim. 3, 7. Aber du möchtest sagen: Wo sind jetzt solche? Wo nimmt man sie? Moses und die Propheten sind jetzt todt, Petrus und Paulus auch todt. Wie ist es gegangen denen, die Moses selbst erwählt hatte zum Regimente, werden sie nicht darnach alle an den Galgen gehangen, 4 Mos. 25, 5., die doch waren Weise, Verständige und Bekannte? Darum so sage ich, daß Regieren und [der] Obrigkeit Stand gar ein schweres Ding ist, und ich wollte, daß einer, der gerne regiert, des Regierens genug hätte. Mit Furcht und Zittern soll man das Regieren angreifen. GOtt will solche dazu haben, die klug, weise und gelehrt sind. Wo will man sie aber nehmen? Müssen es denn so gar Weise, Kluge und Bekannte sein? Man wird sie nicht finden, die so ganz und gar rechtschaffen wären.

43. Ehrgeizige Leute lassen sich dünken, daß sie große Erfahrung haben, und sind die Allerweisesten; hauen denn andere zur Bank, schänden und lästern sie, und wollen alles ausrichten und alle meistern. Da wollen wir noch weit nicht hin. GOtt behüte uns vor den Predigern, die allen Leuten gefallen, und die von jedermann ein gut Zeugniß haben. Wenn mich alle Leute für einen guten Prediger hielten, so wollte ich nimmermehr ein Prediger sein. Ursache ist, daß, welche die Klügsten sind, können die andern behauen und reformiren. Wenn du aber sie recht ansiehst, so ist nicht ein Körnlein Weisheit in ihnen; doch geben sie Zeugniß von großer Weisheit.

44. In einer Stadt müssen auch solche Leute sein, die die andern alle richten und meistern, und sind doch selber die närrischten, oder rechte Narren in der Haut. Wer da will ein Prediger sein, der muß Richter haben, die sonst nichts zu thun haben, denn daß sie uns richten und klügeln; wenn sie es aber ein halb Stündlein versuchen sollten, so ginge es ihnen wie dem Pfeifer, der den Tanz verderbte. Also waren ihrer auch nicht viel, die etwas von Mose hielten, sondern murrten wider ihn, achteten ihn für einen wahnsinnigen Menschen, und setzten sich allezeit wider ihn, ja, wollten ihn und Aaron steinigen und erwürgen, 4 Mose am sechzehnten Capitel, V. 1. 2.

45. Darum so erwähle man zu Aemtern in einer Stadt Leute, die da fromm, geschickt, erfahren, ehrbar, Eheleute sind, die da ihre Kinder ehrlich aufziehen, auch friedliebend sind. Ob sie nicht so klug und weise als Salomon und Moses sind, da liegt nicht Macht an. Es ist genug, daß sie bei vernünftigen Leuten ein vornehmlich Ansehen haben. Dann andere, die da wollen klug sein, sind ruhmreißige, stolze Esel; und wenn es denn kommt, daß man einen Regenten erwählen soll, so nehme man den, welchen man für fromm achtet, und lasse den Esel fahren. Du wirst nicht eitel Salomones, David, Mose und Aaron finden, sondern, hast du etwa einen frommen, ehrlichen Mann, als, auf dem Dorfe oder in der Stadt, so nimm denselbigen redlichen, vernünftigen Mann, der sein Leben fein zugebracht hat, und befehl ihm ein Amt, GOtt wird zu seiner Regierung wohl Gedeihen und Segen geben. Ob solche darnach nicht klug sind als Salomon, so haben sie auch denn nicht Wenig zu regieren. Was ihnen denn fehlt und mangelt, das befehle man unserm HErrn GOtt; und so etwan andere darunter sind, die daran nicht ein Gefallen haben, die lasse man gehen.

46. Wohlan, so erwähle weise Leute auf deinem Dorfe oder Stadt zum Regiment. Du wirst bald sehen, welches ein frommer Bürger ist; man kann es aus seinem äußerlichen Wandel spüren, gleichwie ein Schulmeister eines discipuli ingenium bald erkennen und unterscheiden mag. Ich will bald sehen, was ein Prediger im Schilde führt. Es ist das Stück darum gesetzt, daß Gott will gewehret haben, daß man ja nicht Regenten wähle nach Reichthum, Geld, Gut, nach Freundschaft, Gunst oder Gewalt. Werden wir nicht nach dieser Regel im geistlichen und weltlichen Regiment wählen, so wird GOtt zum Regieren nicht Gnade geben.

47. Aber also siehet und wählet jetzt die Welt nicht. GOtt gebe, er sei ein Ketzer, unweise, Ehebrecher, oder wie er sonst wolle etc., hat er nur sonst große Freundschaft, Gunst, Reichthum, so dringet er durch. Wo bleibt da die Weisheit? O! da fragt niemand nach. Es ist recht so. Im Papstthum hat niemand Acht gehabt, ob ein Bischof gelehrt sei, sondern nur, ob er hohes
\quellseite{836}
Stammes, edel, gewaltig sei. Sollen nun im weltlichen Reiche weise Leute zur Regierung gezogen werden, wie viel mehr soll man im geistlichen Reiche nicht wählen nach dem Adel, Gut, Freundschaft, Gunst und dergleichen! Darauf soll man vornehmlich sehen, ob er ein gottesfürchtiger, vernünftiger, aufrichtiger, redlicher Mann sei. Denn auch ein Handwerksmann nimmt lieber einen Gesellen an, der seines Handwerks redlich ist, denn der reich, großer Leute Kind, und doch daneben ein Bube ist etc. Einen solchen Regenten soll man ehren und fördern, der weise, verständig und bekannt ist, ob er schon arm sei; nicht um Gunst willen. Nun, das ist dies Stück.

\begin{schriftzitat}
V. 16. 17. Und gebot euren Richtern zur selben Zeit [und sprach]:\footnote{„und sprach“ fehlt in der Eislebenschen und in der Erlanger Ausgabe.} Verhöret eure Brüder, und richtet recht zwischen jedermann, und seinem Bruder, und den Fremdlingen. Ihr sollt keine Personen im Gerichte ansehen, sondern sollt dem Kleinen hören wie dem Großen, und vor niemandes Person euch scheuen; denn das Gerichtamt ist des HErrn.
\end{schriftzitat}

48. Der HErr Christus hat die Gnade, die ihm angeboren ist, daß er sich vor niemand darf fürchten noch heucheln. Dahin hat er es gebracht, wie der Pharistäer Knechte bekannten, Matth. 22, 16., daß er ein Lehrer wäre, der von GOtt kommen, und den Weg GOttes recht lehrete, und fragete nach niemand, achtete auch nicht das Ansehen der Personen. Wir Menschen aber haben alle den Gebrechen, daß wir nach dem Ansehen der Personen urtheilen, oder nach Gunst und Furcht richten. Darum befiehlt GOtt hier den Regenten, daß sie keine Person ansehen sollen, weder groß noch klein, und sich vor niemand fürchten. Das sind eitel außerwählte, köstliche Worte, ein güldener Text; [aus] der Ursache gehört zum Regiment nicht ein reicher oder gewaltiger, sondern ein weiser und verständiger Mann.

49. Also haben die Heiden auch gesagt, als der weise Mann Bias: Wenn einer in das Regiment kommt, da sieht man, was es für ein Mann ist. Magistratus ostendit virum. Man erfährt nicht eher, was in einem Manne steckt, er komme denn in das Regiment. Zuvor ist er weise und klug, ja, klüger denn hundert Salomon, in seinem Sinne; ja, er hat in einem Finger wohl dreißig Salomon sitzen, weiß und kann eine ganze Stadt regieren; wenn er aber dazu berufen wird, taugt er weder zu sieden noch zu braten, in einer einigen Sache kann er nicht ein Haarbreit finden, das redlich und dienstlich wäre.

50. Da ist dieser Text vornöthen, daß sich ein Regent rüste und die Landstraße einhergehe, und nicht in Furcht oder Liebe, oder jemandes Person ansehe, sondern strack hindurch gehe und alle Umstände aus den Augen schlage, auch Furcht und Liebe hintan setze, und denke, als lebte kein Mensch auf Erden, der dir was schaden oder dienen könnte. Wie die Sache ist, also richte. Aber da muß ein Mannsherz zu gehören, das da fest und steif stehe, gleich als eine Steinflippe oder Fels im Meere unbeweglich steht, da alle Billigen, Wogen und Wasserwellen anschlagen, und sich an dem Fels abquetschen. Somit, liegt einer im Haß, der beugt bald das Recht; und der nach Liebe richtet, der läßt einen andern genießen der Gunst und Freundschaft etc. Und da muß das Recht eine mächtige Nase gewinnen, da nimmt man denn auch eine Handvoll Gunst, und läßt einem einen Sackvoll Kunst oder Rechtes. So gehört nun ein Mann dazu, der bestehen kann in der Mitte, und kann anhin schleichen lassen Gunst, Haß, Geschrei, Schaden, und frei sitzen\footnote{In der alten Ausgabe Walchs und in der Erlanger: Schaden und Freistiften, und :c.} und derer keines nicht achten. Aber da muß ein Mannsherz sein. Darum ist es vornöthen, daß er Gehülfen habe, die bei ihm sitzen in der Regierung, und auch treulich halten.

51. Aber wer will helfen? Da spricht Moses: „Das Gerichtamt ist des HErrn“; und drunten im 32. Capitel, V. 35., sagt GOtt auch: „Mein ist die Rache, ich will vergelten.“ Das soll euer Trost sein. In einem jeglichen Rathhause sollte dieser Text geschrieben stehen, daß der HErr spricht: Ich nehme mich deß an; das Amt, das man treibt, ist GOttes Amt, und die im Amte sitzen, sitzen an GOttes Statt, und ihr Gericht ist eben, als das GOtt selber vom Himmel richtet; wie St. Paulus zun Römern Cap. 13, 2. auch sagt: „Alle Obrigkeit ist von GOtt, und wer ihr widerstrebet, der widerstrebet GOtt selbst.“ Und Daniel spricht Cap. 2, 21. 4, 22.:

\quellseite{837}
Gott ordnet und verändert die Regimente, und gibt sie, wem er will. Und wo der Trost nicht wäre, was sollte einer machen? Das ist die Mauer, der Wall und die Wehre, die da hält das Regiment, es sei geistlich oder leiblich; sonst wäre es wider den Teufel und die Welt, so dawider stürmen und wüthen, ungehalten. Wie könnte ein Bürgermeister bestehen wider so vieler Leute Haß und Zorn, ohne dies Wort? Darum halten sie das Regiment nicht, sondern der Text, der hier geschrieben steht: „Das Richteramt ist des HErrn.“ GOtt ist das Fundament und die Grundfeste aller Gerichte und Obrigkeit.

52. Dieser Text sollte mir lieber sein, denn viertausend Trabanten, denn er ist erstlich eine Mauer und Wehre der Regimente; es soll sich auch ein Regiment darauf stützen und verlassen. Wohl fest und sperrt sich der Teufel dawider, wie oft geschehen, und wollte es gerne alles über einen Haufen werfen; doch kommt es wieder zu seinem vorigen Stande. Dieser Spruch, daß das Gericht des HErrn sei, hat bisher Könige, Kaiser, Fürsten und Herren geschützt, sonst wären sie lange alle erschlagen. Das sollte nun die Unterthanen schrecken und anhalten,\footnote{Erlanger: angehalten.} daß sie von Herzen ihren Oberherren gehorsam wären, und gedachten: Bist du ihnen ungehorsam, so greiffst du Gott ins Amt und sündigst, nicht wider den Bürgermeister oder deine ordentliche Obrigkeit, sondern wider GOtt. Darum haben die Heiden den Ungehorsam gegen die Obrigkeit geheißen crimen laesae majestatis, daß [er] eine Sünde sei wider\footnote{Erlanger: in.} die Majestät, und haben es nicht unrecht vorgegeben, auch den Aufrührern und muthwilligen Buben, so sich wider die Obrigkeit gesetzt haben, den Kopf hinweggerissen.

53. Zum andern ist der Text auch gesetzt zu Trost der Fürsten und aller frommen Obrigkeit, daß sie wissen, wo ihr Trost sei. Wenn ich ein Fürst wäre, der Text sollte mir lieber sein, denn vier hunderttausend geharnischter Trabanten; welche wir wohl auch müssen haben, aber auf diesen Text sollen wir trauen und uns verlassen, als daß GOtt das Regiment hat zu sich gefasset, und spricht: es sei sein. Und wenn einer das hört, sollte er Hände und Füße lassen fallen, und nicht Unrecht thun im Regimente. Darum haben die Fürsten keinen andern Schutz denn den HErrn, sonst müßten sie verderben; als durch der Bauern Aufruhr der Teufel vorhatte, sie von den Aemtern zu stürzen und umzubringen.

54. Wer nun ein Regent ist, der wisse, daß sein Amt GOttes Amt sei. Darum regiere auch ein jeder also, daß er niemand ansehe, und nicht aus Liebe oder Haß richte, niemand fürchte, vor niemand erschrecke, sondern spreche: „Der mir das Amt befohlen hat, der ist größer denn alle Menschen. Ihr seid zwar Könige, gewaltig, reich etc., aber es ist ein anderer, der sagt: „Das Gericht ist mein.“ Der wird mir stark genug sein. Aller Obrigkeit Trost, Trotz, Schrecken und Vermahnung ist in diesem Texte. Kurz ist’s gesetzt, aber es reicht und greift überall weit um sich. Und wer da soll oder will in der Welt regieren, der mag sich wohl freuen dieses Textes, der ihm sagt, er führe einen göttlichen Stand, Amt und Beruf, sei GOttes Diener, und richte göttliche Werke aus. In Summa, nicht einer allein, sondern alle Herrschaft überall sei in einem seligen Stande. Das ist der Text, den Moses von der Obrigkeit gehandelt hat, und der vor der Auslegung der zehn Gebote hergeht, damit er will das Volk gefasset haben in das geistliche und leibliche Regiment, auf daß sie vor allen Dingen gute Regenten und Oberherren haben, und hernach geschickter sind, die zehn Gebote zu hören.

55. Also haben wir bisher gehört aus seinen eigenen Worten, was Moses für ein Mann sei, wie er dies Buch anfähet, darinnen er will die zehn Gebote erklären, und vor diesem, wie er die Nothdurft beiderlei Regiments erwägt; nämlich, daß er erstlich, ehe denn er die zehn Gebote gibt, bestellet die Personen, Fürsten, Bürgermeister, Rathsherren oder Richter, die das Recht sollen erhalten, und gibt ihnen diese Regel, daß sie sollen weise, kluge [Leute] und bekannt sein im Volke. Und wenn sie hernach zum Amte greifen, daß sie die Augen zuthun, und keine Person ansehen, sie sei gleich gewaltig, reich, großes Standes etc., und mit dem Gerichte schnurgleich durchgehen; wie er weiter auch niemand zu Liebe, Leid, zu Freundschaft oder zu Nachtheil richten soll, sondern wissen, daß Rechte und Gerichte nicht sein sei, sondern ihm befohlen von GOtt, denn es heißt: „Das Richteramt ist des HErrn.“

56. Also hat er unterrichtet und gelehrt, wie
\quellseite{838}
die Häupter geschickt sein sollen, die da solche Gesetze und Ordnung handhaben und handhaben sollen. Denn das muß man auch haben im Regiment, daß die Personen oder Häupter dermaßen geschickt, das ist, weise, klug und bekannt seien; darnach, daß man auch gute Gesetze und Rechte habe. Der beider kann keines mangeln in einem Regiment, und wo eines mangelt, so geht es nicht recht. Habe ich fromme Richter und weise Herren, und dagegen unrechte Gesetze und Ordnungen, so taugt es gar nicht; und wiederum, wo gute Gesetze sind, und niemand ist, der sie treibe und darüber halte, daß es an Richtern und fleißiger Obrigkeit fehlt, da ist es abermals verloren.

57. Dieweil nun Moses das erste Stück gewiß hat, daß er die besten Gesetze gegeben, so ist ihm am meisten daran gelegen, daß er Leute bestelle, die da würden fleißig sein, über den Gesetzen zu halten; darum ordnet er Regenten, ehe denn er die Gesetze wiederholt und auslegt. Und ist zumal eine schöne, feine Ordnung, daß man trachte am ersten nach Leuten und guten Personen, die es thun sollen. Denn was hilft es, daß ein Haus fein angerichtet ist, und ein Edeler darnach hinein komme, und spreche: O, es ist ein alt Haus? Wenn es nicht einen Hausvater hat, der es im baulichen Wesen erhält, so zerfällt es bald; aber ein guter Hausvater kann es in einem Jahre anrichten, und wieder bauen.

58. Also liegt es an Leuten, auch im weltlichen Regimente, die es treiben. Es muß das Recht, und dann auch die executio juris sein. Und das ist die Klage der ganzen Welt, daß zu den Regimenten nicht geschickter Leute genug sind. Es wird in diesem Leben nimmer gebrechen an Geld, Gute, Städten, Gesetzen, sondern an Leuten, die die Gesetze treiben, und die da nachdrücken werden. Also wird es im geistlichen Amte auch gehen. Was hilft es, daß man das Evangelium und alle Bücher voll hat, wenn man nicht daneben christliche, gute, gelehrte und bewährte Prediger habe? Es muß wahrlich ein Bischof sein, der wohl gelehrt sei, und darnach auch fleißig über der richtigen Lehre stehe und halte, dieweil dieselbe unangefochten nicht bleibt.

59. Darum liegt es an dem, daß man gute Gesetze habe, und dann fromme Regenten, die es treiben; da wird eine Stadt und gute Polizei aus. Wo nicht weise Leute sind, da richtet man mit Gesetzen nichts aus. Denn man muß erst den Mann haben, darnach muß man ihm das Schwert an die Seite gürten. Man muß erst den Hausvater haben, ehe man das Haus anrichte. Also hat Mose dem Volke auch zuvor weise Richter und Häupter geordnet und gesetzt, ehe denn\footnote{Erlanger: denn ehe.} er ihnen das Gesetz gibt.

\begin{schriftzitat}
V. 20. 21. Da sprach ich zu euch: Ihr seid an das Gebirge der Amoriter gekommen, das uns der HErr, unser GOtt, geben wird. Siehe da, das Land vor dir, das der HErr, dein GOtt, dir gegeben hat; zeuch hinauf, und nimm es ein, wie der HErr, deiner Väter GOtt, dir geredet hat. Fürchte dich nicht, und laß dir nicht grauen.
\end{schriftzitat}

60. Als die Kinder Israel an die Grenzen des Landes Canaan kamen, zeigte Moses ihnen das Land, und befahl ihnen, sie sollten es einnehmen. Erinnert sie auch GOttes Worts, oder der göttlichen Zusagung und Verheißung, die sie von dem Lande hatten, dadurch ihr Glaube sollte gestärket worden sein, damit sie das Land desto getroster eingenommen hätten. Aber die Kinder von Israel fallen in den leidigen Unglauben, lassen GOttes Zusagung und seine vorige Hülfe und Wohlthaten fahren. Der Teufel reißt ihnen das Wort und Zusagung GOttes aus dem Herzen weg, daß sie nicht wollen glauben, sondern fürchten sich, schicken Kundschafter ins Land, so sich des Weges erkundigten, und des Landes und der Leute Gelegenheit erforschten. Als die nun wieder kamen, und ihnen den Grauen groß machten, da wollten sie nicht fort, sondern fürchteten sich, daß sie würden den Heiden zu schwach sein, und erzürnten GOtt, daß sie wieder zurück müßten, und acht und dreißig Jahr in der Wüste umherziehen, und darinnen verderben.

61. Diesen Schaden that der Unglaube, da die menschliche Vernunft und Weisheit sich von GOttes Wort nicht will leiten und führen lassen, sondern auf ihren Zaum gehen, will es alles besser wissen und machen, denn GOtt; will alles ausrechnen, ob es möglich oder unmöglich sei; steckt darnach mitten in der Furcht, und tröstet doch nicht, sondern zäumt das Pferd am unrechten Orte. Denn Vernunft, menschliche Weisheit und Rathschläge, oder Vertrauen auf menschliche Kräfte müssen fehlen, wenn sie GOtt wollen in sein Regiment greifen; drauf folgen denn
\quellseite{839}
die Strafen, wie an diesem Orte. Es müssen die Kinder von Israel zu diesemmal nicht ins gelobte Land gebracht werden. Sie fallen in GOttes Zorn und Ungnade, welcher seine Verheißung zurückzieht; sie wird also zu Wasser, wie denn im 95. Psalm, V. 11., gesagt wird: er habe ihnen geschworen in seinem Zorn, daß sie nicht sollten in seine Ruhe kommen, fallen bald darauf in Verzweiflung und in Gotteslästerung, daß sie sagen [V. 27.]: „Der HErr ist uns gram“ etc., wissen nichts mehr von den Wunderthaten GOttes, an ihnen zuvor erzeigt. Auch muß Moses entgelten solches Unglaubens der Israeliter, und nicht in das gelobte Land kommen, darum, daß er, als der Heerführer und Hauptmann dieses Volks, am Haderwasser selber auch in Unglauben fiel, um ihrer, der Israeliter, Murrens, mannigfaltiger Versuchung GOttes, Ungeduld und Unglaubens willen.

62. Und zuletzt ist GOtt auch mit leiblicher Strafe hinter den Israelitern her. Denn da die Kinder von Israel diese ihre Sünde des Unglaubens erkannten, bereueten, beklagten, und ließen es ihnen herzlich leid sein, nahmen ihre Wehre und Waffen, und wollten nun die Feinde, die Amoriter, angreifen und schlagen, wider GOttes Verbot, aus unzeitiger Kühnheit, da wurden sie von den Amoriten gejagt und schändlich geschlagen. Das war des Unglaubens Lohn und Trankgeld, welches wir uns zur Warnung mögen lassen gesagt sein, vor dem Unglauben uns zu hüten, und GOttes Wort zu glauben und zu folgen. Denn Glaube und Gottseligkeit durch GOttes Kraft alle Noth wendet, wie denn Moses allhier spricht: „Fürchte dich nicht, laß dir nicht grauen.“ Wiederum aber, Unglaube führt in alles Unglück.

63. Das andere und dritte Capitel in diesem fünften Buch Mosis wollen wir stehen lassen, dieweil es nur Wiederholungen und Repetitionen sind desjenigen, so in den zweien vorhergehenden Büchern gehandelt ist.

\section{Das vierte Capitel}

\begin{schriftzitat}
V. 1. 2. Und nun höre, Israel, die Gebote und Rechte, die ich euch lehre, daß ihr sie thun sollt, auf daß ihr lebet, und hinein kommet, und das Land einnehmet, das euch der HErr, eurer Väter GOtt, gibt. Ihr sollt nichts dazu thun, das ich euch gebiete, und ihr sollt auch nichts davon thun, auf daß ihr bewahren möget die Gebote des HErrn, eures GOttes, die ich euch gebiete.
\end{schriftzitat}

1. Im vierten Capitel dieses Buchs macht Moses eine schöne, herrliche Vorrede über die zehn Gebote. Gleich als wenn ich zu euch spräche: Ich will eine Vorrede machen, oder eine Regel geben, die euch dienen wird zur Vermahnung. Also spricht Moses allhier auch: Sehet zu, daß ihr nicht etwas dazu setzet, oder davon nehmet. Eine solche Sorge hat Moses, und ist auch der Sorge wohl werth. Denn das ist gewiß, wenn GOttes Wort angeht oder weltliches Recht, auch die besten Ordnungen, so hält man es doch nicht durchaus. Es kommen bald Leute darüber, die es überklügeln und meistern. Darum wären diese Worte wohl werth, daß man sie mit güldenen Buchstaben an alle Wände schriebe. Denn Moses will sagen: Ich wollte euch gerne etwas Gutes lehren, aber ihr werdet nicht dabei bleiben; sehet zu, daß ihr nicht davon euch lasset abwenden.

2. Der Teufel ist unseres HErrn GOttes Affe; er hat neben dem gebahnten Wege und der Landstraße des göttlichen Worts allezeit seine Holzwege und Fußsteige, dadurch er die Leute verführt. Als denn jetzt auch geschiehet. Da man hört und weiß, was das Evangelium, was Christus sei, noch legt man zu und ab, wie es einem jeden gefällt, wie die Schwärmer und Rottengeister thun. Also thut der Pabst in der Christenheit auch; der trägt seine Decreta und Decretales auf den Predigtstuhl, und will jedermann, ja, GOtt selbst meistern. Das soll nun nicht sein; denn zu der heiligen Schrift soll man nichts zulegen, noch auch nichts davon nehmen.

3. „Höre, Israel“, thue die Ohren hieher, höre, was ich dich lehre. Siehe einen andern nicht an; merke,\footnote{„merke“ fehlt in der Erlanger.} was ich dich unterrichte. Dies ist eine treffliche Vorrede, daß er die Leute bringe
\quellseite{840}
und erhalte auf dem rechten Wege. Es sollen alle Worte mit großen Buchstaben geschrieben sein. Davon ist dieses kürzlich keine Meinung, und so viel will er sagen: Niemand lehre selbst, noch einen andern, sondern höre zuvor Mich, was ich ihn lehre.

4. „Sitten“ und „Rechte“, praecepta et judicia. Ich habe es\footnote{nämlich an anderen Orten, z. B. Ps. 2, 7.} gerne verdeutscht: praecepta, seine „Weise“, das eine gemeine Gewohnheit im Land ist. Auf hebräisch heißt es Chukka, das man jetzt Ceremonien nennt, zu deutsch: eine Weise, als zum Exempel, daß man das Thor in mancher Stadt vor sieben Uhr nicht aufmacht. Es ist nicht ein Recht, sondern Ceremonie, eine Weise. Dagegen ist es eine Unweise, wenn es ohne Ordnung zugeht, das ist, nicht zu rechter Zeit. Also theilt Moses seine Lehre in zwei Theile, „Weise“ und „Recht“; spricht: Ich will euch lehren, erstlich „Weisen“, wie ihr leben sollet; zum andern „Gerichte“, wie ihr richten sollet.

5. „Weise“ ziehe ich dahin, wie man leben soll im Gottesdienst, oder wie man gegen GOtt sich halten soll. Daher ist im andern Psalm, V. 7., geschrieben: Praedicans praeceptum [id est].\footnote{Von uns hinzugefügt.} narrabo ceremoniam, „Ich will von einer solchen Weise predigen“, wie man leben soll, das ist, an Christum glauben. Summa, Moses will seine Zuhörer eine Weise lehren, wie sie leben sollen gegen GOtt. Darnach „Gericht“, wie man gegen den Menschen sich verhalten solle. „Recht“ ist das weltliche Regiment, das Landrecht, Stadtrecht, darnach die Bürgermeister und Fürsten sollen regieren in äußerlichen Dingen, da es betrifft Haus, Hof \&c. Da gehört „Recht“ zu, als ihr Sachsen den Sachsenspiegel habt. Daher spricht man: Hausrecht, Frauenrecht, Tochterrecht. Was will und bedarf ein Mensch zu diesem Weltleben mehr, denn daß er die Kunst weiß, wie er sich vor GOtt und Menschen halten solle, daß er wisse, „Weise“ und „Recht“ zu vollziehen?

6. Diese zwei Stücke lehrt Moses reichlich. Hiermit zeucht er das Volk an sich, behält sie aber nicht an seiner Person, weist sie von sich zu Christo, daß sie ihn hören, und sehen auf seinen lieben und heiligen Mund. Also sollen wir auch thun, und Mosi hierin folgen, darnach lassen den Papst und die Patres, auch allerlei Geister und Lehrer, reden, was sie wollen und nicht lassen können. Denn an das göttliche Wort sollen wir gebunden sein, das sollen wir hören, und soll ohne GOttes Wort aus seinem eigenen Kopfe niemand etwas lehren. Also bindet Moses sie alle an seinen Mund, und legt also allen Lehrern seine Weise und Exempel vor. Niemand verdammt er und widerspricht alles andere Vorhaben, und sagt, er wolle allein lehren, denn er habe des auch guten Schein und Recht. GOtt habe es ihm befohlen. Nun, diese zwei Stücke sind vonnöthen, nämlich „Weise“ und „Recht“. Im Papstthum haben sie es genannt Kirchengepränge, und dahin deuten wollen; aber diesem Worte ist Unrecht geschehen.

7. „Daß ihr sie thun sollt.“ Das gehört auch dazu. Denn Recht und Weise sind nicht darum uns vorgestellet, als wären es nur Geheimnisse von Weisheit, die man aus den Worten Mosis speculiren müsse, und darnach in Wind hinsprechen und vergeblich reden möge, oder einen Ruhm und Ehre davon haben; wie ihrer viel vor dieser Zeit gedacht haben, und sind derer noch ein großer Theil, die ihnen solches träumen lassen. Es will’s aber GOtt hier nicht so haben, sondern daß man es thue, und dem nach sein Leben anstelle; wie in der ersten zu den Corinthiern [Cap. 4, 20.] Paulus auch spricht: „Das Reich GOttes stehet nicht in Worten, sondern in der Kraft.“ Darum so soll man es thun [Matth. 7, 21.]. Wir rühmen uns alle des Evangelii, aber das Thun macht einen Christen. Also sagt Moses viel mit wenig Worten. Das ist auch die rechte Kunst zu predigen und die Leute zu lehren, sollen sie anders mit Nutzen zuhören und etwas davon tragen.

\subsection{Auf daß du lebest}

8. Alhier setzt er dazu Dräuung und Verheißung. Verheißung, „daß ihr lebet und das Land ererbet“, wenn ihr das Gesetz haltet. Dräuung aber den Uebertretern, wie drunten weiter wird gesagt werden. Er will so viel sagen: Ich lehre euch nicht darum allein, daß ihr es wisset, sondern, daß ihr es thut und davon lebet. Werdet ihr es nun thun, so werdet ihr haben Land und Leute, Glück und Wohlfahrt; wo nicht, so werdet ihr aus dem Lande gestoßen werden, wie hernach folgt [V. 27.]: Wirst du nicht gehorchen, so wird er dich verblasen, wie
\quellseite{841}
der Wind den Staub oder die Spreu verwehet. Diese Worte sind trefflich groß, und haben eine Kraft bei sich, die Leute fortzutreiben und aufzumuntern.

\subsection{Daß ihr hinein kommet, und das Land besitzet}

9. Das ist, daß ihr Haus, Hof, Weib, Kind und allerlei Genieß bekommet. Wiederum, so ihr nicht gehorsam werdet, werdet ihr nicht allein aus dem Lande gestoßen werden, und von Haus und Hof, Weib und Kind ins Elend vertrieben, sondern ihr werdet auch jämmerlich erstochen und ermordet werden. Das ist eine starke Vermahnung, die da fleißig und lustig macht die Zuhörer, daß sie GOttes Wort nicht mit schlafenden Ohren hören, und nimmermehr daran gedenken, daß es GOttes Wort sei, und sie darnach thun sollen. Denn also soll man die Predigt ansehen und treiben, daß die Leute fleißig zuhören. Aber jetzund hören ihrer viel die zehn Gebote und andere Predigten, als wären sie es nicht schuldig zu hören, viel weniger verpflichtet zu thun; als im Pabstthum vorzeiten auch geschah, und noch, da man es nur für ein Menschenwort oder Pfaffengetrieb hält. Aber Moses will allhier sagen: Darum predigt man es, daß du darnach thun sollst, und in alle deinem Leben dich also halten.

\begin{schriftzitat}
V. 2. Ihr sollt nichts dazu thun etc.
\end{schriftzitat}

10. Sehet zu, daß ihr es lasset bleiben bei den Geboten, die ich gebe, und machet sie weder kleiner noch größer. Das ist ein groß, stark Wort, welches man mit großen Buchstaben schreiben sollte, und ist ein Donnerschlag, damit er wehrt allen Rottengeistern und falschen Propheten, die wunderlich sind, und ihre eigenen Köpfe haben. Denn es verdammt dieser Text alle Menschenlehre, stürzt die Ketzer zu Boden, und weissagt von ihrer Herzen Leichtfertigkeit, daß sie alles tadeln. Aber niemand soll ihnen\footnote{Hier haben wir „nicht“ getilgt. In den Ausgaben: nicht nachfolgen.} nachfolgen. Ja, wenn zehntausend Propheten kämen, und brächten das Gesetz und Wort Mosis nicht, da schließe für gewiß, daß es Ketzer sind.

11. Mit diesem Texte hat Moses allerlei Schwärmerey verdammt. Und wiederum, so hat der Text viel Propheten zu Märtyrern gemacht. Denn die Propheten haben aus diesem Texte der Könige Israel Abgötterei gestraft, und sind um dieser Strafpredigt willen alle erschlagen worden (Matth. 23, 37. 1 Kön. 22, 24. Jer. 20, 2.). Denn so ging es damals unter demselben Volke zu: Dieser nahm diese Weise vor, ein anderer jene Weise, und gab dann ein jeder vor, sie wären nicht unrecht, und sprachen: Wer bist du denn? meinst du, daß du den Heiligen Geist allein bei dir hast? Aus diesem Texte habe ich viel Schwärmer darniedergelegt, und Moses hat ihnen hier einen Zügel gesteckt, darüber können sie nicht.

12. Heutiges Tages hauet und sticht dieser Text auch um sich, und schlägt zu Boden alle Secten, allerlei Gerechtigkeiten und Frömmigkeiten. Unter dem Pabst ist mancherlei Weise, als, der Predigermönche, Carthäuser, Minoriten, Laienbrüder etc., da hat ein jeglicher seine sonderne Weise: der fastet den Heiligen, jener einem andern, der läuft zu St. Jakob, jener ins Grimmenthal. Diese Weisen hat er hier alle verboten, und diese allein uns gegeben, nämlich: Glaube an Jesum Christum und liebe ihn.

13. Bei der Weise will man es nicht lassen bleiben, sondern [sie] wollen, daß wir die Weisen der Päpstler annehmen, wo nicht, so wird man drüber verbrannt als ein Ketzer. Aber mit diesen Worten verwerfen wir alle eigene Andacht und menschliche Werke, und sollten ihrer noch so viel drüber erschlagen werden. Darum ist dieser Text der Schrein aller Gottseligen, die zu Unrecht sind erwürgt. Abel ist darum erschlagen, daß er nicht führte eine Weise wie Kain. Das ist geschehen lange zuvor, ehe denn Mose diese Worte gepredigt; wie jetzund auch uns geschieht. Ich lasse geschehen, daß der Kaiser mache „Rechte“, wie er will, aber nicht eine „Weise“ soll er mir in GOttes Sachen vorstellen; denn der Vater im Himmel hat uns gebunden an Christi Mund, und gesprochen Matth. 17, 5.: „Den sollt ihr hören.“ So spricht alhier Moses auch: „Thue nichts dazu“, das ist, du sollst bei dem Texte des göttlichen Worts allein bleiben.

14. Möchte aber einer fragen: Was soll ich denn machen? Jesaias, Jeremias und andere Propheten haben auch ein jeglicher seinen Geist gehabt und viel Bücher geschrieben, die nicht einander gleich sind, und Christus hat viel dazu gesetzt und davon genommen? Darauf antworte also: Das heißt nicht dazu thun, wenn ich Mosis Worte ausstreiche und erkläre; wenn ich nur bei seinem Verstande und Meinung bleibe, so mag
\quellseite{842}
ich es wohl mit viel Worten erklären. Also, Jesaias bleibt bei den zehn Geboten und Worten Mosis; ob er wohl andere Personen hat, zu denen er redet, und andere Worte braucht, daran ist nichts gelegen. Jesaias kann allerdings nicht reden vom Pharao in Egypten, wie Moses; aber eben die Sachen hat er auch gehabt, die er hat müssen strafen, jedoch auf andere Zeit. Darum wirst du in den Propheten nichts Anderes finden, denn das sich reimt auf die zehn Gebote.

15. Das heißt aber dazu thun, wenn einer will über den Verstand etwas Anderes lehren. Zum Exempel, wir lehren von Christo, daß er allein vom Vater uns gemacht sei zur Gerechtigkeit, Heiligung, Erlösung und Versöhnung, 1 Cor. 1, 30., da kommen die Papisten her, und thun zu der Gerechtigkeit auch ihre Frömmigkeit, Heiligkeit und Verdienst; als, daß man müsse zu St. Jakob laufen. Das will sich zusammen nicht vergleichen. „Davon thun“ heißt, etwas Anderes setzen und herauszuzwacken, wie die Sacramentenschwärmer thun; die lassen diese Worte, „das ist mein Leib“, an sich selbst dem Buchstaben nach bleiben, aber nehmen gleichwohl den Verstand hinweg, und deuten es anders; als, daß im Abendmahl nicht sei der wahre Leib und Blut Christi, sondern nur Zeichen und Bedeutung desselben. Das ist den wahrhaftigen Verstand ausgeflössen, und dafür einen andern Verstand den Worten angehängt, denn sie sonst lauten. Es schadet sonst nicht, wenn ich hundert(er)lei Worte führete; allein, daß einerlei und der wahrhaftige Verstand bleibe. Kurzum, man soll keinen andern Verstand bringen, sondern, daß Eine Meinung und Glaube bleibe; ohne das hindert es gar nicht, obschon andere Worte gebraucht werden. Also ist Mosis und der Propheten Lehre gar einerlei, obwohl nicht gleiche Worte sind von ihnen geführt worden. Denn die Propheten haben auch geredet, getrieben durch den Heiligen Geist [2 Petr. 1, 21.].

16. Das ist die Vorrede, die gehört auf solche heilsame Predigt, da er spricht: „Bewahret die Gebote eures GOttes, die ich euch gebiete.“ Das hänget er so oft daran, daß es schier verdrießlich ist zu lesen einem fleischlichen Menschen. Denn Moses gebraucht dieser Worte auch sonst vielmals; aber er weiß sehr wohl, es sei nicht zu viel, wenn er es schon tausend- oder mehrmal sagt, denn es ist hoch vonnöthen. Und er thut es auch darum, daß [es] in der Welt so zugeht, daß immer kommen, die es wollen besser machen. So läßt der Teufel solch Meistern, Klügeln und Ueberklügeln auch nicht; es ist sein Handwerk, daß er im Paradies mit Adam und Eva angefangen [1 Mos. 3, 4. 5.] und bisher getrieben, nämlich, an das Wort und Gebot GOttes einen andern Verstand flicken, und besser machen wollen. Darum muß man nicht nachlässig, sondern wacker sein, und täglich dawider vermahnen.

\subsection{Auf daß ihr haltet x}

17. Ein seltsamer Anhang und Warnung ist das. Wer hätte sie da gesucht? St. Paulus hat freilich auch hieher gesehen, da er spricht zum Tito am 1. Capitel, V. 10. 11.: „Es sind viel widerspenstige und unnütze Schwätzer und Verführer, sonderlich die aus der Beschneidung, welchen man muß das Maul stopfen; die da ganze Häuser verkehren, und lehren, das nicht taugt, um schändliches Gewinnes willen.“ Und bald hernach spricht er, V. 13. 14.: „Um der Sache willen strafe sie scharf, auf daß sie gesund sind im Glauben, und nicht achten auf die jüdischen Fabeln und Menschengebote, welche sich von der Wahrheit abwenden“ etc.

18. Denn das ist der Menschenlehre Art, daß sie vom Glauben und den Geboten Gottes abführt.\footnote{Erlanger: abführen.} Wie denn auch St. Paulus Röm. 16, 17. 18. sagt: „Ich ermahne euch, daß ihr aufsehet auf [die], die da Zertrennung und Aergerniß anrichten, neben der Lehre, die ihr gelernet habt, und weichet von derselben; denn solche dienen nicht dem Herrn Christo, sondern ihrem Bauche, und durch süße Predigt und prächtige Worte verführen sie die unschuldigen Herzen.“ Welches wir bisher wohl erfahren haben unter dem Papstthum. So bald ein Zusatz kommt über GOttes Gebot, so wendet der Mensch sich von den zehn Geboten, und kommt von dem rechten Worte der Wahrheit und den zehn Geboten und derselbigen Verstande, und fällt in Verführung und Irrthum. Also hat man vorzeiten auch wohl gepredigt vom Glauben, aber man ist dabei nicht geblieben, sondern auf die Werke geraten. Da sind kommen die heiligen Antoniter und Einsiedler, die haben gefastet und anderes mehr gethan, und die Leute von Christo und dem Glauben an ihn auf die guten Werke
\quellseite{843}
geführt etc. Das ist daher gekommen, daß sie sich der Möncherei allein beflissen, und gar verirrt sind von der Lehre des Glaubens auf die Menschensatzungen und Leben. Das zeugt gewaltig von GOtt ab.

19. Und so geht es, wenn man nicht Acht hat auf den Unterschied des geistlichen und leiblichen Reichs; da heißt denn „zusetzen“, daß die Leute abgewendet werden von der Wahrheit. Das bekennt hier Moses frei in diesem Texte: Wo ihr davon oder dazu thut, so kommt ihr davon. Wollt ihr es aber behalten, so hütet euch vor dem Zuthun und Abthun, und hütet euch vor den Holzwegen und vor Beinwegen, daß ihr nicht auf einen Affenschwanz geführt werdet. Zwar die tägliche Erfahrung gibt es an sich selbst, daß, wer da sicher wandeln will, der bleibe auf der Landstraße, oder er wird nicht zurecht kommen. Dieses vermahnt Moses fleißig, darum soll man es auch ernstlich halten, daß man nichts zu GOttes Wort thue, noch davon nehme. Darum sollen wir Menschenlehre meiden, und Christi Lehre allein anhangen.

20. Jetzt versteht ihr, wie weit man solle Menschengebote halten, nämlich, daß man sie halte freiwillig, allein in diesem Leben. Denn sie gehören nicht in das andere Gebot, zur Heiligung GOttes Namens. Summa Summarum, GOttes Gebote werden nicht gehalten, man bleibe denn in dem, was sie lehren. Denn sobald man etwas Anderes dabei will setzen, so kommt man gar davon. Und daß solches wahr sei, spricht Moses weiter von denen, so dem BaalPeor angehangen:

\begin{schriftzitat}
V. 3. 4. Eure Augen haben gesehen, was der HErr gethan hat wider den BaalPeor; denn alle, die dem BaalPeor folgeten, hat der HErr, dein GOtt, vertilget unter euch. Aber ihr, die ihr dem HErrn, eurem\footnote{„eurem“ fehlt im Original und in der Erlanger.} Gott, anhänget, lebet alle heute des Tages.
\end{schriftzitat}

21. Als wollte Moses sagen: Seid hierinnen klug, und sehet euch vor, daß ihr zu dem göttlichen Wort, Gebot und Befehl weder zu- noch abthut. Und daß ihr vor diesem Zusatz oder Abzug aufs treulichste möchtet gewarnt sein, so stellet euch dies Exempel vor die Augen, wie GOtt diejenigen vertilgt hat, die ihn verlassen, und dem Abgott BaalPeor angehangen haben.

22. Wir haben nächst gehört, wie Moses das Volk vermahnt in seiner Vorrede, dadurch er will anfangen die zehn Gebote auszulegen. Darum hier, im folgenden Texte, nimmt er sonderlich vor sich das erste Gebot zu handeln, denn aus dem fließen, als aus einer Quelle, und folgen die andern Gebote alle, und richten sich nach dem ersten Gebot. Und [wir] hören vornehmlich, daß er spricht: „Machet euch kein Bildniß“ [2 Mos. 20, 4. 5 Mos. 5, 8.]. Da will er, daß man kein Bild noch Götzen machen soll, sondern daß wir uns straks an das Wort und die Stimme GOttes halten, und wissen, daß wir es hören sollen ohne Abgötterei und Bilder. Darum verbietet er so hart die Abgötzen, und will, daß wir allein am Worte GOttes hangen.

23. Aber was haben unsere Schwärmer jetziger Zeit zu schaffen? Sie sehen nur auf die äußerlichen Bilder; und auf den Wort, „nicht Bilder haben“, fußen sie, und befleißigen sich die Bilder zu stürmen. Aber nach dem andern, nämlich, das Wort GOttes halten, da fragen sie gar nichts nach, und lassen das Wort GOttes anhin fahren.

24. Nun folgt weiter, wie er sie also höchlich vermahnt, daß sie ihnen kein Bild machen, noch auch keine Creatur aufwerfen, ehren oder anbeten sollen. Das ist die Summa des folgenden Textes, da Moses spricht:

\begin{schriftzitat}
V. 15–19. So bewahret nun eure Seelen wohl; denn ihr habt kein Gleichniß gesehen des Tages, da der HErr mit euch redete aus dem Feuer, auf dem Berge Horeb, auf daß ihr euch nicht verderbet, und machet euch irgend ein Bild, das gleich sei einem Manne oder Weibe, oder Vieh auf Erden, oder Vogel unter dem Himmel, oder Gewürm auf dem Lande, oder Fisch im Wasser unter der Erde. Daß du auch nicht deine Augen aufhebest gen Himmel, und sehest die Sonne, und Mond, und die Sterne, das ganze Heer des Himmels, und fallest ab, und betest sie an und dienest ihnen, welche der HErr, dein GOtt, verordnet hat allen Völkern unter dem ganzen Himmel.
\end{schriftzitat}

25. Das ist das erste: „Du sollst dir kein Bild machen.“ Danach das andere, daß man zusehe, und je keine Creatur, welche geschaffen ist, für einen Gott halte, es sei gleich die Sonne, welche doch der vornehmsten Creaturen eine ist,
\quellseite{844}
oder der Mond, oder die Sterne, oder sonst des Himmels Heer; das ist, der ganze Haufe Gestirns zusammengezogen. Denn dies alles ist GOttes Werk, GOtt hat sie geschaffen, sie sind seiner Hände Arbeit. Da sollst du wissen, daß alle diese Creaturen viel zu geringe sind, daß sie sollten dein Trost sein, und daß du ihnen dienen solltest; ich geschweige, daß sie sollten deine Götter sein. Denn was ist Sonne und Mond gegen GOtt, ihren Schöpfer, zu rechnen?

26. Er verdammt diese GOttes Geschöpfe nicht, sondern bestätiget’s vielmehr, daß Sonn und Mond GOttes Creaturen sind, von GOtt geschaffen und geordnet zum Dienst der ganzen Welt, daß sie den Menschen scheinen, zu der Zeit und Jahren dienen sollen. Darum ist es nicht der frommen Sonne und Mondes oder der Sterne am Himmel Schuld, so sich an ihnen etwas zuträgt oder Zeichen erscheinen, sondern an Gott, des Himmels, ihrem und der ganzen Creatur Regenten, ist es alles gelegen. Darum darfst du sie nicht anbeten und ehren als einen Gott, auch ihnen nicht dienen als unserm HErrn GOtt. Der Text ist wahr genug, und redet von unserer Bosheit, und nicht von den\footnote{Erlanger: „ber“.} lieben Creaturen. Und [wir] haben oftmals gesagt, wenn man das will wegreißen, was Ursache zu Sünden gibt, so müßte man gar nichts bleiben lassen; man kann die Creaturen um des Mißbrauches willen nicht wegwerfen. Propter abusum non tollitur substantia. Man müßte die Sonne und Sterne, ja, sonst viel unzählige Creaturen auch wegnehmen.

27. Die Schwärmer meinen, sie haben der Sache geraten, wenn sie die Götzen haben weggethan, darum, daß sie von Leuten sind angebetet worden; aber sie werden es wohl bleiben lassen. Thue die Sonne und den Mond auch weg, denn man betet sie auch an. So müßte man keine Creatur bleiben lassen, wenn man ein Ding sollte wegthun, weil es die Menschen mißbrauchen; sondern man sehe zuvorhin in das Herz, und lehre erstlich, wie man die Götzen im Herzen zerbreche und einreiße, oder herauswerfe, daß man sie nicht anbete. Er redet vom Anbeten; da steht die Sünde innen, daß du sie nicht ansiehst und abfällest, und sie anbetest. Nicht, daß ich wollte, daß man solche Götzenbilder sollte anrichten, wie wir bisher haben gethan, mit welchen man GOtt dienete; ich will solchen Götzen nicht hören; sondern, daß man predigen sollte auf das Anbeten, und daß das Trauen auf die Creaturen möchte nachbleiben. Denn das Dienen und Anbeten soll aufhören; sonst muß man die Creaturen um unseres Mißbrauchs willen nicht wegwerfen.

28. Das Gold ist fein und gut an ihm selbst, noch heißt es ein Mammon; nicht seiner Schuld halben, denn was hat es gethan? Der fromme Gulden hat nicht gesprochen zu dir: Ich bin dein Gott; ja, er würde vielmehr zu dir sagen, wenn er reden könnte: Ich bin dein Knecht. Des falschen Herzens Schuld ist es, daß es das Gold, Gut und Geld also sehr liebt, ihm hold und daran hangt, und ihm dient und einen Gott daraus macht; welches es doch nicht ist, noch sein kann oder will. Also, die fromme Sonne hat noch keinmal gepredigt, man hat es von ihr nicht gehört, daß sie gesagt hätte: Bete mich an und mache einen Gott aus mir. Also hat auch der Mond nicht gesagt; noch sind Leute da gewesen, die sie haben angebetet. Also, die frommen Goldgulden sagen nicht zu dir: Bete mich an, setze deine Hoffnung und Vertrauen auf mich; sondern er läßt sich legen, werfen und ausgeben, wie du willst. Und wenn sie könnten reden, so würden sie sprechen: Ich bin dein armer Knecht, wo du mich hingibst, da gehe ich hin, da muß ich sein. Setzt man nun das Vertrauen auf ihn, daß er Gott sei, das hat er nie begehrt. Also sagen alle Creaturen.

29. Das soll man fleißig predigen, auf daß man die Leute führe aus diesem Stücke, das da heißt, anbeten und dienen, und also den Abgott aus dem Herzen treibe. Denn was sonst externa sind, das wird bald vergehen. Aber das thun unsere Schwärmer und falschen Geister nicht. Darum soll allhier ein rechter Lehrer zusehen, daß er zum Worte erst greife, lehre und unterrichte das Volk. Denn da liegt am allermeisten an, daß man die Bilder nicht anbete, noch für Gott sie halte; wie Moses allhier sagt: Laß sonst das andere alles sein, laß Sonne und Mond sein, Gold und Gut sein, bete es nur nicht an, und halte es nicht für deinen Gott.

30. Es ist nicht der Streit von der Substanz oder Wesen der Creaturen, was die Sonne, Mond und andere Creaturen für Wesen haben, sondern von dem Gebrauch und Mißbrauch. Wir predigen nicht, was das Wesen an ihm
\quellseite{845}
selbst unbelangt, sondern von dem Mißbrauch eines bösen Herzens. Wir wollen nicht die Dinge verändert haben, sondern das verkehrte Herz. Denn Geld ist Geld, die Sonne ist Sonne. Das soll man predigen, nicht, was man an den Creaturen ändern solle; allein wie man den Mißbrauch des schalkhaftigen Herzens wenden möge; denn das sucht man, daß du anders werden solltest, und man nicht die Creaturen anders mache. Denn du sollst auf etwas Anderes deine Hoffnung nicht setzen, sondern allein GOtt sollst du vertrauen; denn, thust du das nicht, so ist der Sache nicht geholfen.

31. Bilder und Abgößen kannst du zwar zerbrechen und Kirchen einreißen, daraus nur ein Stolz wird, dieweil Steine leichtlich können zerbrochen werden; das Herz ist aber darum nicht gebrochen. Das arme Vieh, als Gold, Silber, Holz und Steine, wird nur zerbrochen; aber das Herz wird nicht bewegt, noch unterrichtet, daß es GOtt allein diene. Ich wollte auch wohl leichtlich Bilder stürmen, wenn man des Herzens Sturm ließe vorhergehen, dadurch die Herzen gestürmt und gebrochen würden. Die Creaturen sind uns zu Dienst geschaffen, aber zu Göttern hat sie uns GOtt nicht gesetzt.

32. Und bekennt Moyses allhier selber, daß vorzeiten die Sonne, der Mond und Sterne sind für Götter gehalten und angebetet von den meisten Leuten, und sind die Heiden darum verdammt, daß sie die Creaturen haben angebetet. Die Egypter haben erstlich Kühe und Pferde und andere irdische Creaturen angebetet. Aber darnach haben sie diesen irdischen Gottesdienst verworfen, und der himmlischen Creaturen Gottesdienste aufgerichtet, Sonne und Mond für Götter gehalten, darum, daß sie so schnelles Laufes wären, und in einem Tage eine solche große Reise thun und vollbringen könnten.

33. Dies merket aus dem Capitel, daß Moses so fleißig wehrt und steuert, daß sie nicht sollen die Sonne und Mond anbeten. Denn, wie ich gesagt, es sind nicht grobe Bauern, sondern weise Leute gewesen, welche geglaubt, daß Sonn und Mond Götter wären. Sie sind nicht gewesen wie unsere gottlosen, teuflischen Zauberer und alte Bredin\footnote{Bredin (Bräckin) = Hündinnen (von Bracke, Hund); Schimpfwort gegen die Hexen.} sind, die den Holunderbusch, ja, den Drachen anbeten und mit ihm Verbündniß und Gesellschaft haben. Laß die Creaturen bleiben, wie groß und hoch sie sind, auch Sonn und Mond.

34. Wahr ist es, die Sonne ist eine edle Creatur und herrliches Werck GOttes neben dem Menschen, unter allen Creaturen; Gold und Silber ist nichts gegen der Sonne. Denn wenn wir die Sonne nicht hätten, so könnten wir nichts sehen, es wäre für und für Nacht und alles finster, gleich als wenn wir in einem dunkeln Keller lebten. So die Sonne nicht wäre, so könnte der andern Creaturen Nutz und Brauch nicht bestehen, man würde keiner Tugend und Wahrheit begehren; item, Gold und Silber, Weib und Kinder, und alles andere verachten, und es wäre einer gleich so übel unter der Erde, als daß er lebte. Darum so ist die Sonne das edelste Werck und Geschöpf unter allen, uns Menschen allein ausgenommen. Wenn die Sonne nicht schiene, so würde nichts stehen, weder Haus, Schloß noch Hof, kein Reich auf Erden.

35. Weil nun die Sonne, als die edelste Creatur, nicht soll noch kann GOtt sein, noch angebetet werden, so doch alles erleuchtet, daß es steht; und gewaltig regieret, und ihre Wirkung hat in die untersten Körper: wie viel weniger soll der ohnmächtige Mammon, und andere Creaturen angebetet werden? Darum haben die hohen weisen Leute gedacht: Was ist Geld? Was ist Acker? Wenn nicht Licht ist, wenn die Sonne nicht schiene, so bliebe es für und für Nacht. Wenn sie in zweien oder dreien Tagen nicht hervor blickt, so ist doch alles traurig.

36. Das haben die vernünftigen, weisen Leute, als Plinius und andere, gesehen, und darum gesagt, die Sonne sei Gott, und haben sie angebetet. Und wahrlich, wenn noch die Sonne drei Tage außen bliebe, wir beteten sie auch an. Darum wird hier gepredigt den weisen Leuten, daß sie sich vorsehen, und nicht die Sonne, die edle und herrliche Creatur, anbeten. Und ist doch von vielen sehr weisen Leuten angebetet; als, Plinius ist sehr klug gewesen, und hat dennoch die Sonne für einen Gott gehalten.

\begin{schriftzitat}
V. 20–24. Euch aber hat der HErr angenommen, und aus dem eisernen Ofen, nämlich aus Egypten, geführet, daß ihr sein Erbvölk sollet sein, wie es ist an diesem Tage etc. So hütet euch nun, daß ihr des Bundes des HErrn nicht vergesset, den er mit euch gemacht hat, und
\end{schriftzitat}

\quellseite{846}
nicht Bilder machet einigerlei Gleichniß; wie der HErr, dein GOtt, geboten hat. Denn der HErr, dein GOtt, ist ein verzehrend Feuer, und ein eifriger Gott.

37. Er spricht: „Sehet zu, vergesset eures GOttes nicht.“ An dem Worte „vergessen“ liegt alle Macht, daß man ihn für unsern GOtt halte. Dieses treibt Moses mit hohem Fleiß. Diesen Text aber hat man auch geführt auf die Klostergelübde, und des Nöthigsten ganz und gar vergessen, ja, auf einem Bündlein liegen lassen; nämlich, was uns GOtt gelobt, geschworen und gelobt haben; daher viel Gelübde gekommen sind. Es soll aber sich das Wort „geloben“ ziehen sonderlich auf das erste Gebot, daß wir GOtt gelobt und geschworen haben, daß er sollte unser\footnote{Erlanger: „unsern“.} Gott sein, und wir dagegen sein Volk. Dies Gelübde vermag und bringt mit sich, daß ich wolle auf keine Creatur trauen, denn allein auf GOtt, und sonst nichts annehmen, daß man fürchte, liebe, dem man vertraue und es ehre. Denn er will allein unser GOtt sein.

38. Moses ist ein feiner Doctor und Lehrer, der weiß von Sachen artig zu reden; er spricht: „Hütet euch, daß ihr des Bundes nicht vergesset.“ Als wollte er sagen: Es wird dazu kommen, daß man GOttes vergessen wird, daß gar viel falsche Propheten und irrige Lehrer aufkommen, welche sagen werden: Willst du gen Himmel kommen, so thue dies und das. Aber siehe du zu, und gedenke, was du hast für einen Bund mit GOtt gemacht; dabei bleibe, „denn der HErr, dein GOtt, ist ein verzehrend Feuer“.

39. Denn GOtt ist solcher GOtt, oder, wie er hier sagt, ein solch Feuer, das nicht feiert, sondern frißt und verzehrt, und ein GOtt, der da eifert, 2 Mos. 20, 5. Dieser Spruch ist wohl zu merken, welcher die angeht, so abtrünnig werden vom Glauben; er gehört für die Mammeluchen und Apostaten. Man glaubt es aber nicht, daß GOtt ein verzehrend Feuer ist, und ein eifriger GOtt. Darum ist es gesagt den Mißbräuchlichen, Verhärteten, Verstockten und Ungläubigen, welche nicht den rechten Glauben behalten, und diese Worte für lauter Fabeln achten und verpotten; denen wird er ein Feuer genannt, das verzehrt. Das ist so viel gesagt: Welcher nicht glaubt an sein Wort, und nimmt es nicht an, der ist verderbt in Grund und Boden. Deshalb glaube, und halte dich an GOtt; wo nicht, so bist du schon dahin.

40. Wollte GOtt, daß die Welt das glaubte, und dies Wort für die Wahrheit hielte, daß GOtt ein verzehrend Feuer wäre. Daher so lebt man auch so wild und wüst in der Welt, und thut ein jedermann wider GOttes Gebot, und lassen GOtt schelten und dräuen, wie er will. Denn GOtt wird bei ihnen für kein verzehrend Feuer gehalten; sondern irgend für Stoppeln, für einen Strohhalm und Tröpflein kalt Wassers hält ihn die Welt. Darum geht es auch so durch einander. Dieser stiehlt und raubt, jener betrügt und lügt; der hurt, jener haßt; ein andrer geizt etc. Da scheint es nicht, daß ein solches Wesen verzehrt werden müsse und untergehen, sondern es läßt sich ansehen, als haben sie einen gnädigen GOtt; wie denn der Prophet Jeremias Cap. 12, 1. auch darüber klagt: Quare prosperatur via impiorum? „Warum geht es dem Gottlosen so wohl?“ Und die Erfahrung zeuget es auch, wie man im Sprüchwort sagt: Je ärger Schalk, je besser Glück. Es geht ihnen eine Zeitlang so hinaus.

41. Daran haben sich die heiligen Leute und die Propheten auch sehr geärgert, daß das Gegenspiel in der Welt daher leuchtet, und die Sünder aller Dinge Ueberfluß haben, und treiben den allergrößten Muthwillen, und geht ihnen doch alles hinaus. Also hat St. Paulus das Widerspiel auch gesehen. Da fallen einem denn diese Gedanken ein: Ei, wäre GOtt ein verzehrend, fressend Feuer, eifrig und zornig, wahrlich, er würde das nicht leiden. Nun leidet er es? Ja, vor der Welt scheint es also, als wäre GOtt ein lauter Gähnemann, der das Maul nur aufsperre, oder ein Hahnrei und guter Mann, der einen andern läßt bei dem Weibe schlafen, und stellt sich, als sähe er es nicht.

42. Wider diese Gedanken hat Moses solches geredet: So hütet euch, ihr habt einen GOtt, der ist ein verzehrend Feuer, das ist, ein solcher GOtt, der euch verzehrt und aufräumt, so ihr gottlos seid, eifert, frißt und macht zu Asche und Staub. Er schlingt einen hinein, und hat eine solche Lust daran,\footnote{daß heißt: Er hat ein solches Gefallen an der Gottlosigkeit, daß etc.} daß er aus seinem Eifer und Zorn dazu getrieben wird, die Bösen zu verzehren
\quellseite{847}
. Geht solches einmal an, so läßt er nicht ab. Solches kann man die Leute nicht bereden, sie glauben es auch eher nicht, bis die Erfahrung kommt; so ist es denn zu lange geharrt. Also konnten die Juden auch nicht glauben, daß sie sollten verstoßen werden gen Babel, wiewohl solches die Propheten geweissagt hatten, bis\footnote{Im Original und in der Erlanger: „als“ statt: bis.} daß dasselbige verzehrende Feuer kam, und sie vertilgte.

43. Also, da Christus auch kam und von der letzten Zerstörung Jerusalem prophezeiete, so half es nichts, sie wollten es nicht glauben, bis daß darnach der Glaube ihnen in die Hand kam, da erfahren sie es mit der That. Wenn jetzund hunderttausend Zungen sagten zu unsern Thranen: „Gott ist ein verzehrend Feuer“, so haben sie nur einen Spott daraus, hören nicht, und halten auch unsern HErrn GOtt nur als einen Strohpoppen, der im Hanf den Vögeln zur Abscheu gesteckt wird. Aber wenn er zu seiner Zeit kommen und alles zu Pulver und Asche machen wird, dann werden sie es gewahr werden. Es geht uns heutiges Tages auch also. Wenn wir gleich lange predigen und dräuen, oder selbst hören, daß GOtt ein Feuer ist, fragen wir nicht darnach, glauben es nicht, bis die Pestilenz, gehender\footnote{gehender (gähender) Tod = schneller, plötzlicher Tod, dasselbe wie § 48: jäher (gäher) Tod.} Tod, Krieg und andere Noth und Plage kommen; dann beginnt es wahr zu werden, und sich im Auskehrig zu finden.

44. Darum schrecket uns hiermit Moses, und jagt uns eine Furcht ein; will, daß wir nicht einen andern Gott machen noch annehmen. Denn er kann es nicht leiden, oder du wirst drüber verzehret, und mußt untergehen. Solches bildet er uns als ein guter Redner vor, wenn es sonst helfen wollte. Es sagt’s, wahrlich, kein Schusterknecht, sondern der hohe Prophet Moses, und spricht: „GOtt ist ein verzehrend Feuer“, das ist, es wird nichts ungerochen noch ungestraft bleiben. Wenn sich’s gleich stellen wird, und scheint oder sich ansehen läßt, als habest du Glück und Fortgang in deiner Bosheit, dennoch wird dir endlich dein Recht widerfahren, daß GOtt dir ein verzehrend Feuer ist. Die Gottesfürchtigen glauben es, und darum erfahren sie es nicht. Wiederum, die Gottlosen glauben es nicht, darum müssen sie es erfahren; wie das die Erfahrung und That bezeugt, daß es wahr sei, und sie bald umkommen und zu Boden gehen [Ps. 37, 20.]. Darum hüte dich vor diesem Feuer.

\begin{schriftzitat}
V. 25. 26. Wenn ihr nun Kinder zeuget und Kindeskinder, und im Lande wohnet, und verderbet euch, und machet euch Bilder einigerlei Gleichniß, daß ihr übel thut vor dem HErrn, eurem GOtt, und ihr ihn erzürnet, so rufe ich heutiges Tages über euch zu Zeugen Himmel und Erde, daß ihr werdet bald umkommen von dem Lande.
\end{schriftzitat}

45. Da legt er aus, was vorhin gesagt ist: „Wenn ihr im Lande werdet Kinder aufziehen, [und] Kindeskinder zeugen.“ Da will er sagen: Jetzt, zu meiner Zeit, möchtet ihr GOttes Gebot halten und ein wenig fromm sein; aber nach unserm Tode wird es gehen, wenn wir das Land geleget haben, dann werdet ihr sehen die Abgötterei der Heiden, daß sie werden die Creaturen anbeten, und solcher Abgötterei folgen. Dann wird es auch gewiß kommen, daß das Wort GOttes wird verlöschen, und ihr werdet vertilgt, auch bald ausgerieben und alle gemacht werden, und zerstreuet unter die Heiden. Das meint Moses mit dem verzehrenden Feuer, daß GOtt wird zufahren, und euch aufreiben.

\begin{schriftzitat}
V. 26. 27. Ihr werdet nicht lange drinnen bleiben, sondern werdet vertrieben werden, und der HErr wird euch zerstreuen unter die Völker, und werdet ein geringer Pöbel übrig sein unter den Heiden, dahin euch der HErr treiben wird.
\end{schriftzitat}

46. Das heißt ein verzehrend, fressend Feuer, daß der Herr zufährt, und straft das Land mit Pestilenz, Hunger, Krieg, Theurung; und, daß noch ärger ist, wenn er schickt falsche Lehre und Abgötterei, dann habt ihr es gar mit einander. Denn er sagt [V. 28.]: „Unter den Heiden, dahin dich der HErr treiben wird, daselbst wirst du dienen den Göttern, die Menschenhände Werk sind, Holz und Stein, die weder sehen, noch hören, noch essen, noch riechen.“\footnote{Im Original und in der Erlanger: trinken.} Also verderbt er sie an Leib und Gut; dahin soll es mit den undankbaren Leuten kommen. Ich meine, daß mag aufgerieben heißen, und das möchte ein Feuer sein. Kommt der Türke unter uns einmal (das wir denn nicht erleben werden, wie ich hoffe), so wird er die zwei\footnote{„die zwei“, nämlich GOttes Wort und leiblich Gut. Vgl. § 50 zu Ende.} auch
\quellseite{848}
hinwegnehmen, und falschen Glauben mitbringen. Dann werden wir es lernen, wie GOtt ein verzehrend Feuer sei, das da allmache, und eifere zu beiden Seiten. Denn er will es nicht gut sein lassen, und dräuet uns auch durch den Türken, oder noch wohl andere, die uns näher sind, solche Strafen.

47. So warnt uns nun allhier Moses als ein treuer Prediger, daß wir von GOtt und seinem Wort nicht abfallen, auf daß wir nicht verzehrt werden. Als sollte er sagen: Summa Summarum, lasset ab von eurer Abgötterei, ihr müßt euch ändern und bessern, ihr könnt und mögt nicht vorüber, ihr müßt gar ablassen von der Abgötterei, und allem gottlosen Vorhaben, denn GOtt läßt von seinem Eifer nicht ab. Ich weiß keinen härteren Spruch oder Text in der Bibel denn diesen, daß GOtt ein verzehrend Feuer sei.

48. Ob er sich eine Zeitlang stellt, als sei er es nicht, und sieht durch die Finger; aber er wird es zu seiner Zeit wohl sein,\footnote{Im Original und in der Erlanger: sein werden.} und die Undankbaren zuletzt greulich strafen. Aber wir wollen es nicht glauben. In Wesphland hat man es gesehen, wie großer Hunger ist gewesen, und die Leute des jähen Todes gestorben sind, das da ist erschrecklich zu sagen. Da man sie warnte, glaubten noch achteten sie es auch nicht. Jetzt wird Deutschland seiner Undankbarkeit halben auch vermahnt, aber sie schlagen uns einen Klipp dafür; aber welche es gereuen und betreffen wird, das werden sie sein. Es wird das verzehrende Feuer dermaleins anheben und angehen, mache man da aus, was man will. Wer nicht glauben und hören will, der fühle es denn; man kann ihm anders nicht thun.

49. Die Welt wollte auch nicht glauben zu Noah Zeiten; sie hat aber die Sündfluth gefühlt. Sodoma und Gomorra wollten auch nicht glauben, es kam denn das verzehrende Feuer. Es ist GOtt ein eifriger GOtt, der nicht scherzen will, noch seine Verachtung ungerochen lassen, er straft es alsobald. Das sei nun gesagt den halsstarrigen, harten Köpfen; den Frommen wird GOtt auch anders, das ist, gnädig und barmherzig sein.

\begin{schriftzitat}
V. 28. Daselbst wirst du dienen den Göttern, die Menschenhände Werk sind.
\end{schriftzitat}

50. Das ist die größte Plage, wenn GOtt das Wort wegnimmt. Es ist wohl ein großer Jammer, wenn GOtt uns das leibliche Brod entzieht; aber viel eine größere Marter ist es, wenn er uns des Wortes beraubt [Amos 8, 11.]. Wiewohl es der gemeine grobe Pöbel nicht achtet, und das viel höher hält, wenn sie nur das Brod behalten, und an dem nicht Mangel leiden. Denn sie gedenken also: Wenn das Brod uns genommen wird, so sterben wir Hungers. Aber langsam befinden sie das: Wenn uns aber das Wort genommen wird, so fahren wir in die Hölle, und sind an Leib und Seele ewig verloren. Wahrlich, eines ist leichter und träglicher denn das andere, und besser wäre es, Hungers sterben, denn in die Hölle fahren. Aber hier bräuet Moses und jagt: Werdet ihr GOtt nicht vertrauen, so soll es euch beides genommen werden, das Wort GOttes und das leibliche Brod.

51. Und bald hernach beschreibt er sehr fein solche Götter, sagt, es sind Götter, die weder sehen noch hören, essen und trinken nichts [Jes. 44, 9. Ps. 115, 4. ff.]. Ist es nicht also? Siehe die alten heidnischen Götzen an, welche waren gülden, silbern, Holz und Steine; und der Papisten Götzen zu unserer Zeit: der ist zu St. Jakob gegangen, jener sonst Wallfahrten nachgelaufen; einer hat Gerechtigkeit der Mönchskappe, der andere den hörnenen Stricken zugeeignet. Wahrlich, die Kappe sieht nicht, die gelobte Keuschheit ist nicht, und andere Werke sehen noch hören nicht. Was ist es denn, daß einer auf seine Werke vertrauet, ein anderer auf seine Weisheit und Gerechtigkeit pocht? welche alle weder sehen, hören, essen noch trinken.

52. Die Werke und alles menschliche Vermögen sind nicht GOtt; dennoch trauet man darauf. Das Werk ist auch allerdings dein nicht, sondern GOttes, der es in dir wirkt. Also ist der Mammon und der Bauch auch ein solcher Gott [Phil. 3, 19.]; der thut der Dinge keines, ißt noch trinkt, sieht noch hört, und schmeckt auch nicht; da liegt er im Kasten. Ein Gulden hat auch nicht Augen, Ohren. Aber pfui dich an, daß sich ein Mensch freuen soll, wenn er Gold ansieht, und soll sich seines GOttes und Schöpfers darüber ermögen! Was ist es denn mit dem Golde und Gut? Es hat weder Ohren noch Augen, keine Sinne noch Fühlen; noch vertrauet man darauf. Pfui und aber pfui des
\quellseite{849}
schändlichen Gottes! Es ist doch ein Mensch viel besser, der sieht doch, hört und versteht; dennoch soll ein Mensch so böse sein, und sich auf den ungewissen Mammon verlassen?

53. Wahrlich, sagt Moses, wenn du das erste Gebot verlierst und GOttes Wort nicht hältst, so kommst du letztlich dahin, daß du Gold, Silber und Geld, ja, einen Pfennig, der so leicht ist, daß ihn wohl eine Fliege möchte wegführen, für Gott anbeten sollt, und darauf trauen, ob er schon weder sieht noch hört. So blind, toll und thöricht kann einer werden. Das wird uns auch begegnen, wenn wir von dem Worte GOttes abfallen werden. Siehe, welch ein fleißiger Prediger ist Moses gewesen, daß er die Leute will bei dem ersten Gebot erhalten, damit das Herz wohl lerne dem Worte GOttes glauben, und daß man alle Creaturen lasse\footnote{Erlanger: lassen.} fahren. Es ist eine gemeine Predigt, welche auch die heiligen Propheten getrieben haben, und mehr denn wir, daß man an GOtt und seinem Wort allein hangen soll; und, wiewohl es eine gemeine Predigt ist, achtet doch der gemeine Haufe ihrer nicht.

54. Das fünfte Buch Mosis soll nicht anders gehalten werden, denn für eine Auslegung der zehn Gebote, in welcher sich Mose reichlich und überflüssig erklärt, und einen großen Raum macht, was er will verstanden haben durch die zehn Gebote; und vornehmlich treibt er sehr auf das erste Gebot, daran die ganze Macht liegt. Mit den andern geringen Geboten treibt er nicht so viel Wesens; angesehen, welcher bei dem ersten Gebot bleibt, das ist, im Glauben, daß er den wahrhaftigen GOtt recht erkennt, der hat bald gelernt seinen Namen ehren, und in den andern mag er auch bedeutet werden. Darum geht er mit viel Worten überher, und thut so merkliche Predigten, und gebraucht alle Kunst, lockt, reizt, treibt, dräuet, verheißt, vermahnt, sucht zuwege allerlei Worte und Argumente, damit er es seinen Zuhörern wohl einbilde. Wie ihr denn das habt genugsam gehört, und liegt ihm auch Macht daran, wie man das erste Gebot recht fasse und dabei bleibe.

55. Denn, wie ihr im Alten Testamente leset, so ist aller Unrath und Jammer in dem jüdischen Volke daher gekommen allein um der Abgötterei willen. Darum ist oftmals so viel Volk erschlagen worden; aller Krieg, Theurung, Jammer, Blutvergießen hat sich daher erhoben, nämlich, um Uebertretung des ersten Gebots willen. Und es sind auch alle Propheten darüber getödtet worden, daß sie gestritten haben um das erste Gebot; denn, wo das weg und verloren ist, so ist alles weg und verloren. Gleichwie mir Schriften auch sagen: Wenn wir die Hauptlehre und den Zweck oder das Ziel, daß Christus unser Heiland und HErr sei, nicht haben, daß wir allein durch den Glauben an ihn, ohne alles Verdienst, aus Gnaden gerecht werden, so ist es geschehen.

56. Welcher aus dem Schiffe fällt, der wird untergehen und ersaufen, es geschehe über lang oder kurz, wenn er auch gleich zuvor drei- oder viermal vom Schiffbruch errettet wäre. Darum will er das erste Gebot geglaubt und getrieben haben, daß wir lassen GOtt unsern GOtt sein, und ihm vertrauen, und auf ihn allein bauen, und ihm die Gottheit nicht entziehen und andern zueignen. Aber das ist die größte Plage und ärgste Marter, daß wir die Abgötterei nicht lassen können. Dies Wort steckt uns noch im Mark und Beinen, welches der Teufel zu Eva und Adam sprach [1 Mos. 3, 5.]: „Ihr werdet sein wie die Götter.“ Auch ist es uns in die Natur gepflanzt und angeboren, daß wir auf Abgötter vertrauen und fremde Hülfe suchen. Der Teufel steckt uns in der Haut, der böse Same will nicht heraus, daß wir GOtt seine Gottheit nicht ließen.

57. Aber so lange das erste Gebot steht, so hat der Teufel keinen Zutritt; wenn wir aber im Fleisch sind, so hat der Teufel einen Vortheil und einen Zugang, daß ihm das Fleisch Beifall gibt. Darum soll man auf dies Gebot sehen, als auf das Hauptstück, und Moses treibt es auch um der Ursache willen so fleißig, denn alle Weisheit und Kunst, die man vorbringen kann, ist nichts gegen dem Gebot; was da kann geredet werden, ist alles in dieses Gebot gefasset.

\begin{schriftzitat}
V. 24. Ich der HErr, dein Gott.
\end{schriftzitat}

58. Es sind drei Worte, aber schwerlich zu verstehen. Darum muß man Mosen nicht verdenken; denn er bildet und bläuet sie uns nicht vergebens mit so viel Worten vor. Wir zwar zu dieser Zeit lassen es, GOtt Lob! an unserm Fleiß nicht mangeln, es hilft dennoch, was es helfen mag, solch stetiges Einträuen, Einbläuen und Einstreichen.

\quellseite{850}
\subsection{Ich bin ein eifriger GOtt}

59. Das sind erschreckliche Worte, daß er eifrig ist. Er verderbt, und ist ihm nicht zu wehren. Droben hat er GOtt genannt ein feuriges Feuer, von welchem keine Errettung ist; hier gibt er ihm noch einen Titel, und heißt ihn eifrig, als den, der es thut und will es thun, denn er ist ein Eiferer.

\begin{schriftzitat}
V. 25. Wenn ihr nun Kinder zeuget etc.
\end{schriftzitat}

60. Ich könnte es nicht heftiger reden und machen, denn so ich sage: Werdet ihr von GOtt abfallen und nicht bei ihm bleiben, so ist kein Rath mehr. Was GOtt für ein verzehrend Feuer und eifriger GOtt sei, das erklärt er hier, ruft Himmel und Erde, und alles, das da ist, zu Zeugen an. Er protestirt nicht mit einem Notarien vor Gericht, sondern Himmel, Erde, und alles, wie es steht und geht, soll seine Zeugen und Notarien sein. Es kann nicht größer Zeugniß geführt werden, auf daß er ihnen nur den hohen Artitel tief einbildet: daß dieser allein der GOtt sei, zu dem man sich alles Guten und Gnaden versehen solle; und werde man davon abfallen, so werde kein Rath mehr sein.

61. Also predigen wir auch, daß unser Behelf und Vertrauen soll stehen allein in GOttes Gnaden, und nicht auf unsern Werken. Da rufen wir auch Himmel und Erde zu Zeugen an wider den Unglauben gegen GOtt, [wider] alle Verfolger, Lästerer und Feinde. Nun will man keinen Artitel weniger leiden, denn eben diesen, welcher doch der nöthigste ist; da wollen wir nicht hinan. Wir wollen den Artitel nicht haben, niemand will GOttes Verheißung allein vertrauen. Da hat einer einen Orden erfunden, der andere ist ein Carthäuser worden, einer sonst ein Mönch. Dieses alles ist köstlich Ding gewesen. Allein GOtt vertrauen, und sich Gutes zu ihm versehen, das mag man nicht. Wenn ich reich wäre, und wollte einem Knecht oder einer Magd Haus und Hof, Acker und Gründe und andere Erbgüter aus Gnaden geben, welches ich ihm sonst nicht schuldig wäre, und er spräche: Ich will es nicht aus Gnaden und Barmherzigkeit haben, sondern will es dir bezahlen, will es abverdienen, daß er einen Rechtsspruch dazu hätte: wäre der nicht unsinnig? Also, wenn einer eine Tochter hätte, und gäbe sie mir aus Gnaden, und ich führe zu und spräche: Die Weise gefällt mir nicht, ich will es besser machen, ich will es verdienen, daß ich sie nicht aus Gnaden kriege, sondern mit Recht; und wollte diese Gnade ausschlagen.

62. Eben den Weg nimmt man ihm auch vor. Die Welt will unserm HErrn GOtt mit Recht den Himmel abgenommen, ihm abverdienen und abkaufen, da er doch läßt ausschreien durch die ganze Welt, er wolle es uns umsonst geben, und spricht: Ich will euer GOtt sein, aus Gnaden will ich es euch geben, aus Gnaden und umsonst will ich euch selig machen, wie St. Paulus Eph. 2, 5. sagt; ich will mir es nicht lassen annehmen. So machet nicht andere Götter, erfindet nicht solche Dinge, die ihr selbst thut. Das treibt diese Predigt hier auch und spricht: Hebet nicht an mit euren guten Werken, lasset mich über euch erbarmen. Es ist wohl eine Schande, daß man uns das soll vorwerfen, daß wir es nicht umsonst annehmen, sondern noch verdienen wollen, und noch GOtt zu geben gedenken, dem, der uns alles reichlich darreichen will. Solche Narren sind wir, daß wir geben wollen, da wir nehmen sollen. Der Bettler kommt hier zum gewaltigen Könige, und bettelt also, daß er nicht will das Almosen umsonst von ihm nehmen, sondern ihm irgend vier Heller oder Käuse dafür geben.

63. Also ist die Welt auch toll und thöricht; sie will dem geben, der alles gegeben hat, und von dem sie alles empfangen sollte; und wenn sie dem Nächsten etwas geben soll, da will sie nichts geben, sondern ihm nur nehmen. Das spürt man dabei: Hat einer ein schön Haus, so denkt der andere, ich wollte, daß ich es hätte; seine Aecker und Vieh begehrt er auch etc. Darum sagt hier der Prophet Mose: Werdet ihr es versehen, der HErr ist eifrig. Er wird es euch, auf meine Seele, nicht schenken; wie wir Teutschen sagen. Diese Worte mögen wir uns auch lassen gesagt sein. Denn, werden wir Christum verlieren, so wird es mit uns nicht gut werden. Wir leben zwar sicher, gleich als wäre Christus ein Narr; aber GOtt läßt sich nicht spotten, Gal. 6, 7. Das werden wir erfahren. Er wird nicht immerdar stille sitzen, wie wir meinen.

\begin{schriftzitat}
V. 27. Und der HErr wird euch zerstreuen unter die Völker, und werdet ein geringer Pöbel übrig sein.
\end{schriftzitat}

64. Das geht jetzt auf unsere Zeit auch. Denn, werden wir dies Gebot verlassen, so wird es uns gehen, gleich wie es den Juden gegangen ist;
\quellseite{851}
da wird das verzehrende Feuer kommen und uns vertilgen, die wir solches verachten, wie sie bisher gethan haben; aber sie wird dies Feuer gewiß auch fressen. Unter dem Papstthum sind wir zerstreuet gewesen, und gejagt worden, einer dahin, der andere dorthin; nicht allein leiblich, gleichwie die Juden unter den Heiden zerstreuet waren in Egypten, Persien, Meden, bis er sie durch die Römer gar zu Grunde zerstörte. Wie sie nun daselbst unter den Heiden waren, ließen sie nichtsdestoweniger den rechten GOtt fahren, und beteten der Heiden Götter an, bis auch endlich GOtt hinter ihnen her war, und spielte das Garaus mit ihnen. Also frißt GOtt um sich.

65. Eben dergleichen wird es uns auch noch gehen. Wir werden von GOtt zerstreuet werden durch allerlei Menschensatzungen, und unzählige falsche Lehren und Secten. Jetzt geht solches allbereit daher. Da kommen die Wiedertäufer und die Sacramentischwärmer und -schänder, die aufrührischen Geister, so Weib und Kinder verlassen; item, so die Gottheit Christi verleugnen. Das ist denn das verzehrende, fressige Feuer; denn du hast GOttes Wort nicht geachtet, noch ihm geglaubt, daß er uns gerne und von Herzen bei dem ersten Gebot erhalten wolle. Aber wir wollen die Barmherzigkeit GOttes, in Jesu Christo uns angeboten, nicht annehmen, ein jeder will etwas Sonderliches machen. Darum wird es uns eben also auch gehen, gleichwie es den Juden gegangen ist. Sehet, wie Moses so gar gerne sein Volk wollte behalten im Glauben des ersten Gebots. Wie wir denn auch noch fleißig solches treiben; aber man spottet und lacht unser dazu, daß wir vom Glauben so heftig predigen, so wir doch klärlich und deutlich bisanher aus dem Munde des Mannes GOttes, Mosis, gehört, wie er so gewiß dahin sieht, und alles sein Lehren, Predigen und Vermahnen auf den Glauben richtet, welcher auch vornehmlich in diesem ersten Gebot von allen Menschen erfordert wird.

\subsection[Wenn sie werden zu den Heiden kommen und zerstreuet]{Wenn sie werden zu den Heiden kommen und zerstreuet sein, werden das Reich verlieren und Knechte werden}

66. Dazu ist es auch gekommen mit den Juden; da sie zuvor Herren waren, sind sie darnach aller Welt Knechte worden, und haben anderen müssen unter den Füßen liegen. Also ist es mit uns vor dieser Zeit auch zugegangen, wir waren rechte Knechte unter dem römischen Stuhl. Ein jeglicher Bube, der nur hieß Official, oder Laienbruder, und ihres Gleichen, hat können mit einem Zettel die Könige und Fürsten zwingen. Wenn ihrer einer auf die Kanzel kam, da mußte alles fliehen, weichen, und die Fürsten und Herren vor ihnen sich fürchten; man mußte ihnen gehorsam sein. Jetzt sind wir, GOtt Lob! wieder von ihnen erlöst, und achten sie nicht werth, daß sie unsere Hausknechte sein sollten, da sie doch jenesmal herrichten über uns. Aber wie solches vom größern Theil erkannt und GOtt dafür gedankt wird, sieht jedermann wohl. Diese schändliche, leidige Undankbarkeit wird uns noch heutzutage den Hals brechen. Denn allezeit ist es so gegangen: Wenn man Christum verliert, und vom Glauben irre geht, so muß man dafür zu Lohne in greuliche Knechtschaft letztlich gedeihen. Dieses dräuet Moses seinen Juden hier auch. So wird es, sage ich, mit uns auch dazu kommen, daß uns regieren und ihres Gefallens tribuliren werden heillose Leute, welche nicht werth sind, daß sie uns die Schuhe sollten wischen.

67. Man sehe doch nur darauf, wie sich die evangelischen Fürsten jetzt stellen. Wenn man soll einem Dorfpfarrherrn seinen Zehuten und anderes sein Einkommen reichen, auch die Leute dahin halten, daß sie verdienten Lohn geben, oder ein dreißig Gulden zur Besoldung zulegen, so sagen sie, wir sind geizig, es könne die Pfaffen niemand erfüllen. Sagen auch eines Theils noch herrlich: Die Papisten haben uns betrogen, diese werden uns verführen. Darum will man jetzt Predigern ihre Unterhaltung nicht geben. Wenn man sie soll ernähren, so zählt man ihnen die Bissen Brods in das Maul. Aber es soll kommen, daß sie müssen ihren Verführern mit Haufen geben, wie vorhin, und wiederum unter die Officialen und den Pabst gezwungen werden.

68. Und wenn ich könnte, so wollte ich selbst die päbstliche Gewalt wieder anrichten, und den Pabst aufs höchste erheben über solche Buben. Aber es wird diese Zeit wieder kommen, denn dieser Text wird nicht lügen: „GOtt ist ein verzehrend Feuer, und ein eifriger GOtt.“ Es sollen wieder über sie kommen geistliche Schinder und Welträuber, und werden sie aussaugen mit ihrem Schinden, und doch nichts Gutes lehren. Nun, man lacht jetzt unser, wie vorzeiten Mosis, denn man will nicht glauben, sondern es erfahren, daß GOtt ein verzehrend Feuer sei.

\quellseite{852}
\begin{schriftzitat}
V. 28. Daselbst wirst du Göttern dienen, die Menschenhände Werk sind.
\end{schriftzitat}

69. Das Volk hat Mosen und die Propheten fleißig gelesen, darum haben sie sehr wohl gewußt, daß zu ihnen gesagt worden ist: Das wird dir widerfahren, du wirst Göttern dienen, welche genannt werden Werk [der] Menschenhände, Holz, Stein. Mit diesen Worten beschreibt er praemium impietatis [den Lohn der Gottlosigkeit], will sagen: Das sollst du zu Lohne haben; wenn du von diesem Glauben wirst abfallen und auf ander Ding trauen denn auf GOtt, welcher dir seine Gnade und Barmherzigkeit anbeut, so wirst du denn einen Gott haben, welcher nichts wird sein denn ein Stein und Holz, der nicht göttliche Werke thun oder beweisen, noch dir wird helfen können. Für den einigen GOtt wirst du bekommen unzählig viel Götter.

70. Aber das wird nicht geschehen, denken unsere Scharrhansen, daß über uns der Teufel also gebieten sollte; sprechen: Das wollte GOtt nicht, ich will mich wohl davor hüten, daß mich der Teufel nicht mehr soll dahin führen. Awe ja, es wird ihm sauer werden, du bist gewiß, daß er es nicht solle enden. Es sagt’s allhier Moses, es solle dir widerfahren. Denn wer von diesem höchsten Artikel abfällt, der hat darnach keinen andern Artikel mehr, den er halten und lernen möge. Wie ist es doch möglich? sagen sie. Ei, sehet hinter euch! wenn wir nicht unsere eigene Erfahrung hätten, so ließen wir’s uns nicht überreden.

71. War denn das nicht der Hände Werk anbeten, daß wir uns verließen auf der lieben Heiligen Fürbitte? und die Mönche bauten auf einen grauen Rock? So haben wir für Gottesdienst und heilig Ding gehalten Gabeln, Monstranzen, Kelche, Altäre, Altartücher, Meßgewande, und was nur in Kirchen und Klöstern ist gegeben, daß man theilhaftig würde ihrer guten Werke und Verdienste. Ist das nicht Holz und Steine angebetet? Frage dein Gewissen selbst drum; welcher hat nicht wollen eine eigene Kirche bauen?

72. Ja, man frage sie: Ist denn der Altar GOtt der HErr, oder das Werk, das du thust? Wahrlich, GOttes Gnade und Barmherzigkeit muß etwas Anderes sein, denn ein Werk und Verdienst, das ein Mensch in einem Orden oder sonst thut! Du mußt wahrlich sagen: Dies ist (ehe denn wir geboren worden) in GOttes Gnade und Barmherzigkeit gestanden; und haben doch diese Barmherzigkeit williglich verlassen, und auf unsere nichtigen Werke und Verdienste der todten Heiligen getrauet. Das soll man nicht leiden, denn das heißt: „Du wirst andern Göttern dienen, die Menschenhände Werk sind.“ Das ist zu meiden. Nun, man kann es nicht umgehen, wo man bei dem ersten Gebot nicht bleibt. Wer aber des Artikels fehlt und nicht behält, der muß dahin gerathen, und fällt dahin, er kann sich nicht erhalten.

73. Darum vermahnt Moses mit so viel Worten zum ersten Gebot, und schreckt uns, und spricht: „GOtt ist ein eifriger GOtt, ein verzehrend Feuer“; hängt Dräuung und Verheißung dran. Als sollte er sagen: Willst du Gnade haben? Da findest du sie, GOtt will dir gnädig sein; magst du aber der Gnade nicht, so habe dir Ungnade, GOttes Zorn und ewige Verdammniß, die sollst du mit deiner sauren Arbeit erlangen; wie denn die Hölle manchem saurer wird zu verdienen, denn anderen der Himmel; und mancher verdient mit großer Mühe und schwerem, saurem Gottesdienst und großer Unruhe GOttes Zorn. Das ist je düre und klar genug geredet. Die Menschenhände Werke sind, da wir auf ein Ding trauen, welches nichts ist, denn Menschenhände Werk, und nicht GOttes Gnade.

74. Sehet, was thun unsere neuen Rotten und Schwärmer anders, item, auch die Bilderstürmer, denn daß sie die Leute auf die Werke führen? Das ist nicht ein Christ, der auf GOttes Barmherzigkeit trauet (sagen sie), sondern welcher die Bilder getrost stürmt. Vermerke hiermit die Gerechtigkeit des Glaubens. Item, die Wiedertäufer, was thun sie, was lehren sie? Sie sagen, die Taufe sei nichts; nehmen aus der Taufe rein hinweg die Gnade, daß keine Gnade und Barmherzigkeit GOttes, keine Vergebung der Sünden drinnen sei, sondern nur ein Zeichen, daß du fromm seiest, und du mußt zuvor fromm sein, ehe du getauft wirst etc., oder die Taufe sei ein Zeichen, daß du dieselbige Frömmigkeit habest. Sie sondern die Gnade ab von der Taufe, und lassen mir allda ein bloß äußerlich Zeichen; da ist kein Fünklein der Gnade, sondern sie ist gar heraus geschnitten. Wenn also die Gnade Christi aus der Taufe hinweg ist, so bleibt ein pur Werk.

\quellseite{853}
75. Also, im Sacrament des Abendmahls des HErrn nehmen die Schwärmer heraus die Verheißung, die uns angeboten wird; sagen: Es ist Brod und Wein, wenn du es issest oder trinkest. Da ist die Gnade, so uns darinnen angeboten wird, auch hinweggeschnitten und verleugnet. Denn so lehren sie: Du thust ein gut Werk daran, wenn du allein Christum bekennst; und wenn du das Brod und Wein nur issest und trinkst im Abendmahl, so muß da keine Gnade sein.

76. Also geht es: Wenn einer von dem ersten Gebot abfällt, der richtet bald einen Abgott und ein Werk an, darauf er trauet. Darum sagt Moses: Lieben Kinder, sehet euch wohl vor, bleibet bei GOtt, dem folget nach, sonst ist euch Abgötterei unvermeidlich, ihr müßt in Abgötterei gerathen, ihr könnts euch nicht erwehren; denn die Gnade wird allezeit vom Teufel angefochten, es kann keine Ketzerei die Gnade Gottes leiden.

77.\footnote{Die Nachschriften scheinen uns hier wieder ungenügend zu sein. Wir würden etwa folgenden vermuthen: Wer einem Idolum dient, ist, wie wir Deutschen es nennen, ein Ketzer, ein Abgötzer [d. i. Abgöttischer]; die ersten [Deutschen] haben es genannt: Götzer, die hochdeutschen heißen es Ketzer etc. — Am Ende des Absatzes wird statt: „Götzen“ und Götzer zu lesen sein: Götzern und Ketzern.} Idolum ist, das wir Deutschen nennen, ein Ketzer, ein Abgötzen; die ersten haben es genannt Götzen; die hohen Deutschen heißen es Ketzer, das ist, welcher von dem ersten Gebot und Glauben abfällt und von dem rechten GOtt, und verleugnet den rechten Glauben. Und Mose zeigt uns damit an, daß es unmöglich sei, ohne Abgötterei bleiben. Wenn wir von dem ersten Gebot weichen, dann werden wir zu Götzen und Götzer.

78. Die Schwärmer heutiges Tages treiben auch alle das erste Gebot; sagen: Wir verkündigen auch Gnade und Barmherzigkeit durch Christum, und vermehren nicht den Artikel des ersten Gebots; und sagen, ich, Lutherus, lüge sie an. Aber siehe ihnen drauf: Sie bekennen den gestorbenen Christum, der am Kreuz gehangen und uns selig gemacht, das ist wahr; aber sie leugnen das, dadurch wir ihn bekommen, das ist, das Mittel, den\footnote{Erlanger: der.} Weg, die Brücke und Steig, den brechen sie ein.

79. Die Juden glauben auch, daß ein GOtt sei; aber den Weg, wie man zu Gott komme, nämlich, durch Christum, durch Christi Menschheit, verleugnen sie. Der Türke bekennt auch GOtt; aber er verleugnet den Weg, das Mittel, die Brücke, darauf man zu GOtt kommt, das ist, die Gnade GOttes. Christum wollen sie nicht haben, auch keine Sacramente, dadurch man zu der Gnade kommt.

80. Es ist gleich und geht mit ihnen, als wenn ich einem predigte: Da habe ich einen Schatz, und hielte ihm doch den Schatz nicht vor die Nase, gäbe ihm auch nicht die Schlüssel dazu; was hülfe ihm dieser Schatz? Sie schließen uns den Schatz zu, den sie uns sollten vor die Nase stellen, und führen mich auf einen Affenschwanz; den Zutritt und die Ueberreichung, den Brauch und Besitzung des Schatzes weigert und nimmt man mir.

81. Darum sagen die Schwärmer auch viel von GOtt, von Vergebung der Sünden und der Gnade Gottes, und daß Christus gestorben sei; aber wie ich Christum erlange, und wie die Gnade zu mir kommt, daß ich sie kriege, daß wir zusammen kommen, da sagen sie: Der Geist muß [es] allein thun; führen mich auf den Affenschwanz; sagen, das äußerliche, mündliche Wort, die Taufe und Sacrament, sei kein nütze, und predigen doch von der Gnade. Das heißt, mir den Schatz verkündigen, und fein davon sagen, aber den Schlüssel und die Brücke weggenommen, darauf ich zum Schatze kommen soll.

82. Nun hat es GOtt also geordnet, daß dieser Schatz durch die Taufe, das Sacrament des Abendmahls, und äußerliche Wort uns gegeben und dargereicht wird. Denn das sind die Mittel und Instrumente, dadurch wir zu GOttes Gnade kommen. Das verleugnen sie. Das sage ich darum, daß der Teufel so geschwinde ist, und bekennt diese Worte; aber er verleugnet das Mittel, dadurch wir dazu kommen; das ist, sie leugnen nicht den Schatz, sondern den Brauch und Nutz des Schatzes, sie nehmen und entziehen uns die Weise, Mittel und Wege, wie wir dazu kommen und des Schatzes genießen, und wie wir zur Gnade kommen sollen und mögen.

83. Du mußt, sagen sie, den Geist haben; aber wie ich den Geist haben kann, das wollen sie mir nicht lassen. Nun, wie kann ich den Geist überkommen und glauben, wenn man mir nicht predigt das Wort GOttes und die Sacramente reicht? Ich muß das Mittel haben, denn der Glaube kommt aus dem Gehör, das Gehör
\quellseite{854}
aber durch das mündliche Wort, Röm. 10, 17. Summa Summarum: Es kann keine Rotte aufkommen, sie muß wider das erste Gebot laufen, und an Christum Jesum sich stoßen, und werden in diesem Artikel alle Ketzer in einer Summa gesammelt.

84. Darum lasset uns bleiben bei dem Artikel: „Du sollst nicht andere Götter haben“, und auf diesen Zweck und scopum fleißig Achtung haben. Denn, lassen wir es uns aus den Augen thun, so ist allen Rottengeistern Thor und Thür aufgesperret. Gott hat nie ohne äußerliche Mittel seinen Gottesdienst in der Welt haben wollen. Im alten Testament hat er den Juden eine Weise gegeben, dabei man ihn sollte finden; da war ein gewisser Ort des Tabernakels oder Hütte des Stifts, der Altar, Leuchter, die Leviten, und ließ GOtt sich nicht finden ohne äußerliche Mittel und Weise. Er hat ihnen allewege ein äußerliches Mittel vorgeschlagen, damit sie ihn ja finden sollten; er ließ sie nicht ohne Weise und äußerlich Mittel in der Irre herumgehen.

85. Aber, wie unsere Schwärmer jetzt laufen und lassen von dieser Weise, welche uns Gott gegeben hat im neuen Testament, also verließen die Juden auch dieselbige Weise, und suchten andere Wege. GOtt kann nicht unser GOtt sein, er gebe uns denn etwas Neuerliches, daran wir ihn finden, als, das mündliche Wort und die zwei Sacramente. Wenn ich GOtt nicht ergreife durch äußerliche Dinge, wie kann ich ihn denn antreffen?

86. Darum sind alle Ketzer wider das erste Gebot gewesen, und haben sich daran vergriffen, in allerlei Menschenwerken, und schneiden aus die Verheißung und Gnade GOttes, so darein gesteckt ist; ja, verleugnen GOtt selbst, verwerfen den Nutzen und Brauch, daß man zu der Gnade nicht kommen möge.

\begin{schriftzitat}
V. 29–31. Wenn du aber daselbst den HErrn, deinen GOtt, suchen wirst, so wirst du ihn finden, wo du ihn wirst von ganzem Herzen, und von ganzer Seele suchen. Wenn du geängstiget sein wirst, und dich treffen werden alle diese Dinge in den letzten Tagen, so wirst du dich bekehren zu dem HErrn, deinem GOtt, und seiner Stimme gehorchen. Denn der HErr, dein GOtt, ist ein barmherziger GOtt; er wird dich nicht lassen, noch verderben, er wird auch nicht vergessen des Bundes, den er deinen Vätern geschworen hat.
\end{schriftzitat}

87. Ich wollte gerne den sehen, der so gelehrt wäre, welcher diesen Text könnte den Juden nehmen und wegthun, ausgenommen die Apostel. Es ist ein sehr gewaltiger Text für die Juden. Er hält aber dieses in sich: In allen Ängsten, die dich treffen werden, wenn du von GOtt wirst sein abgefallen, da wirst du schreien und GOtt suchen, und der HErr wird sich erinnern seines Bundes. Da hat ein Jude aus diesem Text: Haben wir gesündigt, und sind von GOtt abgefallen, jetzt suchen wir den HErrn von ganzem Herzen. Darum wird er uns nicht verlassen, und sich unser wieder erbarmen, und gedenken an seinen Bund, so er mit uns gemacht hat. Und scheint dieser Text stark wider das neue Testament [zu] sein.

88. Es ist aber ein gemeiner Text, daß GOtt die nicht will verlassen, welche sich herzlich bekehren, und wieder an ihn glauben: ob sie gleich gefallen sind, er will sie wieder zu Gnaden annehmen. Wie ihr denn seht, daß uns auch geschieht. Obgleich wir die Strafe verdienen, doch, so wir nach der Uebertretung uns wieder bekehren zu GOtt, so kommt die Barmherzigkeit wieder zu uns; wie der Prophet Joel Cap. 2, 12. 13. auch sagt: „Bekehret euch zum HErrn, eurem GOtt, denn er ist gnädig, barmherzig, geduldig, und von großer Güte, und reuet ihn bald der Strafe.“ Daß aber droben [V. 24.] gesagt ist, GOtt sei eifrig und ein verzehrend Feuer, ist darum geschrieben, auf daß man nicht sicher noch vermessen sei, wenn GOtt nicht straft; es wird doch gewißlich letztlich geschehen. Darum soll sich einer nicht darauf vertrösten, daß GOtt barmherzig sei, wie allhier der Text sagt, und drauf hingehen und sündigen; denn du wirst es erfahren, daß er auch ein verzehrend Feuer sei, und wirst du sündigen, so wird er dich auffressen. Wiederum, wenn du bist gefallen, und kannst dir nicht helfen, so wird dich GOtt nicht verlassen, sondern wird dir wieder helfen, daß du wieder seine Gebote haltest.

89. Er ist ein barmherziger GOtt, straft nicht also, daß er gar ausrotte und vertilge, als der Teufel sonst thut, daß keine Hülfe noch Rath mehr da sei. Er schlägt mit Pestilenz, Hunger, theurer Zeit und Krieg, verwüstet Länder und Städte, verderbt Vieh und Menschen; aber GOtt erhält dennoch einen oder zwei, die wiederum eine Stadt ansehen zu bauen. Als, in der Sündfluth behielt er ihm acht Menschen im Kasten Noäh
\quellseite{855}
[1 Mos. 8, 18.], also erhielt er Lot und seine Töchter zu Sodoma [1 Mos. 19, 30.]. Aber diese Predigt wird denen gesagt, welche in der Strafe sind, und wollen sich bekehren und glauben; für dieselben gehört solche tröstliche Verheißung. Den andern aber, so gottlos, verwegen, und ohne alle Beisorge oder Betrachtung zukünftiger Strafen dahin leben, soll man etwas vom Teufel und der Hölle sagen; sie werden von solcher Predigt nur ärger. Also ist eine gemeine Rede, daß GOtt mit Maßen straft, daß er nicht alles zu Grunde verderbe, auf daß man nicht verzweifle; denn die Strafe gehört für die da böse sind, und nicht wollen aufhören zu sündigen.

90. Aber daß die Juden diesen Text anziehen für ihre Erlösung, nachdem GOtt Jerusalem zerstört habe, und sie noch hoffen, daß ihr Königreich wieder solle aufgerichtet werden, da sage du ihnen wieder, daß hier auch geschrieben steht: „Wo du ihn von ganzem Herzen, und von ganzer Seele suchen wirst, so wird GOtt dir ein barmherziger GOtt sein.“ St. Paulus handelt diesen Spruch in der Epistel zum Römern Cap. 11, 31. f. gewaltiglich. Wer könnte auch dies Räthsel sonst auflösen, wenn es nicht St. Paulus hätte aufgelöst?

91. „Wenn du geängstiget sein wirst.“ Er sagt am selben Orte nicht, daß Gott die Juden will ungestraft lassen, wie die Juden meinen, sondern er will sie allzumal hernehmen. So ist auch dieses seine Meinung nicht, daß er sie alle werde zu Gnaden nehmen, und in sein Reich setzen, sondern, wiewohl er sie gestraft hat, so hat er doch ihrer viel angenommen zu Gliedern der christlichen Kirche, nämlich, welche wieder zurecht gekommen sind; wie auch heutiges Tages noch etliche bekehrt werden.

92. Also könnten die Papisten heutiges Tages auch sagen, GOtt verläßt nicht seine Kirche, Christus bleibt wahrlich bei seiner Kirche bis zum Ende der Welt, er läßt sie nicht sinken; wer wollte ihnen das nehmen? Denn Christus hat das selber zugesagt Matthäi am 28. Capitel, V. 20.: „Ich bin bei euch bis ans Ende der Welt.“ Und in Johanne Cap. 14, 18.: „Ich will euch nicht Waisen lassen.“ Darum schließt der Papst und die Bischöfe: Wir werden wohl bleiben, und nicht untergehen, denn wir sind die christliche Kirche. Ja, antwortte du: Wenn man sich von ganzem Herzen und ganzer Seele bekehrt.

93. Derhalben sagt Moses nicht: die da sich begehren in Sünden unthätiglich zu fühlen, wie eine Sau im Koth, und sich wider GOtt setzen, die will er aufrichten, sondern: Die den HErrn suchen von ganzem Herzen und ganzer Seele. Er will keinem Volke, es seien Juden oder Heiden, gnädig sein, es sei denn, daß sie durch Angst und Noth bekehrt werden, das ist, GOtt von ganzem Herzen glauben. Darum kannst du nicht sagen, daß er irgend einem Papisten, oder sonst unbußfertigen, Barmherzigkeit verheißen habe, sondern allein denen, die sich mit ganzem Herzen zum HErrn bekehren, es sind gleich Jude oder Heide, Papist oder böser Christ, Papst oder Bischöfe. Allen, die von Herzen zu ihm seufzen und rufen, und von der Abgötterei und ihrer Hände Werk und andern Sünden abfallen, denen verheißt er die Barmherzigkeit. Darum dürfen die Juden nicht stolzieren, daß sie GOtt werde wiederbringen; denn sie rufen GOttes Barmherzigkeit nicht an, sondern rühmen noch ihre Werke, fleischliche Geburt und Ankunft, daß sie Abrahams Same sind, Joh. 8, 33.

94. Darum ist dieser Spruch zum Schein stark für die Juden und Papisten; GOtt wird sein Volk und Kirche nicht verlassen. Aber sage du, welche sein Volk und Kirche seien, nämlich, die von ganzem Herzen und ganzer Seele den HErrn suchen. Lasse uns die Kirche also beschreiben, daß es sind diejenigen, welche auf nichts trauen, denn allein auf GOttes Barmherzigkeit, und erkennen das erste Gebot; die von ihrem falschen Glauben abtreten. Das heißt GOttes Volk und die christliche Kirche, die auf nichts Anderes denn auf GOttes Gnade und Barmherzigkeit sich verlassen. Die andern mögen sich die Kirche rühmen, oder auch wohl Engel nennen, sie sind es aber darum nicht. Verkleidet sich doch der Teufel auch wohl in einen Engel des Lichts, 2 Cor. 11, 14., er ist es aber darum nicht. Darum mögen sie sich wohl bekehren, und von ihrem falschen Vertrauen ihrer Werke ablassen, denn es ist damit gar nichts. Diesen Spruch habe ich wollen handeln um der Juden und unserer Papisten willen, welche solches stets im Maul haben, und sehen gleichwohl dasjenige nicht, so dabei gesetzt ist von rechtschaffener Bekehrung und Ablassen vom gottlosen Vorhaben.

95. Bisher haben wir gehört, wie Moses so ernstlich gefordert hat, daß die Kinder Israel
\quellseite{856}
Gottes und seiner Gebote nicht vergessen sollen, oder ihnen nicht Götzen und einigerley Bilder, sie sind auch wie sie wollen, machen sollten, denn GOtt sei ein verzehrendes Feuer. Eine solche ernstliche, große, treffliche Vermahnung läßt Moses vorhergehen, ehe er die zehn Gebote aufsagt, auf daß er sie erweckte und ermunterte, damit sie es ihnen einen Ernst sein lassen, auf daß sie nicht meinen, sie hörten etwan eine Fabel, sondern es sei ein rechter Ernst, weil er GOtt nennt ein verzehrend Feuer, ein Feuer, das da frißt, und dazu eifrig, der da nicht nachläßt. Als wollte er sagen: Es muß gefressen und verzehrt und verwüstet werden, welcher die zehn Gebote nicht hält.

96. Das „Feuer“ ist das gewaltigste und heftigste Element unter allen. Was man nicht bezwingen, biegen, zerbrechen noch ändern kann, das greift man mit Feuer an, so wird es schmeidig gemacht und gezwungen. Es kann Silber, Gold, Eisen und alle Metalle zwingen. Darum vergleicht GOtt sich diesem Element. Als wollte er sagen: Ihr könnt mir nicht entlaufen, hebe ich dermaßen an Hausführung zu thun, so kann mir niemand widerstehen. Sündigt nun jemand wider meine Gebote, so will ich ihn wohl finden, und wissen zu strafen.

97. Darnach nennt er ihn „eifrig“, der es nicht lassen will, er muß über seinem Worte halten. Da ist beides beisammen, die Kraft und der Wille, Stärke und Macht, daß er kann und will strafen. Wenn wir solches für die Wahrheit hielten, so würden wir seine Gebote nicht so verachten; aber niemand glaubt es, daß er der sei, bis wir es erfahren.

98. Die zu Jerusalem haben auch nicht geglaubt, daß sie wider GOtt thäten, und daß sie von ihm gestraft sollten werden, bis die Römer kamen und die Stadt zerstörten, daß nicht ein Stein auf dem andern blieb, Luc. 19, 44. Also, ein böser Mensch, der seinem Nächsten Schaden thut, stiehlt und raubt, der glaubt auch nicht, bis der Henker kommt und knüpft ihn an den Galgen. Desgleichen alle anderen Laster, Schande und Sünde, welche wider die zehn Gebote gehen, als, GOttes Verachtung, Fluchen, Mord, Ehebruch; die haben das Urtheil schon über dem Hals, daß dieser Eifer, Zorn und Feuer über sie wird kommen, und sie vertilgen. Es sei denn, daß sie Buße thun, sonst wird keiner in der Welt entlaufen können, er thue was er wolle, er komme denn der Strafe zuvor mit wahrhaftiger Buße; denn so will GOtt mit der Strafe an sich halten und schonen, wie Jeremia am 18. Capitel, V. 7., gesagt wird: Wenn sich ein gottlos Reich, dem ich dräuete, es gar zu verstören, bekehren wird von seinem gottlosen Wesen, so soll mich auch gereuen des Unglücks, das ich über sie gedachte, und will denn dasselbige Reich bauen und aufrichten.

99. Darum soll man von Jugend auf die zehn Gebote lernen und üben, daß man wisse, was man thun oder lassen soll, und was für gute Werke einem jeden in seinem Stande gehören. Denn lebt man gottlos, so wird man der Strafe nicht entfliehen. Entgehet einer aber hier, daß er mit der Strafe verschont wird, so wird er doch sein Gericht anderswo bekommen. Darum, entgehet er an einem Orte dem Staubbesen, so bekommt er doch anderswo einen Strick dagegen. Summa Summarum: Endlich entläufst er doch diesem Feuer nicht, denn GOtt spricht: „Ich bin ein verzehrend Feuer.“ Als wollte er sagen: Nimm dir es nur nicht vor, du seist auch wer du wollest, du sollst meiner Gewalt nicht entlaufen; wenn du auch gleich aus der Welt liefst, so sollst du doch mir nicht entwerden, Ps. 139, 7. ff.

100. Das sehen wir auch in der Erfahrung. Es geht gar wunderlich zu, daß ein Mörder, Ehebrecher, Dieb der Strafe zu Theil wird. Denn GOtt hat die Obrigkeit dazu geordnet und eingesetzt, daß sie solche Sünden strafe, und es kommt oft plumpsweise, daß solche Buben gestraft werden, ehe man sich umsieht. Darum lerne eben so mehr, von Jugend auf GOtt fürchten und seine Gebote halten. Denn du kannst nicht entlaufen der Strafe. Entläufst du der Strafe in diesem Leben, so fährst du doch in die Hölle, und mußt des ewigen Todes sterben, und also deinen verdienten Lohn dennoch empfangen.

101. Aber gottlose Herzen verachten es, und glauben es nicht, daß GOtt so nahe sei, versehen sich gar nicht des Feuers; sehen nicht den Blitz vor dem Donner, das ist, sie bekommen erst die Strafe nach begangener Sünde. Es ist aber in der Wahrheit kein Scherz, es kommt der HErr, ehe denn diejenigen, so GOtt nicht fürchten, sich’s versehen. Unterdeß sieht er zwar durch die Finger, und stellt sich, als wollte er nicht strafen, und sie sündigen auch redlich drauf, gleich als sähe es GOtt nicht, und wisse nicht
\quellseite{857}
drum [Jes. 42, 14. Ps. 50, 21.]. Wohlan, ihr hört allhier, daß GOtt es ungestraft nicht läßt, sie sollen es nicht also hinweg tragen, das Feuer soll über sie kommen. Dieses habe ich wiederholen müssen, um dieser Weise zu reden willen, auf daß ihr wisset, was es sei, wenn die Schrift sagt: „Ein verzehrend Feuer“; item, „eifrig sein“. Darum hat er droben [V. 25. 26.] auch gesagt: „Wenn ihr nun Kinder zeuget und im Lande wohnen, und werdet euch Abgötter machen, und Uebels thun vor dem HErrn, eurem GOtt, und ihn erzürnen, dann werdet ihr nicht lange darinnen wohnen.“

102. Und solches aus dieser Ursache (wie er allhier spricht), denn es wird das Feuer kommen über euch, und ihr werdet schlecht vertilgt werden; so bald ihr Abgötzen machen werdet, so wird euch das Feuer treffen. Und er legt es selber aus, welches das Feuer sei; nicht, daß allezeit vom Himmel falle, und die Abgöttischen, Mörder, Diebe und Schäffte verzehrt, sondern also, daß ihr gewiß seid der Strafen, dieselbigen habt ihr allerdings vor euch. Und sagt ferner: „So rufe ich Himmel und Erde über euch zu Zeugen an, daß ihr bald werdet umkommen von dem Lande, in welches ihr gehet über den Jordan, daß ihr es einnehmet; ihr werdet nicht lange drinnen bleiben, sondern werdet vertilget werden.“ Als sollte er sagen: Daß ihr es je gewiß seid, daß ihr bald werdet umkommen, so wisset, ungestraft wird es nicht bleiben; denn GOtt ist ein Feuer, das verzehrt, frißt und eifert; das ist, er bringt euch um, wie das Feuer ein Haus verzehrt, zu Asche und Staub macht.

103. Da seht ihr, daß er das erste Gebot am allermeisten treibt. Es ist ihm ein Ernst und allein zu thun um die Bilder und Abgötterei des Herzens, und zeigt sonderlich an, und spricht: „Wenn ihr nun Kinder ziehet“, das ist, daß ihr wachst, euch mehrt, und daß es mit euch dahin kommt, daß ihr das Land einnehmt, und GOttes vergesst, der euch aus Egypten geführet hat, und Abgötterei stiftet, und euch richtet nach der Weise der Länder, darinnen ihr wohnt, so wird die Strafe nicht lange außen bleiben.

104. Unter uns Deutschen ist sonst ein Sprüchwort, das lautet also: Mancher verdammt seine Seele um seines Weibes und Kinder willen, oder um des leidigen Guts willen; und es ist schier allzuwahr. Das zeigt hier Moses auch an. Wenn GOtt einem beschert Weib und Kind im Hause, auch andere Nahrung mehr, welches ihm sollte eine Ursache sein, auf daß er desto mehr GOtt vertraute, so mißbraucht er dieselben, und vertraut und glaubt GOtt weniger denn vorhin, da er nichts hatte. Nun, weil er sie hat, fällt er in die Bauchsorge, wie er sie hoch erziehe; arbeitet dahin, daß sie einen höhern Stand möchten haben denn er selbst, oder ein jegliches ja so viel habe, als er hat; denkt nicht, der ihm Leib und Leben gegeben hat, werde auch Futter und Decke geben, Kleidung und Nahrung bescheren [Matth. 6, 31.]. Allobald fällt man in Mißglauben von GOtt, und geräth in den verfluchten Geiz [1 Tim. 6, 9. 10.].

105. Das ist sehr gemein in der Welt; denn ein jeder will seine Kinder versorgen, daß er sich nicht läßt begnügen an der schlechten Nahrung, so ihm GOtt beschert, sondern er will noch höher fahren, denn GOtt ihm gibt; darum muß Abgötterei des Geizes folgen, Eph. 5, 5. Denn, wo Sorge der Nahrung angeht, so folgt denn bald ein Scharren und Kratzen darauf, und der Mammon regt sich weidlich; da ist denn GOtt verloren. Deß könnte ich aus den Historien viel Exempel sagen; aber wir stecken selbst drinnen bis über die Ohren.

106. Darum sagt Moses: GOtt wird dich segnen, er wird dir Kindeskinder geben, daß du sie aufziehest, und zu seinem Dienste, Ehren und Furcht haltest, und er will dir Speise und Nahrung dazu geben [Ps. 128, 3.], allein, daß du GOtt für deinen HErrn GOtt auch haltest. Aber wir meinen nicht anders, denn GOtt habe unser vergessen, und sind also die Kinder eine Ursache und Occasion zum Unglauben, die uns doch GOtt gegeben hat, daß sie uns zum Glauben reizen sollten; dann kommt auch der leidige Mammon, und wird angebetet in allen Winkeln.

107. So ist es nun Mosi alles zu thun nur um das erste Gebot. Thue dasselbige, so hast du die andern alle leichtlich auch gehalten; denn wider das erste Gebot geht es alles. Darum, wenn man GOttes vergisset, ihn verachtet, und nicht fürchtet, man sucht etwas Anderes, darauf man trauet, so ist es bald geschehen. Das ist der vornehmste Gebrechen. Da verläßt sich der auf seinen Fürsten, jener auf einen reichen Freund, der auf seinen Acker: ein jeder sucht einen Gott, darauf er sich gedenkt zu verlassen, dadurch er sich und die Seinen ernähren will, gleich als wäre sein GOtt im Himmel, der es thun könnte, der
\quellseite{858}
da Kinder gäbe und Weiber schaffte, und auch Nahrung und Unterhaltung bescheren würde.

108. Die wahre Abgötterei ist im Herzen, daß man GOtt nicht fürchtet, noch ihm vertraut, sondern das Datum auf andere Dinge setzt. Die äußerlichen Bilder sind nicht die rechte Abgötterei; die äußerliche Abgötterei ist ein grob Ding. Nun, zur Warnung so sagt er: Wisse, das Feuer wird dich aufzehren, und wird dich alsdann deine Sorge und dein Mammon nicht helfen, sondern du wirst ausgerottet werden aus dem Lande, darinnen du wohnst; darüber rufe ich zu Zeugen Himmel und Erde. Er wollte es ihnen ja gerne ins Herz bilden; er dräuet ihnen hart, auf daß er die Leute gottesfürchtig mache, und es ihnen ins Herz pflanze.

109. Und ist zumal fein, daß uns geschiehet nach unserm Unglauben; denn was wir fürchten, das widerfährt uns, Spr. 10, 24. Deß haben wir die Juden zum Exempel; die fürchteten sich, daß, wo sie Christum lebendig ließen, so würden die Römer kommen und ihre Lande zerstören, Joh. 11, 48. Und da sie Christum erwürgt hatten, geschah ihnen also, daß sie und ihre Kinder zu Bettlern wurden.\footnote{Im Original: worden.} Und die noch also scharren und kratzen, und fürchten, wo sie sich allein auf GOtt verließen, so würde es ihnen an der Nahrung mangeln, wollen also selbst ihren Kindern zu großen Ehren und Gütern helfen, sie zu Reichthum und hoch empor bringen, und versorgen: aber ihr Gut verschleißt und verschwindet oft, daß niemand weiß, wo es hinkommt, und es gebeihet ihnen eben letztlich dahin, daß sie Bettler werden mit ihren Kindern. Wir sehen, daß ihrer Viele Schätze sammeln, wenn dieselben in der Erben Hände kommen, so zerrinnen sie, und sind ihrer wenig gebessert. Wie kommt das? Aus diesem Texte ist es klar und offenbar. GOtt will ein Feuer senden, und solche Kinder umbringen, ja, auch solch übel gewonnen Gut verbrennen, wie sonst auch Salomon sagt: Erosio in domo impii. Das Wort steht da, die Erfahrung ist vor unsern Augen, sie sollen dem Feuer nicht entlaufen. Es muß doch wahr werden das Sprüchwort: Uebel genommen, schändlich zerronnen. Item: De male quaesitis non gaudet tertius heres.

110. Wiederum, wo fromme Leute GOtt trauen, und arbeiten nach Gelegenheit ihres Standes, die ernährt GOtt fein, und kommen weiter denn die andern, so reich sind. Wahrlich, die Reichen regieren nicht auf Erden, sondern am meisten kommen empor, welche der HErr hervorzeucht als arme Aschenbrödel; die herrschen in der Welt. Darum singt David auch im 113. Psalm, V. 7. 8.: „Der den Elenden aus dem Staube erhebet, und setzet ihn neben die Fürsten seines Volks.“ So singt es Maria in ihrem Lobgesange auch, Luc. 1, 52., „daß GOtt erhöhe die Niedrigen“. Und St. Petrus in seiner Epistel [1 Petr. 5, 5.] spricht, er gebe den Demüthigen Gnade, wie das die Exempel der heiligen Schrift klärlich ausweisen. Als, Saul, der von den Eseln zum Königreiche genommen ward [1 Sam. 10, 1.]. Und David, hinter den Schafen gefunden, mußte zum Könige über Israel gesalbt werden [1 Sam. 16, 11–13.]. Das kommt alles daher, daß GOtt spricht: Laß mich dich versorgen, ich gebe dir Kinder und Nahrung; willst du nicht, so bleib ein Bettler. Denn, wer mein Wort: Ich bin ein Feuer, nicht will glauben, der erfahre es. Es scheint wohl, als daß andere Leute mit Reichthum und Freundschaft viel können thun und helfen, aber es ist nichts. Glaube mir allein, und fürchte mich, spricht GOtt, so sollst du genug haben; sonst ist es kein Wunder, daß du ein Bettler bleibst, verderben müsse[t?], verzehrt und umgebracht wirst. Wer es nicht will glauben, der fühle es.

\subsection{Mache dir kein Bild}

111. Das ist, traue mir, und keiner andern Creatur. Mißtraue mir nicht, setze dein Herz auf nichts Anderes denn auf mich. Ob du nicht überflüssig hast, liegt nichts dran. Hast du es nicht mit Scheffeln, so hast du es doch mit Löffeln. Bist du nicht ein Bürgermeister oder Edelmann, das schadet dir nicht, bleibe eine Weile noch ein Bürger und Bauer. Traue du GOtt, so wirst du genug haben; suchst du aber einen andern Gott, als, den Mammon, so will ich dich aufressen, und in den Grund verderben mit deinem heillosen Götzen und Kastengott. „Wer Ohren hat zu hören, der höre“ [Matth. 11, 15.].

\begin{schriftzitat}
V. 27. 28. Und der HErr wird euch zerstreuen unter die Völker etc. Da werdet ihr fremden Göttern dienen, welche Menschenhände Werk sind.
\end{schriftzitat}

\quellseite{859}
112. Das wird uns auch widerfahren, und nicht besser werden wir es haben. Nur immer wieder angefangen das Leben, welches unter dem Papstthum gewesen ist. Man hole wieder Prediger, die da Kappen tragen, Platten scheren, härene Hemder anziehen, auf Holzschuhen gehen, und ihre lausigen, stinkenden Werke für einen Gottesdienst rühmen etc. Denn wir wollen nicht hören von guten Werken, so in den zehn Geboten uns befohlen sind; darum laß man uns nur martern und bläuen mit den guten Werken des Papsts und der Mönche; dieselbigen heißen denn solche Götter, die Menschenhände Werk sind. Ist es nicht wahr, die Worte verstehen wir nicht, aber die Erfahrung haben wir? Stelle dir vor die Augen des Papsts, der Fürsten und der Juden Wesen, siehe, was sie für einen Gott haben, und wem sie dienen.

113. Die Türken fasten im Sommer einen ganzen Monat, baden alle Tage, trinken nicht Wein, leben nicht im Hause mit Fressen und Saufen, wie wir Deutschen. Wem dienen sie aber mit solcher Kasteiung und Eingezogenheit, und daß sie also stetig kriegen, Leib und Leben wagen, und ihr Blut vergießen? Ihrem Mahomet, welcher sie verführt hat, damit plagen und martern sie sich. Und das geschieht nicht GOtt zu Ehren, sondern dem Teufel Mahomet, welcher sie diesen falschen, erdichteten Götzendienst also gelehrt hat, als, fasten, nicht Wein trinken, kriegen etc. In diesen Werken üben und zermartern sie sich, und haben nichts davon denn Mühe und Arbeit. Ihr Glaube steht nicht auf GOtt, sondern in Vermessenheit ihrer eigenen Werke, dadurch sie wollen selig werden. Dies ist denn nicht GOtt vertrauet, sondern es sind die Werke [der] Menschenhände, das ist, ein Wesen und Leben, oder Werk, von Menschen erdacht; es heißen Menschenwerke.

114. Also, die Juden gehen mit Werken um, welche nach ihrer Selbstandacht erfunden und erdacht haben; essen nicht Blut, nicht Schweinefleisch. Solches sind eigentlich Menschenwerke, da dienen sie ihren Götzen.

115. Die Papisten mit ihren Mönchen und andern sehr vielen Ordensbrüdern und Anhängern halten ihre Regeln, Gebote und Satzungen, von Menschen erdichtet; es wird ihnen sauer Tag und Nacht. Das muß nun also sein, und kann nicht anders gehen. Denn welcher vom ersten Gebot abfällt, der fällt in ein Wesen, welches von Menschen erdacht ist, er geräth in menschliche Werke. Denn da ist zwischen GOttes Gnade und Ehre, und den Menschenwerken kein Mittel. Wenn du auf GOtt nicht trauest, so fällst du auf ein Werk oder Stand, welchen du erwählst.

116. Also, die Rottengeister fallen vom ersten Gebot und hangen nicht an GOttes Gnade; darum sind sie herunter gepurzelt, haben einen Abgott, das ist der, daß sie das Sacrament zu Brod und Wein gemacht. Da bläuen sie sich mit, und das ist ihr Abgott.

117. Darum, entweder es muß Gott gedienet sein, oder einem Abgott vertraut, oder von GOtt gar abgefallen, und sich auf die Werke und anderes verlassen. Denn da ist es unmöglich, daß ein Mittel wäre zwischen dem Vertrauen und Werken. Es muß einer entweder GOtt vertrauen, oder guten Werken. Das will Mose hier sagen: Werdet ihr nicht GOtt dienen, so werdet ihr der Heiden oder heidnischen Götzen und Menschenwerken dienen, denn man muß GOtt dienen, oder abgöttisch werden, und auf Werke fallen; es läßt sich nicht mitteln.

118. Diesen Wechsel macht hier Moses. Wer das nur glaubte, das wäre ein groß und sehr nützlich Ding; und muß doch entweder geglaubt sein, oder auf ein Werk gefallen und zum Heuchler geworden. Die nun mit Werken umgehen, das sind auch Abgöttische. Jetzt rede ich nicht mit den groben, äußerlichen Mißbräuchen,\footnote{Mißbräucher = einer, der Mißbrauch treibt.} sondern die da wollen fromm sein. Dieses alles ist darum so geistlich verkehrt, auf daß ihr merket, wie es Moses alles zu thun ist um das erste Gebot.

\begin{schriftzitat}
V. 29. Wenn du aber daselbst den HErrn suchen wirst.
\end{schriftzitat}

119. Damit sie aber nicht gänzlich verzweifeln, setzt er ihnen diesen Trost hinzu. Er läßt es nicht bei dem Schrecken und Träuen allein bleiben, daß GOtt ein Feuer sei den Bösen, das ist, der die Ungehorsamen umbringe, die seine Gebote verachten, und ihn nicht fürchten. Denn, wenn er es so ließe bleiben, so möchten die Herzen verzweifeln; darum setzt er es beides. Träuen und Schrecken um der Verstockten willen, welche GOttes Gebot und Lehre in Wind schlagen und verachten; meinen, sie mögen nur frei sündigen, sie werden darum bald nicht gestraft. Wiederum
\quellseite{860}
auch Verheißung denen, welche die Strafe trifft, und das Unglück über sie kommt, daß, wenn sie das Feuer angreift an Leib und Gut, daß sie dann nicht verzweifeln. Da ist es hoch vonnöthen, daß die getröstet werden, derer Herz in Verachtung GOttes gesteckt ist, nun aber sich erkennt und Gnade begehrt.

120. Ein menschlich Herz ist eitel eiserner Stahl, ein Fels und Adamant, wenn es gottlos ist. Wiederum, wenn es erschrickt, durchbrochen, zermalmt und flüchtig ist, so ist nichts Weicheres, Furchtsameres noch Blöderes. Ein Löwe und Pferd sind die freudigsten Thiere, so sie muthig sind; wiederum, wenn sie schüchtern werden und erschrecken, so ist kein blöder Thier unter allen, als sie; man kann sie nicht halten. Je größerer Muth, je verzagter, wenn das Herz entflieht. Also ist auch eines Menschen Herz ungleich gesinnt. Wenn es wohl geht und das Feuer nicht frißt, so ist es stolz und ungelent; wiederum, wenn das Feuer angeht und angreift, so wird es gar verzagt, daß es nicht leicht zu erhalten ist; man kann es nicht wieder aufrichten, wenn es zu Boden sinkt. Darum ist es beides noth zu thun: die Verstockten zu schrecken, daß sie nicht für und für stolzieren und Muthwillen treiben; wiederum, auch die Erschrockenen zu trösten, und süße Worte ihnen zu predigen, auf daß sie nicht verzweifeln, sondern fest halten. Solches thut Moses allhier auch, und spricht:

\begin{schriftzitat}
V. 30. Wenn dich nun treffen werden alle diese Dinge etc.
\end{schriftzitat}

121. Gott hat seine Gebote gegeben, daß sie die Blöden trösten, und wiederum, die Halsstarrigen schrecken sollen. Wir haben aber bisher gehört, wie mancherlei Rede er geführt hat, damit er das erste Gebot treibt, und uns reize und locke, dasselbige zu halten. Mit den andern Geboten hat er nicht so viel Wesens; hier aber hat er sonderlich vor sich genommen die Kinder und Landes Gewohnheit. Das ist auch eine Sache, die da hindert und abzeucht von dem ersten Gebot, und wird doch am andern Orte befohlen, daß sie ihr Haus versorgen sollen. Aber hier, da es das erste Gebot betrifft, sagt GOtt: Kind hin, Kind her, du sollst um deines Weibes oder Kinder willen dir keinen Abgott machen und ihn anbeten, das ist, anderswo Trost suchen, denn bei mir, noch mich verachten. Höher sollt du mich halten, denn deine Kinder; sie sollen dich von mir nicht abwenden, daß du um ihretwillen mein Wort verachtest. Und solches lehrt und gebeut auch der HErr Christus [Matth. 19, 29.], daß man um seinetwillen Haus, Hof, Acker, Weib und Kinder, und alles soll verlassen; denn wer das nicht thun werde, der sei sein nicht werth [Matth. 10, 37.]. Und allhier spricht er: Thut du es aber nicht, und dienst dem Mammon, so wird das Feuer kommen und dich fressen; trauest du mir, so wirst du genug haben.

122. Er trifft uns fein, und weiß wohl, wie es zugeht in der Welt. Will derhalben GOtt allhier sagen: Wenn du Kinder hast, und des Landes Gewohnheit vor dir, und bist genenget unter die Heiden, und hast ihre Sitten und Werke im Angesicht, und man zu dir sagt, wie das deutsche Sprüchwort lautet: Thue wie die andern, so irrst du oder narrst du nicht, so wird es mit dir bald geschehen sein, daß du von dem ersten Gebote wirst abfallen.

123. Die Welt ist voller Aergerniß. Es richtet niemand sein Herz zu GOtt. Alle Menschen auf Erden beten schier den Mammon an, und die wenigsten dienen GOtt nach dem ersten Gebot. Darum geht es also, daß man spricht: Wo mein Rathsherr und Nachbar hinfährt, da fahre ich auch hin; wie es ihm geht, also geht es mir auch; ich sehe niemand, der ihm etwas darum thue. Wer unter den Wölfen ist, spricht die Welt, der muß mit ihnen heulen. Item: Si fueris Romae, Romano vivito more; man muß temporisiren, und sich nach den Leuten richten. Item, man sagt im Papstthum: Ich bleibe bei der Religion meiner Vorfahren; sind die nicht zum Teufel gefahren, so komme ich auch nicht in die Hölle. Und mancher gibt vor, er wolle bei dem großen Haufen bleiben, als, den Päpsten, Cardinälen, Bischöfen, großen Potentaten, Fürsten und Herren; führen die denn in die Hölle, so habe er gute Gesellen, hinter denen wolle er herfahren. Also bringt uns denn das Land, darinnen wir sind, mit seinen Aergernissen dahin, daß wir nach des Landes Exempel auch gebaren, und GOttes drüber vergessen.

124. Davon sagt hier Moses: Wenn du wirst in dein Land kommen, so wirst du GOtt an das Gras schlagen. Wenn das einer thut, so thun es ihrer drei, darnach zehn, zwanzig, ja hundert, bis gar eine Gewohnheit draus werde. Also gewohnt man der Leute und des Wesens auf
\quellseite{861}
Erden, daß einer meine, es müsse nicht anders sein. Sagen denn: Ich sehe, daß diese und jene also thun, und ihre Kinder darauf ziehen; item, ich sehe, daß die auf diesem Berge, oder in jenem Thale oder Wiesen beten, und es gefällt den Leuten wohl, darum wollen wir es auch nachthun.

125. Jetzt ist es eine gemeine Rede, daß man auch also sagt: Wie kann ich ihm thun? will ich unter den Leuten sein, so muß ich thun, was des Landes, dieses Herrn und der Welt Brauch ist. Ich kann nicht ein Besonderes anrichten, ich muß es um meines Nutzes oder meiner Kinder willen thun; wie kann ich sonst unter den Leuten bleiben, und mich enthalten? Aber es heißt: Thue nicht wie die Welt und Leute wollen, sondern wie GOtt will. Darum sagt er allhier: Deine eigenen Kinder, der Welt Sitten, Brauch und Gewohnheit werden dich bewegen und betrügen; darwider wehre dich, feste und hänge dein Herz an mein Gebot, und sage: Ich will nichts ansehen, denn GOttes Gebot. Aber der Teufel, die Hofleute und Weltköpfe mit ihren Krempeln, Weib und Kinder, gute Freunde und Nachbarn reißen uns davon, und machen, daß wir vergessen des Worts GOttes.

126. Dieses sind Ursachen und Warnung genommen von GOtt durch Mosen. Aber man spricht allhier: Was! traue GOtt, du kommst noch wohl über zehn Jahre dazu, daß du fromm werdest, und GOtt fürchtest; ich will mittlerweile thun als mein Nachbar. Aber denke du, GOtt spricht allhier: Hüte dich, daß du mich nicht verachtest, oder das Feuer wird dich verzehren. Thust du das aber, daß du mich mit Ernst suchst, so will ich dein GOtt sein, und dich wiederum ehren, und dir alles Gutes thun. Und wenn du fällst, und an mich wieder von Herzen gedenkst, so will ich dich nicht verlassen. So hoch ich dich verworfen und geniedrigt habe, so hoch will ich dich wieder erheben, und in deine vorige Würdigkeit setzen.

127. Dieses alles redet Moses darum, daß man GOttes nicht vergessen, auch nicht verzweifeln solle. Denn wider die Vermessenen sagt er, GOtt sei ein Feuer; zu den Erschrockenen aber spricht er: Der HErr, dein GOtt, ist ein barmherziger GOtt. Darum soll man auf der Mittelstraße einhergehen, welches heißt, GOtt trauen, so wird man erhalten. Also ist Moses ein rechter Meister und Lehrer mit Auslegen des ersten Gebots; es hat es auch niemand zuvor ihm gleich gethan, wird auch noch keiner kommen, der es ihm nachthun wird.

\section{Das fünfte Capitel}

\begin{schriftzitat}
V. 1. Und Mose rief dem ganzen Israel, und sprach zu ihnen etc.
\end{schriftzitat}

1. Bisher haben wir gehört in dem vierten Capitel den Eingang oder die Vorrede, welche Mose vorher gestellt hat auf die Predigt, die er will von den zehn Geboten thun. Denn dies ist ein Buch, darinnen der rechte Verstand, Nutz und Brauch der zehn Gebote gepredigt wird. Nun, so erzählt er in diesem fünften Capitel den Text der zehn Gebote, und streicht sie heraus durch die folgenden Capitel, und nimmt vor sich eines nach dem andern. Darum so wollen wir den Text hören.

\subsection{Höre, Israel, die Gebote und Rechte etc}

2. Das ist der Text, darauf folgt die Glosse. Denn Moses sagt weiter, daß die Juden gebeten haben um einen Ausleger dieses Texts, denn sie haben selber nicht vermocht zu hören, da GOtt aus dem Feuer mit ihnen redete und die zehn Gebote Mosi gab; da baten sie um Mosen zum Mittler, daß er es ihnen sollte auslegen [2 Mos. 20, 19.]. Das rückt er ihnen allhier auf, daß sie haben um ihn zum Mittler, Dolmetscher oder Ausleger gebeten. Darum setzt Moses allhier dazu [V. 5.]: „Ich stund zur selbigen Zeit zwischen dem HErrn und euch.“ Erstlich sagt er also:

\begin{schriftzitat}
V. 6. Ich bin der HErr, dein GOtt etc.
\end{schriftzitat}

3. Da bindet sie unser HErr GOtt also an, da er sagt: „Ich bin der HErr, der dich aus Egyptenland geführet hat“, siehe, daß du dir nicht einen andern Gott erwählest außer mir. Denn dieses ist allein der rechte GOtt, der sich
\quellseite{862}
hier also nennt. Das habe ich euch oft gesagt und vermahnt: welcher will gewiß wandeln, der befleißige sich, daß er sich halte an GOttes Wort, hange an demselbigen, und mache ihm nicht ein ander Bild, Gedanken noch Träume von einem andern Gott unter GOttes Namen. Mose wußte zuvor, wie es würde kommen mit diesem Volke; wenn sie ins Land Canaan kämen, da würden sie gleich so einen großen Gottesdienst und Religion geistliches Standes finden, also herrlich zugerichtet, der wohl so schön gleißete, als den die Kinder von Israel hätten; dieselbigen Heiden würden GOtt sowohl dienen wollen, als die Kinder von Israel; wie der Gottlosen Gottesdienst gemeiniglich scheinbar und gleißend ist.

4. In Ammon, Moab, Capptern und zu Damasko, haben sie so großen Gottesdienst gehabt als die Israeliter; wie noch auf diesen heutigen Tag das Papstthum schöner gleist, denn unser, der rechten Christen, Gottesdienst. Da nun die Canaaniter sich rühmen würden, daß sie GOtt, der ein Schöpfer des Himmels und Erden ist, sowohl dieneten, als sie; daneben aus den Juden auch würden viel falscher Propheten aufstehen, die da einen falschen Gottesdienst lehreten, als hernach geschehen ist mit des Astharoth, Baals, und anderer Feier und Dienst, daß so viel Götter erwuchsen in Israel, als viel Städte darinnen waren; wie die Propheten solches heftig strafen, und sie auf den Wiesen, Bergen, in Wäldern etc. opferten (wie Jer. 3, 13. angezeigt wird), damit kann es endlich dahin, daß ein jeglicher Narr wollte einen eigenen Gottesdienst haben; darum so kommt Moses allhier der Abgötterei zuvor, und faßt die Juden so enge, und spricht zu einem jeden insonderheit (denn so geht seine Rede zu dem großen Haufen Volks, als wäre es nur Eine Person): Siehe zu, daß du nicht einen eigenen Gottesdienst anrichtest, sondern bleibe bei demjenigen, so ich dir vorlege und vorschreibe [5 Mos. 4, 23.].

5. Unter dem Papstthum, und jetzt, wenn die Schwärmer in die Schrift geraten, meinen sie, eine Lehre sei wie die andere. Das ist nicht wahr. So jetzt einer hat das Neue Testament gelesen, und irgend eine Predigt gethan, flugs rühmt er vom Geist. Daher sind vorzeiten gekommen Pelagianer und Arianer, jetzt unsere Schwärmer, die Sacramentirer und Wiedertäufer, und ihres Gleichen mehr; die führen alle GOttes Namen, und hat ihr Irrthum guten Schein, wie vorzeiten der Juden Abgötterei.

6. Mit der Betrügerei im Papstthum ist es auch so zugegangen. Erstlich kamen die Benedicter, darnach der Predigerorden, Barfüßer, Augustiner, und aus der Barfüßer Secte sind darnach wohl sieben andere Secten worden. Dies hatte alles den Namen, daß es Gottesdienst wäre; zu dem kam des Papstes Götzendienst mit Anrufung der Heiligen, Wallfahrten ins Grimmenthal, zur Eiche, zum Birnbaum,\footnote{Die Wallfahrtsorte der Maria zur Eiche und zu unserer lieben Frauen zum Birnbaum waren in der Nähe von Leipzig (Seidemann, Reformationszeit, S. 118); das Grimmenthal liegt in Sachsen-Meiningen-Hildburghausen im Meinunger Unterlande.} und anderswo zum Galgen zu. Dieses hat alles müssen ein Gottesdienst heißen. Da hat man wollen GOtt dienen. Dieser Name, Gottesdienst, der thut es; wenn der an eine falsche Lehre und Irrthum geschmieret wird, daß man hört GOttes Namen und Wort, solch Ding bewegt den gemeinen Pöbel, der plabt denn zu. Mit diesen Reiher vögeln oder Lockvögeln fähet der Teufel, der Vogelsteller, auf seinem Finkenherde die Leute, und berückt sie fein mit ganzen Beinen.\footnote{„mit ganzen Beinen“ mag wohl von den Gebeinen der Heiligen zu verstehen sein. Sonst muß ein Vogelsteller erst Pfeifen schnitzen aus den Beinen, um Vögel zu fangen. Vgl. das Sprüchwort: „Wer Pfeifen haben will“ etc. am Ende des achten Bandes unserer Ausgabe.}

7. Das weiß Moses wohl, daß das Volk würde die Augen aufsperren, und der Heiden Gottesdienst nachfolgen, darum warnt er sie also ernstlich. Dazu geht es denselbigen Abgöttischen allezeit glückseliger und besser denn denen, die rechtes Glaubens und gottseeligen Lebens sind. Die Christen sitzen in Gefahr Leibes, Guts und Ehre, und sind wie die Schlachtschafe [Ps. 44, 23.], wie denn der 73. Psalm nach der Länge darüber klagt, daß er sei geplagt täglich, und seine Strafe sei alle Morgen da etc. Aber die Abgöttischen sind in Sicherheit, treten Könige und Fürsten mit Füßen, grünen und brüsten sich daher wie ein Lorbeerbaum, wie der 37. Psalm, V. 35., sagt. Das thut der Teufel darum, daß er den letzern solch Glück beschert, auf daß der rechte Gottesdienst solle untergehen. Das macht denn den gemeinen Mann irre, der sperrt Maul und Augen auf, und spricht: Da ist Gefahr und Schmach, dort Ehre und Freiheit; und sieht
\quellseite{863}
nicht den Irrthum darunter. Da geht denn der gemeine Mann dahin, und spricht: Ich will es mit denen halten, und jene lassen alles Unglück haben.

8. Also ist es den Juden auch geschehen. Man sagt noch heutiges Tages: Ei, die Mönche haben also gesungen, viel gebetet, gefastet, und dieses alles GOtt zu Lob und Ehren gethan. Das gefällt dem gemeinen Mann wohl, er kann nicht hinüber, er kann sich nicht enthalten, sondern fällt dahin.

9. Aber was sagt unser HErr GOtt dazu? Ich will des Gottesdienstes ganz und gar nicht, denn es ist eitel Abgötterei: sie haben die Art an sich, daß sie die Leute führen auf ihr eigen Vertrauen. Ein jeder selbsterwählter Gottesdienst hat die Art und Eigenschaft an ihm, ob er sich wohl sonst schmückt unter dem göttlichen Namen, daß der Mensch will selig werden durch sein eigen Werk. Erwählter Gottesdienst führt den Menschen auf sich selbst; wenn man diesen Namen könnte davon reißen, so wäre es nicht Abgötterei.

10. Die Minoriten gedenken: Durch diesen Orden will ich selig werden, und andern auch zur Seligkeit neben mir helfen. Wenn dieser Zusatz nicht wäre, so wäre auch kein Mönch. Also hat ein jeder erdichteter Gottesdienst den Zusatz, Unflat und Koth an ihn hangen, daß ein Mensch darauf vertrauet; darum sind sie alle wider das erste Gebot, welches lehrt GOtt allein fürchten, lieben und auf ihn trauen. Alles, was nun dawider strebt, daß selbige ist Abgötterei. Also sind alle Pharisäer; weil sie ihren Stand rühmen, so fürchten sie GOtt nicht; denn sie meinen, sie dürfen GOttes nicht. Also ist bei den Barfüßermönchen das Fürchten und Glauben an GOtt auch hinweg. Darum so nennt sich der HErr allhier selber, und spricht: Den GOtt sollst du ehren, lieben, fürchten und anbeten, der dich aus Egypten geführet hat. Malt sich also ab, und gibt ein Gebot, daß man einen andern GOtt, der nicht sie hat aus Egypten geführt, unangebetet soll lassen; wo es ein anderer ist, den sollst du nicht ehren, nicht fürchten.

11. Also predige man den Christen auch: Du hast einen Heiland, Christum, der ist für dich gestorben, und vom Tode wieder auferstanden, und hat dich von deinem Jammer erlöst. Wenn du nun einen andern findest, welcher nicht für dich gestorben ist, noch dieses bei dir gethan hat, an den glaube nicht. Wahrlich, St. Franciscus, Benedictus, Bernhardus noch Augustinus sind nicht für dich gestorben; warum willst du den lassen fahren, welcher für dich gestorben, und sein Blut für dich vergossen hat, und denen trauen, welche nicht für dich gestorben sind? Werden wir nun nicht diesen allein suchen, sondern auf etwas Anderes das Vertrauen setzen, so werden wir verloren sein. Denn, wie St. Paulus zum Corinthiern sagt, so ist er uns allein vom Vater gegeben zur Heiligung, Gerechtigkeit und Versöhnung [1 Cor. 1, 30.].

12. Darum sagt Moses: Merket die Götter dabei, und fraget sie, ob sie euch aus Egyptenland geführt haben? Dies soll euer Wahrzeichen sein; aber es ist keiner neben mir. Das ist aber sehr schwer zu halten; denn es ist der Schein da, welcher heißt Gottesdienst. St. Francisci Regel gleist so schön, daß dieselbigen Mönche Christum darüber verleugnen, und St. Francisco nachfolgen. Der Baalspfaffen scheinbares Leben, daß sie sich mit Pfriemen stachen und mit Messerlein ritzen [1 Kön. 18, 28.], hatte viel einen größern Schein der Heiligkeit denn des Propheten Elia Lehre.

13. Darnach schreckt vom rechten Gottesdienst auch ab des Teufels und der Welt Toben und Wüthen wider die Christen, die Verfolgung und das Kreuz, so unsere Lehre mitbringt. Wer Christo will anhangen, dem ist die ganze Welt feind; er ist anders nicht, denn wie die Eule unter den Vögeln. Und also waren dazumal alle umliegenden Völker den Juden auch feind und arg; sie wohnten nicht anders unter den Heiden, denn als ein Schäflein, das mitten unter den Wölfen wäre [Matth. 10, 16.]. Die Heiden waren ihnen feind, daß sie dem GOtt dienten, der sie aus Egypten geführt hatte.

14. Darum schaue darauf, sagt Moses allhier, daß du bei dem reinen und gewissen Worte GOttes bleibest. Wo du das reine Wort GOttes nicht hast, da nimm keinen Gott noch Gottesdienst an, Jes. 8, 20., sondern sage: Ich habe ein gewisses, festes Wort [2 Petr. 1, 19.], nämlich, daß Christus mein Heiland und Mittler [1 Tim. 2, 5.] und das Lamm GOttes sei, so da trägt die Sünde der Welt [Joh. 1, 29.]. Daß aber St. Francisci Regel wahr, und daß er im Himmel sei, für mich GOtt bitte, davon habe ich kein gewiß Wort GOttes, darum will ich auch nicht an ihm hangen. Da sehe ein
\quellseite{864}
jeder zu, daß er in dem Gottesdienst und Stande gefunden werde, welcher in GOttes Wort gefasset ist, und bei demselbigen Gottesdienste bleibe. Denn, wahrlich, sonst kein Gottesdienst beschlossen ist mit GOttes Wort, denn an Christum glauben. Dieser Gottesdienst wird uns auferlegt im neuen Testament, wiewohl er auch im alten Testament ernstlich geboten ist.

15. Denn das erste Gebot ist das Hauptstück unseres ganzen Christenthums; es ist der Brunn des Glaubens, alles Verstandes, Weisheit, Erkenntniß und Gesetze; und alles, was gut ist, das steht im ersten Gebot. Das erste Gebot will alles an sich gebunden haben, und treibt uns das Vertrauen auf Creaturen. Ursache ist diese: Wenn du sollst GOtt von Herzen fürchten und vertrauen, so kannst du nicht den Mammon, Fürsten, oder deine eigene Gerechtigkeit, Frömmigkeit und dergleichen fürchten. Denn diese Worte nehmen alles hinweg: „Ich bin der HErr, dein GOtt“; Ich, und nicht ein anderer. Was sollte ich mich denn vor dem Teufel fürchten, vor den Fürsten, Pabst, Kaiser? sie heißen noch alle nicht „Ich“. Mit dem Wörtlein zeucht er alles an sich, da er ausdrücklich setzt: Ich, der HErr. Der Glaube zeucht alles zu sich.

16. Aus diesem Gebot fließen her, als aus einem Brunn und Quelle, alle Lehre der Propheten und Psalmen; item, alle Flüche, Dräuungen, auch alle Verheißungen. Jeremias am 17., V. 5., spricht: „Verflucht sei, der sich auf Menschen verlassen läßt, und der Fleisch hält für seinen Arm.“ Item, daß der 118. Psalm, V. 8. 9., sagt: „Es ist nicht gut, sich verlassen auf Fürsten“; wie denn derselbige ganze Psalm alles ganz und gar an sich zeucht; will sagen: Was kann dir schaden oder helfen, ohne dieser GOtt? Alle Propheten, und die ganze heilige Schrift des Alten und Neuen Testaments, kommen aus dem ersten Gebot her; denn er heftet es gar alles an sich, und will sagen: Wenn ich dein GOtt bin, warum trauest du nicht auf meine Gütigkeit, und fürchtest oder vertrauest andern mehr denn mir?

17. Und Christus ist hier auch mit hinein gefaßt, wie hernach im 18. Capitel, V. 15., Mose mit deutlichen Worten dahin sie weisen und bescheiden wird, daß sie sollen auf den zukünftigen Heiland Christum warten, als er sagt: „Einen andern Propheten wird euch GOtt aus euren Brüdern erwecken, den sollt ihr hören“; [er] schleußt Christum hinein, daß sie an ihn allein glauben sollen, und auf ihn hoffen und sich gänzlich verlassen. So ein groß Ding ist es um das erste Gebot: „Ich bin der HErr, dein GOtt“; darinnen alles begriffen ist, und alle Gesetze, Lehre und Regimente müssen sich darnach richten und lenken lassen.

18. Also habt ihr erstlich gehört, warum GOtt sich einen HErrn nennt, der sie aus Egypten geführt hat, und sich an diese That bindet; [er] will also heißen; spricht nicht schlechthin: Ich bin dein GOtt, sondern auch: „Der HErr, der dich aus Egypten geführet hat“; auf daß er alle Abgötterei verwerfe, welche nicht heißen: Ich; item, welche dich nicht aus Egypten geführt, oder welche nicht das gewisse GOttes Wort für sich haben. Wenn einer das nicht für sich hat, so ist es Abgötterei, es gleise und scheine, wie herrlich und helle es wolle. Hat es nicht GOttes gewisses Wort für sich, so werfe man es weg als lauter Abgötterei.

19. St. Benedicti und Francisci Regel hat einen guten Schein und schöne Gestalt der Geistlichkeit und Heiligkeit, es ist auch Glück und Gedeihen dabei; der Teufel ist da Schubherr, und thut ihnen keinen Schaden. Aber es hat nicht GOttes Wort, es lehrt nicht GOtt fürchten und an ihn glauben; sondern sie [die Regel] stiftet und macht nur Abgötterei, und lehrt auf Menschen trauen. Dergleichen ist es in allen andern Gottesdiensten. Aber man soll fahren lassen allen successum, auch allen gleisenden Schein, mit allen prächtigen Ceremonien, so nicht GOttes Wort bei sich haben, denn es hilft nicht.

20. Wenn wir aber den rechten GOtt annehmen, da stößt sich's denn. Es scheint unter uns, als hätten wir keine guten Werke, und als wollte es gar nicht fort; aber lasse du GOtt walten, er wird es wohl durchbringen, ob es gleich das Ansehen nicht hat. Was ist das für ein Schein, daß ein Hausvater Kinder auferzeucht, und sie mit GOttes Wort lehrt und unterrichtet? Was ist es auch für ein Schein, daß ein Prediger auf die Kanzel geht, und hat nicht einen Chorrock an? Darum wird er darüber geschändet und gelästert. Man hat Predigens kein Glück oder Gewinst; ja, vielmehr wird man noch dafür von der Welt verfolgt, und muß allerlei darüber leiden.

21. Abgötterei aber hat viel herrlichere Glosen und Schein, wird auch für größere Heiligkeit gehalten, über die rechten Gottesdienste.

\quellseite{865}
Ja, unser GOtt muß ein Narr, schwach und ohnmächtig sein; wie St. Paulus [1 Cor. 1, 23.] sagt. Bei seinem Gottesdienst wird nichts denn Thorheit und Schwachheit gesehen, denn es hat nicht eine gleichende Gestalt, noch groß Glück, wie jenes. GOtt stellt sich, gleich als wäre er krank. Aber die andern Götter heißt man Götter der Weisheit, der Stärke, und was sie thun, das ist eitel Weisheit, Glück und wohlgethan; sie allein auch schön, werden nicht für närrisch gehalten. Also geht man dahin, und mit dem betrügt man die Welt.

22. Darum ist diese Warnung bei den Juden sehr noth gewesen, daß GOtt allhier vermahnt und spricht: Siehe dich vor, ich werde für einen Narren gehalten werden, als ein ungeschickter Gott scheinen, schwach und unweise mich ansehen lassen; die andern Götter aber werden dagegen stark, weise geachtet und gerühmt werden, es wird ihnen wohl darüber gehen. Darum hüte dich, und siehe darauf, ärgere dich nicht daran, weiche darum nicht von mir, denn ihr Glück und Weisheit dürfte euch wohl von mir abziehen; sondern sehet nur das an, ob es der GOtt sei, welcher dich aus Egyptenland geführt hat. Wie wohl sie nun aufs treulichste hier gewiesen werden, sind doch hernach ihrer wenig gewesen, die bei GOtt geblieben und bestanden wären, der sie aus Egyptenland geführt hatte, und seinem Worte geglaubt; denn ihrer viel sind abgewichen. Aber fromme Herzen, wenn sie gesehen haben, daß man hin und wieder gelaufen ist, und GOtt gesucht, so haben sie gesagt: Ich will bei dem GOtt bleiben, der uns aus Egypten geführt hat.

23. Es möchte aber jemand sagen: was dieses uns angehe, sintemal wir nicht aus Egypten geführt sind? und warum sollen wir jetzt glauben an den GOtt, der Israel aus Egypten geführt hat? Darauf antworte du: daß dies nicht rechte Christen sind, die sich dieses GOttes äußern und nicht rühmen wollen; denn es leidet sich nicht, daß wir es wollen allegorice deuten; „aus Egypten“, das ist, aus der Sünde. Es geht uns dieser GOtt auch an. Denn dies Stück, „der dich aus Egypten geführt“, gehört wohl für die Juden; allein er hat es nicht länger gestellt und wollen gelten lassen, denn bis auf den Propheten Christum, davon hernach in diesem Buch am 18. Capitel, V. 15., geschrieben ist, den sollten sie hören.

24. Da hat er die zehn Gebote gestellt bis auf Christum, welchen sie annehmen sollten, und an ihn glauben. Denn Christus ist in das erste Gebot gefasset, er ist der GOtt, der sie aus Egypten geführt hatte, der ihnen das Himmelbrod gab; er war auch der Fels, davon sie in der Wüste tranken, die Wolken- und Feuersäule, die ihnen Tag und Nacht vorgingen. Aber da blieben sie bei dem Namen „aus Egypten“, und wollten diesen Propheten nicht hören. Nun sind wir hinzugekommen, daß wir nicht Mosen mehr hören, wie sie, die Juden, der sie aus Egypten geführt hat, sondern unsern HErrn und Seligmacher, Christum. Auch gibt Mose GOtt den Namen nicht weiter, denn bis auf die Zukunft des Propheten Christi. Der ist nun gekommen, und die Schrift ist erfüllt, und hat GOtt einen andern Namen bekommen, von seinem lieben Sohne Christo, der uns aus der Sünde, Tod, Teufel und Hölle Gewalt geführt und erlöst hat, und uns gebracht zur Gerechtigkeit und ewigem Leben.

25. Also ist Christus auch im ersten Gebot zugleich gefasset, wie GOtt sein Vater. Weil er nun sagt, daß er würde einen Propheten senden, welchen man hören sollte, so folgt, daß man dem Worte, welches er predigen wird in GOttes Namen, glauben müsse, und sich fürchten vor ihm, und ihm zufallen. Darum ist es nicht noth, daß wir glauben und fürchten GOtt in dem Namen, daß er sie aus Egypten geführt hat; also hat er nur eine Zeitlang geheißen; sondern jetzt hat er einen andern Namen, nämlich, daß Christus, wahrhaftiger GOtt, für uns gestorben ist, davon die Propheten auch viel gesagt haben. Als Jeremias: Ich schwöre bei dem lebendigen GOtt, es wird die Zeit kommen, daß man nicht mehr wird sagen: GOtt, der uns aus Egypten geführt hat; sondern man wird den Namen aufheben und sagen: GOtt, der uns aus allen Landen gesammelt hat [Jer. 32, 37.]. Item [Jer. 31, 33.]: „Ich will ihnen mein Gesetz in ihr Herz schreiben“ etc.

26. Dieses sage ich alles darum, daß ihr nicht irre werdet über dem Namen, „der dich aus Egypten geführt hat“. Denn zur selbigen Zeit war es noth, also zu sagen und GOtt [so] zu nennen, ihn [so] auszudrücken, und an dies Wunderwerk anzubinden, um der andern Götter willen im Lande Canaan, die sie selbst erfinden möchten. Darum will er sie abhalten von Abgöttereien
\quellseite{866}
, und sagen: Bete nicht Baal, Ashtaroth noch Moloch an; frage nicht nach ihnen, ob sie schon auch Götter Himmels und Erden genannt werden; sondern deß sollst du dich annehmen, der dich aus Egyptenland geführt hat. Wo findet man aber diesen GOtt, der sie aus Egypten geführt hat? Im Tabernakel zu Jerusalem. Nun, GOtt will den Dienst haben, welchen er selbst hat angerichtet; er will sich nicht von uns nennen lassen, sondern will unter seinem Wort mit Glauben gefunden werden, wie Moses in diesem Buch schreibt. Dabei soll es auch bleiben.

27. So schleichen nun alle Ketzer und Verführer unter dem Schein des göttlichen Namens einher. Denn also sagen die Papisten auch: sie erkennen Christum, und beten ihn an, und ehren GOtt in ihm. Aber wie kann das wahr sein, dieweil der wahrhaftige Christus nicht von ihnen erkannt wird, und sie Christum ihres Gutdünkens träumen, und vorbilden anders, denn er sich selbst hat abgemalt? Darum mißbrauchen sie nur seines Namens. Denn unter dem Namen soll ich bleiben, den er mir selbst hat offenbart; so findet man ihn auch. Das ist nun dieser Jesu Christi, den GOtt gesandt hat in die Welt [Joh. 3, 17.], und er ist auch gekommen, hat gepredigt, Wunder gethan, ist gestorben für unsere Sünde.

28. Dieser rechte Gottesdienst wird durch sein Wort verbracht.\footnote{verbracht = vollbracht.} Wenn man doch GOttes Namen an unsere Werke nicht schmierte, so könnte man niemand betrügen, und wir blieben auch unbetrogen. Man muß sich nicht auf Werke verlassen. Wenn ich St. Augustino zu Gefallen lebte, und trüge aus freiem Muthe, von mir selber, eine schwarze Kappe, wie ich sonst um guter Gesellen willen etwas thun möchte, das müßte man hingehen lassen. Aber wenn du also sagst: Es ist Gottesdienst, wer also lebt, der ist auf der rechten Bahn; das ist Abgötterei. Denn man nimmt den Orden an, daß GOtt damit gedient werde; dadurch wird Christi Leiden und Sterben vergessen, hintan gesetzt, ja, verleugnet.

29. „Aus dem Diensthause“ etc. Er wiederholt oft das alte Wunderzeichen und das Gefängniß, welches sie in Egypten getragen hatten, und zeucht es ihnen zum Gedächtniß hervor, daß sie es nicht vergessen sollen; will, daß sie sollen zurück auf die vorigen Werke sehen, und auf die erlittene Dienstbarkeit. Warum steht dies daselbst geschrieben, daß sie Pharao geplagt hat? Darum, wie vom Pilato im Credo steht, also sollten sie an die Erlösung aus Egypten gedenken. So sollten wir auch rühmen des Herrn Christi Wohltat. Denn, ehe wir Christum erkannt haben, sind wir in Sünden, Tod und Irrthum gesteckt. Diese unsere Irrthümer und Dienstbarkeit sollten wir zu Gemüthe führen, gleichwie die Juden nicht vergessen sollten Egyptenland. Darum hat das Diensthaus und Egyptenland müssen obenan stehen im ersten Gebot, um der Ursache willen, wie der Teufel, Judas, Pilatus, Herodes und Caiphas im Evangelio stehen.

30. Ich kann die Erlösung nicht loben, noch GOtt dafür preisen und danken, es sei denn, daß ich des Teufels, meiner Feinde, Irrthums, Angst und Noth mich erinnere. Also hat er uns jetzt erlöst von des Papsts Gewalt, welche nicht ein Königreich, sondern ein Diensthaus und Tyrannei der armen Gewissen gewesen ist; auf daß wir ja fleißig gedenken der Wohltat, und des Uebels, darinnen wir gewesen sind. Dergestalt war die Erlösung aus Egypten den Juden ein Denkzeichen der Plage, Marter, Aengstigung und Teufels Regiment, als, da die jungen Kinder in das Wasser geworfen, und die Alten mit Frohndiensten beschwert worden. Dieses war wohl den Egyptern ein Scherz und Wohlust, aber den Juden eine Angst, Betrübniß und großes Anliegen. Darum steht ihnen dies Gebot zur Gedächtniß da, daß sie nimmermehr aus der Acht lassen sollten, unter was Feinden sie damals gesteckt, und wie schwach sie im Glauben gewesen wären, auf daß sie nun GOtt um ferneren Trost und Stärke herzlich anrufen sollten.

\begin{schriftzitat}
V. 7. Du sollt keine andern Götter haben neben mir etc.
\end{schriftzitat}

31. In diesem fünften Capitel wird erzählt der Text der zehn Gebote von Worte zu Wort, welche GOtt Mosi gegeben hat. Nun greift folgendes Mose dazu, und will den Text in den nachgesetzten Capiteln auslegen und erklären, wie man ihn verstehen soll. Und wird unter andern anzeigen mancherlei Ursachen, welche machen, daß sie den rechten GOtt verlassen, und falsche Götter suchen, oder zu Abgöttischen werden würden. Darum ist Moses sehr fleißig, daß er uns bei dem einigen GOtt behalte, welcher
\quellseite{867}
sagt: „Ich bin der HErr, dein GOtt“ etc. Nun, das wollen wir auch vor uns nehmen, und davon mit GOttes Hülfe handeln; denn allhier liegt alles, was Weisheit heißt und ist auf Erden. Darum wollen wir fleißig davon reden.

32. Ich habe aber oft und viel gesagt, wiederhole es auch immer wieder, aber nicht ohne wichtige Ursache, daß in geistlichen und Glaubenssachen niemand sich unterwinde, mit GOtt zu handeln durch seine Gedanken. Er soll bleiben bei dem Wort, sonst hebt er es nicht wohl an mit seiner Vernunft.

33. Die Vernunft muß wider ihren Willen bekennen, daß [es] ihr zu hoch ist. Weil’s nun ihr zu hoch ist [so muß sie es bei dem Worte bleiben lassen];\footnote{Von uns hinzugefügt, um Sinn zu geben.} denn trachtet sie darnach aus ihren Kräften, so wird sie zur Närrin darüber. Sie wird es wohl unerlangt und bei dem Wort bleiben lassen müssen; darum muß sie sagen, daß in ihren Augen und Sinnen Thorheit sei, was sie auch nachdenkt. Es ist ihr gleich, als wenn ich wollte mit dem Finger an Himmel reichen, oder die Sonne mit der Hand verdunkeln und verfinstern; das wäre doch unmöglich und ein närrisch Vornehmen. Darum lässet uns nicht GOtt gleich sein, oder auch ihm überklügeln, sondern weit, weit uns ihm unterwerfen.

34. Derhalben sind es gar grobe, große Narren, die da trachten, GOtt zu erkennen aus ihrer Vernunft. Am Worte soll man hangen und darein sich wickeln, wie ein Kind in seine Windeln eingebunden wird, sonst kommen daraus so viel Rotten und Secten, Erden, Aberglauben und Abgötterei in der Welt, daß ein jeder seine Gedanken will ausbreiten als etwas Sonderliches, und ohne GOttes Wort, aus seinen Kräften und Vernunft, GOtt speculiren und erfinden. Es hat Mühe und Arbeit genug, wenn man gleich rein bei dem Worte bleibt, daß man bestehe; was will denn werden, wenn einer irre flattert mit seinen Gedanken und Vernunft? Denn, hätten wir es können treffen mit unserer Vernunft, so wäre es nicht vonnöthen gewesen, daß er sich hätte lassen vom Himmel hören, und das Wort lassen klingen, und alles in das Wort gefasset und gebunden.

35. GOtt hat sonst kein ander Mittel gegeben, ohne sein göttliches Wort, darinnen man allein Christum hören solle. Und niemand weiß, daß der Teufel so ein künstlicher Meister und Geist ist, der sich pflegt zu verstellen in die göttliche Majestät, und also geschicklich der Vernunft sich vorzubilden, auf daß er sie betrüge. Und wenn wir ihm folgen, so geschieht der Vernunft, daß sie meint, es sei eitel Gold, was da gleiche, und alles eitel göttliche Weisheit; aber darnach ist nichts dahinter, und niemand besteht damit. So nun die Vernunft gefangen, und sage: Ich will nichts sehen noch hören, denn was GOttes Wort ist, dabei bleibe ich. Und nun lehrt und treibt die zehn Gebote auch darum, daß es Worte sind, daran wir kleben und haften sollen.

\begin{schriftzitat}
V. 9. Ich bin der HErr, dein GOtt.
\end{schriftzitat}

36. Was GOtt sei, das lernt man leichtlich aus dem Gegenspiel, wenn man bedenkt, was ein Abgott oder ein falscher Gott ist; darnach kann man erkennen den rechten und wahrhaftigen GOtt, als Paulus zum Römern Cap. 1, 18. anzeigt. Alle Welt heißt das einen GOtt, darauf der Mensch trauet in Noth und Anfechtung, darauf er sich tröstet und verläßt, da man alles Gutes will von haben, und der helfen könne.

37. Also haben die Heiden gethan, und etlich den Heyden zum Helfer und Gott gemacht, daß er gut Regiment gäbe; den Marten, in Kriegen zu helfen; darum ihn angebetet. Darnach haben sie aus der Vernunft viel Abgötter gemacht. Die Römer haben viel Götter aufgeworfen, um mancherlei Anliegen und Hülfe willen, die ihnen noth war; daß einer den Leuten hilfe im Kriege; einem andern diese Gewalt gegeben, jenem das zugeschrieben, als, der sollte Korn lassen wachsen, jener zu Wasser helfen im Schiffbruch. So manche Noth, Gut und Nutzung auf Erden war, so manchen Gott hatte man erwählt, bis sie auch Gewächse und Knoblauch zu Göttern gemacht. Davon Augustinus viel schreibt im Buch de civitate Dei.

38. So beschreibt die Vernunft GOtt, daß er sei, was einem Menschen Hülfe thue, ihm nütze und zu Gute gereiche. Daraus spürt man, daß Vernunft nur so viel weiß von GOtt, als Paulus zum Römern am 1. Capitel, V. 19. 20., sagt: „Daß man weiß, daß GOtt ist, das ist ihnen kund gethan, damit, daß GOttes unsichtbares Wesen ist zu merken an den Werken, die er thut an der Welt, die er geschaffen hat.“ „Daß GOtt sei“, das ist, daß GOtt nütze sei, und helfe in Nöthen.

\quellseite{868}
39. Also haben wir unter dem Papstthum auch Götter gemacht. Eine jegliche Krankheit oder Noth hatte einen eigenen Helfer und Gott. Die schwangeren Frauen, wenn sie in Nöthen waren, riefen sie St. Margarethen an, die war ihre Göttin, als die ihnen konnte zu Hülfe kommen, und eine leibliche Frucht geben, die Gefahr und Schmerzen der Geburt zum Besten wenden, und ihnen davon abhelfen.

40. Die Geizigen haben auch einen besonderen Gott, der heißt Mammon. Das bezeugt und überweiset sie ihr eigenes Herz, dieweil sie solche Gedanken haben: Da liegt mein Geld oder Baarschaft; habe ich nicht Essen und Trinken, so wird mir es der Mammon wohl geben; er wird mir helfen, und mich erretten vom Hunger. Darum so setzen sie ihren Trost darauf.

41. Aber es wird der wahrhaftige GOtt in der Schrift [Ps. 9, 10. Jac. 1, 17.] genannt ein Nothhelfer, und ein Geber alles Guten. Und liegt hier die Macht daran, wem diese göttliche Gewalt und Namen rechtlich gebühren, und von uns Menschen gegeben werden sollen.

42. Da haben Adams Kinder beschlossen, daß wir GOtt anrufen, der uns in Nöthen aushelfen solle, und Gutes mittheilen. Das sagt auch St. Paulus, wie droben angezeigt. Aber wenn sich das Ende herzunaht, und die rechten starken Püsse beginnen sich zu finden, da wackelt und wanckt die Vernunft in minore, und spricht: Jupiter, Mars, Venus wird es thun; und allda theilt es sich. Die Vernunft meint, der oder dieser soll es thun. Daher kommen so mancherlei Rotten und Spaltung; wie jetzt auch im Papstthum sind. Einer sagt: Ich weiß, daß GOtt helfen wird. Soll man aber denselben GOtt namhaftig machen, so wird dieses und dergleichen daraus. Der sagt: St. Margaretha wird den schwangern Weibern helfen; jener spricht: St. Erasmus wird der reichen Nothhelfer sein. Da hat man vierzehn Nothhelfer an Einem Orte gehabt. St. Christoffel hat denen helfen sollen, die da in den letzten Zügen liegen. Also gibt ein jeder dem den Namen GOttes, da er sich am meisten Gutes zu versieht, dann wird der Glaube, der Trost, Zuversicht und Vertrauen einem andern gegeben, und nicht dem wahren Gott.

43. Darum sage ich noch einmal, die Vernunft wisse etlichermaßen, daß GOtt könne und solle helfen, aber den rechten GOtt kann sie nicht treffen. Sie mag wohl von ihm reden; aber wer er sei, wo er sei, und wie er helfe, das weiß sie nicht. Die Juden haben den Text zuvor gehabt, welchen die Heiden nicht gewußt: dennoch haben sie im Felde und im Walde, auf Bergen und Höhen Altäre gemacht, geräuchert, Kirchen gebauet, GOtt daselbst angerufen und gesagt: Hier ist GOtt, dort ist GOtt. Das heißt recht, nach GOtt tappen, wie ein Blinder nach der Wand. Also sind wir auch im Papstthum hin und her gelaufen, zu St. Jakob, gen Rom, in das heilige Land, gen Jerusalem, und an andere unzählige Oerter. Da haben die Leute GOtt gesucht. Nichts Anderes thut die Vernunft, wenn sie GOttes Wort nicht hat; sie weiß wohl von ihm zu sagen, und zeugt, daß ein GOtt sei, aber sie kann es nicht recht treffen, denn sie weiß nicht, wer er ist.

44. Dieser Ungewißheit halben muß GOttes Wort uns zu Hülfe kommen, und muß GOtt sich heraus an Tag geben, und sich selber abzirkeln in einem äußerlichen Worte und Zeichen, daß man ihn hören, sehen, greifen, fassen und erkennen möge; sonst ist uns ungerathen. Das siehst du auch in diesem Texte: „Ich bin der HErr, dein GOtt, der ich dich aus Egypten geführet habe.“ Da erzählt er, was GOtt sei, was seine Natur und Eigenschaft sei, nämlich, daß er wohlthue, erlöse aus Gefährlichkeiten, und helfe aus Nöthen und allerlei Widerwärtigkeiten, dieweil er die Israeliten aus ihren Beschwerungen errettet, und aus Egypten geführt, that ihnen noch Gutes, und gab ihnen das Land, welches er ihren Vätern verheißen hatte. [Er] malt sich so hac definitione, oder mit diesen Worten, Werken und Wohlthaten ab, welches sie mit ihren Augen gesehen, und an ihrem Leibe und Person erfahren hatten, da er ihnen aus Egypten geholfen. Daher setzt Mose die Erkenntniß GOttes im ersten Gebot, daß GOtt sei, der allezeit und ewiglich wohlthut und hilft, auf daß, wer Hülfe bedarf, hieher laufe, da wird er Hülfe empfinden. Denn dieser GOtt ist gnädig, barmherzig, und thut wohl denen, die ihn fürchten, in das dritte und vierte Glied. Das haben wir aus der Erfahrung.

45. Es nimmt aber auch die Schrift den Namen GOttes, und theilt den auch mit den Gottseligen, Frommen und allen Gotteskindern, der Obrigkeit, den Fürsten und Richtern, und nennt sie Götter [2 Mos. 22, 28. Ps. 82, 1.]. Solches
\quellseite{869}
geschieht nun darum, daß sie in ihrem Stande andern dienen, und sind die, durch welche GOtt hilft und wohlthut. Also, David und die andern Fürsten sind Götter gewesen, denn sie haben ihren Ländern wohlgethan, ihren Unterthanen geholfen, wenn sie in Nöthen gewesen sind. Darum hat man sie auch angebetet, und ist ihnen göttliche Ehre erzeigt, von wegen des göttlichen Werks, daß sie haben den Leuten wohlgethan und geholfen; wie man noch heutiges Tages die Knie beugt vor der Obrigkeit. Denn ihr Amt ist dahin geordnet, daß man der Armen, Dürftigen Sache schlichten solle, recht richten und Frieden schaffen, daß sie dem rechten GOtt dienen und sich ernähren können.

46. Also sind Prediger, Eltern und Zuchtmeister Götter, gegen ihren Zuhörern, Kindern, Gesinde und Schülern zu rechnen. Denn sie treiben Wercke, welche GOtt eigentlich zugehören, unterweisen sie das Beste, lehren und wehren, helfen und rathen, nachdem es die Noth erfordert; sie geben und thun wohl oder Gutes. Über dies Wort ist nunmehr unkenntlich worden, man verstehet es nicht recht, und hierin wird geirrt, daß, dieweil GOtt wohlthut durch Obrigkeit, Herren und die Creaturen, so plagt das Volck zu, hängt an den Creaturen, und nicht an dem Schöpfer; sie gehen nicht durch sie zum Schöpfer. Daher ist es gekommen, daß die Heiden aus den Königen haben Götter gemacht, und die Juden haben Mosen auch zum Gott machen wollen. Daher ist alle Abgötterei gekommen. Denn man kann es und will es nicht merken, wie das Werk oder die Wohlthat von GOtt komme, und nicht schlecht von der Creatur komme, und nicht schlecht von der Creatur;

tur; ob die wohl ein Mittel ist, dadurch GOtt wirkt, uns hilft und gibt.

47. Aus diesem Ungedichte und Mißverstande ist es anfänglich gekommen, daß die Heiden die Sonne angebetet; denn sie gibt Licht und Nahrung. Etliche haben den Mond und Sterne geehrt. Gleichwie die Zauberinnen noch auf diesen Tag dem Teufel anhangen, mit ihm zu thun und Verbündniß haben. Die nehmen auch Hülfe, das ist, was er bringt und ihnen zuwendet, von ihm an, hofieren und dienen dafür dem leidigen Bösewicht, und nicht dem Schöpfer und Helfer durch die Creatur, welchem man doch billig anhangen sollte.

48. Also siehst du den Ursprung aller Abgöttereien und Ketzereien, daß, weil der Mensch so blind ist und fällt auf die Creaturen, klebt und hängt an denselben, daß sie ihm helfen werden, so doch GOtt in diesem Gebot sich hören läßt, daß er der einige wahrhaftige GOtt sei. Solches aber können wir nicht verstehen, wenn wir nicht bei seinem Worte bleiben. Darum sollen wir in der Landstraße einhergehen, und im göttlichen Worte uns finden lassen.

49. So habt ihr nun gehört: Erstlich, was GOtt sei. Zum andern, daß er nicht könne getroffen werden, man höre denn sein Wort, das ihn abmale. Jetzt wird Moses solches noch besser erklären, wer derselbige GOtt sei, und wer diejenigen sind, die wider ihn streiten, und was da verhindere, oder wie es sich widert,\footnote{„sich widert“ = sich dagegen stellt, Hinderniß verursacht.} daß die Leute nicht bei ihm bleiben.

\section{Das sechste Capitel}

\begin{schriftzitat}
V. 1–5. Dies sind aber die Gebote, daß du den HErrn, deinen GOtt, fürchtest etc. Höre, Israel, der HErr, unser GOtt, ist ein einiger HErr etc.
\end{schriftzitat}

1. Da sehen wir, daß Moses das erste Gebot anhebt meisterlich und fleißig auszutreiben. Also hat er gesagt: „Ich bin der HErr, dein GOtt“ etc. Was heißt das? Darauf antwortet er: daß du dem GOtt und HErrn vertrauest, und ihn fürchtest, daß du thust, was ihm wohlgefällt, daß du nicht anderswo hingehest, seine Gebote und Gesetze nicht übertretest; sondern sollst dich zu ihm versehen, daß er dir alles Gutes thun will, und es soll dir wohlgehen. Also haben wir es auch ausgelegt.

2. „Nicht andere Götter haben“ heißt, GOtt lieben, fürchten, und ihm vertrauen. „Fürchten“ ist, daß ich bei dem einigen GOtt bleibe, fürchte
\quellseite{870}
mich vor ihm, daß ich nicht einen andern Gott suche, oder einen andern ergreife, sondern mein Vertrauen auf diesen setze. Denn welcher einem andern vertrauet, der verachtet den vorigen, er fürchtet ihn nicht, und ist so keck, daß er sich darf nach einem andern Gott umsehen, und darauf seine Zuversicht setzen; so trauet er ihm auch nicht. Und da ist der Glaube mit dem ersten Gebot hinweg.

3. Deß nimm ein Exempel aus dem Papistthum. (Ich wollte nicht gerne, daß man der Greuel vergesse, um Undankbarkeit willen.) Wer sein Vertrauen, Hoffnung und Trost setzt auf St. Margareth, oder einen andern Nothhelfer und Heiligen, der schlägt GOtt in die Schanze, und verachtet ihn aufs äußerste; gedenkt: Wer weiß, was GOtt thut? er kann mir nicht helfen; aber die heilige Jungfrau St. Margaretha wird mir helfen, denn die hat es verdient. Und das heißt GOtt verachten, und auf die Creatur sich begeben, welche heißt St. Margaretha oder St. Barbara, da man ihr in Kindes- und Todesnöthen mehr trauet und glaubt denn GOtt, so doch unser Vertrauen allein auf GOtt sollte stehen, daß er helfen würde. Aber der größte Theil unter den Menschen halten ihn für einen Kloß, darum rufen sie ihn nicht an [Ps. 53, 2.].

4. Ein abgöttischer Krieger ruft St. Barbara an, daß sie ihn nicht lasse ohne das Sacrament sterben; oder sonst St. Marco, daß er ihnen helfe. Er weiß von GOtt nichts. Heißt das nicht GOtt verachten? Man hält ihn nicht dafür, daß er könnte den Kriegern zu-Hülfe kommen, da er doch sagt: „Ich bin der HErr, dein GOtt.“ Ja, der tapfere Held David bekennt im 144. Psalm, V. 1., daß GOtt seine Hände und Fäuste zum Kriege abrichtete. Und du hältst GOtt nicht für den, der dir helfen könnte; aber dagegen St. Margaretha, St. Barbara und St. Marcus, die können dir helfen?

5. Maria, die liebe heilige Jungfrau und Mutter GOttes, ist auch die schändlichste Abgöttin worden, die hat uns auch sollen gnädig sein und in höchsten Nöthen aushelfen. Alle sind wir so geschickt gewesen, daß wir sind von GOtt gefallen, und sie hat sollen unsere gnädige Königin sein; Christus ist nichts gewesen; daß alle Tempel und Altäre sind gemeiniglich in Maria Ehren gestiftet und gebauet worden. Heißt denn das nicht GOtt verachten? Hilft denn GOtt, was darf ich denn Marien Hülfe, oder anderer Heiligen? Sehe ich aber mein Herz auf die Jungfrau Maria, daß sie mir helfen soll und Gutes thun, was darf ich denn GOttes? Er sitzt nur müßig im Rauchloch.

6. Ja, sagt man dann: Man soll gleichwohl die lieben Heiligen ehren, denn sie haben es verdient. Wo steht das geschrieben? Allhier hört du, daß Moses sagt: „Höre, Israel, der HErr, unser GOtt, ist ein einiger HErr.“ Darfst du was? Rufe ihn an. Es ist ohne alle Noth, daß du zu einem andern laufest; denn er läßt dir sagen: „Ich bin dein GOtt“; sagt nicht: Ich will dein GOtt sein, ich werde es noch werden; sondern: Ich bin es allbereit; allein, gedenke und glaube mir, ich will dir wohl helfen.

7. Eben so geht es auch zu mit den Mönchen und Nonnen. Wenn sie kommen in die hohe Abgötterei, so deuten sie: Wir haben drei Gelübde gethan, als, Armuth, Keuschheit und Gehorsam, und haben ihren Orden, Regel, Statut. Diese ihre Werke, so sie darinnen thun, sind ihr Abgott. Denn sie weichen von GOtt ab, fürchten ihn nicht, dürfen nicht seiner Gnaden und Gaben, als, der Vergebung der Sünden, sondern kommen getrollt und wollen selig werden durch ihren Orden, Kappen und Plattenwerck, und dadurch Vergebung der Sünden erlangen. Und dadurch werden sie treulos, fallen von seiner Gnade und Barmherzigkeit ab, welche sie rechtfertigen sollte und aus Gnaden ihnen die Sünde vergeben. Aber sie dürfen es nicht; ihr Stand, Kappe und Abgötterei kann es wohl ausrichten. Das heißt GOtt verachten, nicht fürchten, und einen andern Gott aufrichten. Denn, weil ihr Orden ihnen hilft, so dürfen sie nicht GOttes. Also haben sie GOtt niemals vertrauet noch ihn gefürchtet, und hat Verzweiflung darauf folgen müssen.

8. Darum ist das der rechte Kern und Verstand des ersten Gebots: „GOtt fürchten“, daß man keinen andern suche, und diesem traue, der dir alles Gutes geben will. „Fürcht“ dient dazu, daß sie uns bei ihm behalte, daß wir nicht einen andern Gott erwählen. „Vertrauen“ bringt mit sich alle Hülfe, daß du ihm allein in allen Nöthen in seine Hände begehst, und fest gewiß, daß er das Beste bei dir thun werde, dir Rath und Hülfe schaffen, dieweil er es hat zugesagt und lügt auch nicht, derhalben sonst keinem andern anhängest.

9. So lauten ja seine Worte: „Ich bin dein
\quellseite{871}
„GOtt“, das ist, ein Helfer in der Noth, der gibt alles Gutes. Er kann nicht deutlicher reden, denn daß er sagt: Ich thue dir alles Gutes, und helfe dir aus allenhalben. Der nun den rechten GOtt trifft, der hat das, nämlich, alles Gutes thun, und aus aller Trübsal und Noth helfen, und dich erlösen. Das wird leichtlich gesagt, aber schwerlich geglaubt. Doch haben wir gesagt, daß es GOtt durch die Creatur gibt; darum die Könige, Fürsten, Prediger, Eltern und Christen auch Götter genannt werden.

10. Möchte einer aber sagen: Thue ich denn nicht Unrecht, wenn ich die Könige anbete, und thue einem Fürsten eine Ehre, und den Priestern eine Reverenz mit dem Kniebeugen oder Hutabziehen? Warum sagst du denn, ich thue übel darinn, daß ich die Heiligen anrufe, Mariam anbete? Darauf antworte ich: Wenn du einen Fürsten also ehrst, daß du siehst GOtt durch ihn dir alles Gutes geben, da ist es recht, so thust du wohl. Denn du empfehlest nicht den Frieden und Schutz hier in diesem Lande von Herzog Johann, Churfürsten, ich verlasse mich auch nicht auf ihn; sondern GOtt gibt dir durch diesen Mann, daß du Frieden habest; daß du also nicht bleibest haften an dem, durch welchen es dir geschiehet, sondern kommest zu dem, der dir es durch den Fürsten gibt. Denn Friede ist ein Werk, das GOtt allein zusteht zu geben, und ist nicht eines Fürsten oder anderer Obrigkeit Werk.

11. Also empfähest du vom Pfarrherrn auch das Sacrament, und hörst die Predigt des göttlichen Worts; aber nicht als von ihm, er kann dir es nicht geben, sondern GOtt gibt es durch ihn, er ist das Mittel und Instrument, dadurch dir es GOtt gibt. Denke also, daß der Fürst und der Pfarrherr nicht dein Gott sind, und daß sie das Sacrament, Wort und weltlichen Schutz nicht zwar als das Ihre reichen, sondern GOtt gibt es durch sie. Darum ehre ich auch die Obrigkeit um GOttes willen, der mir solches durch diese Mittel gibt, Röm. 13, 1. Sonst mögen sie ein gutes Jahr haben, [welche] die Fürsten fürchten und ihr Vertrauen auf sie setzen; denn sie müssen zu Scheitern gehen und verflucht sein. Es ist ein wahres Wort, das man pflegt zu sagen: Fürsten Gnade ist wie Aprilenwetter; das währt nicht lange. Daher spricht der 146. Psalm, V. 3.: „Verlasset euch nicht auf Fürsten“; und Jer. 17, 5.: „Verflucht sei, wer sich auf Fürsten verläßt.“ Psalm 146, 3.: „Fürsten sind Menschen, sie können ja nicht helfen.“ Solches bezeugt die heilige Schrift allenthalben, daß, wer sich auf Menschen verläßt, der geht zu Boden.

12. Nun ist es ein ander Ding, Wohlthat empfangen durch einen Menschen, und dem Menschen vertrauen, oder auf ihn sich verlassen. Vom Fürsten, Predigern und von Eltern soll ich es als von Creaturen empfangen, wiewohl mir es GOtt der HErr durch sie gibt; aber sie fürchten und vertrauen auf sie, als wäre kein anderer GOtt, das ist nicht zu thun. Ich soll sie nicht fürchten noch ihnen vertrauen; denn Furcht und Vertrauen gehört alles hinauf zu GOtt. Da soll ich sagen: Was recht ist, dabei will ich bleiben, GOtt gebe, es zürne Fürst, Herr, Vater oder Mutter.

13. Man findet ihrer aber allzuviel, die wissentlich wider GOttes Gebot handeln, fürchten die Fürsten, und thun um ihretwillen, daß sie nimmermehr können verantworten, und sie sonsten wohl ließen ansehen. Diese fürchten nicht GOtt, sondern den Fürsten. Darum so geht es also: Wenn sie einen gnädigen Fürsten haben, sind sie stolz, und kann niemand mit ihnen überein kommen. Ja, unterm Herrn Gott selbst fügen sie nicht ohne Habersack. Ein Fürst und Obrigkeit müssen von wegen ihres Amtes und Befehls Gutes thun, und ihren Unterthanen helfen, darum soll man auf sie nicht vertrauen, noch sie fürchten, oder um ihretwillen wider GOtt thun; von ihnen soll man aber Gutes empfahen als von GOtt, und GOtt durch sie ehren. Darum baue nicht auf Menschen.

14. Daß nun etliche sagen: Was soll man mit den todten Heiligen thun? Soll man sie nicht ehren, wie man sonst die Eltern und Fürsten ehrt? Antwort: Du siehst nicht, daß die Heiligen, als, Maria, St. Margaretha und St. Barbara dir helfen, wie du erfährst, daß dir dein Fürst und Vater hilft; darum darfst du nicht die todten Heiligen anbeten. Der Fürst schützt dich, handhabt dich zu Recht, wie denn die Obrigkeit ist geordnet; GOtt will durch sie seine Gottheit dich sehen lassen, und gibt dir durch sie Landesfrieden, daß sie dich vertreten. Das hast du nicht von den todten Heiligen, die haben deß keinen Befehl, Ordnung, noch dazu einen gesetzten Stand; darum sollst du die Heiligen nicht anrufen, sie nicht fürchten, noch auf sie vertrauen. Dazu ist noch das Allergrößte, Gott hat dich es nicht geheißen. Denn weil du
\quellseite{872}
dich nicht sollst fürchten, noch ihnen vertrauen, welche doch GOtt selber dir geordnet hat zu Nutz und Gedeihen, und daß er durch sie dein GOtt sei, viel weniger will er haben, daß du die Heiligen vorziehen und auf sie hoffen solltest.

15. Durch deine Eltern gibt er dir Leib und Leben, ernährt und versorgt dich; dennoch sollst du ihnen nicht vertrauen. Gebrauchen magst du ihrer nach GOttes Ordnung, aber über GOtt sie nicht achten; denn sie sollen nur ein Mittel sein, dadurch er dich in diese Welt schafft, und eine Zeitlang erhält. Aber die verstorbenen Heiligen sitzen nicht in diesem Stande, thun dir auch der keines; sie sind aus der Welt genommen; wir können ihrer nichts genießen, als sonst der Lebendigen nach GOttes Einsetzung.

16. Darum, wenn ich sage: Der Heilige hat das Amt im Himmel, das erdichte ich aus meinem eigenen Kopfe, ohne GOttes Wort, und der Teufel betrügt mich durch seine Verführung und falschen Schein. Und ist doch dahin das größte Vertrauen gerichtet, daß wir uns besondere und eigene Götter erwählt haben, die wir viel mehr gefürchtet denn unsern HErrn GOtt selber. Also haben wir gefürchtet St. Valentinum, Antonium und Sebastian. Denn, weil es von uns erdichtet ist, so gefällt es uns. Da hat man gesagt: Ja, St. Margaretha hat mir geholfen. Wie, wenn es der Teufel gethan hätte? denn der kann auch einem helfen. Aber das kannst du sehen, und allhier hast du GOttes gewisses Gebot, daß er also diese Stände geordnet hat, daß dir dein Seelsorger, Fürst und Eltern vorstehen, und sie sind von GOtt dazu verordnet, daß er durch sie dir helfen will; darum gebeut er auch, daß man sie ehren solle.

17. So ist es nun ein lauter Aberglaube und Abgötterei, wenn ich Hülfe und Rath suche anderswo, denn bei GOtt. Ich soll zu denen keine Zuflucht haben, da GOtt mir seine Hülfe verordnet hat, ich soll sie auch nicht ehren. Denn darnach folgt Verachtung und Mißtrauen gegen GOtt, daß GOtt so schwerlich zu glauben ist, und dem Teufel man so leichtlich glaubt; denn was GOtt eine Zeitlang aufzeugt, das gibt der Teufel bald. Aber hiervon wollen wir auf eine andere Zeit reden.

18. Auf diesmal lasset uns behalten, daß GOtt sei, der alles Gutes thut, und hilft aus Nöthen; nicht allewege durch sich selbst, sondern es kommen allerlei seine Mittel auch dazu; als, durch seine Engel, Fürsten, Herren, Eltern, Prediger, Christen \&c., vor allen aber durch Christum. Das ist nun wohl recht, daß ich weiß die Mittel, durch welche mir von GOtt Gutes widerfährt; aber ich soll nicht auf sie bauen, sie nicht über GOtt heben, um der Könige, Fürsten und Herren willen Unrecht thun, und um ihrer Ungnade und Zorn willen das Recht lassen; sondern auf GOtt, den Schöpfer, soll ich allein mein Vertrauen setzen.

19. Aus der Ursache ist David angebetet, das ist, geehrt worden, denn er hatte die Hände voll Wohlthaten, durch seine Hände ging dieses großen Volks Israel Regierung, Schutz und Schirm. Darum ward David Gott genannt; nicht, daß er ein wesentlicher GOtt wäre, sondern daß er GOttes Werk trieb, und GOtt durch ihn halb aus Nöthen, that wohl und erlöste das Volk von den Feinden. Darum war nicht weiter auf ihn zu bauen, denn allein sofern er GOttes Diener und Instrument blieb, an seinem Worte und Befehl festhielt.

20. Wir haben im fünften Capitel gehört den Text der zehn Gebote, und im Anfang des sechsten Capitels wird folgen, wie es Mose beginnt zu erklären, nämlich also: „Höre, Israel, der HErr, unser GOtt, ist ein einiger HErr“ \&c. Da hören wir, daß die Auslegung des ersten Gebots sei diese: „Du sollst GOtt, deinen HErrn, lieben von ganzem Herzen, von ganzer Seele, von allem Vermögen.“ Was aber das sei, haben wir oft im Evangelio hören predigen, da denn stets uns vorgehalten worden, daß dies Gebot nicht so leicht und geringe sei, als es scheint, sondern es ist eine Summa und Beschluß aller Weisheit und Verstandes.

21. Denn „lieben von ganzem Herzen“, das steht in den höchsten Affecten, und ist nicht ein schlecht, kalt Werk, das in äußerlichem Wandel oder Uebertretungen nur beruhe, wie es die Juden und Papisten verstehen; als, vor den Bildern und Götzen nicht die Kniee beugen, oder die Hände gegen ihnen aufheben. Wenn sie das nicht gethan haben, so meinen sie, sie haben den rechten GOtt, und an Abgötterei sich mit nichten vergriffen. Item, wenn sie fasten, lange Kleider tragen, ihre Zeiten gebetet, sich äußerlich in allerlei Geberden nach ihrer Weise gehalten, so denken sie: Wir sind heilige Leute, und haben wohl gelebt; wie heutzutage ihrer noch viel gefunden werden, die sich mit solchen
\quellseite{873}
äußerlichen Dingen heilig dünken lassen. Aber hier sagt Mose: Willst du das erste Gebot halten, und wissen, was da sei, nicht andere Götter haben? höre: „Du sollst den HErrn, deinen GOtt, lieb haben von ganzem Herzen“, das ist, daß dir nichts Lieberes sei, denn GOtt, sein Wort und Wille; in Summa, daß wir über ihn nichts im Himmel oder Erden lieben.

22. Wir haben zwei Sonntage nach einander gehört, daß GOtt nicht könne ergriffen werden, denn allein durchs Wort; ohne das könne man ihn nicht sehen noch fühlen [Joh. 14, 23.]. Stellt man sich recht zum Worte, daß man es liebt, und meint es von Herzen, so wird GOtt auch geliebt. Nun, so kann man an uns nicht sehen, fühlen oder erfahren diese Liebe, damit wir GOtt fürchten und lieben, denn wenn man sieht, wie wir uns zum Worte GOttes stellen, oder wie wir uns gehorsamlich gegen der Predigt halten. Wo dir dies Wort, Geschäfte und Ordnung lieber ist denn alles auf Erden, so ist die Sache schlecht; denn ist’s ein Zeichen, daß man GOtt liebe, und so wirst du die Eltern ehren, deinen Nächsten auch lieben, nicht todtschlagen, ehebrechen, stehlen etc. Darum, wenn dir das Wort geliebet über deinen Leib, Leben, und was du sonst hast, so steht die Sache eben wohl; so wirst du deinem Nächsten an alle dem Seinen nirgend Schaden thun, sondern alle anderen Gebote und Werke halten. Derhalben, sind mir die zehn Gebote lieb, so lebe ich darnach, ich lüge und trüge nicht, lasse eher Leib, Leben und alles drüber. [2 Macc. 7, 30. 7, 2. ff.]

23. Aber wo findet man sie? Wenn man sie sollte zählen, würden ihrer sehr wenig sein. Der Teufel, die Welt und unser Fleisch lehren uns viel anders. Wir dürfen um eines losen Hellers willen wohl alle Gebote GOttes, sein Wort und Nächsten in die Schanze schlagen. Denn, heißt das GOttes Gebot und Wort lieben und halten, wenn du übel von deinem Nächsten redest, schändest ihm sein Weib, oder vervortheilst ihn auf dem Markte, und sonst, wo du nur kannst? Den Teufel möchest du wohl lieb haben; ja, einen schäbichten Heller liebst du mehr denn deinen GOtt. Du solltest dich aber eher in einen Zwinger beißen, denn daß du etwas wider GOtt und sein Wort thätest.

24. Also sollte es zwar zugehen, wenn du Liebe hättest zum Wort GOttes: Ehe du deinen Nächsten um einen Heller betrögest, oder ihn ein Wort zuwider redest, du ließest eher Leib, Ehre, Gut und alles fahren, und setzest es dran. Denn ein gotteliger Mensch zeucht GOttes Wort allen Dingen vor, als den edelsten Schatz. Wenn du aber anfingest, solche zu zählen, so würdest du ihrer keinen finden, der GOtt liebt von ganzem Herzen. Denn die Gottlosen pflegen um eines Worts und (wie man sagt) Taubenfußes willen GOtt und den Nächsten hintanzusetzen, zu lästern. So halten die Geizhälse den Mammon für ihren Gott.

25. Nun, dies ist das erste Gebot: „Du sollst nicht andere Götter haben“, das ist, GOtt, deinen HErrn, sollst du lieben, sein göttliches Wort hören und dir gefallen lassen; was sein Wort heißt und verheißt, das laß dir gesagt sein, und dem nach halte dich. Ja, dies sein Wort soll dir das edelste Kleinod auf Erden sein. Dir soll nicht mehr gelieben dein Leib, Leben, Ehre, Gut und alles, was du hast. Aber wir laufen über die Gebote GOttes, wie eine Sau über das Heiligthum, um unserer eigenen Lüste und mannigfaltigen Begierden willen, gleich als wäre nie kein Gesetz von GOtt gegeben. Folgt weiter im Mose:

\begin{schriftzitat}
V. 6. 7. 8. Und die Worte, die ich dir heute gebiete, sollst du zu Herzen nehmen, und sollst sie deinen Kindern schärfen und davon reden, wenn du in deinem Hause sitzest, und auf dem Wege gehest, wenn du dich niederlegest oder aufstehest, und sollst sie binden zum Zeichen auf deine Hand, und sollen dir ein Denkmal vor deinen\footnote{Im Original und in der Erlanger: denen.} Augen sein, und sollst sie über deines Hauses Pfosten schreiben, und an die Thür.
\end{schriftzitat}

26. Siehe, wie heftig er treibt zum ersten Gebot, auf daß ja in unsern Herzen und innerlichen Affecten eitel Glaube und herzliche Liebe gegen GOtt und dem Nächsten brenne und herrsche. Und spricht: „Du sollst diese Gebote schärfen.“ Das ist ein Meister mit Predigen und Auslegen des Gesetzes. Er weiß, daß am ersten Gebot die meiste Macht und Kraft liegt, darum handelt er es auch so gar mit hohem Fleiß. Will sagen: Laß dir diese Gebote, und sonderlich das erste, mit allen Treuen befohlen sein, daß du sie nicht allein aufs Papier klebest, oder in ein Buch fassest, darnach im Winkel
\quellseite{874}
liegen lässet, oder sie dir allein in die Ohren erschallen, vor den Augen gemalt stehen, sondern fasse und schreibe sie in dein Herz, daß du Tag und Nacht daran gedenkest, gerne davon redest und handelst, du arbeitest, stehest oder gehest gleich wo du wollest. Summa: Diese Worte sollen dein höchster, liebster Schatz sein, denn wo dein Schatz ist, da ist auch dein Herz [Matth. 6, 21.]. Du sollst zu jeder Zeit dir sie lassen am meisten angelegen sein.

27. Warum redet aber Mose also? Warum, daß er weiß, daß so viel Leute auf Erden sind, wenn sie sollen die zehn Gebote hören und lernen und ins Leben ziehen, so wissen sie dieselbigen schon allbereit auswendig, und können sie zuvor allzumohl. Darum ist es ihnen eben, als hörten sie eine Fabel und Märlein, wenn man dieselben vorsagt und erzählt; oder, wie eine neue Zeitung vom Türken oder vom Könige aus Frankreich; achten es nicht, meinen auch nicht, daß sie die zehn Gebote etwas angehen, viel weniger sinnen sie darauf, daß sie dieselben ins Herz schließen sollen, darnach thun und sich halten. Sie meinen, es sei genug, wenn sie es zu Zeiten einmal hören, und davon wissen etlichermaßen zu reden, wie wenn einer sonst etwas Seltsames gehört, und kann es nachreden. Und also ist der mehrere und der größte Theil und Haufe der Welt gesinnt, daß sie entweder Ruhm darinnen suchen, und machen nur davon, oder thun etwas zum Schein allein, das ihnen doch nicht ums Herz ist. Solches ist sehr verdrießlich, wenn die Leute so bald flug werden, und des Worts GOttes müde und überdrüssig, gleich als wäre es nur darum zu thun, daß man es bisweilen höre, und lasse es nachmals gut Wetter sein; wiewohl ihrer gar sehr viel auch das Zuhören verachten und aufstehen lassen. Jene, so bald sie es gehört, und es ausgesagt ist, so denken sie nach einem andern, und wenn das andere kommt, suchen sie noch ein anderes, und so fort immerdar; vergessen darnach des Ersten und Nöthigsten.

28. Ich habe oft gesagt, und sage es noch, welcher die zehn Gebote, und sonderlich das erste Gebot recht kann, dem will ich von Herzen gerne zu Füße sitzen, und ihn lassen meinen Doctor sein. Ich halte mich gelehrter, denn die Schwärmer sind, denn sie können nicht die zehn Gebote. Ich kann sie aber, GOtt Lob! Das weiß ich aber gleichwohl, daß die zehn Gebote noch mein Donat, ABC, ja, meine Bibel sind; ich muß noch darin ein Schüler bleiben, ob ich gleich die Bibel oft ausgelesen habe. Aber die Klugen und Naseweisen, wenn sie nur ein Sermonichen können, so meinen sie, sie können es alles, und treiben ihre Hoffahrt mit unmäßigem Waschen; aber den Teufel auf ihren Kopf wissen sie.

29. Es ist ein feindseliges Volk, das so bald gelehrt wird. Denn sie gedenken nicht anders, denn daß sie die zehn Gebote als gemeine Dinge hören, und können davon waschen, plaudern und schwätzen, als sonst von etwas Anderem. Es gilt aber, lieber Gesell, nicht Schnatterns, sondern daß du sie ins Leben und in die That hinein bringest, und gegen jedermann öffentlich beweisest, daß du also gesinnt seist, ehe du wolltest wider einiges Gebot GOttes handeln, du wollest lieber alles auf Erden drüber fahren lassen. Da sei nun Meister, und beweise dich redlich.

30. So schließen wir derhalben, daß kein Mensch auf Erden die Gebote GOttes hält und erfüllt. Ja, auch solche Geister und Wäscher greifen es nicht mit einem Buchstaben, oder dem geringsten Werkt an. Darum ist dieses ihre Art: Auf Einen Tag lernen sie das Gesetz und Evangelium gar aus, daß sie viel können nachweisen, aber im Treffen, in der That und Beweis ist niemand daheim. Es sind nur Federn und Haare.

31. Diesen bösen Geistern und schändlichen Plauderern zuwieder setzt Mose diese Worte und spricht: „Du sollst sie zu Herzen nehmen.“ Das ist, werdet mir nur nicht zu bald klug; denke nicht, daß du es alles getroffen, verschlungen und vorlängst hinweg habest, oder daß es so bald ausgelernt sei, als es gehört wird. Nein; ich lasse mir nicht genügen, daß du es auf die Zunge und in die Ohren fassest, darnach überhin lassest rauschen und schweben. Ins Herz sollst du es fassen und stecken, darinnen soll es dein Trost, Trotz und Behelf sein und bleiben. Darnach strebe, dahin komme.

32. Nach diesem schärfe und webe sie deinen Kindern. Erstlich gedenke für deine Person, daß du sie wohl lernest, nicht allein im Munde, in den Ohren, und auf der Zunge sie tragest, sondern ins Herz senkest, daß dir es von Herzen gefalle, darnach zu leben. Darnach, wenn du diese Gebote im Herzen hast, so fahre fort, daß
\quellseite{875}
du sie deine Kinder auch lehrst. Und sagt eigentlich von schärfen und wetzen; sagt nicht, du sollst sie allein lehren und ihnen vorsagen; sondern mit diesem Worte „schärfen“ zeigt er an, was wir für Gesellen sind. Es will mit schlechter Anweisung und Vermahnung sich nicht thun lassen, sondern getrieben, geübt, angehalten und nachgedrungen sein. Denn unser Herz ist stumpf, verstockt und verrottet, daß es an ihm nicht haftet; es fällt sehr bald vom Worte, darum muß man es immer handeln. Sonst hat es einen Ekel vor dieser Speise, wird ihrer bald überdrüssig; ja, es will immer Meister sein, ehe denn es Schüler worden ist. Darum thut Anhalten stets noth bei denen, die uns Predigern befohlen sind, daß sie es desto besser lernen.

33. Ich kenne selbst etliche, die da meinen, wir dürften keiner Prediger oder Pfarrherren, und man müsse die Pfaffen Gewohnheit und alten Herkommens halben dulden; man könnte die Besoldung und Unkosten, so auf sie jährlich gehen, wohl in andere und bessere Wege gebrauchen. Gleich als wären sie (wie jener sagt) ein necessarium malum. Sonderlich der Adel und etliche Klüglinge sagen: Haben wir doch Bücher, daraus wir es ebensowohl lesen können, als hörten wir es in der Kirche vom Pfaffen. Du liefst den Teufel auf deinen Kopf, der dich denn besessen hat. Wenn unser HErr GOtt gewußt hätte, daß das Predigtamt nicht noth wäre, er wäre ja so weise und klug gewesen, daß er dir es nicht durch Mosen hätte predigen lassen, und wäre nach deinem gottlosen, teuflischen, thörichten Gedanken und Reden ohne Noth gewesen, daß er nachmals das levitische Priesterthum hätte geordnet, und allezeit Propheten ausgesendet, wie er selber sagt [Matth. 23, 34.]. Er würde auch dieser Zeit Prediger und Seelsorger wohl heißen daheim bleiben. So hätte er auch nicht den Eltern dürfen einbinden, daß sie es so ernstlich trieben, und nicht ließen hingehen. Er weiß gar wohl, wo es uns mangelt.

34. Das folgt von Natur: Wenn du dahin kommst, daß du meinst, du könnest nur aufs beste das Evangelium, die zehn Gebote und Worte GOttes, so bist du verloren, und hat der Teufel gewonnen Spiel. Darum, wenn man dieser Lehre überdrüssig wird, und sie uns schmeckt wie eine Neige vom Fasse, da will das Herz was Neues haben, des Vorigen ist es müde und übersatt [2 Tim. 4, 3.]. Dann sagt man gemeiniglich also: O das habe ich vorhin lange gehört, sage mir etwas Anderes her. So will man immerdar eine andere Lehre haben, und stinkt den Leuten die Nose nach Neuerung. Wenn also das Herz GOttes Worts müde ist, es nicht mehr für seinen besten Schatz hält, dann ist die Thür hinten und vorne offen, daß der Teufel einen freien Zutritt hat, und allerlei Irrthum einführen mag [Luc. 11, 26.].

35. Also ist es mit den Schwärmergeistern auch zugegangen, wie zuvor in allen Kirchengeschichten; hiermit sind sie vom Teufel herumgerückt, gefällt und übervortheilt worden, dadurch hat er Platz bekommen. Es war ihnen das Evangelium und der Glaube von Christo nicht ein großer Schatz, sie mußten etwas Neues haben; vor großer Kunst hatten sie einen Ekel vor dem Evangelio, Christus war ihnen zu geringe; sie konnten’s alles; es war ihnen eine verdrießliche Lehre. Wenn ein Kranker nicht mehr Speise mag, sondern es grauet und ekelt ihm davor, so ist er nicht weit vom Tode, er stirbt bald. Also auch, welchem die Himmelsspeise des göttlichen Worts unangenehm ist, und nicht mehr schmeckt, der wird es nicht lange antreiben.

36. Es gedenke niemand, so lange er lebt, daß er das erste Gebot werde auslernen. Weil GOtt selber viel davon hält, und will, man könne nichts Herrlicheres noch Besseres predigen, und Moses sich nicht schämt, daß er immerdar Ein Lied auf Einer Saite fiedelt, das ist, einerlei treibt, so sei es uns auch nicht Schande, ein ewig Ding stets zu lehren und zu hören. Ich sollte mich ja billiger schämen, und einen Ueberdruf daran haben, daß ich immer einerlei vorbläue und lehre, denn ihr, die ihr es sollt von mir anhören. Aber der Heilige Geist und GOtt wird es nicht überdrüssig, einerlei Ding zu lehren, und wir sind also ein feindselig Volk, und verachten alles. Der liebe GOtt behüte uns vor demselbigen Geschmeiß, daß wir nicht auch dahin gerathen, daß wir uns dünken lassen, wir können es gar. Wahrlich, der Text ist solchen sattan(i)amen, ekeln Geistern zuwider hier gesetzt, auf daß niemand sich lasse dünken, er habe es alles und gar ausgelernt.

\begin{schriftzitat}
V. 7. Und davon reden, wenn du in deinem Hause sitzest, oder auf dem Wege gehest, wenn du dich niederlegest, oder aufstehest.
\end{schriftzitat}

\quellseite{876}
37. Siehe, wie Moses dies so fleißig treibt: Du seiest daheim oder auf dem Felde, so sollst du von dem Wort reden; darum sollt du deß\footnote{Im Original und in der Erlanger: dies.} nicht überdrüssig werden. Davon sollt du handeln, wenn du aufstehst, und wenn du schlafen gehst, es in die Hände malen, und über die Thür schreiben, daß du es allezeit vor Augen habest, und immer darnach thuest, und allzeit dir ein Gedächtniß sei. Was ist das? Wie meint er das? Mose wollte dies Gebot nicht allein gern ins Herz zu bläuen und geben, sondern auch ins ganze Leben.

38. Die Juden haben eine Weise aus diesem Texte genommen, davon Matth. 23, 5., daß sie eine Pergamenthaut ums Haupt machten, daran die zehn Gebote geschrieben waren, und schrieben sie auch um die Kleider; gleichwie wir jetzt GOttes Wort predigen, lesen, singen, malen, drucken und schreiben. Dieses war bei den Juden nicht eine böse Weise und Gewohnheit; denn sie wollten GOttes Wort vor den Augen haben, und malten es an allen Orten, auch in den Gärten, und haben solche Weise gewißlich aus diesem Texte genommen. Aber es waren Buben und Schälke, wie man pflegt im Sprüchwort zu sagen: Ein Paternoster tragen am Halse, einen Schalk im Herzen etc. Es ist zwar solches nicht böse, aber es ist nur ein Schein, da man sich gottelig stellt, und mit der That es gar läßt anstehen; darum ist es eine Heuchelei, und Christus ist ihnen auch feind, und straft sie hart darum.

39. Darum meint es Moses also: Du sollst von diesem Gebote predigen und reden im Hause, das ist, was du thust im Hause oder auf dem Felde, oder an welchem Ort du seiest, so sollst du immer daran gedenken, daß du dawider nicht thuest. Bin ich auf dem Markte, daß ich an GOttes Gebot gedenke, daß ich meinem Nächsten keinen Schaden thue; denn GOtt hat mir geboten, ich soll nicht stehlen; wo dir das Gebot gefällt, und von Herzen Gott lieb ist, so stiehlst du nicht. Wenn du dies Gebot, „du sollst nicht stehlen“, auf dem Markte hast, oder auf dem Acker, wo man Güter, und [Aecker]\footnote{Von uns hinzugefügt, um Sinn zu geben.} an seinen Nächsten stoßen hat, da sei ihm gehorsam; hast du GOtt und sein Wort lieb, so wirst du vor allem Betrug, Wucher, Geiz und Hinterlist dich vorsehen, und nicht wider ihn thun. Das heißt davon handeln, daß man das Leben darnach richten soll.

40. Also, wenn du im Hause handelst, bist ein Handwerksmann, Brauer, Schuster, Schneider, Bäcker etc., so gedenke: Ich will also handeln mit meinem Nächsten, daß ich keinem zu nahe greife, noch übertheure, übersehe, betrüge, vervortheile. Denn GOtt hat geboten im siebenten Gebot: Ich soll nicht stehlen, niemand umrücken noch betrügen etc. Siehe ich Gottes Wort und Gebot, so werde ich niemand Unrecht thun. Aber wo sind solche? Christen sollen allezeit so sagen: Ich will mein Leben darnach richten, daß ich wider meinen GOtt nicht sündige, noch meinen Nächsten beleidige. Und diese halten und erfüllen dies Gesetz, die es also ins Leben hinein bringen.

41. Die Pharisäer und Heuchler mag man fahren lassen, die das Gesetz auf die Hüte und an die Röcke schrieben, und nimmermehr an GOttes Gebot gedachten, und ihren Nächsten betrogen. Denn dies Schreiben an die Hüte hilft nichts, wenn du es auch gleich mit Scheidewasser einbeizt,\footnote{Original: einbeißt.} denn ein solcher ist und bleibt doch ein Schalk. Aber ein frommer Christ soll also sagen: Ich will mein Leben, Werk und Geschäft mit GOttes Hülfe also richten, daß ich wider meinen GOtt nicht sündige, und niemand Unrecht thue; daß du gedenkest, GOtt zu lieben, fürchten und vertrauen, und niemand Leid, sondern jedermann Gutes zu thun; da fahre fort, du bist auf rechtem Wege.

42. Welche also die Gebote in ihr Leben treiben und fortsetzen, die schreiben das Zeichen auf ihre Hand. Es ist alles darum zu thun, daß du stets daran gedenkest in allen deinen Worten, Werken und Vorhaben, daß du GOtt fürchtest, ihm trauest und niemand Schaden thust, nicht tödtest, nicht ehebrechest, nicht stiehlest etc., sondern jedermann nützlich seiest. Hebe das an, und greif dein Leben an; treibe es ein Jahr lang, und sage mir es über ein Jahr wieder, was du kannst vom ersten und den andern Geboten allen, so wirst du es gewahr werden. Gedenke, daß du nicht wollest das Deine suchen, deinen Nächsten nicht hintergehen, niemand übertheuern, dann wirst du sehen, was GOtt lieben heißt, und wirst inne werden, daß du noch nicht
\quellseite{877}
gelehrt hast das liebe ABC, auch dich für einen groben Sünder erkennen. Dann würde nicht so viel Diebstahl und böser Thaten sein. Liebtest du GOtt, so würdest du nicht den Mammon anbeten. Dein ganzes Leben zeuget wider dich, daß du GOtt nicht liebest, sondern vielmehr hassest.

\begin{schriftzitat}
V. 9. Und sollst sie über deines Hauses Pfosten schreiben.
\end{schriftzitat}

43. Das ist: Du sollst gedenken, du gehest aus oder ein, du handelst daheim oder draußen mit deinem Nächsten, daß du also lebest, daß du nicht wider deinen GOtt thust. Darum will GOtt diese Gebote allenthalben vor die Augen geschrieben haben.

44. Das ist eine nöthige Vermahnung, die Moses zum ersten Gebot thut. Denn er hat gesehen, daß die Leute aus Hoffahrt sich dünken lassen, wenn sie nur GOttes Wort gehört haben, so können und wissen sie es alles vollkommen, so sie doch ärgerlich leben, und wenig darnach thun. Darum will er sie weisen von dem Gehör ins Herz und Leben, wie sie mit dem Herzen und in ihrem Wandel leben sollen. Dann werden die Leute auch sehen, wie sie dem Nächsten Schaden thun mit Worten und Werken, suchen stets das Ihre, und denkt ein jeder nur auf seine Sache: GOtt gebe, sein Nächster bleibe, wo er wolle. Denn die Welt hat den Mammon lieber denn GOtt.

45. Aber das sind erst die rechten hohen Stücke, wenn GOtt, der HErr, uns nicht allein seine Gebote vorlegt, die wir verachtet, sondern wenn er uns auch angreift, Plagen und Unglück zuschickt, in das Kreuz und Verfolgung wirft, Armuth und Krankheit an Hals hängt, daß du meinst, GOtt sei dein Feind. Dann siehe, ob du in solcher Trübsal und Jammer GOtt lieben könnest, und wie geduldig du darinnen seiest; ob du auch ohne Murren solches leidest, als Hiob, der da noch lobt und dankt in seinem Kreuz, und spricht [Hiob 1, 21. 2, 10.]: „Haben wir Gutes empfangen von der Hand des HErrn, warum wollten wir das Böse auch nicht fürlieb nehmen? (Gelobet sei der Name des HErrn.)“ Da wirst du die rechten Knoten befinden. Wenn du schon deinem Nächsten nichts hast Uebeles gethan, und daß du die zehn Gebote GOttes ein wenig hast gehalten, so wird sich gleichwohl finden, daß dir der Wille GOttes nicht gefällt, der doch dein höchster Trost und Schatz sein sollte. Item, wenn du geschmäht und gehändet wirst, siehe dann, ob du deinen ärgsten Feind auch mögest lieben.

46. Aber das lassen wir jetzt fahren, und bleiben nur bei der gemeinen Weise, daß man diese Gebote nicht schlecht ansehe, als, daß man GOtt allein liebe mit der Zunge, sondern daß man es fasse in das Herz und in das ganze Leben, daß du stets gedenkest, was du nur redest, thust, vorhabest und beginnest, in alle deinen Sachen und Leben, daß du die zehn Gebote nicht übertretest, GOtt und deinen Nächsten nicht erzürnest, damit das Wort GOttes sei dein bester Schatz, und dir am allerliebsten.

\begin{schriftzitat}
V. 10–12. Wenn dich nun der HErr, dein GOtt, in das Land bringen wird, das er deinen Vätern, Abraham, Isaak und Jakob, geschworen hat, dir zu geben; große und feine Städte, die du nicht gebauet hast, und Häuser alles Gutes voll, die du nicht gefüllet hast, und ausgehauene Brunnen, die du nicht ausgehauen hast, und Weinberge und Oelberge etc., so hüte dich, daß du nicht des HErrn vergessest, der dich aus Egyptenland geführet hat.
\end{schriftzitat}

47. Nachdem Moses erklärt hat das erste Gebot, nämlich, „von ganzem Herzen GOtt lieben“, so fährt er weiter, und will, daß wir bei demselbigen Worte bleiben sollen und nicht eine andere Predigt suchen, wir können denn diese recht wohl. Nun will er jetzt auch die Hindernisse und Aergernisse erzählen wider dies Gebot, oder die Sünden, damit man wider das erste Gebot handelt und sich vergreift, und will diese aus dem Wege räumen, daß wir bei dem ersten Gebote bleiben, und solche Stücke meiden und fliehen.

48. Das erste Hinderniß und Aergerniß ist Junker Mammon, der heißt Reichthum und Ueberfluß, darüber Moses harte klagt und prophezeit, der wird dich reißen von deinem GOtt, und von seiner Liebe, daß du fallen wirst auf schöne Häuser, Oel und Korn, Acker und Weinberge, und andere zeitliche Güter, das wird alsdann dein Gott werden, und wirst des HErrn, deines GOttes, darüber vergessen, denn du mit solchen leiblichen Gütern überschüttet wirst [5 Mos. 32, 13–15.], wie sonst aller Welt Gott ist der verfluchte Mammon. Da will nun GOtt allhier im Texte wehren, und ihnen einbinden: Das alles, was jetzt erzählt, soll dir nicht so lieb
\quellseite{878}
sein, als GOtt; denn diesen sollst du über alles lieben, und sollst ehe Haus, Hof, Schlösser, Acker und alles fahren lassen, daß du nur die Gebote GOttes behaltest.

49. Aber diese Worte sind bald geredet: „Du sollst GOtt lieb haben“,\footnote{„lieb haben“ von uns gesetzt (nach der Bibel) statt: „lieben“, um fünf Worte zu machen, wie gleich folgt.} jedoch haben sie einen schweren Anblick; es klingt und klappt solches nicht in unserm Herzen; GOttes Wort gilt nicht in unserm Herzen. Was liegt an den fünf Worten? denkt mancher; um derenwillen kann ich Haus, Acker, Wiesen und andere meine Güter nicht verlassen; ich nehme hundert Gulden für diese Worte: Liebe GOtt, deinen HErrn. Mir liegt an den Worten nicht viel, mein Haus und der Mammon geliebet mir viel mehr. Darum ist Mammon oder der Geiz der erste, der dich hindert an der Liebe GOttes. Daß du seiner vergißest und ihn verachtest, das werden machen die schönen Häuser, Weinberge. Das ist der Mammon. Davon sagt auch Baruch, der Prophet [Cap. 3, 17.], daß Menschen auf Gold sich verlassen. Und David weiß der Menschen Seuche und Krankheit wohl, da er im 62. Psalm, V. 11., sagt: „Fällt euch Reichthum zu, so hänget das Herz nicht daran.“ Darum sagt auch St. Paulus [Col. 3, 5.], daß der Geiz sei eine Abgötterei, und eine Wurzel alles Bösen [1 Tim. 6, 10.]. Darum seid aus diesem Texte gewarnt. Denn es ist der erste Teufel und die erste Anfechtung der Mammon oder Reichthum; die machen, daß man GOtt nicht lieb hat, und daß man vom Gottes Wort abgezogen wird.

50. Geht es nicht in der Welt also zu? Bezeuget es nicht die Erfahrung? Wenn die Kinder erwachsen und groß werden, machen sie einen Unterschied zwischen den Gütern und Eltern. Welche sind es, die lieber ihre Eltern haben denn die Güter? Wie gedenken sie des Schmerzens, Jammers und Mühe, welche die Eltern erlitten haben, daß sie sind erzogen worden, ja, daß sie Ehre, Leib und Gut an sie gesetzt? Da sie nun erzogen sind, wollten sie, daß die Eltern, Vater und Mutter, todt wären, daß sie allein die Güter bekämen. Denn das Gut ist einem viel lieber denn die Eltern. Ja, man findet Frömmichen, die rechten mit den Eltern um die Güter. Sie danken den Eltern also, daß sie dieselbigen gerne todt sähen, damit sie das Patrimonium bekämen und an sich brächten. Da hebt sich denn der Zank, Zorn und Haß gegen den Eltern. So könnt einer sprechen, der Gott liebt und fürchtet: Ehe ich mit dir, lieber Vater, wollte zanken, ehe wollte ich alles Gut fahren lassen. Welche nun gottesfürchtig sind, die thun also.

51. So sehen wir auch, wie oft Brüder untereinander um Guts willen Todfeinde werden. Wer zerreißt also die Liebe unter den Brüdern? Junker Mammon. Aber wo hier brüderliche Liebe wäre, da sprächest du: Ehe ich mit dir zürnen will; ich wollte eher, daß alles Gut in der Elbe schwämme.

52. Also handelt ein Bürger und Bauer auch mit dem andern. Wer kann nun die Kunst, daß die Kinder die Eltern verachten? Wer richtet solche Uneinigkeit zwischen Brüdern an, die unter Einem Herzen gelegen, Eines Fleisches und Blutes sind? Der Gott Mammon, der richtet es alles aus. Also zertritt der Mammon GOttes Gebote, daß sie nicht angesehen werden, weder von Kindern, Schwestern, Brüdern noch Nachbarn, und müssen darüber hintangesetzt werden die natürlichen Eltern, Brüder, und alle guten Freunde; wie man sieht in den Theilungen, da man stracks gedenkt: Hätte ich nur das Haus, die Wiesen, mein Vater, Mutter, Bruder und Schwester wären, wo sie wollten. Siehst du alhier, wie ein gewaltiger Herr und Gott der leidige Mammon sei, daß er auch alle GOttes Liebe und Gebote zunichte macht, und für nichts halten lehrt?

53. Davor warnt uns Moses treulich, und will sagen: Der Mammon wird dich anfechten und reizen, thue die Augen auf, und bleib bei dem einigen GOtt, und gedenke: Pfui dich an, du leidiges Gut, daß ich um deinetwillen sollte Feindschaft, Haß, Zwietracht und Neid haben wider meine Eltern, Bruder und Nachbarn; wenn ich nimmermehr keine Güter hätte, so habe ich doch GOtt, der kann mir sie wohl geben. Ich soll der weltlichen Güter, nach der Lehre St. Pauli, gebrauchen, gleich als hätte ich ihrer nicht, 1 Cor. 7, 30. 31.: „Die da kaufen, als behielten sie es nicht, und die dieser Welt brauchen, als brauchten sie ihrer nicht.“ So laß dir GOtt lieber sein, denn alle Güter auf Erden. Hat er dir Haus, Gärten und Weinberge gegeben, die du doch nicht gepflanzt hast, so traue auf ihn; der dir das gegeben hat, kann dir anderes mehr geben.

54. Wahrlich, was du erlangst von deinen
\quellseite{879}
Eltern, das hast du nicht erworben noch verdient, sondern GOtt hat dir es gegeben durch die Eltern. Das müssen wir ja bekennen. So du aber rathfragst den Mammon: Welches edler und besser sei, GOtt und Eltern, oder Güter? Da sagt er: Güter sind besser denn GOtt und Eltern. Darum vermahnt Moses allhier, daß wir uns GOtt und sein Gebot lieber lassen sein, denn alle Güter auf Erden. Denn wenn dir alles genommen wird, so ist er doch dein GOtt noch. Er kann dir wohl mehr geben. Und wenn er dir gleich nicht mehr geben wollte, so wird er dich doch nicht lassen Hungers sterben, denn er spricht: „Ich bin dein GOtt.“ Es wird gewiß sein, daß dir der HErr, dein GOtt, verheißen hat; allein, hange an seinem Worte, und wisse, daß Güter wohl hernach folgen werden.

55. So sage ich nun: Das erste Aergerniß wider dies Gebot ist der Mammon; der will ein anderer Gott sein, und ist auch ein anderer Gott. (Es ist bei den Juden also gewesen, und auch noch bei uns, daß wir mehr auf das Gut sehen denn auf GOtt und sein Gebot. Das richtet der Mammon an. Denn, wenn du Gott liebest und seine Gebote, so würdest du nicht um Guts willen wider GOtt und sein Wort, auch wider deinen Nächsten dich setzen. Nun aber hanget man an dem Mammon, verachtet und haßt GOtt. Das ist ein gewiß Zeichen, daß man GOtt nicht liebt, dieweil so vorsätzlich wider GOtt und den Nächsten gehandelt wird.

56. Aber es wird uns nicht geschenkt werden, sagt Moses [V. 15.], „denn dein GOtt ist ein eifriger GOtt unter dir“. Darum siehe darauf, daß nicht einmal sein Grimm erzürne, und vertilge dich von der Erde. Da siehst du, daß [es] nicht ein Scherz sei, darfst es nicht in Wind schlagen, noch ein Gespött daraus machen. Denn, halte ich es heute nicht, morgen auch nicht, so werde ich nicht mehr einen GOtt an ihm haben, sondern ein verzehrend, fressendes Feuer, das alles verschlingen wird, wie droben [Cap. 4, 24.] gesagt; das ist, er wird dich aus dem Lande ausrotten und vertilgen, dich bringen um Leib und Leben, und darnach um die Seele, daß er dich mit deinem Reichthum und Gütern ewiglich verdamme.

57. Wie uns denn die Erfahrung solches lehrt mit den Wucherern und Geizhalsen, welche die Leute aussaugen und schinden bis auf den Grat mit ihrem Stehlen und Rauben, darüber GOtt und sein Wort, seine Diener und alle Warnung verachten; so verachtet er sie wieder, daß ihr Reichthum, Geld und Gut auf den andern oder dritten Erben nicht kommt. Denn male quaesit, male perdit, sagt der Wale. Da sie aber als Christen ihnen Gott lieben lieber sein denn den Mammon und die Güter, so hätten sie aller Dinge genug, und wäre Reichthum und die Fülle in ihren Häusern, auch bei ihren Nachkommen, wie der 37. Psalm, V. 25., sagt: „Ich bin jung gewesen, und alt worden, und habe nie gesehen den Gerechten verlassen, noch seinen Samen nach Brod gehen.“ Item [Sprüchw. 10, 22.]: Benedictio Domini divites facit. Summa Summarum: Das ist der erste Teufel und Hinderniß, welches macht, daß wir GOttes und seiner Gebote vergessen. Da sind wir aber gewarnet, daß wir lernen sein Gebot theuer und werth halten, denn alles Gut, so auf Erden ist. GOtt wolle uns vor dem Mammon behüten, Amen!

58. Also ist bisher gehandelt das erste Hinderniß wider das erste Gebot, nämlich, der Mammon, so ein rechter Abgott und ärgster Feind GOttes ist. Und Moses warnt ernstlich die Seinen wider diesen Abgott, daß sie sich sollten vorsehen, daß sie sich den Mammon nicht lassen reiten und bethören, noch ihm dienen, und das zeitliche Gut höher lieben denn GOtt. Darum hat er gesagt, sie sollen die zehn Gebote ins Herz schließen, stets treiben und üben, als eine Lehre, die nicht könne genug geübt, eingebildet und gehandelt werden. Denn der Teufel, die Welt und unser Fleisch sitzen und schaffen allezeit das Widerspiel, stürmen und toben wider dies Gebot. Dieweil nun diese drei nicht müde werden, wider dies Gebot zu fechten und sich dawider zu legen, so sollen wir auch nicht müde noch verdrossen werden, dies Gebot dagegen zu üben und zu treiben.

59. Also bleibt diese Lehre allezeit nöthig im Volk, daß Moses spricht: „Du sollst“ nicht allein „diese Worte schärfen“, sondern, „wenn du im Hause sitzest, auf dem Wege bist“ \&c. Das ist, es ist nicht vonnöthen, daß du es an die Wände mit Buchstaben schmierest, über die Thüre und Pfosten zum ewigen Gedächtniß malest; wie die Juden solche Weise gehabt, sondern in allen Geschäften, Hantierung, Nahrung und Arbeit sollen wir zusehen, daß wir uns nicht an jemand vergreifen, oder ihm Leid thun, sondern also handeln, daß GOtt mehr bei uns vermag, und mehr
\quellseite{880}
Recht bei uns habe denn der Mammon. Denn, wir gehen aus oder ein, daß wir unsere Sachen beschicken, wir kaufen oder verkaufen, es sei unser Wesen, was es wolle, so sehe einer sich vor, daß der Mammon ihn nicht betrüge, daß er demselben mehr anhange denn GOtt. Es soll ein jeder auf seine Hand, Augen und Werke sehen, was er thut, wie sein Wesen und Handel stehe, daß ihn der Mammon nicht verführe zum Argen und Bösen.

60. Der schändliche Mammon ist der größte Gott auf Erden, und es ist gar ein gefährlich Ding, mit Gut umgehen [Matth. 19, 23.]. Die da haben Reichthum und Güter, die trauen gemeiniglich darauf [Luc. 12, 19. 16, 19.]; die nicht haben, die trachten Tag und Nacht darnach, und wollen gerne etwas gewinnen. Wer Reichthum hat, der hat seine Lust und Freude daran; dagegen fürchtet sich jedermann vor Armuth. Es muß ein jeder eine Werbung, Geschäfte, Arbeit und Hantierung haben, und derselben warten; und da ist es denn gefährlich, wo etwas erworben und zusammen getrieben wird. Denn es schlägt gerne dazu, daß man entweder seine Hoffnung und Trost hat aufs Gut, oder Zweifel vom Mangel. Dies Leben schwebt zwischen Mangel und Habe, zwischen Gut und Armuth; daher neigt sich’s an der zweier Enden eines, das ist, auf Vertrauen oder Verzweiflung. Hieraus kommen denn allerlei Bübereien, daß man unter einander schindet und frißt, scharrt und trägt, und daß einer den andern verportheilt und betrügt, wo er nur kann, nach dem gemeinen Sprichwort: Welcher den andern vermag, der steckt ihn in Sack.

\subsection[Große und kleine Städte, die du nicht gesehen hast,]{Große und kleine Städte, die du nicht gesehen hast, Häuser voll Gutes, Wein- und Oelberge etc}

61. Da streicht er es mit mehr Worten und etwas weitläufiger aus; und hier sieht man den emsigen Fleiß Mosis in Ausstreichung dieses ersten Gebots. Erstlich hat er gesagt, und diese Worte gesetzt: „Ich bin der HErr, dein GOtt, du sollst nicht andere Götter haben.“ Dazu thut er eine Erklärung: „Ich bin ein eifriger GOtt, der die Sünde der Väter heimsuchet an den Kindern, bis ins dritte und vierte Glied.“ Item, legt auch dazu eine Verheißung, als: „Denen, die mich lieben und meine Gebote halten, denen thue ich wohl ins tausendste Glied“ [5 Mos. 5, 8–10.].

62. Zuletzt spricht er: Willst du die Meinung wissen, ich will dir es sagen: Wenn du wirst ins Land kommen, und schöne Häuser, Gärten, Weinberge, Oelberge einnehmen und besitzen, Wasser mit Fischen, und was ich dir alles überflüssig und reichlich geben werde ohne deine Arbeit und Verdienst, dann schaue darauf, daß wird der erste Teufel werden, welcher dir wird die Augen blenden, daß du GOttes Furcht aus dem Herzen lässest, GOtt verachtest, auch des ersten Gebots vergessest. Denn du wirst dich hängen an die köstlichen Häuser, Weinberge, Aecker und andere Güter, so ich dir geben werde etc., und wirst einen Gott daraus machen, und dein Vertrauen darauf setzen, und nicht so viel nach mir fragen, sondern mein vergessen.

63. Also ist schon eine Abgötterei angerichtet, da ein Mensch mehr auf den Mammon trauet denn auf GOtt, und fürchtet sich, daß er ihn ja nicht möge verlieren. Moses hat dies nicht vergebens gesagt, und so fleißig gewarnet und vermahnt, denn er hat gewußt, daß der Mammon ein solcher Gott sei. Und hier hören wir mit starkem Beweis darthun, daß unser Wesen, Natur, Vernunft und Wille geneigt ist wider das erste und alle anderen nachfolgenden Gebote. Wir sind solche Gesellen: Wenn es uns wohlgeht, so vergessen wir GOttes, und können ihm nicht danken, und erkennen nicht seine Wohlthaten, daß er uns dazu geholfen hat, können auch nicht zurück sehen auf den Stand, da wir es nicht hatten.

64. Ist aber Gefahr, Trübsal und Noth vorhanden, oder daß es mangeln will, so äuget\footnote{Im Original: eiget; „sich äugen“ (äugen) = sich zeigen.} sich keine Zuversicht, Glaube und Hoffnung in uns; das Gebet liegt danieder, Verzweiflung geht uns zu. Das haben nicht allein die Propheten und Kirchenlehrer, sondern auch die Heiden gesehen. Und viel weiser Leute haben über das Glück und Unglück, über gute und böse Tage geklagt, daß man sich hierin nicht wisse zu schicken. So ist auch ein deutsch Sprichwort: Es müssen starke Beine sein, die gute Tage tragen können. Item: Der Mensch kann alles erleiden, allein gute Tage nicht. Man sagt auch: Wenn dem Esel zu wohl ist, so geht er aufs Eis tanzen, und bricht ein Bein.

65. Der Poet Ovidius spricht sein: Luxuriant animi, rebus plerumque secundis, nec facile est, aequa commoda mente pati. So
\quellseite{881}
kann nun ein Mensch alles tragen, wie groß und schwer es ist, ausgenommen gute Zeit, die ist ihm unmöglich zu tragen. Das merkte Virgilius auch, und machte diese Verse drauf:

\begin{quote}
Nescia mens hominum fati sortisque futurae, et servare modum rebus sublata secundis.
\end{quote}

Ein Mensch, wenn es ihm wohlgeht, kann er sich nicht halten aus eigenen Kräften, vermißt sich, und überhebt sich des Reichthums und Glückes, und geht dahin.

66. Lies die Historien, beide der Christen und Heiden, da wirst du es fein finden. Ueber das so gibt es auch die Erfahrung, daß Gut macht Muth, Muth macht Uebermuth, Uebermuth thut nimmer gut. Das heißt, gute Tage nicht können ertragen. Man sieht es auch in der Haushaltung, und ein jeder in seinem Stande greift es doch mit den Händen, wenn er es nicht wollte sehen in der heiligen Schrift. Wenn ein Knecht allzu einen frommen Herrn hat, ein Kind einen freundlichen Vater, eine Magd eine gütige Frau, so kann man es nicht leiden; je gütiger ihnen geschieht von ihrer Herrschaft, je weniger sie es achten. Also thun auch die Bürger; sie sind allezeit desto muthwilliger, je gelinder der Rath und die Obrigkeit ist.

67. Die Bauern, je gelindere Herren sie haben, je ärger sie sind. Soll jemand mit Glimpf und Güte sie zu was bringen, der muß gar gut Glück haben; darum ist es also: Wenn man einem Bauer fehlt, so verschuldt ihm der Bauch davon. Also ist es mit einem jeden Menschen. Wenn er gute Tage hat, ist nicht wohl mit ihm umzukommen, bis er ein Unglück anrichtet, und GOtt ihm die großen Federn ausrupft; denn das Fütterlein sticht ihn, es ist ihm zu wohl. Davon hat man gar viel Sprüchwörter, die gehören alle hieher, und sind auf diese Meinung gestellt, daß, wenn es einem Menschen glücklich geht, und er sich etwa fühlt, so ist nimmer wohl mit ihm umzugehen; denn er wird muthig, eigensinnig und wunderlich; der Teufel kann nicht mit ihm auskommen.

68. Dieses klagt hier GOtt durch Mosen, und will gleich nach menschlicher Weise so reden und sagen: Ich thue nicht wohl, daß ich dich reich und satt mache, denn, wenn ich dich in das Land bringen werde, so wirst du dich umkehren, und wider mich löcken.\footnote{löcken = frech und muthwillig sein; hüpfen, springen, hinten ausschlagen.} Wie sonst im Mose geschrieben steht [5 Mos. 32, 15.]: „Israel ist dick und fett worden“, darum fähet er an zu löcken, und schlägt mich mit der Ferse. Zu Summa, wir müssen mit dem Kreuze und Leiden gedrückt werden, auf daß wir nicht stolz seien.

69. Darum ist der Welt nichts Besseres denn ein scharf, hart, stark Regiment, wie noch bei den Türken ist, und etliche Junker vom Adel, geschwinde, züglerische Regenten sind, welche den Bauern den Mantel nehmen sammt dem Rocke. Solche groben Tyrannen gehören in die Welt, nicht fromme Obrigkeit, auf daß sie ihre Hoffahrt dämpfen. Denn die Welt läßt sich nicht regieren mit Güte und Sanftmuth, denn der bösen Buben ist allezeit zu viel.

70. Daher siehst du, worüber GOtt klagt, nämlich, wenn er es durch seine väterliche Güte den Leuten läßt wohlgehen, so vergessen sie seiner und aller seiner Wohlthat. Jetzt fehlt dem Deutschlande nichts, denn daß [es] ihm so wohl geht, ein jeder thut, was er will; Deutschland hat gar diese Plage, ihm ist zu wohl, und die Obrigkeit zu gelinde. Darum ist in allen Ständen ein lauter Muthwille, es lebt fast ein jeder seines eigenen Gefallens. Es mangelt, daß die Frösche nicht Störche zu Königen haben; wahrlich, es geht der Esel aufs Eis tanzen. Aber sie werden so lange Muthwillen treiben, bis daß der HErr kommt, und sagt: Du hast mich lange genug verachtet, ich will dich Mores lehren, auf daß du wissest, ich sei noch GOtt. Denn die Erfahrung trifft immerdar mit der heiligen Schrift überein.

71. Derhalben so sagt Mose: Siehe zu, es wird dir nicht geschenkt werden; geschiehet es nicht, so strafe Mosen Lügen. Der fromme Moses läßt es an seinem Fleiß nicht mangeln, wiederholt immer die vorigen Warnungen, und sagt: Siehe dich vor, dein GOtt ist ein eifriger GOtt, daß du seinen Zorn nicht über dich erregest, und er dich ein Stücklein sehen lasse, darauf du nie gedacht hättest. Das reibt er ihnen stets in die Nasen, daß sie ja GOtt fürchten sollen, und vom HErrn sich nicht lassen abwendig machen, seiner auch nicht vergessen, denn er sei ein zorniger GOtt, und strafe bis ins dritte und vierte Glied. Aber die lieben Juden fragten nichts darnach, GOtt mag im ersten Gebot eifrig oder unleidlich sein, wie er wolle.

72. Diese Seuche steckt uns auch zwischen Fell und Fleisch. Wenn wir nur gute Tage haben,
\quellseite{882}
So fragen wir weder nach GOtt, Pfarrherren, Fürsten, noch Vater und Mutter, ja, nach dem Teufel dazu nicht. Aber wenn auch die Stunde kommen wird, daß einmal der Zorn angeht, so wird es heißen: Vertilgt von der Erde! Dann wird es dich jammern, daß dein Weib vor deinen Augen geschändet wird, deine Kinder getödtet, die Eltern erschlagen werden, ja, dir dazu die Haut wohl ausgeklopft wird, und über das der Hals abgeschnitten. Da will man Hülfe, Rettung und Schutz haben. Dann wird es aber zu lange geharrt sein.

73. Darum will dir der HErr hier zu verstehen geben: Wirst du wider mich Muthwillen treiben, so will ich auch unbereit und unwillig sein zu helfen, wenn der Tod, Teufel, Krieg, Pestilenz und alles Unglück mit Haufen kommen, und werde dich auch lassen zappeln und stecken, ja, der Letzte\footnote{Letzte = Valete, Garaus.} mit dir spielen.

74. Es wird dir hier zuvor gesagt, damit du dich habest darnach zu richten, daß, ob du gleich wirst schreien zu GOtt, so wird er dich nicht erhören. Denn allhier steht der Text klar, der es dräuet, GOtt will deinen Muthwillen, wider ihn und deine Herrschaft getrieben, also heimsuchen, daß du selber sagen sollst, er habe dich recht bezahlt und getroffen. In Summa: gute Tage können nicht anders, denn GOttes vergessen, auch folgt darauf Vermessenheit, die geht über, und GOttes Furcht geht unter.

75. Ach! daß ein jeder jetzt in seinem Stande gedächte: Es geht mir, GOtt Lob! wohl; ich will GOtt fürchten, ihm von Herzen dafür danken und dienen, und will daneben dem Meinen getreulich nachkommen, und thun, was mir in meinem Stande auferlegt ist, will demüthig und eingezogen sein; ich weiß, der HErr wird mir Barmherzigkeit widerfahren lassen. Werde ich aber ungehorsam und undankbar sein, so hat mir der HErr gedräuet, er wolle mich von der Erde vertilgen. Wenn man das gedächte, würden nicht so viel Klagen und Plagen sein von der Schinderei, und so gar übermachter Untreue der Leute. Jetzt meint aber ein jeder, unser HErr GOtt sei gestorben, es werde seine Büberei und Bosheit nicht gestraft werden. Also soll ein jeder sich drein schicken, unserm HErrn GOtt gehorsam zu sein, und Gutes zu thun, ehe denn sein Zorn anbrenne.

76. Im türkischen Regimente geht es mit dem Dienstgesinde also zu: Wenn eine Magd nicht will gehorsam sein, verkauft man sie um drei Groschen; dann kauft sie irgend einer, der schlägt in sie, wie du auf eine Kuh oder Esel schlägst, ist mit einer Peitsche hinter ihr her, wenn sie nicht thun will, was sie soll. So wird auch sonst allenthalben ein strenges Regiment gehalten, und geschieht besser Aufsehen denn bei uns Christen. Der Türke hat mehr Gulden, denn unsere Fürsten Heller haben; dennoch thut er Fleiß bei seinen Sachen. Ein jeder unter dem türkischen Dienstgesinde hat sein Abgemessenes an Speise, Trank, Arbeit etc.; verbringt er es nicht, wie er soll, so ist er bald da mit Ruthen und Peitschen; hilft das nicht, so schlägt er mit dem Schwert drein, das Messer folgt bald hernach und hauet ihm den Kopf herab.

77. Aber bei uns, wenn ein Arbeiter oder Gesinde einen Tag oder zwei versäumt, oder seinem Herrn sonst ungehorsam und untreu ist, so hat es kein Gewissen drüber, sondern meint, es thue recht; darum ist jetzt schier niemand Vater und Mutter oder Obrigkeit gehorsam. Aber wird der Teufel einmal unter dich kommen, so wirst du von der Erde vertilgt werden. Darum so muß ein solches türkisches Regiment folgen, das uns zu Trümmern schmettere. Eine solche Herrschaft gehört in die Welt, daß man anders den Leuten zuspreche, gleich als der Türke thut; wenn er einen Finger aufhebt, so geht es von statten, wie er will, in seinem ganzen Reich. Und wo solche Regimente nicht sind, da wird GOttes vergessen, und werden allerlei Stände, als Obrigkeit, ihre Gebote, Eltern, Herren und Frauen, und andere verordnete Personen verachtet.

78. Denn die muthwilligen Leute lassen sich dünken, man muß uns wohl dulden und haben, man kann unser nicht entbehren. Das heißt, ohne GOttes Furcht gelebt, GOttes vergessen, GOtt nicht gescheuet; wir vergessen GOttes, und fürchten ihn nicht vor guten Tagen, und vor dem Mammonsdienste, denn der ist unser Gott. Wenn es aber also wäre, daß ein jeglicher nur ein Stücklein Brods und Trünklein Wassers haben sollte, und doch seine abgemessene Arbeit, als, ein zehn Fuder Mist aufzuladen und hinauf zu führen, und wenn er das nicht thäte, flugs hinter ihm her mit der Staupe, dann würden sie wohl frömmer werden, und GOtt nicht also
\quellseite{883}
verachten. Aber weil das nicht ist (wie wir in glückseliger Zeit dahin nicht kommen mögen), so ist es auch unmöglich, daß sie besser werden. Darum folgt:

\begin{schriftzitat}
V. 15. Daß nicht der Zorn des HErrn, deines GOttes, über dich ergrimme.
\end{schriftzitat}

79. Was hat GOtt den halsstarrigen Juden gethan? Zur Zeit des Königes Ahab in viertelhalb ganzer Jahr ließ er nicht ein Tröpflein regnen [1 Kön. 17, 7.]. Da kam der Zorn auch über sie, da verdorrete das ganze Erdreich. Es ist schrecklich und greulich zu sagen, da war eine große Theurung [1 Kön. 18, 2.], daß da keine Speise für die Menschen, kein Futter für das Vieh war, also, daß man alles Vieh schlachtete. Und wie ging es nachmals in der endlichen Zerstörung Jerusalems, bezahlte sie GOtt nicht redlich? daß sie nach andern Strafen zuletzt Taubendreck und Eselköpfe kauften, einen fast um drei Gulden, welches sie sonst nicht hätten essen dürfen. Die Kinder starben Hungers den Eltern in den Armen; andere sahen, daß ihre Kinder erschlagen wurden; etliche schlachteten ihre Kinder; darnach wurden die Städte geplündert. Also rächte GOtt ihre Undankbarkeit, und strafte ihre Verachtung. Als wollte er sagen: Ihr habt meiner vergessen, nun, ich will euer jetzt wieder vergessen. Und da die Noth vorhanden war, schrieen sie; aber er hörte sie nicht, und half ihnen nicht.

80. Also wird es uns undankbaren Deutschen auch dermaleins widerfahren, daß wir im Unglück und Nöthen schreyen werden, und zum Bescheid bekommen das Wort: Mein Zorn ist ergrimmet, und werden kein Gehör noch Hülffe haben. Denn wir haben gar zu gute Tage, wir wissen nicht (mit Urlaub), wo uns der Hintern steht. Auf meine Seele, es ist nicht zu scherzen; es ist nicht möglich, daß [es] ungebüßt bleibe; die Strafe kann nicht ausbleiben. Es muß über Teutschland eine große Plage kommen; ich fürchte, es werde alles auf einmal kommen, Pestilenz, Krieg, Hunger. Niemand fürchtet GOtt, es ist alles muthwillig; Gesinde, Bauern, Handwerksleute, thut alles, was es will. Niemand straft, ein jeder lebt nach seinem Willen, bescheißt und betrügt den andern. Es ist unmöglich, daß [es] lange so sollte bleiben, denn die guten Tage betrügen uns.

81. Wollte GOtt! daß unter uns etliche wenige solches verstünden, die da ernstlich beten; daß doch die Strafe aufgehoben würde, und, wo sie nicht gar verbliebe, daß doch der Zorn nicht käme zu unsern oder unserer Kinder Zeiten. Denn wenn Krieg und der Türke kommt, so ist Glaube, Predigtamt, Erkenntniß Christi, Sacrament, als, Taufe und Abendmahl, alles dahin; wie dies unter dem Türken zu sehen ist; und werden lauter Bestien aus uns, daß sie uns treiben und jagen wie die Kühe. Dann wirst du müssen thun, was du jetzt nicht thun willst. Und welche nicht nach GOttes Gebot jetzt wollten leben, denen schickte man nur den Türken über den Hals. Wenn das türkische Reich dem Evangelio nicht schadete, so wollte ich ihn über uns wünschen, daß er unser Herr wäre, und uns wohl plagte, die wir jetzt also sichere Geister sind. Aber es ist unmöglich, er vertilgt das Evangelium, und die Kirche GOttes zerstört er.

82. Darum ist besser, wir leiden das Ungemach, Muthwillen der Leute und Verworffenheit der Bauern, bei dem Evangelio, denn daß wir den Türken zum Herrn hätten, ohne das Evangelium, und das Evangelium verloren. Wären wir aber unter dem Türken, so müßte ein jeder seine gemeine Speise, gewiß Biel und Zeit zu essen, und eine bestimmte Arbeit haben, und nicht einen Trunk Weins, sondern Wasser, und sein Stück Brods, und nicht einen Bissen drüber. Du müßtest da nicht zum Bier gehen; wie es denn gut wäre, daß man der Leute Schlampamp und Freßerei und Säuferei wehrete. Wie denn jetzt in der Türken den gefangenen Christen das Essen gewehret ist, und jetzt Hunger leiden, und arbeiten müssen wie die Esel. Wenn du es dann nicht wolltest thun, so käme der Herr mit dem Stricke, bände und geißelte dich. Dann würdest du erfahren, daß ich nicht gelogen hätte; würde ich lügen, so schilt mich. Du würdest dann gerne für den Frieden danken, und GOtt der HErr würde dich sein aus der Bierschenke in des Teufels Namen holen, müßtest dagegen arbeiten, als rauchte dir die Haut. Da würdest du an mich gedenken. Lüge ich denn, so lügt Moses auch; aber er wird nicht lügen werden, darum wird es bei dir auch wahr werden.

83. So thue nun ernstlich, was Gott dir gebeut. Sei nicht ein Knecht des Mammones, vergiß Gottes nicht. Darnach stiehl nicht, raube
\quellseite{884}
nicht, dienen, arbeite treulich, thue dein Amt, was dir befohlen ist, mit Fleiß und Treue. Sei gehorsam deinen Oberherren, so wirst du einen gnädigen GOtt haben, und wirst gesegnet sein, reich werden, es wird dir wohl gehen. Wirst du aber nicht treulich und redlich arbeiten und mit dem Nächsten aufrichtig umgehen, so sollst du vertilgt werden; der Zorn, Grimm und Verfolgung soll dich treffen.

84. Lieber, laß dich die guten Tage nicht betrügen. Aber das ist vergebens, und wird den Tauben geredet. Wie man denn sagt, es sei unmöglich, daß man könne gute Tage ertragen. Darum muß man es mit GOttes Wort einbläuen, uns ein wenig aufrichten, und vermahnen und treiben, daß wir GOtt fürchten, und seiner nicht also vergessen, ob es doch ein wenig helfen wollte, wie es sonst ohne die Predigt, wahrlich, niemand würde achten. Denn wenn der Natur zu wohl ist, und sie ohne GOttes Wort dahin lebt, so geht sie aufs Eis tanzen, und läuft ins Lerchenfeld, da sticht sie das Futter. In Gefährlichkeit aber kann sie sich auch nicht erhalten, da ist nichts denn Zittern, Beben, Fürchten und Verzweifeln.

85. Darum ist die Predigt des göttlichen Worts angerichtet, dadurch man einen Menschen vermahnt zur Furcht GOttes, und bei GOtt zu bleiben, wenn es ihm wohl geht; herwieder, auch nicht zu verzagen, wenn es ihm übel geht, sondern es hält an, im Glauben festiglich zu beharren. Also vermahnt uns auch hier das göttliche Wort, und spricht: Siehe dich vor, hüte dich, wenn du gute Tage hast; denn, weißt du, was es bedeutet? Wenn man die Schweine mästet, so will man sie schlachten, so ist den Säuen der Tod näher, denn wenn sie in der Irre umher laufen; wie es Jeremias, der Prophet, auch sagt [Jer. 12, 3.].

86. Also, wenn GOtt dir auch gute Tage gibt, gibt Schrot\footnote{Schrot = geschrotetes Korn.} und Träber, und du mästest dich mit GOttes Gütern, so gilt es wahrlich deinen Balg. Einer wird kommen, dich schlachten, und Würste aus dir machen, und dich auffreßen, und das Schmeer, Schmalz und Fette verzehren, das ist, dein gesammeltes Gut und Schatz hinwegnehmen.

87. Das ist ein Stück, daß Moses will sagen: Lieber, wenn gute Zeit da ist, und keine Heimsuchung oder anderer Unfall vorhanden, und ihr sitzet in Bierhäusern, saufst, singst und löret, und fragt nicht viel darnach, was [das] Korn gilt, wo man Essen und Trinken nehme. Wie das Gesinde denkt: Der Herr und Frau müssen es wohl schaffen; sie sind ohne Sorgen, und die Tagelöhner arbeiten vier Tage in der Woche. Denselbigen wird hiermit gesagt: Weil es so wohl steht, sehet euch vor vor der Sicherheit und Vergeßung GOttes, werdet nicht muthig und stolz; denn es ist Zeit, daß ihr an GOttes Gebot gedenket, und nicht GOttes vergesset.

88. Wahrlich, du solltest GOtt danken, daß du da mit guter Ruhe sitzest als ein Juncker, und GOtt, der HErr, dir diesen Frieden so lange erhält; deine Schafe treibst du ohne alle Sorge und Fahr aus der Schäferei aus und ein, daß dich die Mörder nicht erschlagen. Aber ja, daß du nur einmal gedächtest: Lieber GOtt, hilf mir, daß ich thue, was du willst haben, daß ich dich liebe und meinem Nächsten diene, ich will niemand betrügen, meine Eier und Käse rechten\footnote{„rechten“, was im Original fehlt, ist von uns aus der alten Ausgabe herübergenommen.} Kaufs geben. Niemand will dieses befinden, und GOtt für solche Gaben danken, der Obrigkeit gehorsam sein, und dem Nächsten dienen; sondern man sagt: Ich will retten, saufen mein Bier, mit Vortheil, Betrug und Ueberseßen das Meine verkaufen; denn es ist mein. Und brauchst also des edlen Friedens und der Gaben GOttes, daß du GOtt schändest und lästerst. Derhalben so wird er auch kommen, und dich in seinem Zorn aus der Erde reißen und vertilgen. Verstehst ihr mich? Ich rede ja deutsch. Ihr habt glückselige Zeit, ich wünsche euch auch Glück zu euren guten Tagen, daß ihr derselbigen recht brauchet. Aber sehet zu, kommt Krieg, da du jetzt zween Pfennige gibst für eine Kanne Bier, dann gäbest du gerne einen Groschen.

89. Du wirst auch nicht alsdann auf dem Markte müßig sitzen, wie jetzt, da man in allem Reichthum und Ueberfluß lebt, und dieser Wohlthat gebraucht, nicht GOtt zu Dank, sondern vielmehr ihm zu Trotz und Schanden, und dem Nächsten zu Schaden; du denkst nicht an den zeitlichen Frieden, wie du GOtt dafür dankest. Also bekümmert sich in einer Haushaltung der Knecht nichts um Essen und Trinken, sorgt nicht für Futter, Kleider und Lohn, achtet es alles für nichts; denn er weiß und verläßt sich drauf,
\quellseite{885}
der Herr muß [es] wohl verschaffen. Denkt nicht, wie gute Tage er habe, läßt einen andern für sich sorgen, ist seiner Herrschaft ungehorsam; meint, er sei darum da, daß er nur esse und trinke; mißbraucht also der guten Tage, und sticht ihn das Futter zu sehr, und macht sich beschissen genug.

90. Aber viel sicherer thun solches die großen Hansen und Gewaltigen. Heißt das nicht GOttes vergessen, wenn er dir Gutes thut, und du weißest nicht, wie du dich übel genug stellen sollst, GOtt zu verdrieß, und dem Nächsten zu Schaden, du gedenkst nicht an den zeitlichen Frieden? Wie, wenn Theurung, Krieg und Unfriede käme, und nähme diesen Schatz des friedlichen, genugsamen Lebens hinweg? Da werden Knechte und Mägde so wohlfeil werden, daß zwanzig, dreißig, vierzig Mägde und Knechte gerne um ein Stück Brod fleißiger dienen, denn sonst jetzund ums Geld. Es geht bereits also an andern Orten, da Theurung ist, da nimmt eine ein Stück Brods, und dient zwei Tage; jetzt dient sie kaum ein halb Jahr so fleißig um großen Sold und Lohn.

91. Derhalben sagt Moses: „Hüte dich, daß du nicht vergesse des Herrn, der dich aus Egypten geführet hat.“ Aber wir sehen es jetzt nicht, es ist außer unsern Augen; aber es wird uns einmal in die Faust kommen; man ringt darnach, und es wird dazu kommen. Wir verfahren alles, und vertrauen allzu sehr auf die guten Tage, da wir uns doch fürchten sollten in den guten Tagen, und nicht zu sicher sein; wiederum, bedenken die Zeit der Noth, daß man alsdann auf GOtt vertrauen und nicht verzweifeln möchte. Denn GOtt wird nicht lügen; wo du wirst deiner guten Tage mißbrauchen, so will er dich vertilgen. Das ist eins, daß wir uns fürchten sollen, und des Bösen nicht vergessen zur glückseligen Zeit; daß du gedenkest, was es würde werden, wenn Krieg oder Theurung käme, und wenn du solltest einen dienen, der dir kaum ein Stück Brod gäbe, und schlüge dich dazu zu Boden. Wiederum, daß man auch auf GOtt hoffe, zur Zeit der Trübsale. Das heißt, als der weise Mann Salomon [Sprüchw. 30, 9.] sagt: Wenn es uns wohl geht, so überheben wir uns, wenn es übel geht, so verzagen wir; dieses hat er aus dem ersten Gebete genommen. Das ist ein Stück unserer Unart: Wenn es der Natur wohl geht, so kann sie sich nicht enthalten, sie vergißt GOttes drüber; und läßt Hände und Füße gehen, wenn es übel steht.

92. Darum muß GOttes Gebot und Wort kommen, uns ermuntern, erwecken und warnen, dir selbst zu gut, daß du nicht zu Schaden kommest, und daß du dich vorsehest, daß, wenn es dir wohl geht, du nicht auf deine gegenwärtigen Gaben und Güter vertrauest, und wider GOtt und deinen Nächsten thuest; sonst, wenn die Predigt nicht wäre, so würde es kommen, daß ein Esel den andern in einen Sack steckte. Also vermähnt uns Salomon gleichergestalt, als Moses gethan hat, daß wir von der Sicherheit nicht betrogen werden, auf daß der Esel nicht aufs Eis tanzen gehe. Darum muß man dem Esel das Futter höher legen, daß er mehrlich\footnote{mehrlich = kaum.} die Haut ertragen kann, und daß man ihm die Rippen zählen möge; Säcke muß man ihm auflegen.

93. Im Buch der Sprüche Sirachs [Cap. 33, 25.] wird gesagt: „Dem Esel gehört dreierlei, Futter, Sack und Stecken; also, einem Knechte gebühret auch Futter, Arbeit, Werk und Ruthen“; das ist, einem jeden, der was zu thun hat, der da Amt oder Dienst hat, er sei Bürgermeister oder Richter, Knecht oder Magd, oder was er sonst für einen Beruf hat, der soll drei Stücke haben. Erstlich, Essen und Trinken soll man ihm geben, aber darnach einen Sack dazu, daß er zu schaffen und zu beschicken habe; daß [es] nicht allein dabei bleibe, daß er esse und trinke, sondern daß er auch arbeite. Und wo er sich sperrt, daß er allein das Futter, und auch nicht den Sack will haben, da gehört dazu ein guter Stecken, und auf die Haut gedroschen, daß man ihm die Striemen zähle, so wird ihm der Kitzel und Lecker vergehen.

94. Ein jeder will nur Futter haben, aber du sollst den Sack auch dazu tragen. Der Esel trägt den Sack nicht gerne, und der Knecht hat auch nicht gerne die Arbeit; sondern das Futter und den Lohn wollen wir alle haben, befleißigen uns alle guter Tage. Es fehlt auch nicht am Futter, ja, wir haben’s zu viel; wir hätten auch Säcke, aber es fehlt am Stecken, es ist kein Ernst, Zorn und Nachdruck da, wie es wohl sollte sein. Der Stecken sollte auch hernach folgen, daß man zuschlüge. Da gehört auf uns der Türke, oder sonst ein greulicher Tyrann; wiewohl unsere Fürsten auch Stecken genug sind.

\quellseite{886}
\begin{schriftzitat}
V. 16. 17. Ihr sollt den HErrn, euren GOtt, nicht versuchen, wie ihr ihn\footnote{„ihn“ fehlt im Original und in der Erlanger.} versuchtet zu Massa; sondern sollt halten die Gebote des HErrn, eures GOttes, und seine Zeugnisse, und seine Rechte, die er geboten hat.
\end{schriftzitat}

95. Das ist das andere Stück. Denn, wie gesagt, wenn es einem zu wohl geht, so fürchtet man GOtt nicht, man vergißt seiner, und wird einer hoffärtig und vermessen. Es gibt keiner auf den andern, meint, man dürfe sein nicht, er dürfe niemand dienen; daher raubt und stiehlt man in der Welt. Wiederum, wenn es übel geht, so kann Fleisch und Blut nichts weniger denn böse Tage leiden. Der Mammon hat zween Arme, damit regiert er die Welt; mit dem rechten Arme regiert er, wenn es wohl geht, da thut ein jeder, was er will; dann fürchtet man den Mammon, und da ist GOttes vergessen. Wenn es übel geht, daß böse Tage vorhanden sind, so regiert er mit dem linken Arm, dann versucht der Mensch GOtt, den HErrn.

96. Nun ist die Versuchung GOttes mancherlei. Allhier aber redet er von der rechten Versuchung GOttes, nämlich, da Israel in der Wüste GOtt versuchte mit dem Himmelbrod; da sie weder Speise noch Trank hatten, gingen sie hin zu Mose, und schalten sich einmal mit ihm, murreten um des Hungers willen wider GOtt und Mosen [2 Mos. 16, 2. 17, 2.]. Wollten schier sagen: Der Teufel hat es euch geheißen, daß ihr uns aus Egypten geführt habt; und meinst du, daß GOtt allhier sei? Sonst hätten wir je zu essen. Wollten derhalben nicht glauben, noch einigen Mangel leiden; sie wollten nicht harren und vertrauen, sondern murrten alsbald, wenn nicht Essen vorhanden war; sie setzten GOtt ein Ziel, Maß, Person, Zeit, Weise: so sollst du mir es geben, durch diesen, auf die Zeit etc.; sie wollten unsern HErrn GOtt gebunden haben, daß er es ihnen nach ihrem Willen und gesetzten Ziel und Zeit geben sollte, und konnten nicht Einen Tag zu GOtt\footnote{„GOtt“ von uns herübergenommen aus der alten Ausgabe statt „Gut“ im Original und in der Erlanger.} harren, und sagen: Ich bedarf ja dies und das, aber ich stelle dir es anheim; wie oder wann du mir es gibst, so will ich es gerne annehmen.

97. Das heißt GOtt versucht, wenn wir ein Ding haben wollen von GOtt, wie wir es ihm abzirkeln, daß er uns geben müsse; oder wir fallen von ihm ab, daß wir auch nicht erwarten können Einen Tag, oder ihm heimstellen, was wir bedürfen, oder ihm vertrauen, und fest auf ihm stehen, wenn es übel geht. Daher kommen im Papstthum Abgöttereien, Wallfahrten, daß wir den Heiligen feiern: also ist man zur Eiche gelaufen, wenn eine Noth hat einen angestoßen, daß die Heiligen helfen sollten. Ein jeder hat Wege gesucht, daß er nichts dürste leiden, und GOtt trauen und von ihm Hülfe gewarten. Das ist GOtt versuchen, wenn mir etwas mangelt, daß er mir es bald geben soll, wenn ich es haben will, und was ich von ihm begehre, und nicht warten oder auf seine Hülfe hoffen wollen.

98. Also thaten die Juden, verfluchten und plagten GOtt wohl in der Wüste, murreten wider GOtt. Was sie nur wollten, mußte er schier thun auf einen gestrackten Termin, als, Wasser, Wachteln und Himmelbrod etc. geben; zielten es ihm ab mit Maß und Zeit, wann er es ihnen geben sollte; aber er nahm sichs nicht an. Also regiert zu jeder Zeit der Mammon auf die zweierlei Weise: Erstlich, wenn es wohl geht, so ist der Mammon Herr, und macht, daß wir GOttes vergessen, macht uns sicher. Wiederum, wenn es übel geht, so fährt er abermals zu, und macht, daß du im Unglück GOtt versuchest, oder ihn wegwerfest, und einen andern annehmest. Also will der Mammon regieren in der Fülle und im Mangel. Sieht man nicht solches vor Augen? Sie wollten es traun gerne in der Faust haben. Haben sie es, so ist GOttes vergessen. Geht es wohl, so taugt es nicht; geht es übel, noch viel weniger. Geht es nun wohl oder übel, so ist ihm nicht zu rathen.

99. Wie soll man ihm denn helfen? Der Mammon steht und sticht. Das erste Gebot steckt in den Nöthen, daß man es nicht halten kann, es gehe wohl oder übel. Wie soll man die Leute fromm machen? Wie soll man doch die Welt regieren? Es gebe doch einer Maß und Weise. Denn geht es wohl, so ist GOttes vergessen; in Gefährlichkeit wird GOtt versucht. Wenn der HErr sagte: Ich will einen gute Tage geben, so wird er mir dienen. Ja, hinter sich, wie die Bauern die Spieße tragen. Gebe ich es ihnen, so danken sie mir nicht; zürne ich aber mit ihnen, so laufen sie von mir, und hängen sich an andere Götter, oder verzweifeln.

100. So geht es: Lache ich, so werden sie
\quellseite{887}
stolz; zürne ich, so verzagen sie. Also kann die Natur die Welt nicht regieren, und können die Gebote GOttes nicht gehalten noch veräußert werden; denn unsere Natur kann es nicht lassen, daß sie nicht sollte stolz werden, wenn es ihr wohl geht; also wiederum, wenn es übel geht, so kann sie auch nicht bestehen und auswarten. Also siehst du, wie GOttes Gebot auf beiden\footnote{Im Original: beider.} Seiten mehrt, und doch gleichwohl verachtet wird. Das wäre ein Mittelweg oder Mittelstraße, wenn sie möchte gefunden werden, daß [es] GOtt nicht zu übel noch wohl ließe gehen, daß nicht eitel Gutes, auch nicht eitel Böses ginge; denn welches geht, so wird GOttes vergessen, er wird verachtet oder versucht. Aber das kann nicht geschehen, darum mag man sich wohl vorsehen.

101. Alle Abgötterei unter und vor dem Papstthum, mit dem Stiften, Wallen gehen etc., kommt her vom Mammon, daß wir gute Tage wollen haben, und uns selber helfen; derselbige wehrt, daß man GOtt nicht vertraue, sondern verachte GOttes Gebot, welches sagt: „Ich bin dein GOtt.“ Aber wenn uns nicht Glücks morgen, oder in dieser Stunde, und gerade auf die Weise, wie mir wollen, geholfen wird, so laufen wir zu einem andern Gott.

102. Siehst du nun die Summa des ersten Gebots, was dasselbige haben will? nämlich, wenn dir es wohl geht, so verachte mich nicht, sondern fürchte mich, spricht GOtt; und wenn dir es übel geht, so verzweifle nicht, sondern glaube an mich und vertraue mir; wenn dir es übel geht, so denke auch an mich. Wir sollen alle bekennen und sagen: O HErr, ich bin nicht werth, daß ich hier mit Weib und Kind und dem Meinen im Frieden sitze, und dies Stück Brod habe; und erkennen, daß es von GOtt käme. Darum, weil es von GOtt ist, so sollten wir ihm dienen und dafür danken, auch, wie der Text lautet, „ihn lieben von ganzem Herzen“, fürchten und vertrauen, und also im Kreuz glauben und geduldig sein, und auf GOttes Barmherzigkeit und gnädige Hülfe warten; denn er will uns helfen. Also, wenn Pestilenz vorhanden ist, wiewohl das böse ist, doch, weil er es uns zugeschickt hat, so wird GOtt auch kommen, und will seines Worts Anhängern helfen. Darum soll niemand in irgend einem Anliegen anderswohin laufen.

103. Also soll dies Gebot auf die Finger geschrieben, und vor Augen gehabt werden in allen Werken, auf daß, wenn es wohl geht, ich nicht stolz werde, auch wiederum nicht hinfalle und verzage in Widerwärtigkeit; denn das sind Hindernisse, so uns von GOtt abwenden. Daran sollen wir nun gedenken, und es ins Leben führen; und weil uns das Fleisch, die Welt und der Teufel damit reizen und plagen, und so stark sind, und uns von GOtt abziehen, darum muß man es stets handeln, und uns zum Worte halten, und nicht lassen von Gott abreißen.\footnote{Anmerkung der Eisl. Ausg.: Diese Zeit über hat Doctor Martinus Luther zu Wittenberg nicht gepredigt, denn er ist zu Marburg in Hessen auf dem Colloquium mit den Sacramentirern gewesen, von dannen um Simonis Judä [28. Oct.] wieder gekommen und hat weiter gelesen. — Das Marburger Gespräch war am 1.—4. October 1529. Am 12. October war Luther auf der Rückreise in Jena; am 17. October schrieb er bereits von Wittenberg aus einen Brief an Myconius (Waldh, alte Ausg., Bd. XXI, 1186.). Deshalb haben wir in der ersten Anmerkung zu dieser Schrift die Wiederaufnahme der Predigten in die Mitte Octobers gesetzt.}

104. Also haben wir die Auslegung des ersten Gebots, welches Moses gar fein erklärt, daß man es verstehen kann, und zeigt alle Hindernisse, Abgötterei und arge List an, welche der Teufel dagegen setzt, daß er uns vom rechten Glauben abführe. Dagegen habt ihr gehört, wie er das allwege schärft und wetzt mit Verheißungen und Dräuungen, dadurch er die schrecke, welche ihn nicht wollen für einen GOtt halten; und wiederum, verheißt er Barmherzigkeit denen, so auf ihn trauen, denn er wolle helfen den Gläubigen.

105. Das erste Aergerniß ist gewesen der leidige Mammon oder zeitliche Gut. Davon hat er diese Worte gesetzt: Hüte dich, wenn du kommst in das gute, volle Land, daß dich die Fülle nicht betrüge, und du nicht deinen Trotz und Hochmuth setzest auf das zeitliche Gut, daß es nicht dein Herz einnehme, und dann nicht gedenkest an den, welcher dich aus Egyptenland, aus dem Diensthause, geführt, und dir dies Land gegeben hat, und du diesen Gütern geholfen; vergessen also deß, durch welchen du dazu gebracht bist.

106. Darum ist der Mammon der erste und vornehmste Abgott, der vom ersten Gebote abführt. Denn das menschliche Herz ist von Natur hochmüthig [Jer. 17, 9.], wenn es weiß, daß es etwas vermag, dadurch sich’s aufhält; denn
\quellseite{888}
Gut macht Muth, und dieser Muth macht einen Abgott. Darum, gewinnst du Güter, so setze deinen Muth nicht darauf. „Fällt dir Reichthum zu, so hänge dein Herz nicht daran“ [Ps. 62, 11.]. Thut ihm aber also: Lasset euch das Gut nicht einen Muth machen, sonst wird GOtt nicht gefürchtet, und wird der Mensch sicher und kühn. Darum spricht er allhier: Denket an den, welcher euch aus dem Elend geführt, daß dich nicht das Gut betrüge, als hättest du es von dir selber; sondern allezeit gedenke, als wärest du noch im Elend. Vergiß nimmermehr des Elends und der Noth, daraus du bist erlöst. Behalte dein Herz also, gleich als hättest du es noch nicht, und als wärest du noch in der Armuth.

107. Also sind wir jetzt auch geführt aus mancherlei Stricken des Gewissens und großen Elends und der Tyrannei des Papstthumes; es ist unsere Hölle und Kerker gewesen. Jetzt sind wir frei, kein Gesetz treibt uns mehr, weder zu feiern noch zu fasten. Da sollen wir uns hüten, daß wir aus dieser Freiheit nicht einen Muth fassen, daß wir stolzirten, gleich als dürften wir GOttes nicht mehr; sondern desto öfter zurück gedenken, und ihm dankbar sein, daß er uns aus diesem Gefängniß erlöst hat. Darum muß man immer wider den Papst predigen; denn wir sollen immerdar gedenken an das Elend und Diensthaus, darin wir so jämmerlich zerplagt und gemartert worden sind, und nicht vergessen der Wohlthaten GOttes; sonst werden wir undankbar, hart und starr, daß wir weder Verheißungen noch Dräuung achten.

\begin{schriftzitat}
V. 15. Denn der HErr, dein GOtt, ist ein eifriger Gott.
\end{schriftzitat}

108. Siehe, wie Moses treibt die Dräuworte in dem ersten Gebot, und wie er so ernstlich vermahnt, hält sie uns stets vor Augen und Ohren; denn er hat immer Sorge und fürchtet sich, daß sie nicht sicher und zu stolz werden, und vergessen GOttes im Glücke. Wie wir uns denn auch müssen besorgen, daß wir uns nicht der Freiheit erheben, und des vergessen, so uns aus den Stricken und Elende geführt hat, sondern sollen gedenken: Der uns aus der Gefängniß erlöst hat, der kann uns auch wieder hinein werfen, und, das noch mehr ist, von der Erde vertilgen, als wären wir Staub. Darum hält er uns GOtt vor als einen Eiferer, und als ein verzehrendes Feuer, der sich wohl weiß zu rächen [Jes. 1, 19. 20.]. Darum so sollen wir uns vor ihm fürchten, und viel lieber glauben denn fühlen.

\begin{schriftzitat}
V. 16. 17. Ihr sollt GOtt, euren HErrn, nicht versuchen, wie ihr ihn versuchtet zu Massa.
\end{schriftzitat}

109. Traue auf GOtt, spricht er, wie das erste Gebot erfordert, daß man an GOtt glaube und ihm vertraue, und nicht auf das Gut; mache dir nicht einen Abgott daraus, daß du mit dem Stolziren wolltest, daß du in dem\footnote{Original: den.} Kasten oder Händen hättest. Thue ihm also auch: So dir es würde übel gehen, als, wenn Krieg, Theurung etc. vorhanden, und mangelte dir an Gewächs der Erde, halte fest, und versuche GOtt nicht. Denn, wie zum öftermal angezogen, so ist das unsere Art, also sind wir gesinnt, daß wir in der Noth verzweifeln; und stolzieren, wenn es uns wohl geht. Im Ueberfluß kann sich des Menschen Herz nicht mäßigen noch enthalten vor Stolz, der Muth wächst uns. Und wiederum, wenn es in Nöthen ist, daß [es] uns übel geht, so will es bald verzagen und verzweifeln, steckt und malt GOtt ein Ziel, Weise und Maß, wie er ihm helfen soll, und will nicht der Zeit, Person noch Stunde erwarten, die ihm gefallen, noch es ihm gänzlich heimstellen. Das heißt GOtt versuchen.

110. Also ein schlüpfrig Ding ist des Menschen Herz: Wenn es wohl geht, so fällt es in Vermessenheit; und wiederum, wenn es übel geht, in Verzweiflung. Darum muß [es] so gemengt sein, und GOtt eine Sorge nach der andern mit einwerfen in die Güter, und sie uns salzen, und den Muth mit einer Furcht dämpfen, auf daß wir im ersten Gebot bleiben. Wenn es wohl geht, oder in der Wohlfahrt, da ist er mit der Sorge hinter uns her, daß du allezeit mußt denken: Jetzt geht es dir wohl, aber GOtt kann es wohl umkehren. Also auch, in der Noth muß man’s auch temperiren mit einer Tröstung, daß wir an GOttes Verheißung hangen, und GOtt nicht versuchen, der da will unser GOtt sein. Also bleibt ein Mensch auf der rechten Landstraße, daß er nicht verzweifelt in der Noth, sondern allezeit stehen bleibt und GOtt lobt, es gehe ihm wohl oder übel, wir gewinnen oder verlieren.

111. Aber der Mammon ist so ein großer und gewaltiger Gott in der Welt, daß, wenn er
\quellseite{889}
vorhanden ist, so macht er die Leute stolz; wenn er nicht da ist, so macht er die Leute verzagt. Darum sucht Moses also heftig wider ihn. Nun, das ist die Auslegung des ersten Gebots, darum soll man GOTT in dem lieben, und nicht andere Götter haben. Das ist, fürchte GOTT im Ueberfluß; wenn dir’s wohl geht, so sei nicht stolz; wenn dir’s übel geht, so vertraue ihm im Unglücke. Laufe nicht an einen andern Ort, und suche nicht andere Götter, sondern glaube, er kann dir aus Nöthen helfen, und dich in Reichthum stärken, auf daß das Herz die Mittelstraße innen behalte, und gleich stehe. Das fordert das erste Gebot.

\begin{schriftzitat}
V. 20. Wenn dich dein Sohn heute oder morgen fragen wird, und sagen: Was sind das für Zeugnisse, Gebote und Rechte?
\end{schriftzitat}

112. Dieses ist eine prosopopoeia.\footnote{prosopopoeia (Personenbedichtung) ist hier, daß den Personen das in den Mund gelegt wird, was GOTT lehrt.} Er befiehlt den Eltern, daß sie gottesfürchtig sollen sein, daß sie lernen GOTT fürchten, und ihm dienen zur glücklichen Zeit, und glauben zur unglücklichen Zeit, und einen reinen Glauben behalten; und das nicht allein für sich, sondern auch für ihre Kinder. Die Eltern sollen der Kinder Schulmeister sein, und die Kinder sollen es förter andere auch lehren; sollen die Nachkömmlinge berichten, daß es stets in frischem Gedächtniß bleibe, wie GOTT sie aus Egypten geführt habe. Sie sollen eitel Schulmeister sein; das will GOTT haben, denn er sagt, daß sie sollen zu ihren Kindern sprechen [V. 21.]: „Wir waren Knechte in Egypten“ etc. Hiermit gibt er ihnen den Befehl, daß sie die Kinder lehren GOTTes Dräuung, und sich vor der Dräuung zu fürchten; wenn es ihnen wohl geht, daß sie GOTT nicht verachten, und zur Zeit der Anfechtung lernen auf GOTT trauen, und seine Verheißung ergreifen, und nicht verzweifeln. Denn er versucht sie, und will sie nicht verlassen. Das heißt recht die zehn Gebote gepredigt. Moses ist ein rechter Meister und Ausleger der zehn Gebote.

\begin{schriftzitat}
V. 25. Und es wird unsere Gerechtigkeit sein vor dem HErrn, so wir halten und thun alle diese Gebote, wie er uns geboten hat.
\end{schriftzitat}

113. Siehe, wie hoch erhebt und rühmt er das erste Gebot; denn er sagt: Wer das hält, der wird vor GOTT gerecht geachtet. Er schreibt der Erfüllung dieses Gebots die Gerechtigkeit zu. Wir lehren also, daß kein Mensch auf Erden durch die Werke gerecht werde, wie sonst die ganze Welt pflegt zu predigen wider die Lehre des Evangelii; darum gehöre allein dazu der Glaube aufs Wort und Verheißung GOTTes [Röm. 3, 24–28. Gal. 2, 16.]. Daher haben wir es im Catechismo gefaßt, daß der Gottesdienst stehe in der Furcht und Liebe GOttes. Ein jeglich Gebot ist gefaßt in die Furcht und Liebe, daß wir sollen GOTT fürchten und lieben. Das ist das erste Gebot, daraus die andern Gebote alle herfließen. Denn die Werke der andern Tafel, als, daß wir nicht ehebrechen, nicht tödten, nicht stehlen, machen uns nicht fromm, können auch von uns nicht geschehen ohne den Glauben; denn wo das erste Gebot nicht ist, da sind die andern Gebote alle nicht, und was von Leuten gethan wird, scheint allein vor der Welt.

114. Aber das erste (denn es geht durch alle Gebote), das lehrt, wie wir gerecht, fromm und heilig werden. Wie denn? Also, wenn du das erste Gebot hältst durch den Glauben und herzliches Vertrauen auf Christum. Darum erfordert dies Gebot allein den Glauben. Dieser Glaube rechtfertigt dich; wie von Abraham 1 Mos. 15, 6. geschrieben steht. Also fordert das erste Gebot das Hauptstück, als, den Glauben; wer das hält, der ist selig. Aber niemand hält es, es werde denn ihm durch den Heiligen Geist gegeben, daß er zum Evangelio gehöre. Darum wird keiner selig, sondern Christus muß hier zuvor sein, daß man an ihn glaube, durch den Glauben wird man gerecht [Röm. 10, 4.]. Darum sagt Moses: Welcher gerecht ist, der hat es daher, daß er es hält und thut. Der aber hält es, welcher von Herzen GOTT vertraut und glaubt, wenn es übel geht; und von Herzen GOTT fürchtet, wenn es wohl geht.

115. Welcher GOTT fürchtet und trauet, der wird keinen totschlagen, niemand Unrecht thun, sondern er befiehlt es GOTT, was ihm Unbilliges widerfährt, und leidet es geduldig um GOTTes willen, und gibt GOTT die Rache heim, der wird es wohl machen. Also wird das fünfte Gebot gehalten, denn es fließt aus dem ersten; denn ein Gläubiger rächt sich darum nicht, daß er sich auf GOTT verläßt und traut ihm, er werde es recht treffen. Also ist das erste Gebot das Mark und der Kern aller andern Gebote. Darum sagt er: Wer den HErrn fürchtet und dient ihm,
\quellseite{890}
der wird Gutes thun, er wird halten und thun diese Gebote. „Halten“ geht aufs erste Gebot, daß man GOtt traue und ihn fürchte. Diesem Halten folgt auch [das] „Thun“, das ist, die andern Gebote werden hernach gehen; das ist, wer GOtt fürchtet, der thut Gutes; also, welcher GOtt trauet, der hat Geduld, er werde ihm helfen.

116. Also ist der nun gerecht, wenn er das hält und thut, das ist, glaubt und beweist den Glauben mit den Werken [Gal. 5, 6. Jac. 2, 14. 17. 22. 24.]. Moses zeigt nur an, wie man leben und die Gebote halten soll; aber es gehört ein anderer Mann dazu, daß wir sie halten, und im Unglück trauen können. Denn ein Mensch, als ein Mensch, muß verzagen im Unglück, und sich erheben im Glück. So sage ich nun, Moses lehre wohl, und heiße, was man thun und halten solle; aber wo man es holen und nehmen soll, lehrt allein das Evangelium, nämlich, daß man an Christum glaube; so wird GOtt Gnade geben, daß man könne den Muth dämpfen, wenn es uns wohl geht, und GOtt glauben möge im Kreuz und in der Anfechtung.

\section{Das siebente Capitel}

\begin{schriftzitat}
V. 1–4. Wenn dich der HErr, dein GOtt, ins Land bringet, darein du kommen wirst, dasselbe einzunehmen, und ausrotten viel Völker vor dir her, die Hethiter, Girgositer, Amoriter, Cananiter, Pheresiter, Hebiter und Jebusiter, sieben Völker, die größer und stärker sind, denn du; und wenn sie der HErr, dein GOtt, vor dir gibt, daß du sie schlägst,\footnote{„daß du sie schlägst“ fehlt im Original und in der Erlanger.} so sollst du sie verbannen, daß du keinen Bund mit ihnen machest, noch ihnen Gunst erzeigest. Und sollst dich mit ihnen nicht befreunden; eure Töchter sollt ihr nicht geben ihren Söhnen, und ihre Töchter sollt ihr nicht nehmen euren Söhnen. Denn sie werden eure Söhne mir abfällig machen, daß sie andern Göttern dienen, so wird denn des HErrn Zorn ergrimmen über euch, und euch bald vertilgen.
\end{schriftzitat}

1. Jetzt geht ein ander Aergerniß an, und treibt Moses wiederum, daß sie sich vor dem Dräuen fürchten, und heißt solch Aergerniß ein Weib und böse Gesellschaft; die bringt manchen in großen Jammer und Noth; denn wie die Gesellen sind, also muß er mitheulen. Geht nun das also zu in äußerlichen, bürgerlichen Dingen, viel mehr wird es geschehen, wenn sich einer gesellt zu denen, welche das erste Gebot GOttes übertreten und dawider toben, und andere Götter suchen und Freunde machen. Denn, gesellen sie sich zu ihnen, so wenden sie sich auch von GOtt. Darum über den Mammon ist die weltliche Gunst; die ist auch ein großer Mauerbrecher und Teufel, wenn man sich gen Hof zu den Reichen und Weisen thut; alle wollen das Hochhinan kommen.

2. Darum verbeut es Moses seinem Volke gar heftig, daß sie sich hüten vor den Leuten, die fremde Götter haben, und sollen nicht mit ihnen Gemeinschaft, Gunst und Freundschaft haben, nicht ihre Töchter nehmen, oder ihre Söhne zur Ehe ihnen geben. Ursache, denn es ist gefährlich Ding um der Welt Gunst; viel wagen darum Weib, Kind, Leib, Leben, und die Seele dazu. Dieweil es denn so ein böse Ding ist, so stellt er ihnen dies andere Aergerniß vor, und vermahnt sie, daß sie sich davor hüten.

3. Wie meint ihr, daß es gangen ist? Die Israeliter haben viel frommer, weiser Leute gesehen unter den Heiden, auch viel junger, schöner Jungfrauen und Frauen, die man alle umbringen und vertilgen sollte; da haben sie aus Barmherzigkeit gedacht: Ei, wir wollen ihnen nichts thun, sondern uns mit ihnen verbinden; und haben denn zu ihnen sich verheirathet, und gesagt: Es wird darum GOtt nicht übel gefallen, er wird uns deßhalb nicht so hart strafen, oder von der Erde vertilgen und ausrotten. So wird es gangen sein; damit haben sie den Abgott in das Haus gebracht, und also ist die Abgötterei geblieben; wie im Buch Josua und der Richter zu sehen ist.

4. Das sollen wir auf uns auch ziehen, es geht also gefährlich auch zu unsern Zeiten. Denn
\quellseite{891}
siehe, wie gar viel sind ihrer, die da verleugnen das Evangelium um der Fürsten und gottlosen Bischöfe und Tyrannen willen, allein, daß sie Gnade und Gunst behalten bei den Menschen. Und ihrer viel nehmen jetzt Weiber, nicht um gotteliges Lebens willen, sondern allein um Reichthums, Gewalt und Freundschaft; und man fragt nichts darnach, ob das Weib oder der Mann Christen sind, wenn sie nur viel Geldes haben. Da muß mancher um eines reichen Weibes willen GOttes Wort fahren lassen. Also hat die ganze Welt diesen Abgott.

5. Darnach ist das noch ärger und gefährlicher, daß man sich zu den Rotten und Secten gesellt, die sich sehr demüthiglich stellen, und bei den Leuten einfließen. Das heißt mit den Cananitern Gemeinschaft halten, das ist, mit denen, die das Wort GOttes verfälschen und fahren lassen. Es ist ein gefährlich und ärgerlich Ding, und ein rechter Abgott in der Welt, also nach Freundschaft, weltlicher Gunst, Reichthum, Gewalt und Weisheit sehen und trachten; denn man wird dadurch betrogen, daß man von GOttes Wort abweicht. Also fallen ihrer viel jetzt zu unsrer Zeit vom Worte, um zeitlicher Ehre und Gunst willen, wie vorzeiten die Juden sich an die großen Fürsten und Könige unter den Heiden gehängt haben.

6. So klebt und hängt man sich auch noch an die großen\footnote{„großen“ fehlt in der Erlanger.} Hansen, und fahren also dahin. Darum spricht Moses allhier: Siehe nicht an einiges Gleisen oder Demuth, Gewalt, Ehre, Reichthum, Gunst, Weisheit oder Freundschaft, sondern allein mich siehe an. Ich will dir es alles geben, und mehr helfen, denn alle Könige, Gewaltige und Weise auf Erden. Aber wiederum, so du mich verlassen wirst, und dich zu ihnen wenden, mit ihnen halten und heucheln, so will ich dir auch mehr Unglück und Plagen schaffen denn ihnen. Fällt derhalben Noth vor, so gedenke an mein Wort; denn mein Wort ist mehr denn dies alles, Reichthum, Gunst und Gewalt der Welt.

7. Also wollte er gerne das Herz rein bei dem Worte GOttes und Glauben behalten. Das kann Fleisch und Blut nicht thun, sondern der Heilige Geist muß den Muth geben, daß einer so fest werde, daß er mehr getraue GOtt, und mehr von seinem Worte halte, denn von allem Gut oder Reichthum, Weisheit, Königen, und Rottengeistern mit ihrer falschen Lehre. Die Vernunft schließt sonst die Augen zu, und verläßt GOttes Wort, geht dahin, und wird bald abgöttisch. Aber ein Gottesfürchtiger, dem der Heilige Geist das Herz stärkt, der hält das Wort theurer und werther denn diese Schätze alle. Was hilft es, wenn du gleich alle Schätze auf Erden hättest, und wenn schon alle Fürsten und die ganze Welt deine Freunde wären, und GOtt wäre wider dich, und dein Feind [Matth. 16, 26.]? Und wiederum, was kann dir mangeln, wenn du GOtt bei dir hast [Röm. 8, 31.]? Wenn du GOttes Wort verlässest, so kann er dich also strafen, daß niemand dich erlösen kann.

8. Das sind die zweierlei Aergernisse, welche Moses erzählt bei der Auslegung des ersten Gebots; nämlich, daß uns erstlich der Mammon anficht; dann ist entweder zu viel oder zu wenig vorhanden, und daß im Ueberfluß wir GOttes vergessen, oder im Mangel GOtt versuchen. Zum andern böse Gesellschaft, Gewalt und Gunst der Welt, auch Rottengeister, die wir meiden und fliehen sollen; denn sie uns abführen vom ersten Gebot und GOttes Wort.

\subsection[Erzählung D. Martin Luthers von der Kanzel, seine Reise]{Erzählung D. Martin Luthers von der Kanzel, seine Reise gen Marburg, und das Gespräch, daselbst mit den Sacramentirern gehalten, belangend}

9. Ihr möchtet vielleicht gerne wissen, was wir zu Marburg ausgerichtet hätten; auf daß ihr nun nicht im Zweifel gelassen werdet, und ihr auch den Lästerern könntet das Maul zustopfen, so will ich dieselbige Handlung euch kürzlich erzählen.

10. Wir haben keine Gefahr erlitten auf dem Wege, und hat GOtt in dem euer Gebet erhört, dafür ihr ihm sollt danken; denn es hat uns besser gegangen, denn wir gehofft hätten. So haben sich unsere Widersacher sehr freundlich und demüthig gegen uns gestellt, mehr denn wir meinten. Denn ich dachte selber, wir würden eitel harte Steine allda finden; aber GOtt gab Gnade, daß sie sich demüthigten, und ob GOtt will, so wird nun ihr hochfährtiges Schreiben wider uns ein Ende haben.

11. De peccato originis sind mir uneinig \&c. Sonst haben sie ein Stück sechs oder sieben willig eingeräumt und eingegangen, als, von
\quellseite{892}
der Kindertaufe, von der Justification, Nutzen und Brauch des Abendmahls, und stimmen in diesen Stücken Gottes Wort mit uns überein, und sind von ihrer vorigen Lehre williglich abgetreten; und geben zu, daß man im Sacrament Glauben und Trost empfah[e?]. Aber, daß allda sei wahrhaftig und leiblich Christi Leib und Blut, das können sie noch nicht glauben. Und haben so viel vermerkt, wenn es bei ihnen stünde, so hätten sie es nachgegeben, und haben ihre Reden also gelautet: Hätte man das Bier wieder im Fasse, so würde man es nicht anzapfen. Aber dieweil sie einen gemessenen Befehl gehabt von den Ihren, so haben sie nicht zurück gekonnt.

12. Brüderschaft aber haben sie von uns begehrt, die haben wir ihnen auf diesmal abgeschlagen, und nicht zusagen können. Denn wenn wir sie für Brüder und Schwestern annähmen, so müßten wir verwilligen in ihre Lehre. Wiewohl man dieses Abschlagen nicht gerne gesehen, und vorgab, man sollte die Liebe gegen ihnen erzeigen, bis sie GOtt auch wieder herzu brächte; denn wir auch unsere Feinde lieben sollten [Matth. 5, 44.]. Nun, wer es übel auslegen will, der mag es thun; wer auch mehr ausrichten kann, der richte es aus. Sie bleiben auf ihrer Meinung, GOtt wolle sie erleuchten; wiewohl sie sich flicken und schmücken, und nachgelassen haben, daß sie nicht verleugnen, daß nicht der wahre Leib und Blut Christi da sei; welches denn lautet, gleich als hielten sie es mit uns.

13. Sie bekennen, daß, die zum Abendmahl gehen, genießen allda wahrhaftig des Leibes und Blutes Christi, aber geistlich, daß sie Christum im Herzen haben. Leiblich zu genießen, das wollen sie nicht zulassen, das haben wir auf ihr Gewissen geschoben. Denn wir haben GOttes Wort und den Text für uns, den sie nicht haben. Darum steht die Sache in einer guten Hoffnung. Ich sage nicht, daß eine brüderliche Einigkeit sei, sondern eine gütige, freundliche Eintracht, daß sie freundlich bei uns suchen, was ihnen fehlt, und wir wieder ihnen dienen. Wo ihr nun werdet fleißig bitten, wird sie auch brüderlich werden.

Ende der Narration oder Berichtes vom Colloquio zu Marburg.

14. In der Auslegung des ersten Gebots hat Moses erzählt das erste Aergerniß, als, den Mammon, wie derselbige ein sehr groß Hinderniß ist, daß man des ersten Gebots nicht achte, verstehe, lerne, thue oder halte. Denn durch den Mammon fällt die Welt dahin, und richtet sich auch wieder auf. Wenn der Mammon vorhanden ist, so ist sie sicher und stolz; wenn er nicht da ist, so verzweifelt sie und versucht Gott.

15. Das andere Aergerniß ist böse Gesellschaft, Freundschaft und Gewalt auf Erden. Die sind auch ein Abgott; denn man mehr darauf trotzt und trauet, denn auf GOtt, da wir doch GOtt alle Dinge vertrauen sollten. Darum will Moses das erste Gebot also erklärt haben, daß wir zusehen, und GOtt mehr vertrauen sollen, denn dem Mammon, Gewalt, allen Gütern und Freundschaft auf Erden. Denn Güter können nicht helfen in den letzten Zügen. Wenn die Pestilenz regieret und der Tod herschleicht, so kann eine kleine Pestilenz uns alle unsere gute Gesellschaft dahin nehmen; was wir uns denn darauf verlassen haben, ist umsonst. Warum wendet sich denn das Herz zu Geld und Gut, oder trotzt auf Menschen? Derhalben verbeut Moses, daß sie sich mit den abgöttischen Heiden nicht sollen befreunden, die um sie her waren, noch Bündniß mit ihnen machen. Und setzt nun die Ursache dazu:

\begin{schriftzitat}
V. 4. Denn sie werden eure Söhne mir abfällig machen, daß sie andern Göttern dienen. So wird des HErrn Zorn ergrimmen über euch, und euch bald vertilgen.
\end{schriftzitat}

16. Er dräuet immer mit den Worten des ersten Gebots, und zieht sie mit ein, die also heißen: „Ich bin ein eifriger GOtt, der die Sünde der Väter heimsucht an den Kindern, bis ins dritte und vierte Glied.“ Das ist, wie er hier sagt: Es hilft kein Mammon, Freund, Gewalt wider GOttes Zorn; wenn der ergrimmet, so reißt er es mit Wurzel und Stamme aus der Erde. Also wiederholt er stets die Worte des ersten Gebots, flicht sie immer mit herein, und behält sie bei der Hand.

\begin{schriftzitat}
V. 5. Sondern also sollt ihr thun mit den Heiden: Ihre Altäre sollt ihr zerbrechen etc., ihre Säulen zerbrechen, ihre Haine abhauen, und ihre Götzen mit Feuer verbrennen.
\end{schriftzitat}

17. Sie sollen nicht Freundschaft mit diesen Heiden machen; und das ist die Ursache: Denn sie werden euch gewiß zu Abgöttischen machen;
\quellseite{893}
wie denn auch geschehen ist. Weil sie das Gebot nicht haben gehalten, und Mosi nicht gefolgt, da haben sie gedacht: Es sind gute, fromme, feine Leute, sie sind auch weise; wie, sollen wir allein klug sein? Darum haben sie ihre Töchter gefreiet, und ihre Götter darnach auch angenommen und ihnen gedient. Diese Gesellschaft zu den Heiden hat auch bei ihnen eine Veränderung der Religion gebracht, damit die Propheten so viel schreien und predigen, und damit große Arbeit haben gehabt.

18. Warum befiehlt er, man soll es gar auswurzeln, zerbrechen und zerpillern, wie im Text steht, daß nicht etwas übrig bleibe, das ihnen Occasion und Ursache gäbe, zur Abgötterei zu fallen; so soll man die Abgötterei zerstören. Denn so lange die Abgötter bleiben stehen, so bleibt immerdar das Aergerniß und ein schönes Bild vor den Augen; denn es ist kein fährlicher noch schädlicher Aergerniß, denn Aberglaube oder ein falscher Gottesdienst. Es gleist kein Gut noch Schatz auf Erden so schön, als ein falscher Gottesdienst; es ist der höchste Schmuck. Geld und schöne junge Leute gleißen und scheinen nicht also sehr, als ein falscher Gottesdienst.

19. Und wo man sich nicht fleißig hält an GOttes Wort und an das erste Gebot, so kann sich kein niemand erwehren. Die Ursache habt ihr oft gehört. Denn ein jeder Mensch ist also gesinnt, daß er lieber hält Menschenlehre und Gesetz, wie schwer und groß sie auch sind, denn GOttes Gebot. Wie viel Fasten und Plagens haben wir gelitten unter dem Papstthum? so viel Geldes gegeben zu den Stiften, Klöstern, Altären, Seelmessen, Vigilien. Das ist uns nicht sauer angekommen noch beschwerlich gewesen. Denn man hat es nicht Abgötterei genannt, sondern es hat einen Namen und Schein gehabt, daß es ein Gottesdienst sei. Es ist der Vernunft nichts schwerer, denn auf GOtt trauen; und wiederum, nichts leichter, denn Superstition und Aberglauben erwählen.

20. Zum andern, daß ein jeder von Natur lieber auf das trauet und bauet, was er selber thut, als, auf unsere Werke; wir sind geneigt, darauf zu fußen; so haben wir gesagt: Also haben wir gelebt, dies und jenes gethan, das sieht Gott an etc. Wenn wir diese zwei Stücke daran hängen, als, daß es erstlich eine heilige Gestalt und schöne Farbe hat, oder schön gleist; darnach, daß der Mensch für sich geneigt ist, daß er gerne etwas thut, darauf er sich verläßt, so ist es bald geschehen, daß wir in Abgötterei fallen und gerathen.

21. Darum ist es sehr schwer, daß wir uns nicht ärgern an GOttes Wort, und dabei bleiben mögen, oder nicht davon abfallen [Matth. 11, 6. Luc. 2, 34.]. Denn das ist die höchste Ketzerey in der Welt. Und was haben jetzt unsere Schwärmer, die Papisten, Wiedertäufer und Bilderstürmer für einen andern Schein und Grund ihrer Lehre, denn die Werke und das falsche Vertrauen? Denn die Werke oder Frömmigkeit und das falsche Vertrauen stets mit eingemischt wird, daß sie meinen, die Heiligkeit und der christliche Glaube stehe in Bilderstürmern; oder, wenn sonst dies oder jenes Werk geschieht; das Vertrauen auf unser Werk schleicht immer mit unter.

22. Diesen Gebrechen versteht keiner; also tief ist es uns eingewurzelt, und also klebt es uns an, es komme denn der Tod oder Sterben. Weil der Mensch lebt und gesund ist, so versteht er es nicht, es kommt ihm nicht in den Sinn; allein in der Todesstunde, wenn der Jammer und Elend herzubringt, und ihn verzweifeln macht, der lehrt es ihn, daß er sage: Ach, HErr GOtt, hilf mir! und daß er sich an das erste Gebot halte. Darum ist GOtt ein GOtt derer, die in höchster Noth und Jammer stecken; sonst ruft man ihn nicht an, man stecke denn in Noth oder im Tode [Ps. 32, 7.]. Wie der Schächer am Kreuz, der da sprach zum andern Schächer: Wir haben unsere Strafen wohl verdient [Luc. 23, 41.], und verschuldet die Pein, darinnen wir stecken; und fiel in das erste Gebot, hing sich daran, vertraute darauf, da er in sich nichts befand von guten Werken; er hatte eitel Sünde. Darum verläßt er sich auch auf nichts, denn auf GOttes Barmherzigkeit, und spricht [V. 42.]: „Gedenke meiner, HErr, wenn du kommst in dein Reich.“ Denn dieser hat gesagt: „Du sollst nicht andere Götter haben, Ich bin der HErr, dein GOtt“; laß mich deinen GOtt sein. Welcher also nicht in Nöthen oder Gefahr ist oder steht, der kann nicht lernen auf GOtt trauen.

23. Also sind jene Heiden feine, treffliche Leute gewesen, und viel frömmer, denn wir unter dem Papstthum; wiewohl wir daselbst auch aufrichteten viel Altäre, stifteten mit großen Unkosten viel Kirchen und Gottesdienste, ordneten
\quellseite{894}
auch viel Ceremonien, und haben auch Poltergeister gehabt, und sie gefragt, und gethan, was die gewollt haben. Die alten Väter sind viel heiliger gewesen, denn sie ihre Kinder, als Söhne und Töchter, geopfert, wie das andere Vieh, und verbrannten sie auf den Altären als ein köstlich, fein Opfer; darum sind es viel feinere Leute gewesen, denn unter dem Papstthum.

24. Woran fehlt es denn, daß sie nicht gottselig gewesen sind, weil sie es allein in GOttes Namen gethan haben, der Himmel und Erde geschaffen hat? Und es hat sollen ein Gottesdienst sein, daß sie es GOtt zu Ehren thäten, und GOtt sollte es ihnen belohnen; denn sie wollten dem rechten, lebendigen GOtt die Kinder opfern. Warum verwirft sie denn GOtt, und sagt, es sei wider das erste Gebot? Wo das erste Gebot nicht regiert, da kann kein Gottesdienst angerichtet werden, und da vermag kein Vertrauen auf GOtt zu sein; sondern da kriecht bald her ein falsch Vertrauen auf das Werk, das du selber thust; wenn das Vertrauen da ist, dann ist es falsch. Denn also gedachten die Heiden: Wenn ich meinen Sohn opfere und verbrenne, so werde ich gewiß Gnade erlangen bei GOtt und selig werden. Also hängt ein solch Herz und Vertrauen an dem Werke, wenn ich meine Kinder opfere und schlachte, GOtt zu Lob und Ehren; das ist, ich verlasse mich auf dies Werk. Solches ist stracks wider das erste Gebot.

25. Nun ist es mein Werk, das ich thue, da verlasse ich [mich] auch auf. Denn, der also gesinnt ist, der darf unseres Herrn Gottes nicht, sondern will so lange wirken, bis so lange er Gnade mit Werken erlange. Erlangt man es nun mit Werken, so macht man sich selber zum Gott, daß man GOttes Gnade nicht bedarf, und unsere Seligkeit kommt dann nicht aus Gnaden. Aber GOttes Gebot sagt, daß GOtt nicht wolle barmherzig sein darum, daß du es von GOtt willst erkaufen; denn er ist dir nicht pflichtig, du wirst ihn nicht dahin bringen. Auch hörst du in diesen Worten, daß er den Juden gebietet, sie sollen solche Götzen weghun, und wird ihnen gebräuet, daß sie sonst sollten vertilgt werden. Wie denn die Heiden sind darob verdammt, und ganz und gar ausgerottet.

26. Also, unter dem Papstthum ist keiner gewesen, wenn er vor seinem Absterben in seinem Testament eine Messe gestiftet hat, der eine andere Meinung gehabt, denn daß er dadurch wollte selig werden und GOttes Gnade mit diesem Werke verdienen; es wäre denn, daß ihm Gott sonst wunderbarerlicher Weise aus der Abgötterei geholfen. Nun ist es nicht so wichtig, wie der vorigen Heiden ihres. Dieselben Heiden haben etwas Größeres daran gewagt, haben es mit einem größern Ernst gethan, als, mit ihrem eigenen Blute, ja, mit ihren leiblichen Kindern, als, ihrem Fleisch und Blut. Unsere Vorältern haben es nur mit Gelde kaufen wollen, etwa mit drei, vier Gulden etc.

27. Dieses ist alles beides wider das erste Gebot gehandelt, wenn man einen andern Gottesdienst anrichtet, und gute Werke thut, darauf die Leute vertrauen; und solcher Gottesdienst heißt denn nicht, „Ich bin der HErr, dein GOtt“, sondern es ist ein solcher Gottesdienst, da der Mensch gedenkt an sein gut Werk, und verläßt sich drauf. Solch Vertrauen ist eitel Gift und Heuchelei, ja, die Verdammniß selbst; denn das Herz vertrauet GOtt nicht, sondern diesem oder jenem gethanen Werke.

28. So ist nun die Summa des ersten Gebots: Vertraue mir allein, und fürchte mich; was hinwider ist, das sei verdammt. Darum, wer sein Kind, Sohn oder Tochter, schlachtet oder verbrennt, der vertrauet nicht ehe GOtt, er habe denn das Kind geopfert; er vertrauet dem Werke mehr denn GOtt. Wahrlich, solch Werk hat scheinbarlich geflissen vor der Vernunft, und ein herrliches Ansehen großer Heiligkeit und Gottesdienstes gehabt. Die Werke sind zu hoch gewesen, und ist dazu ein groß Gespränge geschlagen; die Vernunft hat sich ihrer nicht können erwehren; niemand hat können drüber springen; wie denn die Juden sich auch sehr daran geärgert haben. Lieber, was thaten wir auch, da zu unsrer Zeit der Mönche Werke im Papstthum so schön glänzten, und für so groß gerechnet worden, daß wir uns nicht haben ihrer können erwehren? Wir haben uns daran müssen ärgern, haben nicht können darüber springen.

29. Darum bläuet ihnen Mosis das erste Gebot so fleißig ein, und will, daß man allein auf GOttes Gnade traue, und nirgend anderswo Trost suche. Das Herz soll allein auf GOttes Gnade und Barmherzigkeit stehen und sich verlassen, und auf kein Werk bauen. Ist aber anderswo ein Gottesdienst, Werk, Gesetz und Verdienst, so traue nicht drauf. Derhalben sage ich,
\quellseite{895}
daß diese Anfechtung der Juden nicht ein geringes Aergerniß ist gewesen, weil sie gesehen haben, daß die Heiden so einen gewaltigen Gottesdienst getrieben; wie du siehst im Papstthum bei den Mönchen und Pfaffen einen gleichenden Gottesdienst.

30. Weil denn nun das Aergerniß so groß ist, und der Heiden, Juden und Papisten, und aller Werklehrer Gottesdienste stracks wider das erste Gebot zustürmen und schießen,\footnote{Erlanger: schließen.} darum sagt Moses: Hüte dich, daß du dich mit ihnen nicht befreundest, noch mit ihnen Gemeinschaft habest, sonst werden sie dich verführen; sondern also sollst du mit ihnen thun, ihre Altäre sollst du zerreißen. Man soll nichts mit ihnen zu schaffen haben, auf daß man nicht durch sie verführt werde; darnach soll man es alles zerstören, zerreißen, damit sie Gottesdienst treiben. Die Altäre sind gewesen, wie jetzt unsere Kirchen und Capellen. Darnach haben sie gehabt Säulen, darauf sie gestanden, und aufgerichtete Bilder. Sonderlich sind sie im Gehölze gemein gewesen, wie zu unsrer Zeit zur Eiche, und an andern einsamen Orten mehr, da der Satan hat sein Narrenspiel gehabt. Darum sind auch Haine kleiner Gehölze, als unser Speckhölzlein, da lustige Wieselein angelegen; als, der Poltersberg ist ein rechter Hain, da man dahin kommt zur Kirchmeß, säuft und frißt, und alle Büberei treibt. Darnach hat ein jeder daheim ein Capellchen gehabt, seinen eigenen Ort, da er hat wollen GOtt anbeten, und seiner eigenen Gottesdienste abwarten. Das sollen die Juden alles ausrotten, daß keine Gelegenheit oder Occasion sei oder überbleibe zur Abgötterei; auch soll man alle Oerter und Städte der Abgötterei meiden, und allein dem rechten GOtt dienen.

\begin{schriftzitat}
V. 6. Denn du bist ein heilig Volk, welches GOtt aus allen Völkern erwählet hat.
\end{schriftzitat}

31. Allhier haben die Bilderstürmer einen Text für sich, und ist die Frage daraus: ob wir alle Bilder wegthun, und aus der Welt stoßen sollen, dieweil uns Mose ein heiliges Volk nennt. Item: ob uns dies Gesetz Mosis dazu zwinge? Höre, das erste Gebot lehrt dich, daß du auf kein Ding auf Erden trauen sollst, denn allein auf GOtt; darauf geht es. Darum, welche Bilder aufgerichtet sind, oder dazu gebraucht werden, daß man darauf baue und einen Gottesdienst anrichten will, die reiße weg, denn das erste Gebot wirft sie zuvor hernieder, als, wenn man meint, GOtt helfe um des Bildes willen, wie zur Eiche. Da ist ein kleines Marienbild auf ein Papier gemalt gewesen, darauf ist das Vertrauen gerichtet, Maria hülfe in dem Bilde, und haben in dem Bilde Mariam, und nicht GOtt, angerufen. Das heißt ein abgöttisches Bild. Dergleichen Marienbild ist auch gewesen im GrimmenThal; item, zu Regensburg.

32. Aber die andern Bilder, da man allein sich drinne erfehet vergangener Geschichte und Sachen halben, als in einem Spiegel, das sind Spiegelbilder, die verwerfen wir nicht, denn es sind nicht Bilder des Aberglaubens; sonst dürften wir auch kein Bild auf der Münze haben, und es dürfte eine Jungfrau auch keinen Spiegel haben, darinnen man deß Gestalt und Angesicht schauet, der hinein gucket. Die Ursachenbilder betet man auch nicht an, man setzt sein Vertrauen darauf, sondern es sind Merkbilder.

33. Darum siehe es fleißig an, und unterscheide es. Bilder stürmen heißt nicht, allerlei Bilder abwerfen, sondern diese, wie der Text mitbringt, die angebetet werden als Götzen. Götzenbilder heißen es, da das Herz anhängt, wie die Birne am Baum hängt, dergleichen bei den Heiden viel gewesen, und noch auf diesen Tag im Papstthum sind, die nicht allein angesehen, sondern mit Vertrauen angesehen werden, da ein solch Bild wird für einen Gott gehalten, und Ketzer solche Bilder haben angebetet, da das Vertrauen anklebt, als die zur Eiche und zum Birnbaum gewesen, Sanct Maria; item, St. Niclas’, St. Wolfgangs in Bayern, St. Leonhards; item, des heiligen Kreuzes Bild. Solche mag man wegreißen, doch durch ordentliche Obrigkeit; denn solche Bilder werden nicht allein angesehen, sondern eine Andacht, Vertrauen und Gottesdienst wird daraus gemacht.

34. Aber, daß ich habe ein gemalt Bild an die Wand, daß ich schlecht ansehe ohne Aberglauben, ist mir nicht verboten, sollen auch nicht weggethan werden. Denn warum wollte ich nicht ein gemalt Tuch ansehen, sehe ich doch wohl eine Sau, Kuh oder Hund an? Die Kinderlein müssen hölzerne Pferdlein, Hühnlein, Docken und dergleichen haben; und ich muß auch ein Bild haben, sonst müßte ich auch keinen Spiegel haben. Wer sie nicht will haben, der
\quellseite{896}
mag sie wegthun, ich kann ihrer nicht entrathen, weil doch davon in der Schrift nichts gedacht wird.

35. Da aber ein Bild wird aufgerichtet oder vorgestellt, darauf man ein Vertrauen setzt, das reiß entzwei. Das erste Gebot sei eine Glosse, und gebe einen rechten Verstand den Bildern. Wenn ein Bild aufgerichtet wird, da man sich vor fürchtet, und einen Glauben drauf setzt, das reiße man hinweg; so es aber nicht ein Götze ist, oder Altar, daß man die Knie davor beugt, auch nicht einen Gottesdienst daraus macht, so ist es nicht ein Götze, sondern ein Bild, das du behaltest, und ist recht und gut. Das ist der Unterschied zwischen den Bildern und Götzen.

\begin{schriftzitat}
V. 6–8. Denn du bist ein heilig Volk GOtt, deinem HErrn. Dich hat GOtt, dein HErr, erwählet zum Volke des Eigenthums aus allen Völkern, die auf Erden sind. Nicht hat euch der HErr angenommen und euch erwählet, daß euer mehr wäre, denn aller Völker; denn du bist das wenigste unter allen Völkern; sondern daß er euch geliebet hat.
\end{schriftzitat}

36. Das ist abermals ein Abgott. Als sollte Moses sagen: Wiewohl du ein so heilig, hoch und stark Volk bist, so solltest du wohl eine Abgötterei daraus anrichten, als müßte dich GOtt ansehen um deiner Macht und Stärke willen, und dich zu seinem Volk erwählen. Aber GOtt will allein angesehen sein, und will nicht, daß du setzest auf den Mammon, Freundschaft, Macht, Gesellschaft, noch auf die Menge und Kraft, noch darauf vertrauest; denn das Vertrauen schlägt GOtt allhier zu Boden, und reißt den Abgott herab.

37. Ich wollte, daß man diesen Text merkte, sonderlich wir Deutschen, die wir auch pochen und trotzen auf unsere Macht und Gewalt, wenn man soll wider den Türken ziehen. Darum reißt er hier den Götzen hinweg. Denn das geschieht von Natur, daß, wenn einer gewaltig ist, so ist er ein Pocher, Scharrer und Schnarrer; da ist keine Demuth, kein Vertrauen auf GOtt und seine Barmherzigkeit, sondern allein auf seine Macht und Gewalt pocht und trotzt er, gleich als wären jenseit des Wassers und jenseit des Berges nicht auch viel Leute und viel Fäuste. Wenn solche Pochhansen etwas anfahen, thun sie es nicht im Vertrauen auf GOtt, sondern um ihrer Stärke und Menge willen, gleich als wäre kein GOtt. Sie fragen nichts nach GOtt, sondern nur nach ihrer Gewalt. Also mächtig bin ich, so viel Volks habe ich.

38. Aber siehe du nichts an, das da stehen soll auf deiner Menge und Kraft, gründe dich nicht drauf [1 Sam. 17, 47.]. Nicht daß [es] böse sei, viel Leute haben, ein gewaltiges Fürstenthum haben, Harnisch, Schwert, Büchsen, Geld und Gut haben. Es ist alles gut, Freundschaft, Gewalt und Gehülfen haben, denn es sind GOttes Creaturen; aber daß du dir einen Götzen daraus machest, und ein Vertrauen drauf setzest, und GOtt nicht mehr trauest denn dem Mammon, deinem Vermögen, der Menge und großen Gewalt, das soll nicht sein. Dies ist alles gut; aber hüte dich davor, daß du darauf nicht vertrauest. Diese Dinge sollen nicht der Grund sein unserer Zuversicht, da ich mein Vertrauen auf setze. Also reißt er alles um, darauf ein Mensch und sein Herz trauet, und dem ersten Gebot widerstrebt.

39. Ich bin nicht so gar alt, aber ich habe ihrer viel erlebt, die sich auf ihre Gewalt verließen, und darauf stolzirten, die jetzt im Dreck liegen. Ursache ist die, sie wollten das erste Gebot umstürzen, und gedachten: Wir haben so viel Städte, Schlösser, so viel tausend Dörfer etc., rühmten sich ihres Kriegsvolks, Büchsen, Bündniß und Gesellschaft höher, denn GOttes; darum kann es uns nicht fehlen, wir werden Glück und Sieg haben. Ja wohl! Wofür sitzt GOtt im Himmel? Sie sind drüber zu Scheitern gegangen, oder heim gekommen, haben das Grauen im Nacken mit sich gebracht, und danieder gelegen, sind dennoch nicht klug worden.

40. Jetzt, weil man mit dem Türken kriegt, so jagt er uns auch in das erste Gebot; noch werden etliche sagen: Wir haben es gethan, und ihrer Stärke und Kraft zuschreiben, was GOtt allein durch solche Mittel ausgerichtet. Darum, wenn diese Dinge vorhanden sind, so ist Fleisch und Blut stolz darauf. Ich habe auch dagegen gesehen etliche, so fett und stolz gewesen sind, wenn ihnen die Menge anhing, daß sie Kraft, Zufall und Beistand hatten, etwa ein Glücklein und Sieglein erlangten; aber da sie wieder von ihnen abfielen, sind sie so verdrießliche, schändliche und verzagte, blöde Bösemüther worden, daß es eine Schande ist, daß sie auch Holz hätten auf sich lassen entzwei hauen und spalten. So sieht man, wer ihr Gott gewesen
\quellseite{897}
ist, nämlich, ihre Gewalt, Anhang, Fürstenthum, Büchsen und Schwerter, und nicht GOtt, auf den sie hätten hoffen sollen. Da muß man ihnen denn singen das Liedlein, für solche Scharmützel gehörig Ps. 60, 13.: Vana salus hominis, „Menschen Hülfe ist kein nütze“; item: Superbia praecedit casum, et exaltatio ruinam, „wer zu Grunde gehen soll, der wird zuvor stolz“; ante ruinam exaltatur cor, hoffärtiger und stolzer Muth kommt vor dem Fall, Sprüchw. 16, 18. Item Ps. 116, 11.: Omnis homo mendax, es ist eitel nichtig Ding um menschliche Hülfe.

41. Also nimmt es die Schrift aus dem ersten Gebot; denn deselbigen Gebots Summa ist, allein GOtt vertrauen, und allen andern Creaturen Urlaub geben. Denn, wer sich auf was Anderes verläßt denn auf GOttes Barmherzigkeit, der kann nicht bestehen, er muß fallen, GOtt gebe, er habe zu Hülfe und Freunde den tüchtigsten Kaiser, oder die ganze Welt mit aller ihrer Macht und Gewalt, dennoch kann er es nicht hinaus führen, es wird ihm fehlen; er wird vor dem Siege wohl stolz sein, aber es wird keinen Bestand haben. Damit will Moses die Herzen rein behalten im einigen Vertrauen auf GOtt, daß niemand etwas wisse oder kenne denn GOtt, und sich allein zu ihm halte. Wer etwas Anderes ansieht, der führt es nicht hinaus.

42. Aber das rührt unser Herz nicht ehe, bis daß die höchste Noth und Anfechtung kommt. Wenn einer ihm dann nicht das erste Gebot vorhält, und bedenkt es, da verzweifelt denn Fleisch und Blut, wo GOtt nicht hilft. Aber es hilft bei der Welt nicht, es geht ihr nicht zu Herzen, bis daß sie zu Boden gestürzt werden; so müssen sie denn bekennen, daß dem also sei, wenn sie von aller Kraft und allen guten Freunden verlassen sind, und niemand haben, der ihnen helfe. Wiederum, wer da gedenkt an das Wort: „Ich bin der HErr“, der mag bestehen, der kann sich wieder aufrichten.

43. Darum sagt Mose: GOtt liebt euch, nicht darum, daß euer viel ist. Ja, hütet euch davor, er hat nicht eure Menge angesehen, weil euer so wenig ist. Wie gefällt dir das? Darum bauet nicht darauf. Also, wenn du sollst wider den Türken streiten, sage: Ich baue nicht auf Büchsen, und daß ich habe ein großes Heer, oder viel tausend Mann und Kriegsvolk; denn du, lieber GOtt, hast es bald geschlagen. Ich habe wohl die Kriegsrüstung bei einander, aber ich verlasse mich nicht darauf, sondern in deinem Namen und Vertrauen will ich streiten, denn du bist unser GOtt; und du sprichst, man soll dir allein vertrauen, und darum will ich Büchsen, Rüstung und anderes nur in deinem Namen gebrauchen etc. [Ps. 20, 6. Ps. 10, 14.] Also sollte man hinan gehen.

44. Was ist nun denn die Ursache, darum er euch lieb hat? Darum geschieht es nicht um deiner Gewalt willen, sondern, daß er so gnädig ist, und sich also gegen dir erböten hat, und daß er den Bund hält, den er euren Vätern geredet hat, daß er sich euch erbeut aus lauter Barmherzigkeit und Gnaden, und spricht:

\begin{schriftzitat}
V. 9. So sollst du wissen, daß der HErr, dein GOtt, ist ein treuer GOtt, der den Bund und Barmherzigkeit hält denen, so ihn lieben und seine Gebote halten.
\end{schriftzitat}

45. Darum, was dir widerfährt, das geschieht alles aus Liebe und Gunst; das ist, siehe hinter dich, gedenke an meine vorige Hülfe, und daher vertraue mir allein etc., darum, daß ich mich deiner habe angenommen. Siehe anders nichts an, denn daß ich dich liebe aus lauter Güte und Barmherzigkeit, und danke mir dafür, daß ich so gnädig bin, und erbitte mich gegen dir, daß ich dein GOtt sein will; denn das thue ich aus lauter Gnaden und Gunst.

46. Also reißt GOtt alle Abgötterei um, und tilgt sie aus unserm Herzen. Und ist solches ein feiner Commentarius des ersten Gebots. Erstlich reißt er das falsche Vertrauen und Zuversicht und Superstition um, die in unserm Herzen steckt, und will, daß wir ihm vertrauen von ganzem Herzen, denn er sagt, er sei unser GOtt. Dadurch will er ausrotten den Mammon, Freundschaft, Gewalt und Menge. Denn wenn du diese drei annimmst, so hast du einen falschen und fremden Gott, und kannst auch nicht bestehen, sondern du mußt zu Grunde gehen. Und er wird ihrer noch mehr erzählen, die gleichen Schaden mit den vorigen anrichten; als da ist unsere eigene Gerechtigkeit, welche auch von dem ersten Gebot uns abzieht.

47. Bisher hat Moses gehandelt die drei Feinde, das ist, Abgötter, welche wider das erste Gebot streiten und fechten, als, Mammon oder Reichthum; Gewalt, eigene Macht oder Menge;\footnote{Im Original und in der Erlanger: „meniger Macht“; menige = Menge.}
\quellseite{898}
Macht und Freundschaft der Heiden. Diese drei Dinge oder Götzen machen, daß gar ein großer Haufe der Welt abfällt vom ersten Gebote, und mehr die Menschen darauf vertrauen, denn auf GOtt. Es ist Gunst und Liebe, Reichthum, Freundschaft, eigene Kräfte und Gewalt der Welt Abgott, wie die Erfahrung zeugt. Dawider hat er gesagt: Wenn du bekommst Haus, Hof und anderes, daß du in Ueberfluß, Reichthum und Nahrung sitzest und alles genug hast, so siehe zu, und begib GOtt nicht. Item: Befreunde dich nicht mit den Heiden, und bete ihre Götter nicht an; hüte dich vor böser Gesellschaft, und denke, daß dich GOtt nicht hat um deiner Stärke und Menge willen angenommen etc., sondern aus milder Güte und Barmherzigkeit.

48. Mit welchem allein Mose also das Volk im Worte GOttes, und im ersten Gebot erhalten will, daß sie GOtt anhangen in allen Stücken, und auf seine Creatur sonst trauen, sondern, nach Laut des ersten Gebots, ihre herzliche Zuversicht setzen allein auf GOtt, sein Wort und Verheißung, und alle anderen Creaturen fahren lassen. Welcher aber auf den Mammon trauet und bauet, der wird mit ihm zu Boden gehen; welcher auch trauet auf Freundschaft und Gewalt oder Menge, item, auf seine eigene Stärke und Kraft, der wird mit denselben fallen, und seinen Hals, ja, vielmehr die Seele abstürzen; sintemal derselbigen Abgötter keiner ist, der\footnote{Original: daß.} da könne dem Menschen helfen aus Nöthen.

49. Ich habe aber gesagt, daß Moses viel Worte mit zubringt, und wollte die Leute gerne überreden, daß sie GOttes Worte traueten, und ihre Hoffnung und Zuversicht auf GOttes Zusagung setzten, und hat das erste Gebot gar fleißig getrieben. Denn wir sehen und erfahren es auch selbst, wie es in der Welt zugeht, wie seltsam der Glaube ist; man läßt es wohl predigen, aber es schafft wenig Frucht und Nutzen. Wer Geld und Gut hat, der ist muthig; wer große Freunde hat, ist trotzig; item, der da hat Gewalt und Stärke, der ist kühn. Und wiederum, welcher es nicht hat, sondern schwach, arm und elend ist, der ist blöde und verzagt, kleinlaut und betrübt. Man läßt predigen, singen und sagen, was man will, und bleibt also, wie es ist; doch muß man es sagen um des kleinen Häufleins willen. Die andern glauben es nicht, sondern sie wollen es erfahren, und werden es auch mit ihrem großen Schaden erfahren; als, zur Zeit des Unglücks und Widerwärtigkeit, auch wenn Sterben hinter sie kommt. Da werden diese drei nicht helfen, sie müssen dennoch den Mammon, Geld, Gut, Freundschaft, Macht und Gewalt hinter sich lassen. Ihr aber, welche das Wort lieben, lernet es in der Furcht GOttes, und lasset es euch gesagt sein [V. 7.]: „Nicht hat euch der HErr angenommen und euch erwählet, daß euer mehr wären, denn aller Völker; denn du bist das wenigste unter allen Völkern.“

50. Als wollte er sagen: Wenn er das hätte angesehen, so hätte er ein großer Volk erwählet; du aber bist das geringste Volk unter allen Völkern gewesen; [V. 8.] „sondern, daß er euch geliebet hat, und daß er seinen Eid hielte, den er euren Vätern geschworen hat“ etc. Einen solchen GOtt hast du, der nicht darnach fragt, wie stark, wie groß, reich und gewaltig jemand sei, sondern daß man sein Wort und seine Gnade ansehe, und nichts Anderes. Er will nicht feiern, sondern gefeiert sein. Er fragt nach der Dinge keines. Es gilt nichts bei ihm, denn seine Gnade und Barmherzigkeit, welche er dir verheißt [Ps. 130, 7.]. Er achtet nicht dein Geld oder Freundschaft: er will es nicht ansehen. Es soll bei dir kein Gut, Kraft, Macht noch Vermögen gelten, sondern über alles und allein seine Gnade, wie er es denn allhier aus der Maßen fein ausstreicht mit einem herrlichen Text, der lautet also:

\begin{schriftzitat}
V. 9. 10. Du sollst wissen, daß der HErr, dein GOtt, ein GOtt ist, ein treuer GOtt, der den Bund und Barmherzigkeit hält denen, die ihn lieben und seine Gebote halten, in tausend Glied, und vergilt denen, die ihn hassen, vor seinem Angesicht, daß er sie umbringe, und säumet sich nicht, daß er denen vergelte, vor seinem Angesicht, die ihn hassen.
\end{schriftzitat}

51. Das sind die Worte, die im ersten Gebote stehen. Gnade und Barmherzigkeit hält er, er leuget nicht. Wahrlich, so hast du einen GOtt, der vergelten wird. Er wiederholt diese Verheißung noch einmal, auf daß er es ja wohl ins Herz bilde: „Es ist ein treuer GOtt.“ Und ist überaus nöthig, daß er dies Wort, getreu, dabei setzt. Als wollte er sagen: Daß er verheißt Gnade und Barmherzigkeit, Hülfe und Trost, das glaube; denn es ist treulich geredt,
\quellseite{899}
und soll auch treulich gehalten werden. Denke derhalben, und halte dich an sein Wort, halte fest darob, so wird seine Gnade und Barmherzigkeit treulich kommen [Ps. 111, 8.].

52. Und wiederum, wird er auch nicht fehlen der Ungläubigen und Ungehorsamen, sondern sie ernstlich strafen; wie er denn seinen Eifer und Zorn ihnen allhier auch vorhält. In Summa, er wollte gern das erste Gebot seinen Juden, und sonst allen Menschen, tief eindrücken; darum wiederholt er so oft die Worte des ersten Gebots. Als sollte er sagen: Thue dich nur zu ihm, so hast du einen treuen GOTT; wie ihn auch abmahlt der 86. Psalm, V. 15.: Tu, Domine Deus, miserator et misericors, patiens, et multae misericordiae, et verax. Er hat verheißen Gnade, Hülfe, Wohlthat und Trost allenthalben, und eben in dem Psalm; das hält er auch treulich. Wiederum, wird er mit der Strafe sich nicht säumen gegen den Gottlosen, so ihm nicht glauben, sondern ihn hassen und seine Gebote übertreten. Daran habe ja niemand Zweifel, es wird beides unsäumlich kommen, wie er gesagt hat: Gottes den Gläubigen, harte Strafe den Ungehorsamen.

53. Das treibt Moses mit großem Ernst, so haben auch alle Propheten aus diesem Texte viel Dräuung GOttes genommen; als, daß GOtt im ersten Gebot den Gläubigen Gnade, Barmherzigkeit, Rath, Hülfe, Trost und alles Gutes zusage, verheiße, und getreulich auch halten wolle; und wiederum, ohne Verzug die Gottlosen zu strafen und heimzusuchen gedenke. Aber die Welt glaubt es nicht; sitzt sie im Ueberfluß, in Gütern, und hat ihren Mammon, ihre Freundschaft und Gewalt, so ist es ihr ein Spott, dann achtet man es nicht, sondern sagt: Ei, laß Mosen predigen, so lange er will. Aber höre du, wahrlich, es wird kommen zur Zeit, da man sichs nicht versieht.

54. So lange der Mammon vorhanden ist, Freundschaft, Glück etc. neben anderm, geht man dahin, und Mose schreit sich zu Tode über diesen Worten: „Ich bin ein eifriger GOTT“, „ein verzehrend Feuer“; item, ich scherze nicht. Aber es ist alles ein eitler und todter Buchstabe den Ungläubigen; man fühlt es nicht, so glaubt man es auch nicht. Die Worte sind uns kälter denn Eis, lauten nicht in den Ohren, bis die Erfahrung und die Noth zur Hand kommt, und das Unglück trifft; dann wird man sein erst recht gewahr. Als jetzt mit dem Türken, der vor wenig Tagen\footnote{Die Türken belagerten Wien den 13. September 1529; dies ist gegen Ende October gepredigt.} Wien belagert etc. Indes, wiewohl GOTT schreiet: „Ich bin ein eifriger GOTT“, ich will es, und kann es nicht leiden, das hilft alles nicht. O es gilt kein Predigen bei dem Pöbel. Aber Narren muß man denn mit Kolben laufen. Wer nicht will Worte hören, dem müssen Spieße und Büchsen rathen.

55. So will Moses die Seinen und uns davor gewarnt haben, und sagen: Gedenket, daß ihr unseres HErrn GOTTes Gebote, Sitten und Rechte haltet; verachtet ihr es aber, und macht einen Scherz daraus, so macht er aus euch wiederum einen Spott, gleichwie mit jenen. Wiewohl sie hofften und gedachten, die Noth und Gefahr wäre noch wohl über tausend Jahr von dannen, doch war sie in einer Stunde da. Also geschah es auch mit den aufrührischen Bauern; die fürchteten sich auch vor keinem Unglück, aber in einer Stunde waren sie alle erschlagen. So haben sich die Oesterreicher des Türken auch nicht versehen, sondern gedachten, er wäre tausend Meilen von ihnen; aber in zweien Tagen hatten sie ihn über dem Halse, und fiel das Unglück über sie, daß ihrer viel erschlagen und gefangen weggeführt worden; die Strafe kam ihnen gar zu bei Zeiten. Darum vermahnt Moses uns also fleißig, daß wir diesen dreien Dingen nicht trauen, sondern daß man sich davor zum fleißigsten hüten und vorsehen solle; denn das Unglück kommt ungesäumt.

56. Deß könnte ich viel Exempel sagen, aber ich habe die zwei neulich gesehen: als, wie die Bauern in dem Aufruhr sind geschlagen und aufgerieben und wie jetzt die Türken grausamlich tyrannisirt haben. Vorzeiten sind andere Exempel gewesen: als, Saul, Sodoma und Gomorra, die Sintfluth der ganzen Welt; denn ehe sie sich umsahen, da kam die Strafe. Das will die Welt nicht glauben, sondern meint, GOTT eile nicht also sehr mit der Strafe; bis so lange, daß sie es erfährt und da liegt in der Asche, oder ersäuft und erstochen ist. Darum, wenn nicht zu rathen ist, denn steht auch nicht zu helfen. Die Welt will mit Kolben gelaufet sein.

57. Und allhier siehst du, wie Moses also treulich und ernstlich warnt, daß man sich soll vor GOTT fürchten; denn er diese Worte dazu setzt: Er wird vergelten ungesäumt. Als sollte
\quellseite{900}
er sagen: Sein Zorn wird plötzlich kommen [Ps. 2, 12.], und strafen die Ungläubigen, die jetzt getrost scharren und pochen, daß, sollten sie unsinnig werden, so wird der Zorn GOttes bald hinter ihnen zuschlagen. Was sie dann werden daran gewinnen, da mögen sie die Schuhe mit schmieren. Dieser Text gehört für die Gottlosen, und ist ihre Lection; aber sie kehren sich nicht daran, sie sind verstockt und verblendet, daß sie des ersten Stücks, welches ihnen, den Gottlosen, gesagt ist, nicht wahrnehmen.

58. Das andere Stück: „Er ist ein treuer GOtt, der den Bund und Barmherzigkeit hält.“ Dieses schreibe beides fleißig an, spricht Moses, nämlich, daß GOtt zornig und grausam ist gegen denen, so an ihn nicht glauben, welchen gehört Zorn und Plage. Wiederum, erzeigt er auch Gnade und Barmherzigkeit gegen den Gläubigen. Solches soll nun beiderseits Ernst und die Wahrheit sein, und es gilt uns, die wir müssen mit der Welt einen Sturm ausstehen, oder das Bad mit ihr ausgießen. Da soll zu solcher gefährlichen und besorglichen Zeit dieses unser Trost sein, daß GOtt ein treuer und wahrhaftiger GOtt sei in der Noth, welcher verheißen hat Gnade und Barmherzigkeit, Hülfe und Rath allen Gläubigen, und werde solches auch gewiß halten.

59. Dieses ist ja so noth auszustreichen, wenn Krieg, Theurung, Pestilenz und anderer Unfall kommen, als das Vorige, wie GOtt die Bösen gewißlich strafen wolle. Denn, wenn die Noth und Trübsal daher geht, so sehen wir unsere Schwachheit, und ist dazu ein groß Zagen und Furcht vorhanden; da werden wir gewahr, wie arme, dürftige Leute wir sind. Bald folgt das Hin- und Wiederlaufen, und hätten gerne Hülfe; wollen denn dieses und jenes stiften; hoffen auf den Mammon; rufen alle guten Freunde und Bekannten an. Aber damit ist nirgends kein Bestand, allhier geht denn alle Abgötterei unter. Man wird so weich, blöde und verzagt, daß uns niemand aufrichten noch trösten kann; gleichwie dagegen die Gottlosen so störrig, stöckisch, unbeweglich und hart sind, daß sie niemand kann weich machen, noch in ein Schrecken und Furcht bringen, bis daß sie gar zu Boden liegen; wie denn wiederum die Weichen, Schwachen und Erschrockenen niemand genug kann hart und stark machen.

60. Darum muß man also thun: daß, welche störrig, blöde, Stein, Stöcke und Sauköpfe sind, die weder GOtt noch den Teufel fürchten, die muß man erschrecken mit ernstlichen Träuen. Wiederum, blöde, erschrockene und verzagte Leute soll man mit diesen lieblichen Trostworten aufrichten und stärken, daß man also zu ihnen sage: Du hast einen getreuen GOtt, der dir Gnade und Barmherzigkeit verheißen hat, und schwört dir noch dazu, er wolle dein GOtt sein. Denn also spricht er im Propheten Ezechiel [Cap. 18, 23. 32. 33, 11.]: „Als wahr ich GOtt bin und lebe, so will ich nicht den Tod des Sünders, sondern daß er sich bekehre und lebe.“

61. Derhalben so scherzt er mit dieser Verheißung nicht; es ist solche Zusagung ihm nicht ein gering und leicht Wort, sondern ein lauterer Ernst, und es soll den Gläubigen gehalten werden; er setzet ehe dran Himmel und Erde, ja, es müßte eher Himmel und Erde vergehen, denn daß sein Wort nicht sollte erfüllt werden [Luc. 21, 33.]. So setze nun deine Zuversicht hierauf, und sei gewiß, GOtt sei ein getreuer GOtt, der nicht lügt, scherzt noch schimpft, sondern der den Bund und Barmherzigkeit hält. „Er beweiset Barmherzigkeit (gegen denen, die seine Gebote halten) ins tausendste Glied.“ Das wirst du gewißlich erfahren.

62. Also thun wir aber: Wenn wir in Armuth stecken, so sehen wir nirgend hin, denn nur wo Geld, oder nicht Geld ist. Wenn Schmach und Schande vorhanden ist, bald sind unsere Augen dahin gerichtet, wo Ehre, oder keine ist. In Todesnöthen, Pestilenzzeiten und in Kriegsläuften sieht man, wo Gesundheit, und wo Friede, oder kein Friede ist. Dahin stehen alle unsere Gedanken. Wir können unsere Herzen nicht erheben, daß sie diese Worte allhier ansehen möchten. Aus dieser Ursache sagt Moses: Wenn du willst GOtt trauen, und nicht fremde Götter haben, sondern das erste Gebot halten, und du kommst darüber in Noth, Jammer, oder in Gefahr Leibes und Lebens, so gedenke daran, daß du nicht sehest, wo menschliche Hülfe und Rath, wo Geld und Gut, wo Freunde oder Feinde, Leben oder Tod bleibe, sondern hieher siehe auf GOttes Wort, höre, was GOtt dir verheißen habe. Denn, wenn schon etliche Noth, Tod, Krieg, Krankheit und Theurung vorhanden ist, dennoch solltest du diese Worte dir nicht nehmen lassen, die allhier gesagt werden: „Ich bin der HErr, ich will meinen Bund treulich halten.“

\quellseite{901}
63. Aus diesem Texte sind gar viel Psalmen gemacht, die alle auf diesen Spruch Mosis sehen, da David auch rühmt, GOtt sei getreu, gerecht und ein Vergelter [Ps. 31, 24.], item, ein Helfer in der Noth; als, im 145. Psalm, V. 18.: „Der HErr ist nahe bei alle denen, die ihn anrufen, die ihn in der Wahrheit anrufen.“ Item [Ps. 34, 20.]: „Der Gerechte muß viel leiden, aber der HErr hilft ihm aus dem allem.“ Item, der 44. Psalm sagt, V. 23.: „Wir sind geachtet als Schlachtschafe“; wie solches St. Paulus zum Römern [Cap. 8, 36.] auch anzeucht, und darauf spricht [V. 37.]: „Aber in dem allen überwinden wir weit, um des willen, der uns geliebet hat.“ Die Psalmen klagen immerdar unsere Noth und Jammer, aber allezeit ist die Zusagung dabei gesetzt: GOtt hilft und rettet, daß also die Psalmen eine rechte Auslegung sind des ersten Gebots. Denn der Bund ist fest, daß GOtt getreu ist [2 Tim. 2, 19.], er will halten, was er zugesagt hat. Wiederum sagen die Psalmen auch, wie GOtt die Ungerechten richten wird, und sie zu Schanden machen. Als, im 91. Psalm, V. 8., wird gesagt: „Du wirst sehen, wie es dem Gottlosen vergelten wird.“ Gott kann wohl diese Kunst, daß, was die Gottlosen zusammen tragen und sammeln, das es soll ein anderer auffressen; ja, um Leib und Gut, Leute und Lande sollen sie kommen, denn sie sündigen wider das erste Gebot. Davor hat uns bisher Mosis gar treulich gewarnt.

64. Das ist nun das Hauptstück der christlichen Lehre, daß wir lernen auf GOtt trauen, ihn nicht verachten, und darauf stolzieren, wenn die drei Dinge vorhanden sind, als, Gut, Freundschaft und Gewalt, daß man sein Vertrauen auf die Creaturen setze, sicher und verwegen des Gelds und Guts halben sei, gleich als dürfte man unseres Herrn Gottes nicht. Wiederum, daß man auch nicht verzage, wenn diese drei nicht da sind.

65. Darum thue man nicht, wie etliche Schelmen unter den Bürgern und Bauern thun. Wenn die Geld und Gut haben, so sind sie muthig und stolz, werden darüber Mammonisten und abgöttische Leute, verleugnen Gott, stolzieren wider ihren GOtt, und treten das erste Gebot mit Füßen; sagen: Hätte ich nur Geldes genug, was wollte ich klagen? Hinter diese will GOtt kommen, ehe sie sichs versehen, und sie unterdrücken, und will GOtt sich da nicht säumen, sie sollen Leib und Gut, Seelen und alles verlieren. Dagegen die Frommen, die da erschrecken vor GOtt, sollen nicht verzweifeln in ihrer Trübsal. Als, ich und meines Gleichen haben nicht viel guter Tage, darum zu uns wird gesagt: es solle nicht Noth haben, GOtt sei getreu, und halte seine Barmherzigkeit und Bund. Wie denn die Verheißung lautet: „Ich will Barmherzigkeit erzeigen“ etc. Ei, wenn wir den Worten nur glauben könnten, und gewiß möchten sein, daß es GOttes Wort wäre, wie sollte unser Herz und Gewissen doch darüber so stolz werden; wie würde ich wider alle Creaturen so muthig sein; da könnte ich sagen: Dieweil ich einen gnädigen GOtt habe, der mir Gnade zusagt, wer ist reicher, denn ich? Denn, wem der etwas zusagt, dem darf er es halten.

66. Wenn du nun solches glaubtest, daß dir GOtt so geneigt sei, daß er dich nicht verlassen wolle in Armuth, Schande und Tod; item, alle Creaturen dir günstig und geneigt sind um seinetwillen, so wollte ich gerne sehen, welcher mir wollte etwas Böses thun! Denn ich könnte den Mammon, alle Freundschaft, Tyrannen, den Türken, oder den König von Frankreich ansehen, gleich als eine Laus oder Fliege. Niemand ist reicher denn er; niemand frömmer, niemand wüßte mehr denn er; denn es ist alles nichts geredet, gegen der Gnade und Barmherzigkeit, die ich an GOtt habe. Ei der Teufel! wer wäre stärker denn ich? Diese Stärke von der Gnade und Barmherzigkeit GOttes sollte man also erheben, daß man auch den Teufel und Türken nichts dagegen achtete.

67. Aber man hört diese Worte allein, und hält sie, gleich als hätte sie nur irgend ein Pfarrherr, Capellan oder Prediger geredet; man glaubt es nicht, man hält es nicht für GOttes, des Allmächtigen, Wort. Aber welcher es glaubt, der hat es alles [Marc. 9, 23.], und hat es in der Faust. Es soll uns aber auch sehr schrecken, daß GOtt daneben sagt: Trauest du mir nicht, so wisse, daß ich zürne, und diese Worte oder diese Träumung sollen dich also treiben, daß dir Himmel und Erde soll zu enge werden, und immerdar bedenken müssest, der Himmel werde einfallen, und das Meer werde dich ersäufen, und die Erde mit dir untergehen. Aber wenn man dies gleich lange sagt, daß GOtt zürnt, so fragt man nichts darnach. Wenn irgend ein Fürst zürnte, oder daß man sagte, der Fürste
\quellseite{902}
käme, da wäre große Furcht; aber GOtt gilt bei dieser Art gar nichts, der da spricht: er zürne mit uns etc. Also sehen wir in der Erfahrung und an diesem täglichen Leben, wie gar arme Leute wir sind, die wir uns nicht fürchten vor dem Worte GOttes, trauen auch nicht darauf. Weil wir denn nicht glauben dem Worte, so müssen wir es mit der Erfahrung lernen, wie diesmal durch den Türken, daß er seinen Zorn über uns ausschüttet, auf daß wir sehen, wie es GOtt kein Scherz sei.

68. Moses wollte gerne die Worte also backen, brauen und kochen, daß sie uns doch eingingen. Denn wenn wir es glaubten, so gedachten wir, so oft wir diese Worte hörten: O der Himmel hängt gar voll Feigen! Diese Worte wären unser süßester Trost, vor allem Mammon und Freundschaft und Gewalt der Welt, denn wir wüßten, daß Himmel und Erde unser wäre. Wiederum, wenn wir GOtt nicht vertraueten, und dann hörten, daß GOtt spricht: Ich will vergelten und ungesäumt strafen, würde man davor erschrecken und sagen: O! zürnt GOtt, so ist er mehr zu fürchten denn hundert Türken.\footnote{„Türken“ steht hier für das türkische Reich oder den türkischen Kaiser.} Diese Worte sollten mir viel erschrecklicher sein, denn wenn hundert Türken oder Kaiser kämen. Ist es nun nicht Schande und Sünde, daß man soll diese drei Dinge, als, Mammon, Freundschaft und Tyrannen mehr fürchten denn GOtt? und daß sich einer der Creatur und eines Pfenniges mehr soll freuen, und den Türken und böse Zeit mehr fürchten denn GOtt? Ein rauschendes Blatt soll stärker sein denn GOtt mit allen seinen Dräuungen; und ein rother Geselle, ein Gulden, erfreuet dich mehr denn alle GOttes Verheißungen?

69. Wahrlich, wir sollten uns selber fluchen, uns anspeien und verdammen, daß wir sollen so verkehrte Herzen in uns haben, die GOtt nicht vertrauen noch fürchten. Es möchte uns wohl gereuen, daß wir leben, dieweil wir also unserm HErrn GOtt ungehorsam sind, und ein rauschendes Blatt mehr, denn ihn fürchten, und einen rothen Gulden mehr lieben. Solche Untugend steckt in unserm Herzen. Dieses wollte Moses gerne den Leuten einbilden (wie ich es denn auch gerne thun wollte), und spricht: GOtt läßt es wahrlich an Gnade und Barmherzigkeit nicht fehlen; er will Güter genug haben, die er gibt denen, die ihn fürchten. Dagegen so will er auch nicht scherzen, und gar genug Plagen und Unglück verschaffen, so er dir zuschicken will, wenn du ihm nicht gehorchst und deines Eigensinnes lebst.

\begin{quote}
2. 12. 13. Und wenn ihr diese Rechte höret, und haltet sie, und darnach thut; so wird der HErr, dein GOtt, auch halten den Bund und Barmherzigkeit, die er deinen Vätern geschworen hat, und wird dich lieben, und segnen, und mehren, und wird die Frucht deines Leibes segnen, und die Frucht deines Landes, dein Getreide, Most und Oel, die Frucht deiner Kühe, und die Frucht deiner Schafe auf dem Lande etc.
\end{quote}

70. Jetzt wird er sich selber erklären, was da sei „Gnade und Barmherzigkeit GOttes“; item, was da sei „der Zorn GOttes“. Denn er hat gesagt, GOtt sei ein getreuer GOtt, welcher Gnade und Barmherzigkeit hält, und dir Gutes thut. Aber was thut GOtt, wenn er Gnade und Barmherzigkeit erzeiget? Darauf antwortet Moses: „Er wird dich lieben, und segnen, und mehren“ etc. Das wird der Segen, die Gnade oder Barmherzigkeit sein, daß er dir wird Weib und Kind geben; und dieses alles wird er auch dir bewahren vor dem Satan, Pestilenz, Wetter, und vor allem Unglück und Uebel, daß dein Geschlecht lange währe ins tausendste Glied. Er will alles in seine Hut nehmen.

71. „Er wird die Frucht deines Leibes segnen“, das ist, deine Kinder, daß die gesund seien, und du einen guten Muth habest, deine Kinder ohne Gebrechen. Darnach soll auch die Frucht des Landes gesegnet sein, als, dein Getreide, dein Most und Oel etc., und folgt darauf: „Gesegnet wirst du sein unter aller Völkern.“ Da siehst du, was da heiße Gnade und Barmherzigkeit, daß er will sagen: Wirst du GOtt vertrauen, so sorge nicht, wie du ernähret werdest, wie dein Haus und Hof, Weib und Kind erhalten werde, da laß mich dafür sorgen; vertraue nur mir, und laß mir die Sorge. Wie denn auch der 55. Psalm, V. 23., sagt: „Wirf dein Anliegen auf den HErrn, er wird’s wohl machen.“ Item, St. Petrus in seiner Epistel spricht [1 Petr. 5, 7.]: „Alle eure Sorge werfet auf ihn, denn er sorget für euch.“ Es haben diese Sprüche David und St. Petrus aus diesen Texte genommen.

\quellseite{903}
72. Und GOtt will allhier durch Mosen sagen: Ich habe noch so viel Engel, und so viel Gewalt, daß ich dich segnen könne, und bewahren vor dem Teufel und vor allem Unglück, auch aus aller Trübsal und Noth helfen. Wir wollten aber gerne einen solchen Gott haben, welcher uns gar kein Leid ließe widerfahren. Da wisse du, es liegt nichts daran, wenn dich die Welt gleich haßt und verfolgt; er will doch dich lieben. Will dich die Welt geringern, so will er dich mehren. Er will ja so stark sein (wo er anders nicht stärker ist), als dein Feind sein mag, und will dich wohl erretten.

\begin{schriftzitat}
V. 14. Gesegnet wirst du sein über alle Völker.
\end{schriftzitat}

73. Du sollst ein selig und gediegen Volk sein. Denn die andern Völker trauen nicht auf mich; und darum, daß du mir trauest, so sollst du auch genug haben, und kein Böses soll unter dir sein, und sollst mein reichlich genießen.

74. Ist nun diesem also, wie kommt es aber, daß GOtt gleichwohl oft seinen Heiligen so viel Plagen und Seuchen zuschickt, und Johannem den Täufer des schändlichsten Todes sterben läßt [Matth. 14, 10.]? Stephan, denn Hiob schickt er viel Leidens zu [Hiob 1, 12. ff.]? Darauf antworte ich: Das schadet alles nicht; läßt er sie leiden, so sind sie darum unverlassen; er läßt sie nicht stecken; er läßt sie in Tod kommen, und bringt sie wieder zum Leben; er steckt sie in Schande, und erhebt sie wieder zu Ehren. Dagegen wenn die Gottlosen allhier in Ehren leben, darnach sterben sollen, so ist alle ihre Freude, Ehre und Herrlichkeit aus [Hiob 21, 17. ff.], denn sie haben nichts mehr, denn nur dies Leben, und sie sollen dennoch dieses Lebens nicht so viel gebrauchen, daß sie so fröhlich sollten sein, als ein Christ ist. Der Christen Trotz und Trost ist, daß sie allhier in diesem Leben sollen unverlassen sein. Die Christen sollen im Mangel und Noth ja so wohl leben, als die Gottlosen in dem höchsten Ueberfluß, denn sie können ihrer Armuth und Güterlein mit besserem Gewissen gebrauchen. Ob ein Christ gleich nicht so viel hat als ein Reicher, so liegt nichts daran, denn Ueberfluß ist uns nicht vonnöthen [Sprichw. 15, 16.]. Trauest du aber GOtt, so wirst du alles genug haben, und wirst nicht Hungers sterben.

75. Und laß es gleich sein, daß ein Gottloser mehr habe, denn du, so mag es ihm doch nicht so wohl zugute kommen, als dir. Darum ist es so viel gesagt: Was du bedarfst, will ich dir alles geben. Wo du aber nicht GOtt fürchtest, so wird er dich lassen Geld, Häuser, Ehre und Gut suchen und sammeln, und lassen gewaltig werden; aber zuletzt wird es alles wieder gehen, wie es hergekommen. Die Kinder werden dir sterben und plötzlich ausgerottet werden; es wird ein anderer dein erworben Gut verzehren, dein Haus, Acker, Wiesen und Weib an sich bringen, wenn du wirst mit Schanden, Schaden und bösem Gewissen davon sein. Darum wäre es besser, daß du mir trauest; dann würdest du gesegnet sein in dieser Welt, und dort im andern und ewigen Leben.

\begin{schriftzitat}
V. 14–16. Gesegnet wirst du sein über alle Völker.
\end{schriftzitat}

Es wird niemand unter dir unfruchtbar sein, noch unter deinem Viehe. Der HErr wird von dir thun alle Krankheit, und wird keine böse Seuche der Egypter dir auflegen, die du erfahren hast, und wird sie allen deinen Hassern auflegen. Du wirst alle Völker fressen, die der HErr, dein GOtt, dir geben wird.

76. Moses will haben, daß ein jeglicher das erste Gebot GOttes wohl lerne halten, denn es soll kein Mensch sein, der sein Herz, Zuversicht oder Vertrauen setze auf etwas Anderes, denn allein auf GOtt und sein Wort. Und das Wörtlein „Herz“ wird darum hieher gesetzt, auf daß man GOtt fürchte, und vor seinem Wort sich entsetze oder erschrecke. Das Herz soll an GOtt mit Glauben und Furcht allein hangen; darum hat er auch verheißen, er wolle dein GOtt sein, Gnade und Barmherzigkeit erzeigen alle denen, die ihn lieben und seine Gebote halten. Damit will er unsere Herzen reizen und locken, ihm zu glauben.

77. Er ist aber nicht allein barmherzig, sondern auch eifrig, und ein verzehrendes Feuer gegen den Verächtern seiner göttlichen Gebote [Sir. 16, 12. 13.], auf daß ja reine Furcht in dem Herzen sei und bleibe gegen GOtt, und wir ihm unterworfen seien mit aller Furcht und Vertrauen, und sonst nichts Anderes vorziehen. Nun hat aber Moses, wie droben [Cap. 6, § 48 ff.] gemeldet ist, unter die Abgötterei auch diese gerechnet, nämlich den Mammon, der die Herzen der Menschen auch gefangen nimmt, daß sie das erste Gebot verlassen, und nichts darnach fragen, daß GOtt spricht: Ich will dein GOtt
\quellseite{904}
sein, und dir Gnade und Barmherzigkeit erzeigen. Darnach hat er noch von einem andern Abgott geredet [Cap. 7, § 36 ff.], der da genannt wird: auf Stärke, Kraft, Gewalt und Herrschaft sich verlassen. Den dritten Abgott [Cap. 7, § 15 ff.] hat er geheißen: große Freundschaft. Diese drei Abgötter sind in der Welt, und die Menschen stolziren mehr darauf, denn daß sie sich auf GOtt verlassen sollten.

78. Das ist nun Mosis Glosse, die er über das erste Gebot macht, daß da sei das erste Gebot übertreten, wenn der Mensch sein Herz jetzt auf den leidigen Mammon, auf Macht und Gewalt, oder auf Freundschaft, und nicht hangen bleibt an den Verheißungen des göttlichen Worts, daß uns GOtt will ernähren, schützen und erhalten. Diese Sünde achtet nun die Welt gering, straft sie auch nicht, oder kann sie nicht strafen. Sie richtet allein Mord, Ehebruch, Hurerei und Diebstahl, aber der heillose Aberglaube bleibt ungestraft vor der Welt. Und weil die Sünde nicht gestraft wird, wie andere, so scheint es nicht eine große Sünde sein; sie bleibt zwar klein im Ansehen, und wird ungeachtet, aber die Sünden, so ernstlich und hart gestraft werden, dieselbigen werden auch mehr geachtet und viel eher vermieden.

79. Das Gegenspiel aber hören wir allhier, daß es viel eine ärgere Sünde sei, bauen und trotzen auf den Mammon, denn der Diebstahl; es reißt auch tiefer ein. Mord und Diebstahl schneißt und frißt nicht also inn sich. Der Schuldige, als der Mörder und Todtschläger, wird allein gestraft, seine Freunde und Geschlecht gehen frei und ledig aus; aber um der Abgötterei willen pflegt GOtt oft ein ganz Geschlecht, eine Stadt, Land und Leute, ein Königreich und Kaiserthum auszureuten. Derhalben ein Christ halte das für die höchsten Sünden, die man in der Welt für die geringsten achtet. Die Welt hält das für ehrliche Freude, daß sie sich auf Geld und Gut verläßt, und daß ein Fürst trotzt auf sein Fürstenthum. Dieses hält man für keine Sünde, sondern man setzt dieselbigen Abgöttischen obenan. Und das ist doch anders nichts, als den Teufel an GOttes Statt setzen, GOtt seine Gottheit, Autorität und Majestät nehmen wollen, und ihn lästern.

80. Derhalben so gebraucht Moses allhier viel Worte, auf daß er uns das erste Gebot ja beibringe, denn es will solch Gebot sonst nicht ins Herz hinein gehen, und es ist auch ein gefährliches Gebot. Denn wer Geld hat, der ist stolz; wer da viel guter Freunde hat, der verläßt sich auch darauf. Wer aber arm ist, der wird kleinmüthig, und hilft bei ihm kein Trost, der ihn aufrichtete. Also ist die Welt gar voller Abgötterei. Darum muß GOtt Noth halben dazu thun, und den Türken oder Pestilenz, theure Zeit und andere Strafen schicken, auf daß man ihn lerne fürchten. Wie er denn um dieser Sünde willen den Lucifer mit seinen Engeln auch gestraft, und vom Himmel gestürzt hat, daß er wider das erste Gebot sündigte, und wider GOtt und seine Majestät sich auflehnte, und mehr sein wollte, denn ihm gebührte.

81. So soll man nun zu unserm HErrn GOtt alles Gutes sich versehen, und unsere Herzen nicht anderswo mit der Zuversicht und Vertrauen hin richten. Wider GOtt soll man auch nicht stolziren, gleich als hätten wir einen bessern Behelf, und dürften dieses GOttes nirgend zu; wie denn dieses eigentlich aller Gottesverächter Gedancken sind. Auch soll man nicht verzagen oder verzweifeln, wo es übel und sonderlich zugeht, aus [der] Ursache,\footnote{Im Original: aus Ursachen.} denn GOtt selber heißt dich, nur kühnlich deines Herzens Vertrauen auf seine Verheißung setzen. Darum ist es dir, wahrlich, bedenklich, daß du einen andern solltest aufwerfen für GOtt, oder sonst irgend Hilfe suchen; denn er will dir genug geben [Job. 10, 11.], und allen Unfall von dir nehmen, das ist, er will dein gnädiger GOtt sein. Widerfährt dir ein Schaden, so will er ihn wenden, und dir Erstattung und alles nach Nothdurft verschaffen.

82. Es sind fürwahr eitel herrliche Verheißungen in diesem Texte, welche für die gehören, so GOttes Gebot halten. Denn was sie nicht haben, das will ihnen GOtt geben; und was er ihnen gegeben hat, das will er ihnen schützen und erhalten. Und aus diesem Text, wie ich droben [§ 63] auch angezogen, sind viel trefflicher Psalmen gemacht, welche David, wie ein Bienlein den edlen Saft, daraus gesogen. So haben die Propheten ohne Zweifel so viel schöner Predigten daraus auch genommen, und ihre Zuhörer aus diesen lieblichen Quellen gelabt und getränket.

\begin{schriftzitat}
V. 17–19. Wirst du aber in deinem Herzen sagen: Dieses Volk ist mehr, denn ich bin, wie
\end{schriftzitat}

\quellseite{905}
kann ich sie vertreiben? so fürchte dich nicht vor ihnen. Gedenke, was der HErr, dein GOtt, Pharao und allen Egyptern gethan hat, durch große Versuchung, die du mit Augen gesehen hast, und durch Zeichen und Wunder, durch eine mächtige Hand, und ausgestreckten Arm, damit dich der HErr, dein GOtt, ausführete. Also wird der HErr, dein GOtt, allen Völkern thun, vor denen du dich fürchtest.\footnote{Original und Erlanger: denen, die dich fürchten.}

83. Dies ist eine Objection, da irgend ein kleinmüthiger, verzagter Mann möchte in diese Anfechtung fallen, oder ihm sein Herz durch den Teufel versucht und blöde gemacht werden, daß er sagte: Dieses Volks ist viel mehr, denn daß wir geringen, wenigen und schwachen Leute ihm sollten können widerstehen, vor ihnen bleiben, oder sie schlagen und aufreiben; es wird unmöglich sein, daß wir also sehr sollten von GOtt gestärkt werden, denn unser ist wenig, jener ist eine große Menge und Haufe; und möchten darüber in eine andere schwere Sünde wider das erste Gebot fallen, die da heißt diffidentia, daß man an GOttes Hülfe verzagt, Hände und Füße gleiten läßt, GOtt mißvertrauet, und an seiner Hülfe verzweifelt.

84. Welches Mißtrauen daher kommt, daß wir schwach sind, und nicht sehen oder fühlen unsere äußerliche Stärke, Macht und Kraft; da besorgen wir, es sei denn gar aus mit uns, GOtt werde uns nicht helfen können. In diesem Spital lagen die Kinder von Israel hier krank, und verkündigten sich zuvor damit auch sehr, da sie Kundschaft ins gelobte Land schickten, und des Landes Gelegenheit, Festungen, starke Leute und Riesen erfordern ließen [4 Mos. 13, 19. ff. 14, 1. ff.]. Dadurch sie also sehr erschreckt worden, daß sie nicht wollten fortziehen; meinten, sie wären den Riesen zu schwach, sie würden sie nicht überwinden können, GOtt würde sie von dieser Heiden Hand und Gewalt nicht erretten mögen.

85. Dieser Objection und Einrede Fleisches und Bluts, oder Versuchung des Teufels, begegnet hiermit Moses, und lehrt, wie GOttes Verheißungen erfüllt werden, nicht durch unsere Kräfte, sondern durch seine göttliche Macht. Will damit zu Boden schlagen die Vermessenheit und das Vertrauen auf eigene Kräfte und menschliche Stärke und Macht, unserm Mißtrauen und Verzweiflung steuern und wehren. Denn wenn die Kinder von Israel auf sich und ihre Gewalt oder Vermögen hätten sehen sollen, so wären sie gegen den Heiden, als, Cananitern, Jebusitern, Pheresitern, Philistern, Ammonitern, und andern gewaltigen, mächtigen Ländern und Völkern, gleich als nichts gewesen, welche doch die Kinder von Israel alle vertilgen sollten, und hätten GOttes Gebot müssen lassen zu Boden fallen, und sich an diese Heiden nicht treiben, sondern sie zufrieden lassen.

86. Aber Moses weist sie auf GOttes Gegenwärtigkeit, welcher treulich zugesagt hatte, daß er wolle bei ihnen sein, und ihnen helfen wider alle diese Heiden und ihre große Macht. Wie denn GOtt desshalben auch hier den Kindern von Israel eine tröstliche Zusagung und Verheißung thut. Und daß man solches desto mehr glaube, so führt Moses die Kinder von Israel auf die vorigen Wunderthaten GOttes, wie er ihnen Hülfe hat verheißen, und dieselbige auch gewißlich ihnen widerfahren lassen. Darum er denn das merkliche Exempel göttlicher Hülfe, Gnade und Barmherzigkeit ihnen vor die Augen vorstellt, und sie daran erinnert, wie er den König Pharao und alle Kriegsmacht der Egypter ins rothe Meer gestürzt hatte, und die Kinder von Israel aus ihrem Zwange und Tyrannei frei gemacht [2 Mos. 14, 7. ff.], welches allein ein göttliches Werk war, und sonst unmöglich gewesen, daß sie ihr Lebenlang des Pharao und der Egypter von ihrem Halse wären los worden. Mit diesem Exempel der vorigen göttlichen Hülfe sollten sie sich trösten; denn der GOtt, der damals geholfen hätte, der könnte nochmals ihnen helfen.

87. Also rüstet und stärkt sich auch David: Da der Philister Goliath an ihn setzte, und er durch GOttes Hülfe zuvor einen Bären und Löwen zerrissen hatte, erinnert er sich auch derselben Wohlthat und Hülfe GOttes, tröstet sich damit und sprach: Der GOtt, der mir zuvor vom Bären und Löwen, den grausamen Thieren, geholfen hat, der wird mich auch von diesem unbeschnittenen Philister erretten [1 Sam. 17, 37.]. Also ermahnt sich auch St. Paulus zum Corinth ern, daß der GOtt, so die Todten lebendig macht, ihn in Asia, da er todtkrank war gewesen, wieder gesund hätte gemacht durch die Fürbitte der Christen, und GOtt werde ihn noch ferner erlösen [2 Cor. 1, 9–11.].

88. Denn man soll der vorigen Mirakel und Hülfe GOttes nicht vergessen, dieweil es sind
\quellseite{906}
Memoriale, Erinnerung und Denkzettel, wie GOtt Handreichung uns gethan, und noch helfen könne und wolle alle denen, die bei ihm Hülfe suchen und im Glauben der Hülfe erwarten. Darum spricht auch David im Psalter [Ps. 119, 52.]: Memor ero mirabilium tuorum, das ist: „Ich gedenke an deine Thaten von Alters her“ etc. Und damit sollen wir unsern Glauben auch stärken, daß wir auf GOttes Wort und Verheißung, und auf die vorigen Thaten GOttes und Exempel seiner Hülfe sehen.

89. Sonderlich aber sollen wir in dieser Tentation uns GOttes Wort und Thaten vor die Augen stellen, wenn uns diese Noth auch zu Händen stößt. Denn da geschieht es gemeiniglich, daß unsere Natur und Vernunft in Anfechtungen pflegen die Contraria gegen einander zu halten [Ps. 42, 6. 43, 2.]. Solches ist der Vernunft Kunst, die wil alles messen und abwägen, daß sie bedenkt des Landes Gelegenheit, oder die Landart, die Menge des Volks. Also haben die Juden die Zahlpfennige auf die Linien gelegt, und überschlagen, wie viel der Cananiter wären, und wie eine geringe Anzahl ihrer, der Israeliter, wäre. Da haben sie also gerechnet: Dort ist viel Volks und große Gewalt, allhier aber wenig Volks und gering Vermögen. Das kann der Teufel, daß er uns ein Exemplar vor die Augen mache, und wir denn unsere Unkräfte und der Widerwärtigen Kraft messen und rechnen; und weil wir in der Waage und in dem Messen sind, da vergessen wir, daß GOtt allhier sagt: „Ich bin der HErr, dein GOtt.“

90. Also thun wir: Wenn wir in Nöthen sind, so sehen wir allein auf die Gefahr des Leibes, auf die Nahrung, Speise, Trank, auf die Krankheit, Armuth, oder was uns sonst beschwert. Da machen wir denn die Rechnung, daß das Wasser gar zu groß sei, und unsere Kräfte zu gering, und messen dieses alles gegen unsern Kräften. Wenn ich denn meine Unkräfte ansehe, und der Gegengewalt und der Kraft in vorstehender Noth inne werde, so bin ich schon verloren und gottlos worden, und frage nichts mehr nach GOtt. Also geht es zu: Wenn mich der Teufel vom Worte bringt, da stecke ich, und der Teufel hat dann gewonnen Spiel, und macht, daß ich entweder verzweifle, oder sonst irgend hinfliehe und Hülfe suche; denn ich kann da nicht bestehen. Dieweil ich denn Worte GOttes nicht glaube, so hat mich denn der Teufel gar bald übereilt, und auf sein Eis oder Schlupfriges geführt, daß ich zu Boden falle und darnieder liege. Da bleibet es wohl bei, daß man in der Verfolgung hält unser Unvermögen und Schwachheit gegen unserer Feinde Macht und Gewalt.

91. Die Weiber thun auch also in Kindesnöthen: Wenn sie in der Geburt sind, da haben sie schon ausgerechnet, daß Mutter und Kind bleiben werden; denn die Mutter sei schwach, so sind auch große Unkräfte des Leibes vorhanden. Auf das kräftige Uebel sehen sie allein, und verzweifeln dann, oder geloben das Kind irgend zu einem Heiligen, haben Zuflucht zu einem Abgott; als, zu St. Margaretha, die hat in solchen Nöthen helfen sollen. Auch haben ihrer viel ein Gelübde gethan: Bleibt der Sohn am Leben in der Geburt, so soll er eine geistliche Person werden, und das Mägdlein soll eine Nonne werden. Niemand hat allda mit herzlichem Vertrauen diese Kindesnoth unserm HErrn GOtt befohlen, oder an das erste Gebot gedacht, daß er einen gnädigen GOtt hätte, und daß er Barmherzigkeit thun wollte in tausend Glied denen, die ihn fürchten. Da sollte ich sagen: Ich vertraue mich dir ganz und gar, mit allen den Meinen, in deine Gnade, denn du hast gesagt, du wollst mein GOtt sein, und mir Gnade und Barmherzigkeit erzeigen, darauf will ich mich verlassen.

92. Derhalben so ist das nur des Teufels Kunst, daß er uns die Schwachheit und die Noth gar zu groß vor die Augen bildet, daß wir drüber GOttes Worts und Gebots gar vergessen, und letztlich verzweifeln müssen. Dahin hat er es mit uns geführet.

93. Aber wir sollten auch wiederum die Kunst lernen, daß wir sagten: Lieber GOtt, du hast mir nicht allein geboten, wenn mir es übel geht, daß ich soll dich anrufen und bei dir Hülfe suchen, sondern du hast mir auch verheißen und zugesagt, du willst mein gnädiger GOtt, mein lieber Vater, Helfer und Schutzherr sein, und bei mir in Nöthen halten, deine Kraft und Gewalt an mir beweisen. Das willst du, daß ich dir es glauben und gänzlich vertrauen soll; denn dein Gebot stehet allda, und fordert von mir den Glauben, daß ich mich dir gar befehlen soll. Und ich könnte dir kein größeres Leid thun, denn wenn ich deinen Verheißungen nicht glauben wollte, in welchen du mir zusagst, daß du wollest mein GOtt sein und mir helfen. Glaube
\quellseite{907}
ich dir nun, so wirst du mich ohne Zweifel nicht verlassen, und werde wohl in deiner Gnade bleiben, und alles Ja sein, was ich von dir bitten und bedürfen werde. Denn dies Wort und Zusagung, so er allhier gibt, zeucht GOtt nicht wieder zurück; es gereuet ihn nicht. Und sollte es kommen, daß GOtt über seinem Wort und Verheißungen nicht hielte, und ich in der Noth bliebe stecken, so ist es gewiß meines Unglaubens Schuld, daß ich GOtt nicht vertrauet habe, sondern in Kindesnöthen hat mir St. Margaretha helfen sollen, in den Zahnwehtagen St. Apollonia, in Kriegsläuften der Ritter St. Georg, in Schiffbruch St. Nickas, und so fortan; habe also meinen Glauben, Zuversicht und Vertrauen auf die Heiligen gesetzt, und bin schändlich betrogen worden.

94. So sollen wir uns nun hüten vor dem Unglauben und Verzweiflung, auch die Fahr und Noth aus den Augen setzen, aber dagegen GOttes Wort und Verheißung, auch die Exempel göttlicher Hülfe uns vorbilden. Denn, sehen wir allein auf die Gefährlichkeit und Noth, so hat der Teufel mit einer so langen Elle die Noth ausgemessen, daß man sie nicht ertragen kann, sondern man muß drüber verzagen. Also hat er den Propheten auch die Noth und Angst vor die Nase gehalten, und sie daran riechen lassen, damit sie sich denn bekümmern, sehr drüber klagen und heulen. Aber, je größer sich eine Noth aufbläst, je mehr sollst du diese Verheißung in diesem Texte ergreifen, und dem Teufel sie wieder ins Angesicht stoßen, da GOtt spricht [V. 21.]: „Laß dir nicht grauen vor ihnen, denn der HErr, dein GOtt, ist unter dir, der große und erschreckliche GOtt.“

95. Deß stellt Moses ihnen ein gar feines Exempel vor, als, ihre eigene Erlösung aus dem Diensthause Egypten, da sie mit trockenen Füßen durch das rothe Meer gegangen waren, als nichts denn der gewisse Tod vor ihren Augen war [2 Mos. 14, 21.]. Denn so stand es mit ihnen: Hinten war der Feind, Pharao mit großer Kriegsmacht und Heereskraft; vorne war das Meer, darüber war keine Brücke gebauet, noch einiger Furt, dadurch sie setzen möchten, sie hätten darinnen alle ersaufen müssen; auf beiden Seiten waren hohe Berge, über diese Berge konnte man nicht kommen; Himmel und Erde war da verschlossen; da war keine andere Rechnung, denn der gewisse Tod. Darum sahen sie auch allda ihre Kraft, die so groß war, daß sie gar verzagten, und meinten, nun wären sie verloren. Schrieen derhalben und sprachen [2 Mos. 14, 11.]: „Waren nicht Gräber in Egypten, daß wir darinnen gestorben wären“ und konnten die Gefährlichkeit, als, die Berge, das Meer und die Kriegsrüstung Pharaonis nicht überschlagen; es war das Facit eitel Tod in ihren Augen.

96. Aber da Moses sich hing an das erste Gebot, das also lautet: „Ich bin der HErr, dein GOtt“, und betete fleißig, schlug auch auf GOttes Befehl ins Meer, da theilt sich das Meer von einander, und wird eine herrliche, große, weite Landstraße hindurch, so breit, als fast von Wittenberg gegen Magdeburg sein möchte, fast ein zwölf oder dreizehn Meilen Wegs lang. Derhalben will alhier Moses sagen: Du hast es in der That erfahren, daß, was dein GOtt verheißt, will er auch thun, und er kann es wohl thun; darum läßt er diese Worte erschallen: „Ich bin der HErr, dein GOtt.“ So zweifle nicht, er kann dir Gesundheit, Nahrung, Kleider, Haus, Hof, Kinder, und alle andere Nothdurft geben; es mangelt ihm an der Macht und Vermögen nicht, er kann es wohl thun [Luc. 1, 37.]; auch ist der Wille dabei, denn er hat dir’s verheißen. Warum willst du denn nicht in aller Noth ihm glauben und vertrauen, daß er dir könne und wolle thun und geben, was dir fehlt? Derhalben sollen wir ihm auch vertrauen, bei seiner höchsten Ungnade, Strafe und Pein.

97. An dieses Mirakel und sonderliche Exempel der Erlösung aus Egypten sollen wir auch täglich gedenken; denn viel Psalmen dies göttliche Werk zurück holen und oft melden; so ziehen es alle Propheten auch an, und sonderlich Jesaias gedenkt dieser Erlösung oft. Denn wenn eine Noth ist vorgefallen, so haben sie daran gedacht, wie GOtt ihre Voreltern hat aus Egypten geführt, und daher geschlossen, daß er ihnen nochmals auch würde helfen. Derhalben will Moses lehren, daß sie sich hüten, damit sie in Anfechtungen und Nöthen nicht selbst aufwerfen, wie groß auch die Gefährlichkeit sei; auch nicht vornehmlich nachdenken, wie viel Feinde du hast, und wie wenig deiner sind; sondern stelle dir vor die Augen das göttliche Wort, die Verheißung und schönen Exempel, die werden dich weisen, wie Gott\footnote{„Gott“ von uns herübergenommen aus der alten Ausgabe Walchs, statt: „ich“ im Original.} helfen könne und wolle.

\quellseite{908}
98. In diesen Hülfspiegel siehe, und lasse den Nothspiegel fahren, sonst möchte er dir Schaden bringen. Es ist ja ein tröstlicher Spiegel, daß GOtt in seinem Worte sich dir vor die Augen hängt und stellt, und spricht: „Ich bin dein GOtt“, ich will dir helfen; hoffe nur auf mich. Und wenn dir sein Wort nicht genug wäre, daß du dich an dasselbige könntest halten, da er dir zusagt, er wolle es thun, und allen helfen, so siehe doch auf die Exempel und bewiesenen Thaten GOttes, da er es an andern wahr gemacht und gethan hat, darauf wir uns haben zu trösten und zu verlassen.

99. Also sollen wir Deutschen heutzutage auch thun. Wir haben zu Feinden den Türken, Pabst, Kaiser, und sonst andere Tyrannen, die alle wider uns sind; so haben wir nichts denn eitel große Gefährlichkeit vor Augen. Wenn wir es nun überschlagen nach den Zahlpfennigen, so ist der Türke, Kaiser und andere Feinde uns zu mächtig, und wir sind gegen ihnen, gleich als wenn Einer gegen Tausenden gestellt würde. Dieser Spiegel ist denn eitel, umsonst und verloren, welchen uns der Teufel vorstellt. Noch müssen wir die Augen von diesem Spiegel wegwenden, und gedenken, gleich als wäre kein Türke, Pabst oder anderer Tyrann, und dagegen den andern Spiegel des Lebens, nämlich das göttliche Wort und die Verheißung hervor ziehen, und uns größer und stärker lassen sein, denn tausend Türken und Päpste; ja, stärker lassen sein, denn die ganze Welt ist.

100. Denn was ist die Welt gegen GOtt? Ist sie nicht, wie der Prophet Jeremias [Cap. 40, 15.] sagt, gleich als ein Tröpflein Wassers, so noch am Eimer klebt, wenn das Wasser ist ausgegossen? Dies Tröpflein ist noch lange kein Strom; die Luft hat solch Tröpflein Wassers gar bald aufgeleckt. Wenn wir dies Wort und Verheißung GOttes nur ergreifen könnten (so wir anders rechte Christen sind), so würden wir empfinden GOttes Hülfe gegen uns. Es würde nichts Gewisseres sein, denn daß wir den Türken in Kürze schlagen würden; dafür sollte uns nicht leid sein, wir wollten den Türken wohl dämpfen. Denn GOtt hat uns ja verheißen, er will unser GOtt, und mitten unter uns sein etc. Wen wollte dieser herzliche Trost, diese güldenen Worte nicht lustig und fröhlich, muthig und trotzig machen, wider alle Feinde, ja, wider den Teufel und die Pforten der Hölle selbst [Matth. 16,18.]?

101. David im 18. Psalm hat auch auf diese Worte gesehen, da er spricht, V. 30.: „Mit dir kann ich Kriegsvolk zerschmeißen, und mit meinem GOtt über die Mauern springen.“ Denn wo der große und erschreckliche GOtt ist, warum sollte da nicht auch Glück und Sieg wider alle Feinde sein? Sollte da nicht Einer tausend jagen und schlagen können? Darum spricht er auch im 46. Psalm, V. 4–6.: „Ob wohl das Meer wüthete und wallet, und von seinem Ungestüm die Berge einfielen, dennoch soll die Stadt GOttes bleiben mit ihren Brünnlein, da die heiligen Wohnungen des Höchsten sind. GOtt ist bei ihr drinnen, darum wird sie wohl bleiben, GOtt hilft ihr frühe“ etc. Daher spricht auch St. Johannes in seiner ersten Epistel am 4.\footnote{Erlanger: „am 3. Capitel.“} Capitel, V. 4.: „Der in uns wohnet ist größer, denn der in der Welt wohnet.“ Und St. Paulus zum Römern am 8. Capitel, V. 31.: „Ist GOtt mit uns, wer kann wider uns sein?“ Damit sollten wir uns allezeit trösten. Denn es ist kein Zweifel dran, daß wir Christen mit diesem Troste und Glauben an diese Worte dem Türken, Pabst, Kaiser und allen Verfolgern bisher gewehrt haben, und wir wollten wohl ferner vor ihnen sicher bleiben, wenn wir GOtt glaubten. Aber wir gehen vor diesen Worten und Verheißung vorüber, und sind nicht solche Leute, die drauf baueten. Darum, wenn auch ein Unglück uns begegnet, daß wir drinnen stecken bleiben, und uns keine Hülfe widerfährt, so ist es kein Wunder; uns geschieht auch daran nicht Unrecht; darum sind wir gottlos und ungläubig.

102. Diesen Trost hatte auch der Prophet Elisa [2 Kön. 6, 14. ff.]; da er von den Syrern in der Stadt Dothan belagert ward, und sein Knecht heraußer ging und die Feinde vor der Stadt sahe, lief der Knecht hinein und schrie: O wir sind alle verloren! und wollte gar verzweifeln. Aber der Prophet Elisa that gleich, als sähe und wüßte er von keinem Feinde, und antwortete also seinem Knecht: Es sind ihrer mehr bei uns, denn bei ihnen. Bat aber GOtt, daß er dem Knecht wollte seine Augen aufthun. Da solches geschah, sah er eitel feurige Wagen und Rosse auf den Bergen. Es waren alle Berge voller Engel, die auf des Elisa Seite standen, die Feinde von der Stadt abzutreiben,
\quellseite{909}
damit dem Propheten Elisa sein Leid widerführe, und sie ihm (nach gemeinem Sprüchwort) die Rüben unversehret ließen.

103. Eben dies zeigt auch an die ritterliche Schlacht, 2 Kön. 19, 34. 35., die Gott am Senacherib vor Jerusalem that. Da er Jerusalem gar fressen wollte und zu einem Steinhaufen machen, wußte er die Stadt wohl vor ihm zu erhalten, und er verlor darüber in Einer Nacht bei hundert und fünf\footnote{Im Original: vier.} und achtzig tausend Mann, die des Morgens gefunden worden im Lager erschlagen. Wenn noch der Glaube bei uns wäre, schimpfs würden die großen Hunnen und Scharhanien zu Boden liegen; aber dieweil wir entweder zaghaftig sind, oder auf unsere Kräfte stolzieren und pochen, da ist es kein Wunder, daß uns zuweilen ein Rad über ein Bein geht, und wir in der Noth stecken bleiben. Blieben wir aber in diesem Glauben bestehen, so würde sich bald Hülfe und Erlösung finden.

104. Und ob uns gleich eine Noth und Trübsal in solchem Glauben über ginge, da wir doch GOtt für unsern HErrn halten, und also drüber geschlagen und getödtet würden, so könnte uns doch nicht besser geschehen, noch wir besser sterben; die Haut wäre theuer genug verkauft. Denn, errettete mich GOtt nicht hier zeitlich, so erlösete er mich doch von dem ewigen Tode. GOtt will die Seinen, die da glauben, erretten; und ob sie der bösen Nachbarschaft zuweilen müssen entgelten (denn ein Nachbar dem andern einen Brandschaden zu leiden schuldig ist, und sonst einer oft dasjenige bezahlen muß, das er nicht gekauft noch geraubt hat), so haben wir Gläubigen dennoch diesen Vortheil, daß, ob wir gleich mit der gottlosen Welt erschlagen werden und umkommen, so fährt sie doch zum Teufel in Abgrund der Hölle, mir aber ins ewige Leben, da wir bei GOtt ewig bleiben wollen.

105. Darum sollten wir uns nicht also in Gefährlichkeit und Unfall fürchten, sondern auf GOttes Wort und seiner Hülfe Beweis sehen, und darauf uns verlassen; denn es sind lebendige und kräftige Worte und Thaten, und zeugen genugsam, wie GOtt will aus allem Anliegen helfen, auch mitten aus dem Tode; allein, daß man sich deß auch tröste, und ungezweyfelt darauf verlasse. Darum setze dich in der Noth ja nicht wider GOttes Wort; richte auch nicht einen Abgott an, auf den du mehr denn auf GOtt wagest; gedenke doch hieran, was der HErr, dein GOtt, dem Pharaoni in Egypten gethan hat. Hast du nun an diesem mächtigen Werke und unwidersprechlichen Hülfe nicht genug, so schaue doch sonst seine gewisse Verheißung an, was er dir von seiner allmächtigen Gnade zusagt. Denn in diesen beiden Stücken, als, im Worte und Exempeln, findest du Trost und Hülfe; allein, daß du dein Herz auf GOtt gründest und gar dahin wagest.

106. Aus dieser Ursache nimmt ihm Moses so viel Raum und Zeit dazu, gebraucht auch so mancherlei auserlesene Worte, daß er solches vorbringe und uns berede, nämlich, daß wir zu jeder Zeit, wenn Unglück vorfällt, sollen die Gefahr und Noth aus den Augen thun, und uns allein hängen an GOttes Wort, und mehr GOttes Verheißung vertrauen, denn alles Unglück und Unfall fürchten, und sonst keinen andern GOtt fürchten noch suchen. Denn an diesem sollst du einen solchen GOtt haben, der seine Hände wird dir unterlegen. Er wird auch seinen Engeln befehlen über dir, daß sie dich auf deinen Wegen behüten, wie der 91. Psalm, V. 11., solche Verheißung uns gibt, und sonst der ganze Psalter nichts ist, denn eine Auslegung und Exempel des ersten Gebots.

107. So sollen wir nun allen Fleiß ankehren, damit wir das erste Gebot wohl lernen, und also unser Herz sich recht in alle Sachen wisse zu schicken, dann wird das andere von ihm selbst folgen, nämlich, daß es uns auch werde wohl gehen, und wir die Kraft dieser Zusagungen erfahren.

\begin{schriftzitat}
V. 20. Dazu wird der HErr, dein GOtt, Hornisse unter sie senden, bis umgebracht werde, was übrig ist, und sich verbirgt vor dir. Laß dir nicht grauen vor ihnen.
\end{schriftzitat}

108. Damit will er sie deß bereden, ihnen ein Herz einsprechen, und sagen: Fürchte dich allein vor GOtt, und sonst vor niemand; denn der HErr, dein GOtt, ist bei dir. Darum sind es noch immer Auslegungen des ersten Gebots, daß GOtt will sagen: Ich bin dein GOtt, es soll denen wohl gehen, die mich lieben. Darum, vertraust du Gott, so will er deine Feinde bald ausrotten, sollte er auch gleich Hornisse unter sie schicken, die sie zu Tode stechen müßten. In Summa, er will aus aller Beschwerung dir helfen, sei nur unverzagt.

\quellseite{910}
\begin{schriftzitat}
V. 22–24. Er, der HErr, dein GOtt, wird diese Leute ausrotten vor dir, einzeln nach einander. Du kannst sie nicht eilends vertilgen, auf daß sich nicht wider dich mehren die Thiere auf dem Felde. Der HErr, dein GOtt, wird sie vor dir geben, und wird sie mit großer Schlacht erschlagen, bis er sie vertilge, und wird dir ihre Könige in deine Hände geben, und sollst ihren Namen umbringen unter dem Himmel c.
\end{schriftzitat}

109. Er beharrt noch auf der Erklärung des ersten Gebots, und setzt lauter herrliche Tröstung, daß sie GOtt stark genug rüsten werde, damit sie den Feinden obliegen und sie überwinden mögen; denn er ist ihr GOtt. Glauben sie an ihn, so wird er unter ihre Gewalt die Feinde wohl zwingen [Ps. 47, 4.], allein, daß sie an ihm halten, und ihn lassen ihren GOtt sein.

110. Auf die Weise würden auch wir zu rechten Rittern geschlagen, und alle Feinde verachten können. Aber wir vertrauen ins Teufels Namen dem Zeitlichen mehr, denn GOtt, und setzen unser Herz mehr auf menschliche Rathschläge, auf Creaturen und gegenwärtige Hülfe, denn auf den lebendigen GOtt, der uns allhier solche gewisse Verheißungen gibt. Darum nehmen wir solche Verheißung nicht anders an, denn als hätte es nicht GOtt, sondern etwa ein Gaukler, Lotterbube, oder ein loser Holzhipler\footnote{Siehe Walch, St. Louiser Ausgabe, Bd. XX, 2395. s. v. Holzhipler.} geredet. Das Herz ist so böse, daß es straks nach dem Mammon, nach einem Schatz, oder Geld und zeitlichem Gut trachtet; und wenn es dasselbige hat, so stolzirt und pocht es darauf, will obenaus und nirgend an. Hat es aber solches nicht, und Armuth drückt es, so will es verzagen, wird so furchtsam, daß es Holz ließe auf sich entzwei hauen und scheitern,\footnote{scheitern = zu Scheiten spalten.} vergißt Gottes gar.

111. Was man um GOttes und des Nächsten willen in der Welt dran setze, deß sind die Gergesener ein Vorbild [Matth. 8, 31. ff.], welche gar übel-zufrieden sind, daß die Teufel, aus dem besessenen Menschen getrieben, in ihre Säue gefahren waren, stießen Christum aus ihren Grenzen, haben lieber ihre Säue denn den HErrn Christum. Aber also sollen sich halten und selber an Tag geben alle, die GOttes Wort verachten, und auf die Creaturen ihr Vertrauen setzen. Wer ihm läßt solchen Trost aus dem Herzen reißen, und will sich durch diese Worte nicht bewegen lassen, der ist nicht werth, daß er von GOtt regieret werde, sondern mag den Teufel zum Regenten haben. Wer nicht will GOttes sein, der mag des Teufels bleiben.

112. Fürwahr, es sind gewaltige Aergernisse und Abgötter, so da fechten und streiten wider das erste Gebot, davon Moses bisanher geredet hat, und an dem Stück fleißig gearbeitet, daß er es seinen Juden ins Herz hinein brächte, und daß demnach jedermann seinen Glauben und Vertrauen auf GOttes Gnade setzte. Der Teufel hat die Welt also klug und so weise gemacht, daß sie meint, sie wisse alles, und weiß doch in der Wahrheit nichts. Also treibe ich auch immer an diesem Stücke, daß wir erstlich lernen, Christum wohl erkennen, und Glauben an GOtt haben. Denn wenn dieser Artikel fället, so liegts alles; wenn er aber bleibt und erhalten wird, so ist alles erhalten. Ich bin ein Doctor der heiligen Schrift, jedoch lerne ich noch immer daran, und weiß es dennoch nicht alles. Aber die Welt weiß es alles. Solches hat Moses wohl gesehen, daß überaus viel daran gelegen, und es ein groß schwer Ding sei, daß man lerne, daß GOtt unser HErr ist, und daß man ihm vertraue; denn man kann es nicht auslernen. Es wird wohl mit kurzen und wenigen Worten geredet, aber es begreift viel in sich.

113. Und hat Moses bisanher befohlen den Juden, daß sie vor der Menge der Heiden nicht erschrecken, die sie sollten umbringen und vertilgen, sondern gedenken, was GOtt zuvor bei ihnen gethan hätte, und stellt ihnen vor sein Wort, und die vorigen Thaten oder Wunderwerke, daß sie getrost wären und frisch dran gingen, denn GOtt würde ihnen nicht lügen, oder seine vielfältigen Verheißungen in Born fallen lassen.

\begin{schriftzitat}
V. 25. Die Bilder ihrer Götter sollt du mit Feuer verbrennen, und sollt nicht begehren des Silbers oder Goldes, das daran ist, oder zu dir nehmen, daß du dich nicht drinnen verstrickst; denn solches ist dem HErrn, deinem GOtt, ein Greuel c.
\end{schriftzitat}

114. Er heißt allhier nicht allein, die Bilder abthun, sondern auch, nicht schmelzen, noch zu sich in seinen Schatz reißen und ziehen. Item, sie gänzlich verbrennen. Ob sie es nun gehalten
\quellseite{911}
haben, oder nicht, das lasse ich an seinem Ort und Würden bestehen. Aber dieser Text hat viel Occasion gegeben den Schwärmern und Bilderstürmern, die für eine Vollkommenheit und Gerechtigkeit vor GOtt gehalten haben das Werk, Bilder verbrennen, oder ruiniren mit den Bildern und sie einreißen, und haben durch solche Werke wollen heilige Leute und gute Christen sein. Aber wenn es dieses gelten sollte, so wäre der Türke der Allerheiligste, und hätte große Gerechtigkeit vor GOtt; denn er reißt gar viel Bilder, Schlösser und Städte ein. Darum soll man sich damit nicht brüsten, oder dies Bilderstürmen für eine Kunst ausgeben.

115. Davon habe ich nun viel geschrieben im Buche wider D. Carlstadt, vom Bilderstürmen;\footnote{Gemeint ist der Abschnitt „von dem Bilderstürmen“ in Luthers Schrift „wider die himmlischen Propheten“, Walch, St. Louiser Ausgabe, Bd. XX, 138 ff.} darinnen habe ich mich hart wider dasselbige Werk gelegt. Und ist das meine Meinung: Ein Christ solle anfangen von der vollkommenlichen Gerechtigkeit des Glaubens, daß er sich verlasse auf das Wort GOttes, und das erste Gebot vom Herzen lerne halten, daß er GOtt liebe, vertraue, und ihn über alle Dinge fürchte. Wer das thut, der wird alle silbernen und güldenen Götzen leichtlich verachten. Denn das äußerliche Bilderstürmen ist nur eine Larve, dadurch man zeitliche Ehre erjagen will, und sich damit rühmen, gleich als hätte man es damit gar wohl ausgerichtet.

116. Aber du hast darum GOtt noch nicht vertrauet, du bist auch noch nicht gottselig und ein Christ, ob du gleich Klöster, Kirchen und Altäre eingerissen hast, denn das könnte jedermann wohl thun. Aber darum liegt’s vornehmlich GOtt und uns Menschen, daß wir das erste Gebot recht halten lernen, welches der ganzen Welt gegeben ist; nicht allein den Juden, sondern auch uns Heiden, daß wir GOtt vertrauen, und ihn erkennen für unsern GOtt, der uns helfe in allen Nöthen und Aengsten, auch daß wir ihn fürchten in allem Glück und Unglück. Der Glaube soll gegen GOtt hervorleuchten im Kreuz und Trübsal. Die Furcht GOttes aber sollen wir gebrauchen, wenn es uns wohl geht. Wer also das erste Gebot wohl betrachtet, der würde der Bilder leichtlich vergessen; aber leichtfertige Herzen fahren also heraus, wie die Bilderstürmer gethan haben.

117. So lerne du erstlich wohl vertrauen unserm HErrn GOtt in allen Nöthen und Gefährlichkeiten, beide geistlich und leiblich, in Anfechtungen Leibes und der Seele, daß du nicht von GOtt ablässest, und anderswo nach Hülfe dich umthuest, oder verzweifelst, als Kain, Saul und Judas. Wiederum, wenn dir es wohl geht, nicht stolz noch frech werdest, und GOttes darüber vergessest, gleich als hättest du Leib, Leben, Gedeihen und alles von dir selber.

118. In Summa: Es soll einem Christenmenschen frei stehen, Bilder haben, und nicht haben. Denn im 22. Capitel des Evangelisten Matthäi, V. 20., lesen wir, daß der HErr Christus der Juden Münze hat angegriffen, und [sich] nicht gescheuet, in die Hände sie zu nehmen, ob sie ihn wohl dadurch versuchten; fragt: welches sei des Groschens oder Guldens Gepräge, oder Bildniß und Ueberschrift? rührt des Kaisers Bildniß an; verbeut nicht diese Bilder, sondern bestätigt sie vielmehr; sonst dürfte kein Christ die Münze von Gold und Silber anrühren. Derhalben so ist es keine Sünde, Bilder haben. Aber solche Bilder haben, daß man Abgötter daraus mache, und das Vertrauen auf etwas Anderes setze, denn auf GOtt und das erste Gebot, solche Bilder soll man zerstören.

119. Wie wir denn viel Bilder der Heiligen haben abgebrochen, und wäre zu wünschen, daß wir viel silberne Bilder hätten, wir wollten sie wohl christlich brauchen. Es will GOtt im neuen Testament nicht viel fragen nach den Bildern, sondern fordert das Herz. Denn die in dem Worte des Glaubens wohl unterrichtet sind, können leichtlich aller Abgötterei und Bilder vergessen, wenn sie auch mitten unter den Bildern stecken. Wo man aber GOttes Wort und Gebot aus den Augen setzt, so bekommt man allerlei Bildnisse und Götzen. Und obgleich nicht Bilder vorhanden wären, so hängt man sich doch an den Mammon, Reichtum, Ueberfluß, an die große Menge, auch an gute Gesellschaft oder Bündnisse,\footnote{Erlanger: „Bildniß“.} wie denn Moses solches alles bisanher ernstlich verboten hat.

120. Die Juden hatten aus diesem und andern mehr Texten ein sonderlich Gebot, die Bilder zu verbrennen und sich vor Bildnissen zu hüten, auf daß GOtt ihnen steuerte, damit sie
\quellseite{912}
sich nicht an denselben vergriffen. Gleichwie wir noch unsern Kindern wehren, daß sie nicht Lichter oder Messer angreifen; nicht um des Lichts oder der Messer willen, sondern des Kindes halben, daß es sich nicht mit dem Lichte brenne, oder mit dem Messer steche oder schneide. Also sollten die Juden die Bilder meiden, auf daß sie GOttes Gebiete recht brauchten im Geist, und nicht in Abgötterei fielen, und sich an dem ersten Gebot versündigten.

\section{Das achte Capitel}

\begin{schriftzitat}
V. 1. Alle die Gebote, die ich dir heute gebiete, sollt ihr halten, daß ihr darnach thut, auf daß ihr lebet und gemehret werdet, und einnehmet das Land, das der HErr euren Vätern geschworen hat etc.
\end{schriftzitat}

1. Dies Capitel ist eine neue Predigt von dem ersten Gebote. Denn Moses, wie wir wissen, hat gar viel Predigten gethan über dies erste Gebot. Er will aber in diesem Capitel uns vermahnen und treiben, daß wir täglich und fleißig gedenken an das erste Gebot, auf daß wir GOttes nicht vergessen. Befiehlt und bindet uns also ein die stetige Betrachtung und fleißige Handlung des ersten Gebots, auf daß wir es mit der That und Werken befinden, wie wir dadurch leben, und vielfältig sollen genährt werden.

2. Zudem will er ihren Glauben stärken durch die Einbildung der vorigen Wohlthaten GOttes, auf daß sie ihm vertraueten, und ja von ihm in Anfechtungen nicht abfielen. Darum spricht er, sie sollen zurück sehen, und bedenken, daß er ihr GOtt gewesen wäre, da er sie vierzig Jahre lang in der Wüste geführt, geleitet, ernährt und erhalten hätte; an welchen Orten gar nichts zu essen und zu trinken war, da weder Bäcker, Brauer, Schuster oder Schneider wohnten; noch hatten sie von GOtt Essen, Trinken, Kleider und alle Nothdurft. An diesem herrlichen Werke sollten sie lernen und erkennen, daß er nochmals ihr GOtt sein würde, und wenn sie ihm vertraueten, so würden sie Hülle und Fülle, und alle Genüge haben.

3. Solch Exempel der göttlichen Sorge und Liebe für sein Volk stellt er jedermann hier vor die Augen, auf daß man durch die Erfahrung doch versucht und gedrungen würde, GOtt zu glauben. So ist nun die Verheißung noch heutzutage vorhanden, darinnen sich der HErr genugsam erbeut, und allenthalben ausrufen läßt, daß er will unser GOtt sein; so thut er auch noch täglich solche Werke, welche uns zum Glauben an ihn gewaltig reizen sollten; aber an dir fehlt es, daß du nicht glaubst. Das Wort hält dir, und sieht vor Augen das Werk an der ganzen Creatur, welches dich billig vermögen, und dahin treiben sollte, zu glauben, daß er auch dir wolle Gutes thun.

4. Er setzt aber in diesem Capitel eine Sünde, so wider das erste Gebot geht, welche so heißt, daß man in Armuth, Hunger, Kummer, Widerwärtigkeit und Trübsal das Herz von GOtt abwendet, und auf ihn sich nicht verläßt. Wie man denn, auf der andern Seite, in Reichthum, Ueberfluß und glückseliger Zeit, GOttes auch vergißt. Und ist solches alles von Mose dahin gerichtet, daß er gerne wollte steuern unserm Unglauben und Verzweiflung, daß wir lernten den Trost, der im 33. Psalm, V. 19., geschrieben steht, daß der HErr die Seinen ernähre in der Theurung, und lasse sie nicht Hungers sterben, gleichwie er die Israeliter in der Wüste erhalten hat.

5. Lernet derhalben auch ihr auf zukünftige Noth, wie man solchen Hunger, Mangel und Anfechtung ansehen und überwinden soll, nämlich, daß es nur eine Demüthigung und Versuchung GOttes sei. Nicht, daß er die angefochtenen oder nothleidenden Christen verlassen wollte, sondern er will erfahren, ob sie auch an ihn glauben, oder nicht. Darum sagt Mose im Text [V. 2. 3.]: „Der HErr, dein GOtt, hat dich vierzig Jahre geleitet in der Wüste, gedemüthigt, und versucht, daß kund würde, was in deinem Herzen wäre, ob du seine Gebote halten würdest, oder nicht. Er demüthiget dich und lässet dich hungern, kommt aber zu rechter Zeit, und speiset dich mit Manna.“

\quellseite{913}
6. Da die Israeliter in Egypten waren, Essen und Trinken genug hatten; auch da sie aus Egypten zogen, und ihre Tücher voller Zeuges mittrugen, stellten sie sich fromm und heilig, und man konnte nicht gewahr werden, ob sie GOtt von Herzen liebten, glaubten und vertrauten, dieweil der Sack noch voll war. Aber da der Sack ausgeleert, und nicht mehr Mehl, Teig und Brod vorhanden war, gar bald ließen sie sich merken, daß sie GOtt bisher des Bauchs halben gedient hatten. Denn da gedachten sie, was sie gethan, und daß sie es mit dem Auszuge aus Egypten nicht recht getroffen hätten: daß sie die Zwiebeln und Fleischtöpfe verlassen, und nun in der Wüste weder zu beißen noch zu broden hätten. Murrten derhalben wider Moses; wie wir im 16. Capitel\footnote{Im Original und in der Erlanger: „im 13. Cap.“} des 2. Buchs Mosis (V. 2. 3.) gehört haben.

7. Also ist es um ein ungläubiges Herz gethan, das vertrauet GOtt nicht länger, denn weil leibliche Güter vorhanden sind; werden aber dieselben entzogen, so weicht es von GOtt ab, und lästert ihn. Darauf redet hier Moses: Daß der HErr dich hungern ließ, das sollte dir eine Demüthigung sein, GOtt um Speise, Nahrung und Leibes Nothdurft anzurufen, und solches von ihm gewiß gewärtig zu sein. GOtt will dein Schulmeister sein, und dich lehren das erste Gebot verstehen, daß dein Herz auf ihn sehe und ihm vertraue, wenn nichts vorhanden ist, und sich ansehen läßt, als wollte er uns gar verlassen.

8. Denn GOtt hat die Art und Weise: Ehe denn er den Seinen hilft, so läßt er sie zuvor erst in Leiden und Anfechtung kommen; darnach, zu seiner Zeit, hilft er ihnen desto wunderbarerlicher heraus. Alhier, in der Wüste, kann er die Kinder Israel vierzig Jahre lang speisen und tränken, auf daß man sehe, GOtt könne viel mehr Rath schaffen und bei allen Sachen thun, denn alle Welt; er kann auch mehr, denn alles Gold, Silber, Geld, Gut, Speise, Trank und Kleider.

9. Denn das muß mir ein hübscher Schneider und Schuster sein, der auf vierzig Jahre Kleider und Schuhe machen kann; und diesen GOtt lasse man einen Speisemeister oder Hausvater sein, der hat einen solchen Kornboden oder Söller, daß er aus der Luft oder steinichten Erbe alsbald Korn und Brod, auch aus den Felsen Wasser hervorbringen kann. Ps. 65, 10. ff. Ps. 145, 15. 16. Ps. 147, 9. ff. Wenn ich sollte Brod suchen nach der Vernunft, so suchte ich es bei dem Bäcker, und Korn auf den Schlössern; so bringt es GOtt aus der ledigen Luft. Da kein Handwerk noch Vorrath ist, muß Brod und alle Genüge sein. Das ist ein Meisterstück, welches man wohl merken soll. Es könnte solches GOtt noch alle Augenblick beweisen. Jes. 50, 2. 3.

10. Daher sehen wir, daß alle Gläubigen, so GOtt vertrauen, und sein Wort lieben, müssen genug haben; obgleich keine leiblichen Mittel von Essen und Trinken vorhanden wären, so weiß GOtt doch wohl Hülfe zu schicken; wie er den Eliam in der Wüste durch die Raben speiste, und hernach durch die Wittwe zu Sarepta ernähren ließ [1 Kön. 17, 4. 15.], und jemals die Juden durch eine neue Speise, als das Manna, ernähret sind worden, darauf sie nie gedacht haben, auch mußten ihre Kleider nicht veralten noch zerreißen.

11. Diesen Trost sollten wir unserm Abgott, dem Bauche, vor die Nase halten, daß wir lernten das erste Gebot, GOtt und seinem Worte und Verheißung in der Mangelszeit vertrauen. Es will aber die ganze Welt das erste Gebot GOttes sehr wohl verstehen; aber wenn sie aufsagen sollten, wüßten sie gar nichts davon, da sie doch zuvor stracks Doctor und Meister sein wollten. Ich bin nun lange ein Prediger, Doctor und alter Kämpfer; aber wenn ich das erste Gebot aufsagen soll, so habe ich es nehrlich bei einem Seiden-Faden ergriffen und gefaßt. Der Teufel kann meine Zuhörer bald ein anderes lehren, wenn Noth und Armuth vorfällt.

12. Wie denn die Israeliter zu Mosen und Aaron in der Wüste auch sagten: Der Teufel hat es euch geheißen, daß ihr uns aus Egypten führet. Also verlösch ihr Glaube. Da der Abgott Bauch den Sack nicht mehr voll Mehls sah, flugs war aller Predigt und des ersten Gebots vergessen. Aber Moses berichtet sie hier der Ursache: GOtt ließ dich mangeln, auf daß du sein Kind und Schüler würdest, und lerntest das erste Gebot recht verstehen, und sagen: Obgleich der Sack ledig worden, daß Mehl heraus ist, und der Teig aufgefressen, so glaube ich ihm doch; denn er hat mir gesagt: „Ich bin dein GOtt“, du sollst genug haben; obgleich kein
\quellseite{914}
Sack, Teig, Mehl etc. vorhanden ist, so habe ich unsern HErrn GOtt noch bei mir zum Vorrath.

13. Die Israeliter aber sagten viel anders: der Teufel, und nicht GOtt, hätte sie aus Egypten gebracht. Also thun wir auch, murren ohne Ursache wider GOtt, unangesehen, daß er zugesagt: Ich will dein GOtt sein, und dich ernähren; oder, bin ich nicht so kräftig und mächtig, daß ich dir zehn Gulden, einen Scheffel Mehls und Korns geben könnte? Das macht alles, daß niemand das erste Gebot recht auffassen kann, niemand denkt ihm auch nach, daß er es lerne, oder mit Ernst sich darum annehmen wollte.

14. Es verstehet aber das erste Gebot niemand recht, er komme denn zuvor in Mangel und Noth, daß sich es ansehen läßt, als vergesse GOtt deiner, als wollte er dich stecken lassen. Da schreiet du denn: Ach GOtt! wo soll ich mich hinhalten? wo soll ich Trost und Hülfe erlangen? Soll ich sehen in den Keller, auf den Boden, in die Küche, in den Kasten, und meinen Schatz hervorziehen? habe ich doch keinen. Da kommt das erste Gebot recht, und schreiet dir und allen Menschen zu: „Ich bin der HErr, dein GOtt.“ Dieser Hausvater wird dir zu essen und zu trinken geben.

15. Es muß aus dem ersten Gebot gebacken und gesponnen werden. Sollst du einen Faden an die Haut, und ein Stück Brods ins Maul haben, so muß dir es das erste Gebot geben, sonst wirst du nimmermehr etwas bekommen; verlache es gleich so höhnisch, als du willst, du gottloser Bube! Das Wort: „Ich bin dein GOtt“, ich habe dir geschworen, Gutes zu thun, das gibt den Hungrigen Brod und Mehl in Sack, gibt auch Gesundheit in Krankheit, und alles, was man sonst weiter bedarf; gleichwie allhier von dem Himmel, aus der lebigen Luft, Brod herabfällt. Da hätten die Israeliter das Brod nicht in der Luft gesucht, sondern irgend bei einem reichen Bäcker; noch konnte es GOtt daher geben und verschaffen. Sollte er denn nicht noch solches oder dergleichen vermögen? Pfui unseres leidigen Unglaubens!

16. An dies Werk gedenke, denn es kommt mit dem Wort und Gebot GOttes überein, da gesagt wird: „Ich bin dein GOtt.“ Denn da steht die That, daß er Gutes thut und Brod gibt, welches niemand von dem Orte her hätte gewärtig sein können. Aber die lieben Juden waren GOtt dafür wenig genug dankbar; noch gibt er es umsonst und aus Gnaden, auf daß solch Werk bei jedermann bekannt sei, und das Wort allezeit wahr bleibe, daß er unser GOtt sei, und wir ihm vertrauen. Wer hätte damals denken sollen, daß er aus den Wolken würde Brod geben? Und dennoch hat es GOtt gethan, auf daß nicht einer sagte: Ich habe es GOtt vorgeschlagen, habe ihn darum gebeten und darauf bedacht, daß er es mir aus der Luft gebe, sondern er thut es aus lauter Güte, auf daß dir kund werde, der Mensch lebe nicht allein von dem Brod etc.

17. Diesen Spruch zeucht der HErr Christus Matthäi am 4. Capitel, V. 4., auch an, und begegnet dem Teufel damit, der ihn mit der Bauchsorge plagte, da er in der Wüste vierzig Tage und Nächte gefastet; hätte Christum auch gerne in Mißtrauen an GOtt getrieben. Da antwortet der HErr Christus: „Der Mensch lebet nicht allein vom Brod, sondern von einem jeglichen Wort, so aus dem Munde GOttes gehet“; als wollte er sagen: Wenn gleich diese Steine eitel Brod wären, so lebt doch der Mensch nicht davon, sondern von GOttes Wort.

18. Ist nun kein Brod da, so halte sich der Mensch nichtdestoweniger auf das Wort GOttes, und ergreife solches durch den Glauben, dann wird zu rechter Zeit wohl Brod kommen, sollte es auch aus der Luft herab geschickt werden; und ehe Trinken mangeln sollte, müßte Wasser aus den Felsen heraus fließen. Diese herrliche Kunst lehrt Moses, dadurch man in Noth und Kummer möge genug haben. Wenn man auch gleich nichts mehr denn nur das bloße Wort GOttes hätte, das also lautet: „Ich bin der HErr, dein GOtt“, dennoch müßte aus diesem Wort Speise, Trank, Brod, Korn, Wolle, Leib und Leben, und alle Nothdurft werden.

19. Diese Kunst lerne du auch, auf daß du GOtt vertrauest. Wenn du gleich nichts mehr hättest denn allein das erste Gebot, so wird doch das andere alles daraus herkommen, sollten gleich die Wolken Brod regnen, oder aus den Felsen zu trinken fließen. Denn wer das erste Gebot hat, der hat alles, und die Welt muß bekennen, wenn sie recht reden will, daß sie alles von dem GOtt habe, was sie hat, der im ersten Gebot spricht: Ich will dein GOtt sein.

20. Solches erste Gebot lerne der Ursache halben auch wohl, daß GOtt sich darinne erbeut, er will dein GOtt und Helfer sein, und dich ernähren, wie er es gewaltiglich in diesem Exempel
\quellseite{915}
mit den Kindern von Israel, mit dem Manna gethan hat, und will es noch an alle den Seinen beweisen. Denn also sagt der 37. Psalm, V. 19.: „Er will sie in der Theurung ernähren.“ Item, David spricht im selbigen Psalm, V. 25.: „er sei jung gewesen und alt worden, habe noch nie gesehen den Gerechten verlassen, oder seinen Samen nach Brod gehen.“ So bezeugen solches auch andere reichliche Verheißungen GOttes, als im 34. Psalm, V. 10., da gesagt wird: „Die den HErrn fürchten, werden nicht Mangel haben an irgend einem Gut.“ Und St. Petrus in seiner Epistel sagt: „Alle eure Sorge werfet auf den HErrn, denn er sorget für euch“ [1 Petr. 5, 7.].

21. Der HErr Christus im Evangelio Matthäi am 6. Capitel, V. 28. ff., hält uns auch vor gar seine Exempel der göttlichen Sorge für die Seinen, da er sagt, wie er kleide die Lilien auf dem Felde, die nicht nähen oder spinnen, und ernähre die Vögel unter dem Himmel, die nicht säen oder pflanzen; viel mehr werde er seine Christen speisen, ernähren, erhalten und schützen, die viel mehr sind, denn die Lilien und Sperlinge. Aber Mosis Exempel das ist auch trefflich und wichtig, daß GOtt aus der Lufft Himmelbrod gegeben, und Wasser aus den Steinen habe fließen lassen, damit zu speisen und zu tränken vierzig Jahr lang die Kinder von Israel. Und es dient überaus wohl, unsern Glauben damit aufzurichten, und das erste Gebot fleißig zu studiren und zu lernen, daß wir bei dem lieben Worte GOttes bleiben, und mit dem ernsten, christlichen Gebet anhalten; denn so wolle GOtt noch allen Gläubigen helfen, und geben, was sie in diesem Leben bedürfen.

22. Unsere eigene Erfahrung stimmt mit diesem Mirakel und Gebot GOttes auch überein, Dies sollte uns je bewegen, daß wir GOtt desto lieber vertraueten. Denn der GOtt (welcher diese hohen Wunderwerke hat können thun, als, Kleider und Schuhe in die vierzig Jahr lang ganz behalten, nicht sich vertragen und vernützen lassen, ob sie wohl gleich täglich die am Leibe getragen; item, der von oben herab vierzig Jahr lang Brod gegeben, damit er so viel hunderttausend Menschen täglich hat können speisen und versorgen), der ist auch unser GOtt. Warum wollen wir ihm denn auch nicht herzlich uns selber und alle das Unsere zu seinen väterlichen Händen übergeben?

23. Diese Worte, die Moses allhier gebraucht, als: „Gedenke alle der Wege, dadurch dich der HErr geleitet hat diese vierzig Jahre“; item: „Vergiß sein nicht, und erkenne, daß der HErr dich gezogen hat, wie ein Mann seinen Sohn zücht“, die gelten auch mir und dir. Wir sollen diesen Text wohl merken, und unser Lebenlang den uns lassen befohlen sein. Insbesondere aber sollen wir [uns] diesen Spruch: „der Mensch lebe nicht allein vom Brod, sondern von einem jeglichen Wort, so aus dem Munde GOttes gehet“, eigentlich einbilden, den theuer und werth halten; denn er befiehlt uns, das mündliche Wort oder Predigtamt hoch und groß zu halten, fleißig zu hören und zu lernen.

24. Wir sollen nicht gedenken, daß es daran genug sei, die zehn Gebote und das göttliche Wort einmal oder zwei aushören, sondern man muß es für und für hören und handeln. Denn geht es doch mit dem äußerlichen, leiblichen Brod also zu, daß, wenn du heute issest, so lässest du es dabei nicht bleiben, sondern morgen issest du wieder, und treibst solch Essen von Tag zu Tage, für und für, auf daß der Leib seine Speise, Nahrung und Unterhalt davon habe. Also laß es auch nicht an dem genug sein, daß du die zehn Gebote und GOttes Wort einmal gehört und gelernt hast, sondern du mußt dein Lebenlang daran lernen, denn man kann sie nimmermehr genug wissen noch auslernen.

25. Daher gebeut auch St. Paulus [Col. 1, 11.], daß man wachsen und zunehmen soll im Erkenntniß des HErrn Christi, welches geschehen muß durch die tägliche Hörung und Handlung des göttlichen Worts. So lehrt auch der erste Psalm, V. 2., daß man am Gesetze des HErrn Lust und Freude haben solle Tag und Nacht. Und Joh. 5, 39. sagt der HErr Christus: „Suchet in der Schrift, denn die zeuget von mir“; will, daß man GOttes Wort fleißig forschen und betrachten solle. Man muß sein nicht satt, müde noch überdrüssig werden, dasselbige fahren lassen, oder unter die Bank stecken, sondern immerdar hören, treiben und wieder auf die Bahn bringen.

26. Wir hören aber allhier von einem schönen Encomio, Preis und Lobe, so Moses dem göttlichen Wort und Geboten gibt und zuschreibt, nämlich, daß es das Leben gebe in dieser Welt, und auch in jenem Leben; sie erhalten einen Menschen zeitlich und ewiglich, Ps. 107, 9. 20.

\quellseite{916}
Weish. 16, 12. Davon St. Paulus auch zum Timotheo redet: „Die Gottseligkeit (das ist, wenn einer GOtt glaubt), die hat Verheißung des gegenwärtigen und zukünftigen Lebens“ [1 Tim. 4, 8.]. Denn wer sein Vertrauen auf GOtt stellt, der hat die Verheißung, daß GOtt ihm will beistehen und ernähren in dieser Welt, und im ewigen Leben selig machen. Denn Mosis sagt ausdrücklich, daß GOttes Wort macht den Menschen zeitlich leben, stärkt ihn in diesem und aufs zukünftige Leben.

27. In diesen Worten redet Moses solches alles vom äußerlichen, mündlichen Wort, wider die Secten und Schwärmgeister, welche das Predigtamt verachten. Denn soll man solches vor die Augen halten. Es hatten die Juden das äußerliche Wort GOttes, das also lautet: Ich will dein GOtt sein, darum sollst du mir glauben; solches war eine gewisse Verheißung, auch was die leibliche Nahrung anbelangt; das hielt ihnen Moses treulich genug vor, welcher dieses Wort führt, und durch seine Predigt überkamen sie Speise. Darum steht dies Exempel auch dabei, daß die Kinder von Israel mit Hunger in der Wüste versucht sind; aber dieweil sie dem Worte GOttes angehangen, und an GOtt fest hielten, so mußte Brod kommen, und wurden also in der Wüste reichlich versorgt, gespeiset und befriedet.

28. Dieses Mirakel geschieht noch täglich an den Christen, daß sie nicht müssen Mangel leiden in leiblichen Nöthen; denn St. Paulus [1 Tim. 4, 8.] spricht, sie sollen allhier in diesem Leben genug haben. Ob sie gleich oft ein wenig mangeln müssen, so ist es doch nur eine Versuchung, und sie geschieht darum, daß man solle zum Wort GOttes laufen, und wissen, daß das erste Gebot lebendig macht. Denn, wenn wir nicht versucht würden, so gedächten wir nicht an dies Gebot, wir lernten auch nimmermehr GOttes Wort. Das erste Gebot ist ein Brunn des Lebens; darum, in Leibesnoth muß entweder Brod kommen, oder an seine Statt andere Speise, die uns ernähre; gleichwie in der Wüste das Manna ohne ihr Zuthun, ohne ihr Werk, vom Himmel herab und aus der Luft kam, dahin es niemand gesäet hatte. Auch mußte das Wasser aus den Steinklüften entspringen; ihre Schuhe, und alles, hielt fest und währte wohl.

29. Will nun jemand wissen und Bericht davon haben, was solches alles unverändert behalten habe? dem antwortet hier Moses selber: Das Wort hat es gethan, darauf sie vertraueten. Davon mußten ihre alten Kleider nicht zerreißen, sondern so fest halten, als wären sie nur jetzt neu gemacht, denn das Wort erhielt sie, das unter sie Mose predigte: „Ich will dein GOtt sein;“ bei welchem Wort sie fest hielten. Und dies Wunderwerk, neben den angehefteten Verheißungen, sollten wir auch fleißig beherzigen und uns lernen nütze machen, denn wir haben eben dasselbige Wort und Zusagung GOttes auch, so steht uns diese That GOttes in seinem Wort vor Augen gemalt. Aber da hängt es, und will mit uns auch nicht fort; welches Moses sehr klagt von den Kindern Israel, daß sie es nicht lange haben angetrieben; es hat die Andacht und Gottseligkeit mit ihnen (wie man sagt) einen Tanz und eine hohe Messe gewährt.

\begin{schriftzitat}
V. 5. 6. ff. So erkennest du je in deinem Herzen, daß der HErr, dein GOtt, dich gezogen hat, wie ein Mann seinen Sohn zucht. So halte nun die Gebote des HErrn, deines GOttes, daß du in seinen Wegen wandelst, und fürchtest ihn; denn der HErr, dein GOtt, führet dich in ein gut Land, ein Land, da Bäche und Brunnen und Tiefen innen sind etc.
\end{schriftzitat}

Diesen Text hat Lutherus gelesen bis zum Ende des achten Capitels.

30. Moses ist gar ein fleißiger Prediger, der immerdar feilt und treibt an den Juden, daß sie das erste Gebot lernen mögen, und hat diese Predigt nicht an Einem Tage gethan. Er sagt aber erstlich: „GOtt hat dich gezogen, wie ein Mann seinem Sohn zeucht.“ Es ist GOtt ganz väterlich mit ihnen umgegangen, er hat sie zärtlich erzogen. Ein Vater säumt zuweilen sein Kind, und hält es hart, schilt und straft es; bald aber hat der Vater zum Ende, er gibt ihm wieder darauf einen Apfel oder Birne; nimmt es ihm auch wohl wieder, will sehen, ob das Kind ihn als seinen Vater liebe, und ob es Gutes sich zu ihm versiehe: er spielt mit ihm, nimmt dem Kinde und gibt es ihm wieder, will es also an sich gewöhnen, daß es eine kindliche Zuversicht zum Vater gewinne.

31. Also thut GOtt auch mit uns. Er nimmt uns oft unsere Habe und Gut, Gesundheit und Liebsten, auf daß, wenn wir in Hunger, Armuth, Krankheit, Betrübniß und Herzleid stecken, zu ihm laufen, ihn flehen und bitten. Denn Versuchung
\quellseite{917}
schickt er uns zu, nicht als ein Feind, oder als unser Henker und Stockmeister, sondern als unser lieber Vater. Er will, daß du rufest und klopfest und sagest zu ihm: Gib, lieber Vater! und solches begehrest mit herzlichem Sehnen und Verlangen. Dahin müßte den Juden ihre Wüste und beschwerliches Reisen auch gereichen, und daß, wenn sie, aus der Wüste geführt, die Gefährlichkeit und Noth ausgestanden hätten, und in Reichthum und Ueberfluß säßen, sie alsdann nicht sicher würden, Gottes dabei vergessen, und andern Göttern nachliefen, sondern GOtt für solche Wohlthaten dankten, und erkenneten, daß sie es alles von ihm hätten; wiederum auch, in vorfallenden Anstößen nicht anderswo Hülfe suchten, welche sie von GOtt sollten erwarten, und ihm in seine Hände sehen. Also kann Gott sein mit uns umgehen, daß wir in Mangel und Ueberfluß uns recht gegen ihn müssen zu halten; und so muß er uns, als ein rechter Vater, nach seiner Hand und Weise ziehen und gewöhnen.

32. Man sehe aber gleichwohl daneben an das gottlose Wesen der Juden, und zwar noch heutiges Tages der ganzen Welt. Die Juden meinten, sie hätten alles durch ihre Kräfte; darum, wenn es ihnen wohl ging, und sie alles genug hatten, von Stund an vergaßen sie unsers HErrn GOttes. Wir befinden solche Undankbarkeit auch noch heutiges Tages wider das göttliche Wort. In was Jammer und Blindheit sind wir doch unter dem Papstthum gewesen, da wir alle unsere Habe und Gut auf denselbigen falschen Gottesdienst wendeten. Wahrlich, wenn wir solche Noth noch recht erkenneten, es wäre keiner so arm und unvermögend, der nicht seinen Rock und Kühe, oder was er sonst hätte, dahingäbe, daß er dieses Zwanges möchte los sein.

33. Aber wir haben es jetzt gar vergessen. Sollten wir jetzt einem Pfarrherrn, oder sonst armen Leuten, einen Scheffel Korn zu der Zeit des Evangelii geben, so müßten wir verarmen. Jedermann geizt und scharrt, und gibt niemand nichts. Die Leute müssen sein jetzt zu sagen: Ich habe es mit meiner blutsauren Arbeit und Fleiß erworben, sollte ich es so hingeben? rauben, tragen, stehlen gleichwohl immer zu sich, gedenken nicht daran, daß wir alles von GOtt haben. Das Leben haben wir je nicht von uns selber; so ist uns der Tod täglich auf dem Halse. Unser Leben hängt in eines Andern Hand und Gewalt; noch thut ein jeder, was er will, gleich als wäre er ein Herr seines Lebens und Guts, fürchten weder GOtt noch seinen Zorn.

34. Darum muß GOtt kommen, und machen, daß wir dieses Lebens müde werden, denn die Leute sind also verstockt, und härter denn die Adamanten. Man findet ihrer noch wohl, wenn man ihnen dräuet, die noch wünschen und wollen, daß der Türke kommen möchte, hoffen auf ihn. Aber, ihr ohnmächtigen Bauern, ihr dürft auf ihn nicht hoffen, ob ihr wohl vorgebt: Ei, er macht uns alle frei, daß wir nicht also Zins, Schabung, Geheiß und Tribut geben dürfen. Ja, wenn der Türke kommt, schlägt er dir den Kopf ab, er hauet dich mitten von einander, führt dich mit Weib und Kind gefangen hinweg. Also macht er alles frei. Aber es ist nicht besser, denn mit euch Alten immer zum Teufel zu. Wir predigen noch um der jungen Leute willen, welcher Herzen noch nicht vergiftet sind, ob GOtt denen wollte Gnade geben, daß sich dieselben besserten. Mit den Alten ist es verloren, die geizen und wuchern, leben dahin, als hätten sie das Leben von sich selber, und als wären sie Herren über Leib und Leben.

35. Wie denn die jetzige Welt anders nichts thut, denn daß Oberherr und Unterthan, Groß- und Kleinhaus, Geld und Gut sammeln, und in GOttes Verachtung stecken, denn Mammon als ihrem Gott dienen, und in leibiger Sicherheit liegen [2 Cor. 4, 4.]. Denn diesen Dank und Lohn hat GOtt davon, wenn er der Welt alles genug und überflüssig gibt, daß man denn nach GOtt nichts fragt; niemand glaubt an ihn, oder vertraut ihm, ruft und betet ihn an, ja, vergessen seiner. Wir haben einen andern Gott, nämlich, unsere Arbeit und unser Handwerk. Ein jeder gedenkt: Ich bin ein reicher Bürger, ein reicher Bauer, ich bedarf niemandes. Darum mag ihm GOtt wohl ein ander Volk schicken.\footnote{Das wird wohl meinen: Darum mag GOtt sich ein ander Volk zurichten.}

36. Wenn GOtt ein Land segnet und reich macht, so vergißt man alsbald seiner. Also thun Bürger, Bauern und Edelleute, ja, jedermann; sie sind über Einen Leisten alle geschlagen, und thun alle nach der Regel, die da steht im Liede Mosis [5 Mos. 32, 15.]: Dilectus meus incrassatus et impinguatus, recalcitravit, da das Volk fett war, da löcketen sie wider GOtt. Das ist die
\quellseite{918}
Danksagung für die göttlichen Wohlthaten gegen uns, und die feine, löbliche, ehrliche Tugend, daß, wenn man einen aus dem Schlamm und Dreck gezogen hat, daß derselbige ihm (mit Urlaub) darnach wieder aufs Maul hofiret, oder wieder in den Dreck stößt. Also vergißt man GOttes und des Nächsten, man hängt allein auf den Creaturen und Gaben Gottes.

37. Aber GOtt wird diese Undankbarkeit auch rächen und strafen, daß allerlei GOttes Plagen darauf folgen. Denn, vergessen wir GOttes, unseres HErrn, so müssen wir verderben; sind wir verhärtet und verstockt, so hat er höllisch Feuer, und schwere Hämmer genug, und also viel, daß er dir das Herz wohl kann weich machen und schmelzen [Jer. 23, 29.]. Sonst will GOtt Vater und unser gnädiger HErr sein, und mit uns handeln, als ein Vater seinen Sohn oder Kind hält. Er will dein getreuer lieber GOtt sein, allein, verachte ihn nicht, und sei nicht geizig, vergiß sein nicht [Hebr. 13, 5.]. Denn wo du das thun würdest, als gnädig, freundlich und väterlich er sonst gewesen, also zornig und greulich würde er sich gegen dir erzeigen [Sir. 5, 7.]. Er würde dich durch den Teufel strafen lassen, und führen in Irrthum, auch Leib, Gut und Seele an dir verderben, und durch des Teufels List und Betrug vertilgen.

38. Also warnt und bräuet Mosis, auf daß ein jeglicher sich vorsehe vor ihm selber, und daß er ihm selber nicht traue, und GOttes (wenn es ihm wohl geht) nicht vergesse, oder auf seine Kraft und Macht trotze. Denn wenn GOtt den Mangel aufhebt, und Güter gibt, und man wird GOtt undankbar, so folgt bald darauf GOttes Zorn, Pestilenz, theure Zeit und Krieg. Welcher nur dahin trachtet, wie er den Pfennig suchen soll, so vergißt er GOttes bald, und wird Gottes Feind.

39. Darum halte sich ein jeder hieher, und bete fleißig, erkenne die empfangenen Gaben, daß sein Häuslein, Höflein, Acker, Weib und Kindlein, Geld und Gut und Vieh Gottes sei, daß er es gegeben habe, und nicht von seinen Händen herkomme; durch seine Kräfte hat er es nicht erlangt. Also bleibe in der Furcht GOttes, daß du sagest: Lieber Vater, es ist dein; meine Hände haben es nicht erworben, deine Zusagung und Verheißung hat es gegeben. Ich sehe solches wohl mit meinen Augen, aber ich habe es von dir, und nicht durch meine Hände.

40. Aber bei den Gottlosen ist keine Furcht GOttes, da dankt man GOtt nicht, sondern der Abgott, Vermessenheit, ist allda, daß man kein Vertrauen auf GOtt hat; ja, man gedenkt: Da habe ich meinen Schatz, mein Korn und mein Bier im Hause. Item, mein Handwerk und Ackerwerk, das habe ich erworben; ein anderer schicke es ihm auch, ich will mich darauf wohl ernähren [Luc. 12, 19.]. Das ist die Danksagung und Lob oder Vertrauen GOttes. So sammelt man auch Schätze mit des Nächsten Schaden, trotzt und pocht denn auf die Güter, jauchzt darüber, gedenkt GOttes nicht, hilft andern Leuten auch nicht, stolzirt, raubt, stiehlt, lügt und betrügt andere, und meint, seine Habe und Güter habe er von seiner sauern Arbeit.

41. Aber ein Gottloser denkt: Was ich von Gütern habe, das habe ich nicht von mir; ich habe mein Leben auch nicht von den Eltern, wenn nicht ein anderer Werkmeister dazu gekommen wäre; denn vor hundert Jahren hatte ich weder Haut noch Haar. Aber, lieber GOtt, es ist alles dein, mein Reichthum, meine Ehre und Kraft. Und dieweil GOtt so viel Gutes an mir beweist, er verheißt mit Worten, und bezeuget es mit der That, daß er mein GOtt sei, derhalben sollte ich sagen: Ich will ihn nicht erzürnen, ich will thun, was er haben will; warum wollte ich stehlen? warum wollte ich meinen Nächsten betrügen? GOtt hat noch mehr, denn er je vergeben hat; warum wollte ich häßlich mit den Leuten umgehen? Ich werde genug haben; denn er sagt allhier, er will meine Kraft, mein Reichthum, mein Geld und Gut sein, und will mich ernähren. Er will mein Procurator und Speisemeister sein, gleichwie er in der Wüste, die vierzig Jahre über, der Kinder von Israel Procurator gewesen ist. Der GOtt thut es allein. Ein solcher Mensch muß fromm sein inwendig und auswendig; GOtt hat ihn auch lieb, und er ist anderen nütze, und lebt ihm selber allein nicht.

42. Aus diesem vielfältigen Innhalten könnt ihr sehen, was Mosis für ein reicher Prediger ist. Er treibt fleißig das erste Gebot, das Herz vor GOtt anzurichten, daß [es] ihm glaube, und vertraue auf sein Wort und Verkündigung. Also sollten wir auch ansehen sein Gebot, seine Verheißung, und diese beschriebenen Thaten oder Exempel.

43. Ueber das alles sollten wir Christen gedenken
\quellseite{919}
, daß GOtt seinen Sohn, den HErrn Christum Jesum, in die Welt geschickt hat, und lassen für uns kreuzigen. Davon sollte man täglich predigen, denn dies Werk ist höher denn aller Welt Schätze und Güter. Aber wir lassen es auch vorüber rauschen, und fragen nichts darnach, so doch diese Liebe GOttes, in Christo uns bewiesen, das reichlichste Werk und höchste Exempel der Gnade und Barmherzigkeit GOttes gegen uns ist, das Himmel und Erde erfüllt. Denn wir haben im neuen Testamente viel ein größeres Liebeszeichen GOttes gegen uns, denn jene im alten Testamente.

\begin{schriftzitat}
V. 11. So bitte dich nun, daß du des HErrn, deines GOttes, nicht vergessest, damit, daß du seine Gebote, und Gesetze und Rechte, die ich dir heute gebiete, nicht haltest etc.
\end{schriftzitat}

44. Moses treibt ganz fleißig das erste Gebot, und spricht: Ob wohl die Welt darüber möchte thöricht werden, dennoch alles, was wir haben, das ist deß, der da sagt, er sei unser GOtt. Hat man es nicht alles von GOtt? Denn in dem, daß er seine Güter und Gaben in die ganze Welt ausgestreuet hat, zeigt er an, daß dies erste Gebot erfüllt sei. Aber wir gebrauchen sie also, gleich als hätten wir sie nicht von GOtt. So will Moses auch nicht nachlassen, solches einzubläuen, und dies unebene Lieblein uns zu singen, nämlich: Hüte dich, daß du nicht denkest, du habest es durch dich selber, dein Handwerk, deine Sorge, dein Vieh, sondern wisse, du habest es von GOtt, aus seiner Hand, Gnade und Barmherzigkeit. Darum siehe dich wohl vor, dieser Sünde halben, und vor dem Abgott, Vertrauenheit; daß du nicht meinst, du hättest alles von dir. Nein, du hast es nicht durch deinen Witz, hast es auch nicht durch deine Mühe und Arbeit, Pferde und Kühe; sondern durch GOttes Wort wirst du ernährt, und darum wird dir es gegeben, nicht, daß du arbeitest, sondern daß er dein GOtt sein will. Das erkenne, und gebrauche des Seinen mit Danksagung, so wird er dich mehr segnen; wo nicht, so wirst du mit deinem Gut kein Glück noch Heil haben.

45. Summa, alle Güter auf Erden sind Werke des ersten Gebots, dadurch GOtt beweist, er sei unser GOtt. Ob du gleich GOtt dafür nicht lobst, noch ihn erkennst, so kann ich es, will Moses sagen, doch nicht lassen, ich muß dir es in die Ohren treiben, und sagen, wer dir Gutes thue; du erkennest nun solches, oder nicht, wird zu seiner Zeit\footnote{„Zeit“ fehlt im Original.} mit dir darum auch geredet werden. Darum, wenn du es erkennest, daß du es von mir hättest, so würdest du es allenthalben genießen; wenn du es aber nicht erkennst, und nicht darum dankst, so wirst du es wiederum entgelten müssen. Denn wenn GOtt nicht mehr dein GOtt ist, so wirst du auch nicht mehr für den Seinen gehalten sein.

46. Also spricht denn GOtt: Ich will mein erstes Gebot auch aufheben, ich bin lange genug dein GOtt gewesen; willst du mich aber nun nicht mehr für deinen GOtt halten, so sollst du auch nicht mehr mein Volk sein; denn so kommt der Türke, oder sonst andere Staupen, dieselben lehren dich recht Mores. Denn wer nicht glaubt, daß alles, was wir haben und heißen, von GOtt komme und gegeben werde, sondern er vertraut den Creaturen, dem wird GOtt die Kolbe laufen, mit seinem großen Schaden.

47. Derhalben so lernet nun, dieweil ihr Zeit und Raum habt, daß ihr zum Verstande und Erfahrung des ersten Gebots kommen möget, und in seinem Wege GOttes vergesset, von dem ihr alles genug habt; auch nicht verzweifelt, wenn ihr nichts habt; denn bei ihm allein findet man alle Schätze, und da sollen wir, was einem jeden mangelt, suchen. So sollen wir nun ihn auch fürchten, wenn etwas vorhanden ist, und stehe ein Gottseliger fest auf den Verheißungen und den Exempeln, und wisse, daß GOtt ihn lieb habe, und werde ihn nicht verlassen.

48. Das ist die lange und herrliche Auslegung des ersten Gebots, da Moses auch erzählt hat alle Hinderungen, so uns vom ersten Gebot abziehen. Denn am ersten Gebot ist es alles gelegen, daß ein Mensch auf nichts achte, und lasse ihm außer GOtt nichts gelieben; alsdann kann derselbige nicht verhindert werden. Denn GOtt ist auch unser GOtt, und ein HErr und GOtt des alten und neuen Testamentes; derhalben so wird er uns auch versorgen hier in diesem und auch im zukünftigen Leben, Amen.

\quellseite{920}
\section{Das neunte Capitel}

\begin{schriftzitat}
V. 1. Höre, Israel, du wirst heute über den Jordan gehen, daß du einkommest einzunehmen die Völker, die größer und stärker sind, denn du, große Städte, vermauert bis in den Himmel etc.
\end{schriftzitat}

1. Im Anfang dieses fünften Buchs Mosis (Vorrede § 2. 3.) habe ich gesagt, daß es nichts Anderes sei, denn eine schöne, köstliche Predigt, oder Auslegung der zehn Gebote GOttes, und sonderlich des ersten und vornehmsten Gebots, welches er durch viel Capitel in die Länge und Breite ausstreicht, mit hohem Fleiß und vortrefflichem Ernst handelt und treibt, als ein theurer Meister und Prediger, auf daß man je lerne, die zehn Gebote theuer und groß achten, und sehe, wie keine bessere, nützlichere oder nöthigere Lehre und Predigt, auch keine höhere Weisheit oder Kunst auf Erden zu finden sei.

2. Nun hat er bisher, das erste Gebot wohl zu erklären und einzubilden, mancherlei Anfechtung und Hinderniß erzählt, so dabider laufen in der Welt auf allen Seiten, nämlich der große Abgott Mammon, Geld und Gut, Freundschaft und Gunst, Gewalt und Ehre, und dergleichen mehr; auf daß er uns warne, daß wir uns wohl vorsehen und hüten, daß niemand sein Herz anderswohin setze denn auf das Wort, das er uns gegeben und in dies Gebot gefaßt hat, er wolle unser Gott sein. Denn es liegt alle Macht daran, daß die Zuversicht oder Vertrauen des Herzens rein bleibe, und niemand nichts Anderes sehe noch wisse, denn was sein GOtt sagt, redet oder gebeut.

3. Aber unter allen Hindernissen oder Aergernissen und Abgöttereien ist keine größere noch schädlichere denn die, so man heißt eigene Gerechtigkeit oder Heiligkeit, dadurch man auf unser Werk und Verdienst trotzt. Doch ist es nicht ein grober, gemeiner Abgott, wie die andern, so in der Welt regieren unter dem Pöbel, sondern der allersubtilste, schönste und schalkhaftigste Teufel, der allein die höchsten, feinsten Leute bezaubert, allermeist aber die rechten Christen plagt und anficht, auch so fest anhängt und klebt, daß sein niemand kann los werden, so lange wir dies Fleisch am Halse tragen.

4. Denn ich habe oft gesagt, wie uns natürlich angeboren ist das leibige, schändliche Uebel, daß man Heiligkeit in Werken sucht und vermißt bei GOtt zu verdienen, daß auch die allerhöchsten Christen ihre größte Mühe und Arbeit damit haben. St. Paulus, achte ich, wird fast der allergößten Heiligen einer sein, dergleichen auch der Prophet David; aber wenn sie aufs höchste gekommen sind, haben sie vor diesem Abgott nicht können sicher sein, wie sie selber immerdar klagen. Denn Paulus bekennt frei, GOtt habe ihm müssen einen Engel des Satans aus der Hölle heraus schicken, der ihn spießete und recht wohl zermarterte, allein darum, daß er sich der großen Erleuchtung nicht überhöbe, als wäre er deshalb besser und heiliger vor GOtt, denn andere [2 Cor. 12, 7.].

5. Siehe, der treffliche Mann im Geist und Glauben steht in so großer Fahr, daß er möchte in solche Hoffart fallen, daß er sich selber kitzelte, und vor GOtt käme mit seiner eigenen Heiligkeit, und mit solchem Register: So viel habe ich gethan, gelitten, gepredigt, bekehrt etc., daß kein Rath war, denn Unglück zu mehren und zu widerstehen, denn damit, daß er einen solchen Teufel am Halse mußte haben, der ihn so plagte, daß er mußte in der Erkenntniß bleiben, daß, was er wäre, wäre aus lauter Gnaden, um keines Werks noch Verdienstes willen. Also war das Loch, sich selber zu brüsten, verrannt.

6. Das bleibt nun aller Christen Klage bis an jüngsten Tag, wie es so sauer und schwer wird, dieser Anfechtung zu widerstehen. Darum hat auch Moses diesen Abgott nicht umsonst zuletzt gesetzt, wie er denn auch der letzte ist. Und wenn gleich die andern alle überwunden und gedämpft sind, so bleibt doch dieser ungetödtet bis in die Grube. Denn, obwohl jene auch gewaltig regieren, kann man sie doch alle schwächen und stürzen: den Mammon mit Armuth, Ehre und Gewalt mit Schande und Unglück etc. So kann man auch noch Weisheit zu Thorheit und Schanden machen. Aber wo die Heiligkeit regiert, läßt sie sich mit keiner menschlichen Kunst oder Kräften schwächen noch dämpfen. Denn, ob man sie gleich will angreifen und verdammen, so rühmt sie sich nur desto mehr ihres Leidens
\quellseite{921}
und Unschuld, macht sich zur Märtyrin um GOttes und der Gerechtigkeit willen. Summa, es ist unsere alte, dicke, zähe und dehnische\footnote{dehnisch = dehnbar.} Haut, die sich nicht will abziehen lassen, wie man sich mit ihr bläuet und schlägt.

7. Darüber haben auch etliche fromme Väter geklagt in der Wüste, und einer unter ihnen hat ein Gleichniß gegeben, und gesagt: es gemahne ihn dieses Abgotts wie einer Zwiebel, welche, so man eine Haut abgeschält, hat sie immer eine andere darunter, so lange etwas davon bleibt; also ist dieser Unflat auch; wenn er gleich in Einem Stück einmal oder zwei überwunden ist, daß GOttes Gnade den Ruhm und die Ehre behält, so kommt er doch immerdar wieder, klebt und hängt sich an alle Werke, daß niemand davor kann unbeschmeißt bleiben, er habe denn auch irgend einen Teufel, der ihn mit Füßen trete, daß er seine Heiligkeit könne aufbringen.

8. Darum ist es eine hohe Predigt um das erste Gebot, sonderlich in diesem Stücke, daß ein Mensch dahin komme, daß er sich seiner Heiligkeit in Werken vor GOtt anmaße. Wiewohl es für den gemeinen Haufen, hoffe ich, nicht groß noth ist, denn er steht nicht nach großer Frömmigkeit oder Heiligkeit, bleibt bei keinem Abgott Mammon, und andern gemeinen Götzen, daß sie so hoch nicht kommen können. Wenig sind ihrer, die auf große Kunst oder Weisheit trotzen, welches auch subtil und hoch ist; aber keiner thut so großen Schaden, als dieser; denn er will nirgend denn bei den Allerbesten sein, und weil er der schönste ist, sucht er auch die schönste Kirche, die andern läßt er wohl zufrieden. So hüten sie sich auch fleißig genug davor; denn wer nicht gute Werke thut, kriegt keine große Anfechtung davon.

9. Das ist nun der rechte Teufel und große Greuel, bisher unter dem Papstthum aufs allerhöchste geehrt, in allen Kirchen, Stiften und Klöstern, auf allen Kanzeln und hohen Schulen ausgeschrien, dazu so viel Gottesdienst gestiftet und aufgerichtet, und der Welt Güter mit Haufen geschlagen sind; ja, so weit eingerissen, daß kein reicher Mann gestorben ist anders, denn der Meinung, weil er nichts Besseres hinter sich lassen, auch nichts Anderes mitnehmen möchte denn gute Werke und Gottesdienst, müßte er sein Testament machen, und so viel zur Kirche bescheiden, damit GOtt gedient und seiner gedacht würde. Summa, der Abgott ist zeitlich angegangen, ja, der erste gewesen im Anfang der Welt; aber sonderlich im Schwange gangen und öffentlich aufgeworfen als der höchste Gottesdienst in der Christenheit, nun länger denn tausend Jahr, da man nichts Anderes gelehrt noch gewußt hat, denn unsere Zuversicht und Seligkeit stellen auf eigene Werke.

10. Darum kommt nun Moses zuvor, daß er diesen Abgott stürze und verdamme, braucht allerlei Weise, daß er keine Juden davon schrecke und reiße; wiewohl es doch auch wenig bei ihnen geholfen hat; denn sie hernach doch so tief in die Abgötterei gefallen sind, als wir bisher sind gewesen. Erstlich aber fähret er an, sie zu warnen, wenn sie ins Land kommen, daß sie nicht zufahren und sagen, sie haben es von GOtt verdient, als wäre ihre Heiligkeit so groß gewesen, daß er ihnen das Land hätte müssen geben, und jene heraustreiben und vertilgen.

11. Merke aber, wie eben und meisterlich er ihre Schalkheit und Heuchelei trifft und abmalt, damit sie einen Schein der Heiligkeit machen. Denn das thun sie dennoch; die frommen Schälke, daß sie GOtt die Ehre geben, daß er es gethan habe; sind nicht so grob, daß sie ihnen selber das Werk zumessen, aber im Herzen spielen sie GOtt dieses mit: Er muß sich von ihnen lassen zum Narren und Götzen machen, und den Namen führen, daß er ihre Gerechtigkeit ansehen müsse, und das Land darum geben, will er anders seine Ehre behalten. Sie aber wollen das Beste thun, daß die Schuld und Ursache des Werks nicht GOttes, sondern ihres Verdienstes sei, und soll dennoch GOttes Ehre heißen, wie gesagt, daß er so fromm sei und gebe den Frommen, was sie verdienen, so es doch wahrhaftig nichts Anderes ist, denn GOtt seine Ehre geraubt, und ihnen selber zugemessen. Denn GOttes Ehre ist allein die, daß er aus lauter Gnade jedermann gebe, helfe, wohlthue und selig mache.

12. Das ist nun gewißlich der eigenen Gerechtigkeit und aller falschen Heiligen Eigenschaft, daß sie große Ehre GOttes vorgeben, und sich trefflich demüthig stellen, sagen, sie sind arme Sünder, GOttes Gnade nicht werth, wollen nicht heilig noch fromm heißen etc., ist aber im Grund nichts anders, denn eine verzweifelte, zweifältige Hoffahrt und Vermessenheit, so nicht
\quellseite{922}
eher vor GOtt will kommen, sie sind denn zuvor rein und heilig, noch seine Gnade annehmen, sie haben es denn vorhin verdient, daß nicht er, sondern sie den Ruhm haben. [Sie] will nichts von ihm geschenkt nehmen, sondern ihm zuvor geben, daß er ihr wieder vergelten müsse, als der unser Schuldner sei. Das heißt aus GOtt einen Treudler gemacht, der unserer Heiligkeit und Verdienst nicht entrathen\footnote{Im Original: gerathen.} könne, und keinen Dank dazu habe, daß er uns etwas gebe, es sei denn wohl verdient. Siehe, das ist eine Tugend der löblichen Frömmigkeit.

13. Die andere ist eben derselbigen gleich, daß sie vor GOtt kommt, und schmückt sich mit einem Werke oder zwei, die sie aufbracht hat, und will damit zudecken alle Sünde und Schande, so sie je begangen hat. Als wir unter dem Papstthum gesehen haben: Wenn jemand dreißig oder vierzig Jahr gelebt hat in eitel Büberei und Untugend, und darnach niedergefallen, ein Testament gemacht, oder so viel Messen gestiftet, denn sollte Gott willkommen heißen, und alle bösen Stücke, sein Lebenlang begangen, mit einem Gulden oder zehn lassen bezahlt sein. Ist das nicht eine teuflische Vermessenheit, ja, ein unleidlicher Hohn und Spott der göttlichen Majestät, daß er, alle Bosheit hintan gesetzt und geschwiegen, um eines Guldens willen sollte alles lassen geschenkt und geschlichtet sein, und den Himmel dazu geben? Ja, er wird ihnen einen Himmel geben, darinnen die Engel mit Keulen laufen (wie die Leute sagen).

14. Solche schändliche Untugend der falschen Heiligkeit will Moses auch aufdecken, wie sie unter einem kleinen Schein unzählige Schande und Laster birgt, und dennoch unverschämt darf vor GOtt treten, Rechnung mit ihm machen, und rühmen, als sei eitel Heiligkeit da, daß er keine Sünde mehr sehen, und alles Zornes vergessen müsse. Darum fährt er hernach fort, erzählt und hält ihnen vor, daß sie ansehen sollen, was sie wider GOtt gethan haben die ganzen vierzig Jahre über, ehe sie ins Land kommen sind. Als sollte er sagen: Lieber, wenn es soll Rechnens und Gleichaufhebens gelten mit GOtt, so mußt du ja so viel Werke dagegen legen. Nun hältst du ihn für einen solchen Pozen,\footnote{Pozen = Vogelscheuche, Popanz.} der deine Untugend nicht sehen, oder je dazu schweigen solle, und deines ohnmächtigen Geldes froh werden; ja, auch nicht dagegen rechnen die Güter, die er dir gegeben hat, Leib und Leben, Haus und Hof, Sonne und Tag, und alles, was Himmel und Erde vermag oder trägt; Summa, aller seiner Wohltat schweigen, welcher keine so gering ist, sie überwöge alle deine Werke, so du ihm geben kannst.

15. Sprichst du aber: Ja, GOtt ist gnädig, fordert und rechnet nicht so strenge, sondern verschont und schenkt, wie hoch wir ihn vorhin erzürnt haben. Antwort: So müsse dich auch ein gut Jahr angehen! Kannst du nun glauben, daß er dir umsonst vergebe so viel Sünde und Untugend, darinnen du so lange gelegen bist, dazu dir Gott allerlei Wohltat beweiset, der du nie keiner bist werth worden, und meinst, daß er dich nicht könne oder wolle gen Himmel nehmen ohne deine zehn Gulden, oder was du ihm geben magst? Was ist das anders, denn daß du entweder ein Narr bist, daß du dich vermissest, mit einem Gulden unzählige Sünde abzulegen, so du nie kein gut Werk gethan hast bis in den Tod, oder GOtt, als ein verzeihender Bube, für einen Narren hältst, der sich wie ein Kind mit einem Zahlpfennig äffen oder täuschen lasse, daß er deine Bosheit nicht sehe?

16. Das ist nun alles der bösen Natur Blindheit, daß, wo sie nur einen Flecken sieht, der ein wenig rein und weiß ist, wenn sie einmal ein gut Werk aufbringt, das sie wohlgethan heißt, kann sie ihr es so nütze machen, sich spiegeln und kitzeln, daß sie GOtt solle für den größten Heiligen ansehen. Welches ist eben die Klugheit, die der Strauß an sich hat, welcher, wo er unter einen Strauch kommt, daß er den Kopf bedecken kann, meint, er habe den ganzen Leib verborgen, daß man ihn nicht sehen könne, und wie die jungen Kinder die Augen zuthun, und meinen, wenn sie nicht sehen, solle sie auch niemand sehen.

17. Darum will nun Moses lehren und schließen: „Wer vor GOtt kommen will, der nehme ihm nur die Weise nicht vor, daß er seine Werke hervorbringe, und auf Rechten troße, sondern ein ander Register ansehe, und rechne, wie viel ihm GOtt gegeben und gethan habe ohne Verdienst, auf daß er lerne, daß alle unsere Gerechtigkeit nichts und verdammt sei, und keine andere Weise oder Weg zur Gerechtigkeit sei, denn solches wissen und bekennen, daß wir vor
\quellseite{923}
GOtt nicht fromm noch gerecht sind, sondern was er uns gibt, aus lauter Gnade empfangen müssen, kein Werk dafür gefordert noch angesehen. Wie der Prophet im Psalter, Ps. 143, 2., spricht: „Gehe nicht ins Gericht mit deinem Knechte, denn vor dir ist kein lebendiger Mensch gerecht.“ Item, Ps. 130, 4.: „Bei dir ist nichts denn Vergebung, daß man dich fürchte.“ Ein solch Herz, das sich selbst verleugnet, keine Frömmigkeit noch Verdienst bei sich findet, ist die rechte Frömmigkeit vor GOtt. Und um solches Glaubens und Bekenntnisses willen heißt es fromm und gerecht, daß es\footnote{Im Original und in der Erlanger: er.} frei heraus beichtet: Lieber HErr, ich weiß von keiner Frömmigkeit in meinem Leibe und in meinem Leben, sondern das ist mein Trost und Trotz vor dir, daß du einem armen Sünder gerne gibst, und alle Sünde verzeihest aus lauter bloßer Gnade.

18. Solches hat uns auch Christus selber vorgemalt im Evangelio Lucä am achtzehnten, V. 11. 12. Der Pharisäer kommt und tritt vor GOtt mit seinen Werken: „Ich danke dir, HErr, daß ich nicht bin wie andere Leute, auch nicht wie dieser Zöllner; ich faste zwier in der Woche, und gebe den Zehnten von allem, das ich habe.“ O ein trefflicher Heiliger, den niemand strafen konnte, ohne mit dem schändlichen Abgott der\footnote{Im Original und der Erlanger: „und“ statt: der. Vgl. im vorigen Capitel § 40 und § 44.} Vermessenheit, daß er kommt getrollt, rühmt, und dankt nicht von der Gnade, sondern von seinem Fasten und guten Werken, als dürfte er keiner Gnade.

19. Daneben steht jener arme Zöllner, hat kein Werk noch Verdienst, weiß nichts aufzubringen, das GOtt ansehen solle, ohne, daß er frei heraus bekennt, er sei ein Sünder, komme eben darum, und bitte Gnade, suche und begehre nichts, denn von GOtt zu nehmen. Ueber diesen spricht Christus das Urtheil, daß er gerecht und gottgefällig sei heimgegangen um dieses Glaubens willen; der andere aber, mit seiner Heiligkeit und vielen Werken, sei verdammt, und nicht werth geachtet, dem Zöllner das Wasser zu reichen.

20. Siehe, daran fehlt es nun uns allen, gleichwie es den Juden gefehlt hat. Denn das ist, wie gesagt, unsere alte Haut, durch Fleisch und Blut, Mark und Bein gewachsen; denn ein jeglicher ja gerne wollte etwas aufbringen, das vor GOtt gelten müßte, darauf er möchte ruhen und fußen, und rühmen, daß [es] sein eigen wäre, und nicht müßte nackt und bloß vor GOtt stehen, und seine Schande vor ihn bringen [1 Mos. 3, 11.]. Darum sind alle geistlichen Orden aufgebracht, alle Kirchen und Gottesdienste aufgerichtet, alle Testamente und Seelgeräthe gestiftet, daß jedermann gehofft hat, GOtt soll es ansehen, und so köstlich achten, als sie meinen, daß er uns darüber müßte gen Himmel heben und nehmen.

21. Daher kommt es auch, wenn man solches verwirft, und dawider lehrt, so will sobald folgen, daß jedermann schreiet: Wohlan, so wollen wir keine guten Werke thun; und fahren flugs den Holzweg, entweder GOtt zum Narren und Götzen gemacht mit unsern Werken, oder gar nichts gethan. Denn\footnote{„Dem“ von uns gesetzt statt: „Denn“ im Original.} kann man nicht wehren, hier kann man nicht halten. Auf der Mittelstraße will niemand bleiben, daß wir uns schlecht verleugnen vor Gott, als die nichts suchten, denn lauter angebotene Gnade und Verheißung zu empfangen, darnach aber hingehen, und thäten, was wir könnten, auch frei umsonst, um keines Verdienstes oder eigenen Genieß willen. Man sagt es, predigt es, und hört es wohl, es kann es aber niemand, wenn man es können und beweisen soll. Ich muß für meinen Theil bekennen, und achte, es soll niemand so heilig und voll Geistes sein, er wird es auch müssen bekennen. Denn, weil es alle Heiligen gefühlt und geklagt haben, werden wir's freilich auch nicht überhoben sein. Ich fühle es wohl, wenn es Treffens gilt, wie sich die Natur wehrt, sucht und um sich greift, wollte ja gerne ein Werk finden, das sie könnte GOtt aufrücken, und sprechen: Das habe ich dennoch gethan, so viel gepredigt und ausgerichtet. Und ob ich gleich weiß, daß [es] nicht gilt noch gelten soll, noch kann ich es nicht lassen, noch dahin kommen, daß ich mich so bloß könnte ergeben auf seine Gnade.

22. Ist es aber nicht ein elender Jammer, daß ein Mensch nicht so viel vermögen soll, wenn ihm GOtt selber anbeut, und spricht: Lieber, da hast du zwei Wege vor dir, nimm,
\quellseite{924}
und wähle doch einen: Willst du lieber meine Gnade und ewige Seligkeit umsonst haben, geschenkt und heimgetragen, ohne alle deine Kost und Arbeit; oder mit deinen Werken darnach stehen zu verdienen, und doch nicht erlangen? Ja, [er] vermahnt und lockt dazu zur Gnade, und dräuet dagegen bei ewigen Zorn und Strafe, wo wir es nicht annehmen. Wie soll oder kann er mehr thun, wenn der keines hilft, daß er seine Gnade so reichlich anbeut und ausschüttet, reizt und gebeut, vermahnt und schreckt, dazu straft und drein schlägt? er wolle selber gerne das Recht lassen fahren, und einen Strich dadurch machen: so wollen wir Recht ohne Gnade haben, und dennoch auch ein eigen Abgöttlein mit uns führen.

23. Solch elend Ding ist es um unser Leben, daß wir den Schatz nicht annehmen, daran wir sollten Leib und Leben wenden, keine Kost noch Mühe sparen, und laufen, daß wir möchten Blut schwitzen, wenn es irgend an einem Orte zu finden wäre. Nun er von ihm selber kommt, und den Schatz vor die Thüre bringt, dazu so herzlich und väterlich vermahnt, so ernstlich gebeut, so schrecklich dräuet; so mögen wir sein nicht. Sollten wir uns doch selber anspeien und feind werden, daß wir so heillose Leute sind, die wider die angebotene Gnade auf Recht trotzen, das ist, Zorn und Ungnade auf uns laden.

24. Denn was ist es anders, Recht suchen und fordern, denn so viel gesagt: Gehe mit mir ins Gericht, und gib mir, was ich verdient habe? Summa: Ich darf deiner Gnade nirgend zu; doch daneben aus den Augen gethan alle Sünde und Bosheit, als hättest du GOtt keinmal erzürnt. Siehe, das wollte Moses gerne niederlegen, warnen und wehren, daß man ja eben auf diesen Abgott aufsehe, und sich davor hüte. Darum lerne, wer lernen kann, daß er beileibe kein Verdienst noch eigene Gerechtigkeit suche, ja, mit keinem Werke zu schaffen habe, daß er vor GOtt wolle bringen, sondern laufe und fliehe davor, als vor dem Teufel selber.

25. So spricht du: Wie geht es denn zu, daß so viel Sprüche und Verheißungen in der Schrift stehen, daß die, so fromm sind und gute Werke thun, sollen beide, hier und dort, ihren Lohn empfahen? als [Luc. 6, 38.]: „Gebet, so soll euch wieder gegeben werden“ etc. Antwort: Da liegt es nun an, daß man wisse, was vor GOtt fromm sein und gute Werke thun heiße; davon ich sonst genug gesagt habe; nämlich daß heißt es, nicht, daß du fromm gethan, und sprechest: HErr, das habe ich gethan, schreibe an, und vergilt es; sondern also sagest von Herzen: HErr, ich nehme alle deine Güter, Wohlthat und Gnade als ein Sünder und verzweifelter Mensch, wie ich gehe und stehe des ewigen Zorns und höllischen Feuers werth, wenn du solltest nach Recht und Verdienst mit ihm handeln. Aber ich sehe nicht meine Sünde, noch was ich verdient habe, sondern dein Wort und ernstlich Gebot an, daß du heißest, vermahnest und dräuest, daß niemand sein Werk vor dich bringe, etwas zu verdienen, sondern aus väterlicher Güte Vergebung der Sünde und allerlei Wohlthat empfache, und in der reinen Zuversicht deiner Gnade stehe und bleibe.

26. Wo nun solche Zuversicht im Herzen ist, sollen auch die Werke, so darin geschehen, gut und gottgefällig heißen, und um derselbigen willen die Verheißung erlangen. Denn das ist die Verheißung: Wer das erste Gebot hält, und in der Gnade steht, soll auch sein ganzes Leben, und was er thut, angenehm und wohlgethan sein; denn außer der Gnade ist so bald ein Abgott daraus gemacht, unter falschem Schein guter Werke. Wer aber diesen Verstand hat, kann nicht darauf trauen, noch sie zum Abgott machen, sondern bleibt fein rein und richtig im Glauben, und thut die Werke in und aus der Gnade. Darum heißen sie auch rechtschaffene gute Werke. Siehe, also redet die Schrift davon im 112. Psalm, V. 1. 2.: „Wohl dem, der den HErrn fürchtet, der große Lust hat zu seinen Geboten, des Same wird gewaltig sein auf Erden, das Geschlecht der Frommen wird gesegnet sein.“ Also auch im 128. Psalm, V. 1. 2., sagt der Heilige Geist eben also: „Wohl dem, der den HErrn fürchtet, und auf seinen Wegen gehet. Du wirst dich nähren deiner Hände Arbeit; wohl dir, du hast es gut.“ Als sollte er sagen, an dem Stücke liegt es gar, das muß vor und in allen Werken da sein; denn die Furcht leidet nicht, daß man auf die Werke traue, oder seinen Trost lasse sein. Wo aber nun des Herzens Zuversicht recht ist auf GOttes Gnade, läßt ihm GOtt die Werke, aus solchem Herzen gethan, so wohl gefallen, daß er sie mit allerlei Segen und Wohlthat begnadet, wie die jetzt angezogenen Psalmen anzeigen.

\quellseite{925}
\section{Der andere Theil}

27. Aus dem allen sehen wir nun, warum Moses mit solchem Fleiß und Ernst auf das erste Gebot treibt, und so treulich warnt vor diesem Abgott, als vor dem schädlichsten Gift. Er braucht aber starke Wehre dawider, zeigt und gibt kräftige Arznei, damit man ihr vorkomme und widerstehe.

28. Zum ersten damit, daß er anzeigt, und schließet, Gott habe sie um keiner Heiligkeit, auch nicht um jener Bosheit willen ins Land gebracht, sondern um seiner Verheißung willen, welche er ihnen vorhält, und spricht [V. 5.]: Auf daß er das Wort aufrichte, das der HErr geschworen hat euren Vätern etc., das ist, wie er sonst oft sagt, um seines Namens willen; daher auch der 25. Psalm, V. 11., bittet: „Um deines Namens willen, HErr, sei gnädig meiner Missethat.“ Wo steht nun der Name? Eben in dem ersten Gebot, da er sich verbindet, und der ganzen Welt von sich predigen und anbieten läßt, er wolle unser GOtt sein, und alles Gutes geben. Damit ist je dies abgesagt: Deine Heiligkeit gilt nichts, deine Frömmigkeit verdient nichts, sondern das Gebot, darin sein Name eingeschrieben, seine Verheißung gebunden und gefaßt ist, das gilt, das thut es, das ist die einige Ursache, warum er euch wohl thut und das Land gibt.

29. Willst du nun wissen, woher du habest Leib und Seele, Geld, Gut, Ehre, Kunst, Weisheit, Gewalt etc., so siehe mir dies Gebot an, das wird dir zeigen, daß du nicht dazu geboren, noch durch dein Glück, Klugheit, Arbeit, oder Mühe erworben oder verdient hast, sondern alles daher, daß er verheißen hat, ehe du je etwas gethan hattest, er wolle dein GOtt sein; daher hast du alles, was du vermagst, bis auf den letzten Heller. Solches ist immer noth zu treiben, wie denn hier Moses thut, daß man doch lerne, dies Gebot recht ansehen, und so erkennen und ausbreiten, daß man sich und alle göttlichen Güter darein ziehe und fasse, ja, auch den HErrn Christum, und das ganze neue Testament, welches gar aus dieser Verheißung quillt und fließt, und endlich darinne bleibt, als in dem solche Verheißung erfüllt und bestätigt ist, daß er unser gnädiger Vater sei, und durch Christum alle Sünde vergebe, von Sünde und Tode erlöse, alle seine Güter schenke, und das ewige Leben gebe.

30. Zum andern, daß ihnen Moses solches wohl einbläue, und ihren Hochmuth lege, sähet er an, ihre Legende zu beschreiben, macht ein lang Register, darinnen er erzählt, und ihnen vor die Nase stellt alles, was sie begangen haben, weil sie sind in der Wüste gewesen; auf daß sie ja sehen und greifen sollen, daß ihnen GOtt nichts um ihres Verdienstes, sondern allein seiner verheißenen Gnade willen gebe. Und ist kurz die Summa des ganzen Capitels: Was sollte er dir geben um deiner Frömmigkeit willen? Bist du doch von Anfang bis auf die Stunde ein halsstarrig, widerspenstig Volk gewesen, das ohne Unterlaß GOtt erbittert und erzürnt hat, daß, wo er sollte eurem Verdienste nach handeln, sollte er euch längst alle auf einen Haufen vertilgt haben.

31. Siehe, da steht unser aller Ruhm und Preis, eine schöne, löbliche Tugend, die der elenden Creatur wohl ansteht gegen ihrem GOtt und HErrn. Wenn uns GOtt alles Gutes thut, mehr, denn wir wünschen können, daß wir ihm zu Dank nichts denn den verfluchten Ungehorsam und halsstarrigen Kopf geben. Wie denn die Welt jetzt allenthalben voll ist halsstarriger, unbändiger Leute, Knechte, Mägde, Bauern, Bürger, Adel, wie die groben, krummen und ungelehrten Blöcke\footnote{Im Original: ungelehrten Blöcher.} und Klötze, die nirgend zu dienen, denn im Ziegelofen mit Feuer zu bemältigen. Lasset ihnen nicht sagen, wehren noch steuern, wie man ihnen dräuet mit Feuer, Wasser, Pestilenz und allem Unglück, als hätten sie Eisen und Stahl im Halse, das niemand brechen kann, bis [es] der Henker mit dem Rad entzwei stoße, oder der Tod strecke.

32. Aber viel härter und halsstarriger ist das Volk, so in den Wahn geräth, daß es heilig und GOttes Volk sei, wie sich die Juden rühmen [Matth. 19, 20.]. Da werden Leute aus, so hart und steif, die weder GOtt noch die Welt tragen kann, an welchen alles verloren ist, wie man mit ihnen umgeht, man warne, dräue, bitte, oder flehe. Wie diese Historie genugsam anzeigt und malt, wie sie die vierzig Jahre so oft geplagt und geschlagen worden, und täglich so viel grausamer, schrecklicher Exempel des Zorns GOttes vor Augen hatten, davor ihr Herz sollte zerschmolzen sein; noch war es von Stunden an vergessen und verachtet. Das ist die Frucht der
\quellseite{926}
großen Heiligkeit, durch eigene Werke gesucht und gerühmt, welche, wie oben [§ 15] gesagt, mit einem Gulden oder zwei gegen GOtt pochen, sich brüsten und trotzen darf, und sich ein Werk oder zwei so aufblasen, stolz und steif machen lassen, als müsse er ihre Untugend und Abgötterei zu großem Verdienst rechnen.

33. Siehe, solches will ihnen Moses eintreiben und einbinden, daß sie sich daran spiegeln und immer vor Augen haben sollen; darum treibt er so viel Worte, macht es groß und greulich, daß sie davor sollten erzittern, wie sie von dem Zuge an, da sie aus Egypten gegangen sind, stets ungehorsam und widerspenstig sind gewesen [2 Mos. 32, 8. 9.]. Und zieht sonderlich das Stück an, daß sie gethan hatten, so bald sie aus dem Lande in die Wüste kamen, an dem Berge Sinai, da GOtt mit ihnen einen Bund machte, und die zehn Gebote gab, nämlich, wie sie die Abgötterei anrichteten mit dem güldenen Kalbe, und GOtt so hoch erzürneten, daß er sie alle auf einmal vertilgen wollte, daß auch Moses mußte das Kalb mit Feuer verbrennen, und zu Pulver machen, darnach den Staub aufs Wasser werfen, und ihnen zu trinken geben (wie man liest 2 Mos. 32, 20. ff.), daß sie ihren eigenen Greuel mußten saufen. Siehe, der Sünde will er nicht geschwiegen noch vergessen haben, schreibt sie ins Buch, daß sie [es] etlichemal im Jahr öffentlich lesen und hören mußten, treibt auch eben die Worte: Gedenke es, und vergiß nicht etc. Daß sie es je nicht hinter sich werfen, und unter die Bank legen, noch aus der Acht lassen, als habe es GOtt vergessen, sondern immer vor den Ohren klingen lassen, und nichts denn ihre Sünde und Schande hören.

34. Das lasse ihm nun jedermann auch zum Exempel gesagt sein, ich sowohl als andere, wider diese Anfechtung; denn es ist doch, wie gesagt, kein besserer Rath noch Kunst dawider, denn daß ein jeglicher in seinen Busen greife, da wird er ein Register finden, welches ihm gar viel Anderes wird sagen denn seine Heiligkeit, daß du mußt den Pfauenschwanz niederschlagen, und sagen: HErr, ich will gerne aller Werke und Verdienste vergessen, daß ich nur mag zur Gnade kommen. Wie der arme Zöllner [Luc. 18, 11. ff.], welcher läßt den Pharisäer sein Register lesen, und eitel gute Werke rühmen; er aber weiß kein Werk noch Heiligkeit, bringt aber ein ander Register, das heißt: Mihi peccatori, darin steht nichts denn eitel Schuld und Sünde; solcher Anblick lehrt ihn sein, daß er an ihm verzweifele, keine eigene Heiligkeit vornehme, und vom Verdienst der Werke zu der Gnade laufe.

35. Darum regieret auch GOtt mit seinen Heiligen also, daß sie dennoch immer Sünde am Halse tragen, ja, auch solche grobe Knoten auf ihnen behalten, daß sie müssen in der Furcht und Demuth bleiben, und nicht können zu sicher oder vermessen werden. Christus wollte St. Petrum und Paulum zu hohen Aposteln haben, aber zuvor wohl anlaufen und tief genug fallen lassen, daß dieser verleugnet [Matth. 26, 74.], und abtrünnig an ihn wird; jener verfolgt [Apost. 9, 4.], lästert und verdammt ihn aufs alleräußerste; allein darum, daß er ihnen einen Kittel an den Hals binde, daß sie ihres Unglücks nicht vergessen könnten, sonst wären sie auch in dies verzweifelte Unglück geraten. Nun aber müssen sie hienieden bleiben, daß sie sich keiner Heiligkeit vermessen, und keinen Sünder verachten dürfen. Darum, wer der schändlichen Vermessenheit entlaufen will, der thue auch also, daß er hinter sich sehe, und rechne, wie er sein Leben vorhin zugebracht hat. Wie ich und andere müssen sagen: So lange bin ich ein Bube gewesen, [habe] mit falscher Lehre und Abgötterei unzählig viel Seelen verführt, meinen HErrn täglich gelästert und gekreuzigt etc., so wird er wohl zum Kreuz lernen kriechen, und froh werden, daß er mag als ein Sünder zu Gnaden kommen.

36. Solches lehrt auch das Buch Jesu Sirach am 5., V. 5.: Noli esse sine metu de propitiato, sei nicht sicher und ohne Furcht der Sünden halben, die schon vergeben sind. Warum das? Eben darum, daß hier Moses auch lehrt dem Unglück zu wehren. Denn, so bald man aus der Acht läßt, und anfähet zu vergessen, was wir gewesen sind, so ist auch vergessen der Gnade, die uns geschenkt ist. Weil man das Unglück nicht mehr ansieht, das uns sollte demüthigen, daraus folgt so bald, daß das Evangelium nimmer schmeckt, und darnach allerlei greuliche Untugend mit Haufen wieder eingehen, Unzucht, Stolz und Hochmuth, Verachtung des Nächsten etc., und wird hernach vielmehr ärger, wie Christus [Matth. 12, 45. Luc. 11, 26.] sagt, denn es zuvor war; wie wir jetzt täglich allzu viel sehen und erfahren.

37. Das hat St. Petrus auch wohl gesehen,
\quellseite{927}
da er von solchen sagt 2 Petr. 1, 9.:\footnote{Im Original: „2 Pet. 2.“} „Bei welchem solches nicht ist, der ist blind, und siehet nichts, und vergisset der Reinigung seiner vorigen Sünde“ etc., das ist, es wird ein solcher Mensch daraus, der dahingehet, als der starrblind wird, und mit offenen Augen nichts mehr sieht, noch achtet oder fühlt, wird sicher und vermessen, als hätte er nie kein Wasser betrübt,\footnote{betrübt = getrübt.} und könne nicht mehr irren. Es heißt aber also: Daß deine Sünden nimmer Sünden und Vergebung sind, ist nicht deiner gelben Haare Schuld, sondern der Gnade und Vergebung; wenn du aber so vergessen wirst, und nicht solches siehst noch achtest, wie viel es GOtt gestanden, und wie theuren Schatz er an dich gewandt hat, [10] gibst du wiederum dem Teufel Raum, daß er dich stürze in Blindheit, und allerlei Sünde und Schande. Daher auch jetzt so viel unserer Rotten von dem Evangelio gefallen sind, allein darum, daß sie sich dünken lassen, sie haben es nun gar, sind satt und sicher, und nicht mehr denken, daß sie zuvor auch in der Hölle und tiefen Blindheit gelegen sind, oder vor dem Netz und Garn gewesen.

38. Darum, auf daß wir auch nicht in dies verfluchte Uebel gerathen, lasset uns mit Fleiß an dem ersten Gebot lernen, und stets damit umgehen. Lernen, sage ich, also, daß wir Schüler bleiben unser Lebenlang, und uns hüten vor dem Ueberdruß und schädlichen Wahn, als hätten wir genug gehört, und könnten es nun wohl; denn wir haben doch keine Ruhe, auch vor allen andern Abgöttern und Aufsetzungen, aber viel weniger vor diesem Eigendünkel der Gerechtigkeit. Ich bekenne mich dazu für der meisten einen, denn ich habe ja Fleiß gethan, so thut auch GOtt große Gnade, und hilft getrost wehren, noch kann ich es nicht dahin bringen, daß ich des Götzen los würde, und muß, so lange ich lebe, des ersten Gebots Schüler bleiben. Alle anderen Leute auf Erden können es allzu wohl, nicht allein die zehn Gebote, sondern noch weit darüber, so viel Orden, und Stände, und Werke zur Uebermaß; meinen, der Heilige Geist habe viel höher Ding müssen lehren und setzen in Conciliis, als wären sie weit über alle zehn Gebote gefahren, so doch nie keiner gewesen, oder noch ist, der das geringste Gebot recht angesehen habe, oder das wenigste Theil verstehe.

\begin{schriftzitat}
V. 8–10. Denn im Horeb erzürnet ihr den HErrn, also, daß er vor Zorn euch vertilgen wollte, da ich auf den Berg gegangen war, die steinernen Tafeln zu empfangen, die Tafeln des Bundes, den der HErr mit euch machte, und ich vierzig Tage und Nächte auf dem Berge blieb, und kein Brod aß, und kein Wasser trank, und mir der HErr die zwo steinernen Tafeln gab, mit dem Finger GOttes geschrieben, und darauf alle Worte, wie der HErr mit euch aus dem Feuer auf dem Berge geredt hatte, am Tage der Versammlung.
\end{schriftzitat}

39. Es ist euch zum öfternmal angezeigt, wie Mosis hat viele Worte gebraucht in der Auslegung des ersten Gebots, welches er treibt überaus fleißig, gleich als ein Schulmeister seinen Schülern die Lection mit großem Fleiß vorliest und einbläuet, auf daß sie es wohl fassen und lernen; oder, wie eine Mutter ihrem Kinde den Brei einstreicht, daß es den zu sich nehme, und ernährt werde; also will er, daß das Volk Israel durch diese vielen Worte und Predigten das erste Gebot auch wohl einbilden, fleißig lernen, und darnach leben sollen. Und wollte sie gerne bringen zu ihrer selber Erkenntniß, daß sie sich bei der Nase ziehen sollten, und hinter sich gedenken, wie sie es in der Wüste getrieben hätten, spricht: Ich habe um euretwillen vierzig Tage lang in der Wüste gefastet, große Arbeit und Mühe gehabt, daß ich euch die Gebote GOttes holete und brächte vom Berge Sinai, welche Mühe ich verlieren mußte; diese Arbeit war umsonst gethan, denn ihr erzürnetet unsern Herrn GOtt mit dem güldenen Kalbe.

40. Diese Tugend liest er ihnen vor in diesem Register, rückt ihnen allezeit auf, daß sie ein halsstarrig und hartnäckig Volk, abgöttisch und ungehorsam sind gewesen, und daß es ihn viel Mühe gekostet habe, sie bei Gott einzuheitdingen,\footnote{einheitdingen = durch Fürsprache wieder zu Gnaden bringen.} und wieder zu Gnaden zu bringen. Dies Register sollen sie lesen, in diesen Spiegel sollen sie sehen. Er legt es ihnen geschrieben vor, auf daß sie es nicht vergessen sollten, sondern durch ihre Legende gedemüthigt würden. Und müssen ihre Sünde und Laster ihnen nun dazu dienen, daß sie hübscher fromm sein möchten.

41. Wir sollen auch unsere begangenen Missethaten nicht aus dem Gedächtniß allerdings kommen
\quellseite{928}
lassen, sondern zu unserer Demüthigung daran gedenken. Also pflege ich ihn auch zu thun. Denn ich, Martinus Luther, bin auch lange ein Schalk gewesen, und im Klosterleben gesteckt, habe die Leute verführt, und ich kann die Seelen, die ich verführt habe, nicht bezahlen. Dies Register und diese Untugend muß ich noch stets vor Augen haben, auf daß sie mit predigen von meinen guten Werken und eigener Gerechtigkeit, darauf ich im Papstthum mich auch verlassen habe, und nun daran gedenken möge, daß ich GOtt nicht mehr schände und lästere. Und wiewohl ich jetzt kein Lästerer noch Schänder GOttes mehr bin, denn ich verführe mit meiner Predigt das Volk nicht, so opfere und kreuzige ich auch Christum nicht mehr in der Messe, sondern ich preise und lobe die Barmherzigkeit GOttes; dennoch, so ist mein Leben also gethan, daß ich GOttes Gnade und Barmherzigkeit bedarf, mein Leben ist unvollkömmlich, und meine Werke also untüchtig, daß ich darauf gegen GOtt nicht darf trotzen, noch darum den Himmel bitten und begehren. Ich muß mit David sagen, Ps. 143, 2.: „HErr, gehe nicht ins Gericht mit deinem Knechte“ etc., mein Leben ist zu schwach und gebrechlich, und mein Glaube zu gering, so ist die Liebe etwas kalt in mir. Ich bedarf des Mittlers und Gnadenthrons, des HErrn Christi Jesu. Unter diesen Deckel muß ich kriechen, und ihn bitten um Vergebung der Sünden; er muß desjenigen, so noch sündlich und gebrechlich an mir ist, vergessen, daselbstige zudecken, und mir es zu Gnaden kommen lassen. Wer aber sonst bei GOtt wird recht suchen wollen, der wird Rechts genug finden.

\begin{schriftzitat}
V. 11. Und nach den vierzig Tagen und Nächten gab mir der HErr die zwo steinernen Tafeln des Bundes etc.
\end{schriftzitat}

42. Das sind die Tugenden, darauf sie sehen sollen; dies Verzeichniß sollen sie anschauen, und herunter sitzen, hinter sich sehen, und gewahr werden, wer sie gewesen sind. Sie sollen dies Gemälde nicht unter die Bank stecken, sondern vor Augen stellen, und ja ansehen, auf der Kanzel davon predigen, und daraus ihnen lassen aufrücken ihre Untugend, und stets vor die Nase halten lassen, wie sie es von Tugend auf getrieben, und was für böse Stücke sie wider GOtt begangen haben, auf daß ihr Herz davor erschrecken möchte, und das Volk sagen: „HErr, handele nicht mit uns nach unsern Sünden, vergilt uns nicht nach unsern Missethaten“, Ps. 103, 10., sondern sei uns armen Sündern gnädig, wir begehren Gnade und nicht Recht [Dan. 9, 18. 19.]; also würden sie fein demüthig bleiben. Darum, wenn wir den alten Schelmen nicht dämpfen, so vergessen wir auch der Gnade GOttes. Daher rückt Moses den Kindern von Israel immerdar vor die Nase, daß er die Tafeln habe entzwei gebrochen, die GOtt selber geschrieben hatte. Er macht ihre Sünde groß und schwer, daß sie schier verzweifeln mußten, denn er spricht, daß er über ihre Sünde erzürnt sei, und die zwo Tafeln der zehn Gebote zerbrochen habe.

43. Sie möchten aber sagen: Ei, obwohl unsere Väter gesündigt haben, so wollen wir nicht also handeln und sündigen. Ja, sie machten es noch ärger, denn ihre Väter, nachdem sie in das Land eingessen waren. Wir haben auch unter dem Pabstthum mit Abgötterei es arg genug gemacht, und auch also Götzenkälber gemacht, mehr denn sie. Das Volk Israel hat nicht das gegossene Kalb angebetet, sondern sie haben den GOtt Israel bei diesem Kalbe anbeten wollen; unter dem Bilde des Kalbes haben sie GOtt dienen wollen, GOtt zu Ehren haben sie es aufgerichtet. Denn sie haben gewußt, daß im alten Testamente die Altväter und Patriarchen, als, Abraham, Isaak und Jakob, unserm HErrn GOtt Kälber geopfert, darum sind sie zugefahren, haben es verbessern wollen, und dem lebendigen GOtt ein gülden Kalb auch aufgerichtet, damit haben sie dem GOtt dienen wollen, der sie aus Egypten geführt hatte; dem Teufel haben sie gewißlich nicht dienen wollen.

44. Gleichwie wir heutzutage mit unsern scheinbarlichen und gleisenden Werken Gott auch ehren und dienen wollen, und ihn uns gnädig machen. Ich bin für meine Person ins Kloster nicht gelaufen, daß ich dem Teufel dienen wollte, sondern, daß mit meinem Gehorsam, Keuschheit und Armuth ich den Himmel verdienete. Das heißt, ein Kalb aufrichten, und das Kalb anbeten. Aber diese, der Juden Sünde mit dem Kalbe, ist noch eine güldene und köstliche Sünde gewesen gegen unserm Abgott und Sünde. Und wenn wir auch wären allda gewesen, so hätten wir vor großer Andacht und Heiligkeit das Kalb gefressen. Sie haben gewißlich zu Gemüth gezogen, und vor sich im Gedächtniß gehabt die rühmlichen Exempel ihrer
\quellseite{929}
Vorfahren, daß die Patriarchen Kälber zu ihren Opfern gebraucht haben, darum wollten sie ihren Vätern nachfolgen, und unsern HErrn GOtt in diesem Bilde ehren, und der ihnen dieses hat sollen umstoßen und ausreden, der hat müssen keck sein.

45. Wir sind alle so gesinnt und geartet, daß wir unserer Sünden bald vergessen; aber Moses zeucht es wieder an, und sagt zum Volke: Gedenke daran, wie du dich gehalten, und was du mich für Mühe gekostet hast, daß ich deine Sünde abwandte, und euch mit GOtt durch mein Gebet wieder versöhnte, seinen grimmigen und entbrannten Zorn aufhielt.

\begin{schriftzitat}
V. 25.
\end{schriftzitat}

Und fiel vor den HErrn, wie zuerst, vierzig Tage und vierzig Nächte, und aß kein Brod, und trank kein Wasser etc.

46. Im andern Buch Mosis [Cap. 32, 20.] habt ihr gehört, wie Moses das güldene Kalb zu Pulver verbrannt, und Wasser in den Staub gegossen, auch den Juden, die das Kalb angebetet hatten, den Staub zu trinken gegeben habe, damit sie sähen, welch einen schwachen, losen und ohnmächtigen Gott sie gedient und angebetet hätten. Das that Moses ihnen zur großen Schande, daß sie ihren Gott mußten fressen, und in sich saufen; er macht ihren Gottesdienst gar zu Staub und zu Asche, vertilgt ihn gänzlich, daß er nicht ein Stücklein einer Erbse groß von diesem Kalbe hätte lassen überbleiben. Daraus wir lernen, daß man nicht einen eigenen Gottesdienst anrichte, noch darauf troste und sich verlasse; denn eigene Gerechtigkeit ist große Sünde.

47. Und daß Moses den Staub auch nicht hat lassen übrig bleiben, damit zeigt er an den großen Unwillen und ernsten Zorn GOttes darüber, wenn wir eigene Gottesdienste anrichten.

48. Wir sollen noch auf diesen Tag auch also thun, und nicht aufhören, des Papsts Abgötterei und falsche Gottesdienste und Mißbräuche zu zerstören. Wir müssen dem Papste und seinem Reiche fluchen, und dasselbige lästern und schänden, und das Maul nicht zuthun, sondern ohne Aufhören dawider predigen. Denn etliche geben jetzt vor, wir können anders nichts, denn den Papst und die Seinen verdammen, schelten und lästern. Ja, das kann nicht anders sein; denn, alsbald man der Irrthümer vergißt, so ist die Gnade GOttes auch vergessen, und wird die angebotene Gnade verachtet; darum muß man immerdar daran gedenken, und dawider predigen. Denn GOtt ist von Herzen den eigenen, erwählten Andachten und Gottesdiensten feind, und will, daß man nicht ein Stäublein davon lasse überbleiben, sondern das Herz rein behalte am Glauben und Vertrauen GOttes, der also sagt: „Ich bin der HErr, dein GOtt.“ Und sollen wissen, daß GOtt uns alles gibt; denn er hat verheißen, er will unser GOtt sein, und will seine Gläubigen alle hier und dort erhalten.

49. Also ist Moses ein feiner Lehrer. Er hat das erste Gebot wohl ausgelegt, und die Leute zu ihrer selber Erkenntniß geführt, und die hoffärtigen, vermessenen Geister gedemüthigt, über das auch allerlei Untugenden vorgeworfen, damit sie etwas Anderes, denn das verheißene Land, verdient hätten. Eben wie wir für das liebe Evangelium uns auch halten und verdienen, daß jetzt unsere Evangelischen siebenmal ärger werden, denn sie zuvor gewesen. Denn nachdem wir das Evangelium gelernt haben, so stehlen, lügen, trügen, fressen und saufen wir, und treiben allerlei Laster. Da Ein Teufel ist bei uns ausgetrieben worden, sind ihrer nun sieben ärgere wieder in uns gefahren, wie das jetzt an Fürsten, Herren, Edelleuten, Bürgern und Bauern zu sehen, wie sie jetzt thun, und sich ohne alle Scheu, ungeachtet GOtt und seine Dräuung, verhalten.\footnote{Anmerkung der Eislebenischen Ausgabe: „Ende dieser Predigten am 4. Sonntage des Advents [den 19. December] Anno 1529. Es hat D. M. Luth. nicht mehr in diesem Buche gepredigt, denn der Reichstag zu Augsburg Anno 1530 drauf gefolget ist.“}
